\documentclass[11pt,a4paper]{article}
\usepackage[margin=25mm,headheight=14pt]{geometry}
\usepackage[T1]{fontenc}
\usepackage{lmodern,microtype,amsmath,amssymb,amsthm,mathtools}
\usepackage{booktabs,longtable,array,multirow,enumitem,graphicx,xcolor,fancyhdr,xurl,needspace}
\usepackage[colorlinks=true,linkcolor=blue!45!black,citecolor=blue!45!black,urlcolor=blue!55!black]{hyperref}
\hypersetup{pdftitle={Proofs, Counterexamples, and Structural Reductions for Written on the Wall II},pdfauthor={KEIO},pdfsubject={An integrated account of research Rounds 1--22},pdfkeywords={graph conjectures, domination, annihilation, induced paths, traceability, degree sequences, reproducible certificates}}
\pagestyle{fancy}\fancyhf{}\lhead{\small WoWII: consolidated results of Rounds 1--22}
\rhead{\small KEIO}\cfoot{\thepage}
\setlength{\parindent}{0pt}\setlength{\parskip}{5pt}
\setlength{\emergencystretch}{2em}
\setcounter{tocdepth}{1}
\setcounter{secnumdepth}{3}
\setlist{nosep,topsep=4pt,leftmargin=20pt}
\newtheorem{theorem}{Theorem}[section]
\newtheorem{lemma}[theorem]{Lemma}
\newtheorem{proposition}[theorem]{Proposition}
\newtheorem{corollary}[theorem]{Corollary}
\newtheorem{conjecture}[theorem]{Conjecture}
\theoremstyle{definition}\newtheorem{definition}[theorem]{Definition}
\theoremstyle{remark}\newtheorem{remark}[theorem]{Remark}
\newtheorem{example}[theorem]{Example}
\newcommand{\res}{\operatorname{res}}
\newcommand{\rad}{\operatorname{rad}}
\newcommand{\diam}{\operatorname{diam}}
\newcommand{\im}{\operatorname{im}}
\newcommand{\ceil}[1]{\left\lceil#1\right\rceil}
\newcommand{\floor}[1]{\left\lfloor#1\right\rfloor}
\newcommand{\NN}{\mathbb N}
\newcommand{\ZZ}{\mathbb Z}
\newcommand{\RR}{\mathbb R}

\begin{document}
% Integrated source fragment: frontmatter.tex
\title{Proofs, Counterexamples, and Structural Reductions\\
for Written on the Wall II\\[6pt]
\large An integrated account of Rounds 1--22 with structural extensions}
\author{KEIO}
\date{26 October 2026\\[3pt]\small Revised integrated edition}
\maketitle

\begin{abstract}
We consolidate the mathematical results in twenty-two rounds of work on
the graph-conjecture catalogue \emph{Written on the Wall II}. Complete
arguments establish the archived statements 434d, 434e, and 298, and
explicit counterexample families refute 394, 392e, 392f, 395a, 382c, 269,
302, 171, 382e, 154, and the specified historical form of 354. The latter
uses coefficient one-half and a union of open neighborhoods. These
fourteen source-qualified conclusions were reached by Round 14; later
rounds strengthen partial results without increasing that count.
This revised edition also proves triangle-free counterexample families
for 394, its validity on bipartite graphs, and an exact bound using an
odd-cycle transversal, together with further independent-domination
corollaries. Further restricted results concern 40, 61, 165, 199, and
384b. An eleven-vertex counterexample refutes the lower-rank quartile
reading of 389a, while the upper-rank quartile-one case is proved;
neither result is counted as an unqualified resolution of 389a.
For statement 133 we prove the original unrounded inequality whenever
average local independence is at most five, with larger ranges at
radii two and three. The proof combines density bounds, radius
classifications, and closed finite extension certificates. For 384b we
show that any failure has a connected representative of order at most
37, with between 11 and 33 vertices of degree at least three. For 199
we establish several split-graph traceability cases and derive strict
Hall and capacity conditions for any residual obstruction; cycle
elimination remains unresolved. Further results concern induced
forests and residue, triangle-free graphs, spanning-tree leaves,
claw-free graphs, even-distance traceability, and degree frequencies.
The paper integrates complete analytic arguments and explicitly
identified finite certificates. Established inputs and adapted
constructions are credited, and the accompanying package records
source integrity, verification programs, and their outputs. No claim
of worldwide priority, formal proof-assistant certification of the
whole paper, or unrestricted resolution of the remaining conjectures
is made.
\end{abstract}

\smallskip
\noindent\textbf{Keywords.} graph conjectures; domination; annihilation
number; induced paths; traceability; degree sequences; computer-assisted
proof; reproducible certificates.

\clearpage
\tableofcontents
\clearpage

\section{Introduction and exact results ledger}
\label{sec:introduction}
\emph{Written on the Wall II} (WoWII) is the graph-conjecture catalogue
associated with E. DeLaVi\~na and Graffiti.pc. Its formulas combine
classical graph invariants with less standard quantities such as
degree-class frequencies, radial-circle orders, and neighborhood
statistics. The precise definitions matter: changing a set's open
neighborhood to its external neighborhood, interpreting a frequency
as a matching number, or rounding an average can change the statement
being tested. We retain the archived formulations and give each
needed definition at its point of use \cite{src:catalogue,src:definitions}.

This paper was prepared from five supplied archive bundles covering
Rounds 1--22. The bundles contain 51 main \LaTeX{} manuscripts together
with proof notes, exact examples, checkers, and certificates. The
extracted source files used here were matched against the archive
bytes by SHA-256. The proofs are reorganized by mathematical theme:
common constructions are proved once, later statements replace
superseded intermediate bounds, and previously separate dependencies
are connected by internal references. Auxiliary notes with proved
statements are included even when they were not separate main papers.
Appendix~\ref{app:provenance} gives the source map and reproduction scope.

\subsection{Fourteen complete conclusions about original formulations}
The affirmative conclusions are the following. Write $i$ for independent
domination, $d_*$ for degeneracy, $B$ for the periphery, $C$ for the center,
$\rho=\max_x d(x,C)$, and $a$ for the annihilation number.
The notation is defined fully in the proof sections.
\begin{center}
\begin{tabular}{>{\raggedright\arraybackslash}p{0.08\linewidth}>{\raggedright\arraybackslash}p{0.60\linewidth}>{\raggedright\arraybackslash}p{0.23\linewidth}}
\toprule
Item&Established inequality and strengthening&Proof\\\midrule
434d&$i(G)\le d_*(\overline G)-c(G[N[B]])+2$.
More generally, $i(G)+p_4(G)-2\le d_*(\overline G)$ for a
maximum distance-four packing $p_4$.&
Theorem~\ref{early:thm:434d-packing};
Corollary~\ref{early:cor:434d}\\[5pt]
434e&$i(G)\le d_*(\overline G)-\rho(G)+2$.
The quantity $\rho$ is maximum distance \emph{to the center}.&
Theorem~\ref{early:thm:434e}\\[5pt]
298&$\gamma_t(G)\le a(G)+\chi_{C_4}(G)$.
The stronger inequality $\gamma_t\le a$ holds whenever a cycle
of length divisible by four occurs.&
Theorem~\ref{early:thm:4k-cycle};
Corollary~\ref{early:cor:298}\\\bottomrule
\end{tabular}
\end{center}
Here $\chi_{C_4}=1$ when there is no ordinary four-cycle, and zero
otherwise. The proof of 298 includes the antecedent annihilation
and total-domination argument with explicit attribution; the global
$a+1$ bound is not presented as a new contribution.

The eleven refutations are witnessed by the following explicit graphs.
Every row is supported by a general construction and an exact invariant
calculation in the indicated section. The displayed order need not be
the smallest possible order of a counterexample.
\begin{center}
\begin{tabular}{rrlp{0.49\linewidth}}
\toprule
Item&Order&Reversed comparison&Location and qualification\\\midrule
394&15&$9>8$&Corollary~\ref{early:cor:394}\\
392e&20&$11>10$&Corollary~\ref{early:cor:392ef}\\
392f&39&$22>21$&Corollary~\ref{early:cor:392ef}\\
395a&33&$24>23$&Proposition~\ref{early:prop:395a}; a unique-mode
witness also has order 34\\
382c&24&$16>15$&Corollary~\ref{early:cor:382c}; an integer-valued
violation avoids rounding ambiguity\\
269&268&$4<5$&Theorem~\ref{early:thm:269-counterexamples}; minimum
only within the specified family\\
302&40&$8>7$&Corollary~\ref{early:cor:302}\\
171&407&$53<54$&Subsection~\ref{late:171}; periphery distances and
even-distance counts use the original conventions\\
382e&21&$4>3$&Subsection~\ref{late:382e}; all Maxine runs have the
same output size\\
154&20&$13<14$&Subsection~\ref{late:154}; the ratio of radial-circle
orders, distinct from item 155\\
354&36&$14<15$&Subsection~\ref{late:354}; coefficient $1/2$ and open
neighborhood union, using a credited Feit construction\\\bottomrule
\end{tabular}
\end{center}
The count is therefore three proofs and eleven refutations of specified
original formulations. It was reached in Round 14. It counts mathematical
conclusions established in this source corpus, not independently certified
new discoveries, journal acceptances, or catalogue updates. Previously
known statements merely used or reproduced, including the separately
recorded items 299, 412d, and 310, are not added to it. Rounds 15--22
add substantial theorems but no further complete original-item conclusion.

\subsection{Structural extensions in the revised edition}
Section~\ref{r11td:extension} allows arbitrary internal graphs of minimum
degree at least two in the 394 spoke construction, gives triangle-free
counterexamples on 15 and 19 vertices and corresponding infinite families,
proves 394 on bipartite graphs, and derives an odd-cycle-transversal bound
attained by both new families. The independent-domination section also
extends the packing inequality to nonempty graphs without isolated vertices
and proves $\diam(G)\leq2(d_*(\overline G)+1-i(G))+3$ for connected
nontrivial graphs. These are stronger structural results for already
counted entries. They do not increase the fourteen-item count.

\subsection{The strongest retained partial results}
The following table records the terminal claims. Their conditions are
part of the conclusions, rather than qualifications to be dropped when
counting resolved items.
\begin{longtable}{>{\raggedright\arraybackslash}p{0.075\linewidth}>{\raggedright\arraybackslash}p{0.64\linewidth}>{\raggedright\arraybackslash}p{0.19\linewidth}}
\toprule
Item&Proved scope&Location\\\midrule
\endfirsthead
\toprule Item&Proved scope&Location\\\midrule
\endhead
40&Graphs with a nonempty vertex cover whose every vertex has at least
two pendant neighbors.&Proposition~\ref{r12support:forty}\\[5pt]
61&Triangle-free graphs with diameter at most six; triangle-free
graphs with feedback vertex number at most one; and triangle-free
graphs whose ordinary matching number is $\lceil\diam/2\rceil$.
The last two regimes give $f-\res\ge\lceil\diam/2\rceil$.
In addition, all support-graph expansions with private leaves in
Theorem~\ref{r12support:main}, without a triangle-free assumption.&
\S\ref{late:equality}, \ref{late:apex}, \ref{late:geodesic},
\ref{r12support:section}\\[5pt]
133&The original unrounded local-independence inequality for
$\ell\le5$, and additionally $r=2,\ell\le7$ and $r=3,\ell\le6$.
Sharp density statements, radius classifications, and four- and
eleven-core theorems are proved on the way.&
Theorem~\ref{local:final}\\[5pt]
155&All nontrivial connected graphs whose maximum spanning-tree
leaf number is at most three.&Theorem~\ref{late:leaves:main}\\[5pt]
165&Maximum degree at most two; or independent maximum-degree set
$M$, $3\le\Delta\le5$, with every vertex outside $M$ of degree
at most two. The latter class satisfies $L\ge2+|N(M)\setminus M|/4$.&
\S\ref{r12leaf:sec165}\\[5pt]
177, 180&All connected claw-free graphs in the original domain.&
Section~\ref{late:claw}\\[5pt]
189&Every deficient independent set forces an explicitly traceable
graph of order $2d_2+1$; every possible counterexample has a perfect
2-matching. Infinite example families exclude several automatic
stronger reductions.&Section~\ref{sec:even-distance}\\[5pt]
199&Split-graph cases include $\kappa\le3$, the stated maximum-column
bound, balanced sides and their padding range, $|C|\le2\kappa$,
equal sides, and Hall equality. The residual case satisfies strict
Hall, an exact degree-two factor condition, strict capacity, and
directed alternating reach. The revision also proves
$\alpha\le\kappa+3$ under the original premise, using the
Tan--Hung classification for the last boundary.&
\S\ref{sec:split-small}--\ref{sec:split-capacity},
\ref{r12:199:gap-three}\\[5pt]
384b&The inequality holds when the number $h$ of degree-at-least-three
vertices is at most 10 or at least 34. If a counterexample exists,
a connected representative has order at most 37 and $11\le h\le33$.
At $h=33$, a counterexample has exactly two cubic vertices with triangle
counts two and three, and the stated attachment restrictions.&
Theorem~\ref{paired:terminal}; \S\ref{r12:384h33}\\[5pt]
389a&Sufficient conditions in terms of degree frequency, the number
of low-degree vertices, and a specified first-quartile convention.
The upper-rank quartile-one case holds. An explicit eleven-vertex
counterexample refutes the lower-rank reading; the source ambiguity
is retained and no unqualified closure is counted.&
\S\ref{late:389a}, \ref{revision:389q1}, \ref{r12:389a:rounding}\\\bottomrule
\end{longtable}
The residue material also includes exact forest equality, matching-padding
and weak-majorization consequences, inheritance under split composition,
and the corresponding unigraph result. These supporting theorems are
included with proofs and are not separate original-item closures.

\subsection{Notation and logical distinctions}
All graphs are finite, simple, and undirected. A connected graph's
order and size are $n$ and $m$. The neighborhood $N(v)$ is open.
Unless a section explicitly says otherwise, $N(A)=\bigcup_{a\in A}N(a)$
is the union of open neighborhoods and can intersect $A$; it does not
automatically mean $N(A)\setminus A$. The closed neighborhood is $N[A]$.
Distance is measured in the ambient graph. The center minimizes vertex
eccentricity, and the periphery maximizes it. The historical catalogue
term ``boundary'' is interpreted according to that periphery definition
where the archived formulation requires it.

We distinguish ordinary domination $\gamma$, total domination $\gamma_t$,
and two-domination $\gamma_2$. An independent dominating set is an
inclusion-maximal independent set, and $i$ minimizes its size.
The ordinary matching number is denoted by $\nu$ or $\mu$ in the
locally defined settings; induced and acyclic matchings are separately
defined. The residue $\res$ comes from the Havel--Hakimi elimination
process. The largest induced forest has order $f$, while $\tau$ in
the traceability sections is the largest \emph{connected} induced-tree
order. The longest induced path $p$ is counted in vertices; diameter
is counted in edges.

For Conjecture 133, the local-independence sum is denoted by $\Lambda$
and its unrounded mean by $\ell$. This avoids confusing that sum with
spanning-tree leaves or low-degree vertex sets. Notation used only in
a construction, such as its parameter $r$ or its low-degree set $L$,
is explicitly reset there. No proof passes between these meanings
without a new definition.

\subsection{Proof organization and verification}
The affirmative independent-domination results precede the shared
two-domination modules and the total-domination arguments. The remaining
counterexample families and structural leaf and residue results follow.
The longer partial programs for 133, 384b, and 199 are presented as
connected proof chains, and the even-distance reduction for 189 closes
the mathematical development.

An analytic proof is distinguished from a finite supplement to it.
For the $P_8$ and $P_9$ core classifications and the exact finite kernels
and rational inequalities in 384b, the specified certificate is itself
a proof dependency. The accompanying verifier checks its complete
coverage and exact arithmetic. Construction programs, canonical-labeling
heuristics, and exploratory bounded searches are not silently promoted
to proof dependencies. In particular, the order-37 result for 384b is
an existence-of-a-representative reduction; it is neither an assertion
that every counterexample has that order bound nor a completed census
through that order. For 199, a factor containing cycles is explicitly
distinguished from a Hamiltonian path.


% Integrated source fragment: closed_early.tex
% Complete, audited consolidation of the ten original items settled in R01--R07.
% All local labels and bibliography keys have the prefix early:.

\section{Independent domination, complement degeneracy, and the center}
\label{early:sec:independent}

We begin with the two affirmative results numbered 434d and 434e.  For a
graph $F$, let $i(F)$ be the minimum cardinality of an inclusion-maximal
independent set, and put $i(\varnothing)=0$.  For a nonempty graph $F$ write
\[
 d_*(F)=\max_{\varnothing\ne X\subseteq V(F)}\delta(F[X]).
\]
Thus $d_*$ is degeneracy itself.  In particular, the catalogue's
Szekeres--Wilf quantity in these two statements is not degeneracy plus one.
For connected $G$, write $C(G)$ and $B(G)$ for its center and periphery, and
define
\[
 \rho(G)=\max_{x\in V(G)} d_G(x,C(G)),\qquad
 d_G(x,C(G))=\min_{c\in C(G)}d_G(x,c).
\]
The set eccentricity $\rho(G)$ is different from the radius.  It is zero
when every vertex is central, and always satisfies $\rho(G)\leq\rad(G)$.
For a vertex set $A$, $N[A]$ denotes its closed neighborhood, and
$\operatorname{comp}(F)$ is the number of connected components of $F$.

With these conventions, the two original inequalities, stated for connected
graphs of order greater than three, are
\begin{align}
 i(G)&\leq d_*(\overline G)-
       \operatorname{comp}(G[N[B(G)]])+2,\label{early:eq:434d-original}\\
 i(G)&\leq d_*(\overline G)-\rho(G)+2.
       \label{early:eq:434e-original}
\end{align}
These conventions agree with the preserved catalogue and with the
independent-domination paper of DeLaVi\~na, Pepper, and
Waller~\cite{src:catalogue,early:dpw-independent}.  Both inequalities will be
proved for all connected graphs of order at least two.

\subsection{Distance-four packings: statement 434d}

Let $p_4(G)$ be the maximum cardinality of a vertex set whose distinct
members have pairwise distance at least four.

\begin{theorem}[Packing strengthening of 434d]
\label{early:thm:434d-packing}
Every connected graph $G$ of order at least two satisfies
\[
 i(G)+p_4(G)-2\leq d_*(\overline G).
\]
\end{theorem}
\begin{proof}
Put $H=\overline G$ and $k=i(G)$.  Every maximal clique of $H$ has at
least $k$ vertices.  If $p_4(G)=1$, a maximal clique immediately gives
$d_*(H)\geq k-1$, which is the required inequality.

Now take $S=\{s_1,\ldots,s_c\}$ with pairwise distances at least four,
where $c\geq2$.  For each $j$, choose a neighbor $t_j$ of $s_j$ in $G$.
The closed neighborhoods $N_G[s_j]$ are disjoint, and there are no
$G$-edges between two of them: any such edge would give a path of length
at most three between the corresponding members of $S$.  Consequently
$T=\{t_1,\ldots,t_c\}$ and $S$ are cliques of $H$, $s_it_j$ is an
$H$-edge for $i\ne j$, and $s_jt_j$ is not an $H$-edge.

For a clique $K$ of $H$, put
\[
 C_H(K)=\{x:K\subseteq N_H[x]\},\qquad
 U=\bigcup_{j=1}^c C_H(T\setminus\{t_j\}).
\]
In particular $S\subseteq U$.  We show that every vertex of $H[U]$ has
degree at least $k+c-2$.  Fix $x\in U$.  There is an index $j$ such that
\[
 T\setminus\{t_j\}\subseteq N_H[x],\qquad
 S\setminus\{s_j\}\subseteq N_H(x).
\]
Indeed, if $x\in N_G[s_j]$ for some $j$, the disjointness and absence of
edges between these closed neighborhoods give both inclusions for that
index.  If $x$ lies outside their union, it is $H$-adjacent to all of $S$,
and any index witnessing $x\in U$ suffices.

Extend the clique $\{x\}\cup(T\setminus\{t_j\})$ to a maximal clique
$Q$ of $H$.  Then $|Q|\geq k$ and $Q\subseteq U$.  None of the vertices
$s_i$ with $i\ne j$ belongs to $Q$, because $t_i\in Q$ and $s_it_i$ is
not an $H$-edge.  Thus $x$ has the disjoint sets $Q\setminus\{x\}$ and
$S\setminus\{s_j\}$ among its neighbors in $H[U]$.  This gives
$d_{H[U]}(x)\geq(k-1)+(c-1)$.  The definition of degeneracy now proves
the result with $c=p_4(G)$.
\end{proof}

\begin{corollary}[Statement 434d]
\label{early:cor:434d}
For every nonempty $A\subseteq V(G)$ in a connected graph of order at
least two,
\[
 i(G)\leq d_*(\overline G)-
 \operatorname{comp}(G[N[A]])+2.
\]
In particular \eqref{early:eq:434d-original} holds.
\end{corollary}
\begin{proof}
Choose one vertex of $A$ from each component of $G[N[A]]$.  Each component
contains a member of $A$.  Representatives from different components
cannot be at distance at most three: a path of length at most three
between two vertices of $A$ lies entirely in $N[A]$.  They therefore
form a distance-four packing.  Apply Theorem~\ref{early:thm:434d-packing}
and then take $A=B(G)$.
\end{proof}

\begin{example}[Sharpness]
\label{early:ex:434d-sharp}
Let $c\geq2$, and subdivide every edge of $K_{1,c}$ once.  Denote its
vertices by $z,t_1,s_1,\ldots,t_c,s_c$, with edges $zt_j,t_js_j$.
Every independent dominating set either omits $z$ and uses at least one
of $t_j,s_j$ for each $j$, or contains $z$ and all $s_j$.  The set of all
$t_j$ is independent and dominating, so $i(G)=c$.  The leaves $s_j$ are
pairwise at distance four, while no other vertex can coexist in a
distance-four packing with another vertex; hence $p_4(G)=c$.
In the complement, deleting $z$ leaves the complement of $cK_2$, a
$(2c-2)$-regular graph.  The complement degree of $z$ is $c\leq2c-2$,
so this also proves $d_*(\overline G)=2c-2$.  The periphery consists of
the $s_j$, and its closed neighborhood induces $cK_2$ in $G$.  Thus both
the packing bound and statement 434d hold with equality.
\end{example}

\subsection{Minimal vertex covers and distance to the center}

For a graph $F$, define
\[
 q(F)=\max\{|Q|:Q\text{ is an inclusion-minimal vertex cover of }F\}.
\]
Minimality is by inclusion, not by cardinality.  Complements interchange
minimal vertex covers and maximal independent sets, so
\begin{equation}
 q(F)=|V(F)|-i(F),\qquad q(\varnothing)=0.
 \label{early:eq:upper-cover}
\end{equation}

\begin{lemma}
\label{early:lem:boundary-cover}
If $b$ is a nonisolated vertex of $F$, some minimal vertex cover of $F$
contains $b$.
\end{lemma}
\begin{proof}
Extend a neighbor of $b$ to a maximal independent set.  Its complement
is the required cover.
\end{proof}

\begin{lemma}
\label{early:lem:cover-tree}
If $F$ is connected and $Q\ne\varnothing$ is a vertex cover, there is a
tree $T\subseteq F$ containing $Q$, all of whose leaves belong to $Q$,
such that $\diam(T)\leq2|Q|-2$.  A one-vertex tree is allowed.
\end{lemma}
\begin{proof}
Choose a tree containing $Q$ with the fewest vertices.  Deleting a leaf
outside $Q$ would contradict this choice.  Vertices outside $Q$ are
independent.  A longest path therefore has endpoints in $Q$, at most
$|Q|$ vertices in $Q$, and at most $|Q|-1$ vertices outside $Q$.
\end{proof}

\begin{theorem}[A local cover bound]
\label{early:thm:center-local}
Let $G$ be connected of order at least two, let $v\in V(G)$, and put
$R=G-N_G[v]$.  Then
\begin{equation}
 \rho(G)\leq\max\{2,q(R)+1\}.
 \label{early:eq:center-local}
\end{equation}
If $q(R)>0$, then for every prescribed $u\in V(G)$ there is a tree
subgraph $T\subseteq G$ whose vertex set dominates $G$, which contains
$u$, and which satisfies $\diam(T)\leq2q(R)+2$.
\end{theorem}
\begin{proof}
Put $A=N_G(v)$ and $q_R=q(R)$.  If $q_R=0$, $R$ is independent.  Each
vertex of $R$ has a neighbor in $A$, by connectivity, so
$\operatorname{ecc}_G(v)\leq2$ and $\rho(G)\leq\rad(G)\leq2$.

Suppose $q_R\geq1$.  Let $F_1,\ldots,F_s$ be the nontrivial components
of $R$.  Choose in each $F_j$ a vertex $b_j$ adjacent to some $a_j\in A$,
and a minimal vertex cover $Q_j$ containing $b_j$ by
Lemma~\ref{early:lem:boundary-cover}.  Their union is a minimal vertex
cover of $R$.  If $q_j=|Q_j|$ and $q=\sum_jq_j$, then $1\leq q\leq q_R$.
Choose a tree $T_j$ as in Lemma~\ref{early:lem:cover-tree}.  Form $T_0$
from the star with center $v$ and leaves $A$, all the $T_j$, and the
single attachment edges $a_jb_j$.  This is a tree, even when some $a_j$
coincide.  Its vertex set dominates $G$: it contains $v,A$, and the covers
$Q_j$, while an isolated vertex of $R$ has a neighbor in $A$.

Two vertices in distinct $T_j,T_\ell$ have distance at most
$(2q_j-2)+4+(2q_\ell-2)\leq2q$.  A vertex of $T_j$ is at distance at
most $2q_j+1$ from any vertex of the star.  Hence
$\diam(T_0)\leq2q+1$.  If $u\notin T_0$, it lies in $R$; attach it as
a single leaf to a neighbor in $\bigcup_jQ_j$ if it is nonisolated in
$R$, and to a neighbor in $A$ otherwise.  The resulting dominating tree
$T$ contains $u$ and has diameter at most $2q_R+2$.

A center $z$ of $T$ has tree eccentricity at most $q_R+1$.  Because $T$
dominates $G$,
\[
 \operatorname{ecc}_G(z)\leq q_R+2,
 \qquad d_G(u,z)\leq q_R+1.
\]
Thus $\rad(G)\leq q_R+2$.  If the radius is at most $q_R+1$, then the
same is true of $\rho(G)$.  In the remaining integer case
$\rad(G)=q_R+2$, the vertex $z$ is a center of $G$.  Since $u$ was
arbitrary, every vertex is within $q_R+1$ of a center, proving the result.
\end{proof}

\begin{corollary}
\label{early:cor:center-degree}
Every connected graph $G$ of order $n\geq2$ satisfies
\[
 \rho(G)\leq\max\{2,n-\Delta(G)-i(G)+1\}
 =\max\{2,\delta(\overline G)-i(G)+2\}.
\]
\end{corollary}
\begin{proof}
Choose $v$ of maximum degree.  If $R=G-N_G[v]$, adjoining $v$ to an
independent dominating set of $R$ gives one of $G$.  Therefore
$i(R)\geq i(G)-1$ and
$q(R)=n-\Delta(G)-1-i(R)\leq n-\Delta(G)-i(G)$.  Apply
Theorem~\ref{early:thm:center-local}.
\end{proof}

\subsection{The equality boundary and statement 434e}

The following boundary argument was already present in the supplied
baseline~\cite{early:baseline}.  We include its proof because it is needed
to complete the unrestricted result.  The local cover theorem above is the
ingredient that treats the remaining positive-slack cases.

\begin{lemma}
\label{early:lem:nested}
Suppose every maximal clique of a $(k-1)$-degenerate graph $H$ has order
$k$, and $G=\overline H$ is connected of order at least two.  Put
\[
 M=\{m:d_H(m)=k-1\},\qquad W=V(G)\setminus N_G[M].
\]
The sets $N_G(w)$, for $w\in W$, are linearly ordered by inclusion.
Consequently, if $W\ne\varnothing$, its members have a common neighbor.
\end{lemma}
\begin{proof}
Every vertex of $H$ lies in a $k$-clique.  Thus $\delta(H)\geq k-1$,
with equality by degeneracy, and $M\ne\varnothing$.  For $m\in M$,
$N_H[m]$ is a $k$-clique.  The set $W$ is disjoint from $M$, is a clique
of $H$, and every $m\in M$ is $H$-adjacent to all of $W$.

For nonempty $P\subseteq W$, let
\[
 L=\{x\notin P:P\subseteq N_H(x)\},\quad
 Z=V(H)\setminus(P\cup L),\quad
 L_0=\{x\in L:N_H(x)\cap Z=\varnothing\}.
\]
We have $M\subseteq L_0$: a neighbor in $Z$ of a simplicial $m\in M$
would be adjacent to every vertex of $P$, contrary to its membership in
$Z$.  Suppose each member of $P$ has an $H$-neighbor in $Z$.  In a
degeneracy ordering of $H$, let $x$ be the first vertex outside $L_0$.
If $x\in Z$, it has no earlier neighbor and, as $x\notin M$, still has
degree at least $k$, a contradiction.  Otherwise
$x\in P\cup(L\setminus L_0)$ and has a neighbor $z\in Z$.  Extend
$xz$ to a maximal clique $K$, of order $k$.  None of $K$ belongs to
$L_0$, because its vertices are adjacent to $z$; all of $K$ therefore
survives.  The remaining closed neighborhood of $x$ has at most $k$
vertices and contains $K$, so it equals $K$.  It also contains all of
$P$.  This forces $z$ to be adjacent to all of $P$, a contradiction.

Some $p\in P$ consequently has no $H$-neighbor in $Z$.  Since $P$ is
independent in $G$, we have $Z=\bigcup_{w\in P}N_G(w)$, and hence
$N_G(p)=\bigcup_{w\in P}N_G(w)$.  Apply this conclusion to every pair
in $W$ to obtain pairwise comparability.  Connectivity makes every such
neighborhood nonempty, and a smallest one provides a common neighbor.
\end{proof}

\begin{theorem}[Previously established zero-slack boundary]
\label{early:thm:zero-slack}
If $G$ is connected of order at least two and
$i(G)=d_*(\overline G)+1$, then $C(G)$ is a dominating set and
$\rho(G)\leq1$.
\end{theorem}
\begin{proof}
Put $H=\overline G$ and $k=i(G)$.  Every maximal clique of $H$ has at
least $k$ vertices, while $(k-1)$-degeneracy forbids a larger clique.
If $k=1$, $G$ is complete.  Otherwise apply
Lemma~\ref{early:lem:nested}.

For each $m\in M$, the vertices nonadjacent to $m$ in $G$ form an
independent set.  Each such vertex has a neighbor, which must itself be
adjacent to $m$.  Thus $\operatorname{ecc}_G(m)\leq2$.  There is no
universal vertex because $i(G)\geq2$, so $\rad(G)=2$ and
$M\subseteq C(G)$.  If $W=\varnothing$, these centers dominate $G$.
Otherwise let $y$ be a common neighbor of $W$.  For any $m\in M$ and
$w\in W$, the edge $wy$ forces $my\in E(G)$, by the independence just
noted; also $y\ne m$.  Hence $y$ is adjacent to all of $M$, is at
distance at most two from every vertex of $N_G[M]$, and is adjacent to
every vertex of $W$.  Therefore $y\in C(G)$.  The centers $M\cup\{y\}$
dominate $G$.
\end{proof}

\begin{theorem}[Statement 434e]
\label{early:thm:434e}
Every connected graph $G$ of order at least two satisfies
\[
 i(G)\leq d_*(\overline G)-\rho(G)+2.
\]
\end{theorem}
\begin{proof}
Write $d=d_*(\overline G)$ and $h=d+1-i(G)$.  A maximal clique of
$\overline G$ has at least $i(G)$ vertices, so $h\geq0$.  If $h=0$,
Theorem~\ref{early:thm:zero-slack} gives $\rho(G)\leq1=h+1$.
If $h\geq1$, Corollary~\ref{early:cor:center-degree} gives
\[
 \rho(G)\leq\max\{2,\delta(\overline G)-i(G)+2\}
 \leq\max\{2,h+1\}=h+1.
\]
Rearranging proves the assertion.
\end{proof}

\begin{example}[Sharpness and the two eccentricities]
For $P_7$, an independent dominating set in positions $1,4,7$ has size
three, while two closed neighborhoods cover at most six vertices.  Thus
$i(P_7)=3$.  Its complement has minimum degree four; deleting the middle
vertex first leaves the complement of $2P_3$, of maximum degree four.
Consequently $d_*(\overline{P_7})=4$.  The unique center has eccentricity
three, so $\rho(P_7)=3$ and statement 434e is sharp with positive slack.

For the subdivided stars of Example~\ref{early:ex:434d-sharp}, taking
$v=z$ gives $q(G-N[v])=0$ but $\rho(G)=2$.  The constant two in
\eqref{early:eq:center-local} is therefore necessary.  Radius cannot
replace $\rho$ in statement 434e either.  The tree with vertex set
$\{0,\ldots,6\}$ and edges
$01,13,14,26,36,45$ has $i=3$: $\{1,2,5\}$ is an independent
dominating set.  A dominating pair would have to contain vertex $1$,
since every other closed neighborhood has size at most three.  But
$1$ dominates neither $2$ nor $5$, and no second vertex can dominate
both, as their distance is five.  Its complement has
degeneracy three, as witnessed below by its minimum degree and above by
the deletion order $1,6,4$, followed by the four remaining vertices.
Its center is $\{1,3\}$, its radius is three, and $\rho=2$.
\end{example}


\subsection{Further consequences of the packing and cover arguments}
The following two consequences extend the preceding arguments. They do
not add catalogue identifiers to the fourteen-item ledger.
\begin{corollary}[Disconnected graphs without isolated vertices]
Let $G$ be a nonempty finite simple graph without isolated vertices, and put the
distance between vertices in different components equal to infinity.
Then
\[
 i(G)+p_4(G)-2\leq d_*(\overline G).
\]
Consequently, for every nonempty $A\subseteq V(G)$,
\[
 i(G)\leq d_*(\overline G)-c(G[N[A]])+2.
\]
\end{corollary}
\begin{proof}
The complement-core proof of Theorem~\ref{early:thm:434d-packing} uses
connectivity only to choose a neighbor of each packing vertex. The absence
of isolated vertices supplies exactly that condition. Distinct closed
neighborhoods of distance-four-separated vertices are still disjoint and
anticomplete, including vertices in different components. All clique
extensions and degree counts therefore apply unchanged. Selecting one
member of $A$ from each component of $G[N[A]]$ proves the second assertion.
\end{proof}

\begin{corollary}[Diameter controlled by the degeneracy slack]
For every connected graph $G$ of order at least two, put
$h=d_*(\overline G)+1-i(G)$. Then
\[
 \operatorname{diam}(G)\leq 2h+3.
\]
\end{corollary}
\begin{proof}
The packing bound gives $p_4(G)\leq h+1$. If $h=0$, no two vertices
are at distance at least four, giving the result. Suppose $h\geq1$.
Choose a maximum-degree vertex $v$ and let $R=G-N[v]$.
Corollary~\ref{early:cor:center-degree} gives
$q(R)\leq n-\Delta(G)-i(G)\leq h$.
If $q(R)=0$, then $\operatorname{ecc}_G(v)\leq2$, so the diameter
is at most four. Otherwise, the base dominating tree $T_0$ in the proof of Theorem~\ref{early:thm:center-local}
has diameter at most $2q(R)+1$. Each graph vertex belongs to $T_0$ or
has a neighbor in $T_0$, and consequently
$\operatorname{diam}(G)\leq\operatorname{diam}(T_0)+2
\leq2q(R)+3\leq2h+3$.
\end{proof}

\section{Rooted constructions for two-domination counterexamples}
\label{early:sec:two-domination}

A two-dominating set $D$ requires every vertex outside $D$ to have at
least two neighbors in $D$; there is no condition on vertices of $D$.
Its minimum size is $\gamma_2(G)$.  We use $\nu(G)$ for the matching
number, corresponding to $\mu$ in the original rows.  Put
\[
 A_\delta(G)=\{v:d_G(v)=\delta(G)\},\qquad
 A_2(G)=\{v:d_G(v)=2\},\qquad
 \lambda_G(v)=\alpha(G[N_G(v)]).
\]
In particular, $A_2$ here means degree exactly two.  The constructions
below refute statements 394, 392e, 392f, 395a, and 382c under their
original definitions~\cite{src:catalogue,early:glossary,early:dpw-two}.
All five refutations have unbounded additive deficit.  The stated small
witnesses are explicit examples, not assertions of global minimal order.

\subsection{A clique module and statement 394}

Let
$M_{\min\lambda}(G)=\{v:\lambda_G(v)=\min_x\lambda_G(x)\}$.
The original statement 394 is
\begin{equation}
 \gamma_2(G)\leq \nu(G)+|E(G[A_2(G)])|+|M_{\min\lambda}(G)|.
 \label{early:eq:394-original}
\end{equation}
The final set is denoted by a capital $M$ in the original row; reading
it as a second matching-number symbol changes the proposed inequality.

\begin{definition}
\label{early:def:394-family}
Fix $r\geq2$, $k\geq3$, and a connected graph $H$ on
$\{c_1,\ldots,c_r\}$.  For each $i$, introduce a $k$-clique
$Q_i=\{q_{i1},\ldots,q_{ik}\}$ and an independent set
$T_i=\{t_{i1},\ldots,t_{ik}\}$.  Add the edges
$c_it_{ij}$ and $t_{ij}q_{ij}$ for all $i,j$.  Finally attach a leaf $p$
to $c_1$.  There are no other edges except the edges of $H$ and the
cliques $Q_i$.  Write $X(H,k)$ for this graph and $H^+$ for $H$ with
the leaf $p$ attached.
\end{definition}

\begin{proposition}
\label{early:prop:394-invariants}
The graph $X=X(H,k)$ is connected and satisfies
\begin{align*}
 |V(X)|&=r(2k+1)+1,&
 |E(X)|&=|E(H)|+r\bigl(2k+\tbinom{k}{2}\bigr)+1,\\
 \gamma_2(X)&=r(k+1)+1,&
 \nu(X)&=rk+\nu(H^+),\\
 E(X[A_2(X)])&=\varnothing,&
 M_{\min\lambda}(X)&=\{p\}.
\end{align*}
\end{proposition}
\begin{proof}
The order, size, and connectivity follow directly from the construction.
Every two-dominating set contains $p$.  Consider the vertices
$\{c_i\}\cup T_i\cup Q_i$.  If $c_i$ is selected, each pair
$\{t_{ij},q_{ij}\}$ contributes at least one selected vertex, since an
unselected $t_{ij}$ needs both its neighbors selected.  This gives at
least $k+1$ vertices in the module.  If $c_i$ is unselected, all of
$T_i$ must be selected.  At least one member of $Q_i$ must also be
selected, since otherwise each vertex of $Q_i$ has only its one selected
neighbor in $T_i$.  This again gives $k+1$.  Conversely, $p$, all the
centers $c_i$, and all the vertices in all the $Q_i$ form a
two-dominating set, proving the formula for $\gamma_2$.

The $rk$ edges $t_{ij}q_{ij}$ together with a maximum matching of $H^+$
give a matching of size $rk+\nu(H^+)$.  In any matching, the edges
outside $H^+$ use the $2k+1$ vertices of one module and therefore
contribute at most $k$ per module.  The remaining matching edges belong
to $H^+$, proving the reverse bound.

The only degree-two vertices are those in the $T_i$, and they form an
independent set.  The local independence numbers are
\[
 \lambda_X(p)=1,\quad \lambda_X(t_{ij})=2,\quad
 \lambda_X(q_{ij})=2,\quad \lambda_X(c_i)\geq k.
\]
For $q_{ij}$ its neighborhood is a clique of order $k-1$ together with
the isolated vertex $t_{ij}$; for $c_i$ the set $T_i$ is independent.
These observations prove the last assertion.
\end{proof}

\begin{corollary}[Refutation of 394]
\label{early:cor:394}
Every graph in Definition~\ref{early:def:394-family} violates
\eqref{early:eq:394-original}, with deficit $r-\nu(H^+)>0$.
For $r=2$, $k=3$, and $H=K_2$, it has order $15$, size $20$,
$\gamma_2=9$, and proposed right side $7+0+1=8$.
\end{corollary}
\begin{proof}
Subtract the right side from the left using
Proposition~\ref{early:prop:394-invariants}.  Since $H^+$ has $r+1$
vertices, $\nu(H^+)\leq\floor{(r+1)/2}$, and the deficit is at least
$\floor{r/2}$.  Alternatively choose $H$ to be the star centered at
$c_1$ and fix $k=3$.  Then $H^+$ is also a star, so
$\nu(H^+)=1$.  This yields graphs of order $7r+1$ with
$\gamma_2=4r+1$ and proposed right side $3r+2$, an unbounded deficit
$r-1$.
\end{proof}

\subsection{One nine-vertex module for 392e, 392f, and 395a}

The next three counterexamples use the same module; consolidating it
avoids repeating its local two-domination and matching arguments.

\begin{definition}
\label{early:def:nine-module}
Let $F$ have vertices $c,t_1,t_2,q_1,q_2,w_1,w_2,w_3,w_4$ and edges
\[
 ct_1,ct_2,t_1q_1,t_2q_2,
 w_1w_2,w_2w_3,w_3w_4,w_4w_1,
 q_1w_1,q_1w_3,q_2w_2,q_2w_4.
\]
The vertex $c$ is its root.  When copies of $F$ are used, only roots
receive neighbors outside their own copies.
\end{definition}

\begin{lemma}[Local cost of the module]
\label{early:lem:nine-module}
In any graph containing such a copy of $F$, every two-dominating set
contains at least five of its nine vertices.  The set
$\{c,q_1,q_2,w_1,w_3\}$ two-dominates the module.  The module also has
a matching of size four avoiding its root.
\end{lemma}
\begin{proof}
Put $U=\{w_1,w_3\}$ and $W=\{w_2,w_4\}$.  At least two of the four
$w$-vertices must be selected.  With none selected each has at most its
one $q$-neighbor selected; with just one selected, the other vertex in
the same part $U$ or $W$ has the same obstruction.

If $c$ is selected, each pair $\{t_j,q_j\}$ contributes at least one
selected vertex, so there are at least five in all.  If $c$ is not
selected, both $t_j$ are forced.  A selection of at most four module
vertices would therefore consist of the two $t_j$ and exactly two
$w$-vertices, with no $q_j$.  The two unselected $q_j$ then require at
least one selected vertex in each of $U,W$.  An unselected $w$-vertex
would consequently see just one selected vertex in the opposite part
and no selected $q$-neighbor, a contradiction.

The displayed five-set works by direct inspection of the unselected
vertices.  Finally
$t_1q_1,t_2q_2,w_1w_2,w_3w_4$ is the asserted matching.
\end{proof}

The original statements 392e and 392f are, respectively,
\begin{align}
 \gamma_2(G)&\leq \nu(G)+\nu(G[N(A_2(G))])+|A_\delta(G)|,
 \label{early:eq:392e-original}\\
 \gamma_2(G)&\leq \nu(G)+\nu(G[N(A_2(G))])+
 \bigl|\{v:|N(v)\cap N(A_\delta(G))|=1\}\bigr|.
 \label{early:eq:392f-original}
\end{align}
Neighborhoods in these formulas are open neighborhood unions.

\begin{definition}
\label{early:def:392-families}
Take $r\geq2$ disjoint rooted copies $F_i$ of $F$, and join their roots
$c_i$ to a new vertex $x$.  For $G_r$, add the single edge $xp$ to a
new leaf $p$.  For $J_r$, instead add the path $xsp$ on two new
vertices $s,p$.
\end{definition}

\begin{proposition}
\label{early:prop:392-table}
The two families have the following invariants, where
$B_1=\{v:|N(v)\cap N(A_\delta)|=1\}$:
\[
\begin{array}{c|cc}
 &G_r&J_r\\\hline
 |V|&9r+2&9r+3\\
 |E|&13r+1&13r+2\\
 \gamma_2&5r+1&5r+2\\
 \nu&4r+1&4r+2\\
 |A_\delta|&1&1\\
 \nu(G[N(A_2)])&0&1\\
 |B_1|&r+1&2
\end{array}
\]
\end{proposition}
\begin{proof}
The leaf $p$ is forced in a two-dominating set.  Each module contributes
at least five by Lemma~\ref{early:lem:nine-module}.  For $G_r$, use $p$
and the indicated five-set in every module; the unselected hub sees all
selected roots, and the lower bound is attained.  In $J_r$, the path
requires at least one of $x,s$ in addition to $p$: otherwise $s$ has
only one selected neighbor.  Choosing $p,x$ and the five-sets attains
$5r+2$.

Each module contributes at most four matching edges.  The edges outside
the modules all meet $x$ in $G_r$, and all meet $\{x,s\}$ in $J_r$.
Thus the matching upper bounds are $4r+1$ and $4r+2$.  They are
attained by the root-avoiding four-matchings, together with $xp$ in
$G_r$, or $xc_1,sp$ in $J_r$.

The unique minimum-degree vertex is $p$.  The degree-two vertices of
$G_r$ are exactly the $t$-vertices.  Their open neighborhood consists
of roots and $q$-vertices, an independent set.  In $J_r$, the vertex
$s$ also has degree two.  Its neighbors add $x,p$ to this set, whose
induced graph is a star with center $x$ and leaves $c_i$, together
with isolated vertices.  Its matching number is one.  Finally
$N(A_\delta)=\{x\}$ for $G_r$, giving $B_1=N(x)$ of size $r+1$;
for $J_r$ it is $\{s\}$, giving $B_1=\{x,p\}$.
\end{proof}

\begin{corollary}[Refutations of 392e and 392f]
\label{early:cor:392ef}
Statement 392e fails on $G_r$ for every $r\geq2$, with deficit $r-1$,
and on $J_r$ for every $r\geq3$, with deficit $r-2$.
Statement 392f fails on $J_r$ for every $r\geq4$, with deficit $r-3$.
The displayed deficit formulas remain valid at the smaller defined values
of $r$, but a zero or negative deficit is not a counterexample.
In particular, $G_2$ has order $20$, size $27$, and
$\gamma_2=11>9+0+1=10$.  The graph $J_4$ has order $39$, size $54$,
and $\gamma_2=22>18+1+2=21$ for 392f.
\end{corollary}
\begin{proof}
Substitute the invariant values in Proposition~\ref{early:prop:392-table}.
The deficits are unbounded as $r$ increases.
\end{proof}

For statement 395a, let the \emph{mode-min degree} be the smallest degree
whose frequency is maximum, and let $\mathcal M(G)$ be its degree
class.  The proposed inequality is
\begin{equation}
 \gamma_2(G)\leq |\mathcal M(G)|+
 \nu(G-A_\delta(G))+\nu(G[N(A_2(G))]).
 \label{early:eq:395a-original}
\end{equation}
The cardinality of the mode class is independent of the choice among
tied modal degrees; the following construction also supplies witnesses
with a unique modal degree.

\begin{proposition}[Refutation of 395a]
\label{early:prop:395a}
Take $r\geq2$ copies of $F$, join their roots to a hub $x$, and attach
$L\geq7r$ leaves to $x$.  Denote the resulting graph by $K_{r,L}$.
Then
\begin{align*}
 |V(K_{r,L})|&=9r+1+L,& |E(K_{r,L})|&=13r+L,\\
 \gamma_2(K_{r,L})&=5r+L,&
 |\mathcal M(K_{r,L})|&=L,\\
 \nu(K_{r,L}-A_\delta)&=4r+1,&
 \nu(K_{r,L}[N(A_2)])&=0.
\end{align*}
Consequently the deficit in \eqref{early:eq:395a-original} is $r-1$.
\end{proposition}
\begin{proof}
All $L$ leaves are forced, and each module costs five.  Conversely the
leaves and the five-sets from Lemma~\ref{early:lem:nine-module} work,
including at the unselected hub.  The degree frequencies are
\[
\begin{array}{c|rrrr}
 \text{degree}&1&2&3&L+r\\\hline
 \text{frequency}&L&2r&7r&1.
\end{array}
\]
Here $A_\delta$ is the leaf set and the mode-min class is that same
set.  The mode is unique if $L>7r$.  Deleting the leaves leaves $r$
modules and the star edges at $x$.  At most one star edge belongs to a
matching, so its matching number is at most $4r+1$, attained by $xc_1$
and the four root-avoiding edges in each module.  As before, the open
neighborhood of the degree-two vertices consists of independent roots
and $q$-vertices, giving the final zero.  Subtraction proves the deficit.
\end{proof}

For $(r,L)=(2,14)$ this gives order $33$, size $40$, and
$\gamma_2=24>14+9=23$.  The choice $(2,15)$ gives the unique-mode
witness of order $34$, size $41$, and $\gamma_2=25>15+9=24$.
Taking $L=7r$ gives order $16r+1$ and arbitrarily large deficit.

\subsection{Triangle coronas and statement 382c}

If $A\ne\varnothing$, define its average vertex eccentricity by
\[
 \overline{\operatorname{ecc}}_G(A)
 =\frac{1}{|A|}\sum_{v\in A}\operatorname{ecc}_G(v).
\]
This averages eccentricities of individual vertices; it is neither
set eccentricity nor average pairwise distance.  Statement 382c is
\begin{equation}
 \gamma_2(G)\leq
 \frac{\alpha(G)\,\overline{\operatorname{ecc}}_G(A_\delta(G))
       +|B(G)|}{2}.
 \label{early:eq:382c-original}
\end{equation}
The outer brackets in the catalogue serve as grouping symbols.  We give
an integer-valued violation as well, so the refutation does not depend
on a possible rounding interpretation.

\begin{proposition}[Exact reduction for a triangle corona]
\label{early:prop:triangle-corona}
Let $H$ be connected on $m\geq2$ vertices.  Form $G=H\circ K_2$ by
attaching to each vertex $z$ a private pair $p_z,q_z$, adjacent to one
another and to $z$.  Then
\begin{align*}
 \gamma_2(G)&=2m,&\alpha(G)&=m,\\
 A_\delta(G)&=\{p_z,q_z:z\in V(H)\},&
 B(G)&=\{p_z,q_z:z\in B(H)\},\\
 \operatorname{ecc}_G(z)&=\operatorname{ecc}_H(z)+1,&
 \operatorname{ecc}_G(p_z)=\operatorname{ecc}_G(q_z)
 &=\operatorname{ecc}_H(z)+2.
\end{align*}
For these graphs, \eqref{early:eq:382c-original} is equivalent to
\begin{equation}
 2m\leq\sum_{z\in V(H)}\operatorname{ecc}_H(z)+2|B(H)|.
 \label{early:eq:corona-reduction}
\end{equation}
\end{proposition}
\begin{proof}
In each private triangle, a two-dominating set must select at least two
vertices: if at most one were selected, an unselected private vertex
would have fewer than two selected neighbors.  Selecting all private
vertices attains $2m$.  An independent set has at most one vertex in
each private triangle, and selecting one private vertex for each $z$
attains $m$.  Private vertices have degree two, while core vertices
have degree at least three.  Distances between different attached
triangles pass through their core vertices.  The eccentricity formulas
follow, as does the description of the periphery.  In particular
\[
 \overline{\operatorname{ecc}}_G(A_\delta)
 =2+\frac1m\sum_z\operatorname{ecc}_H(z),\qquad
 |B(G)|=2|B(H)|.
\]
Substituting these values and $\gamma_2=2m$, $\alpha=m$ gives
\eqref{early:eq:corona-reduction}.
\end{proof}

\begin{corollary}[Refutation of 382c]
\label{early:cor:382c}
Let $H_r=K_r\vee 2K_1$ for $r\geq1$, where the two vertices outside
the clique are nonadjacent, and put $L_r=H_r\circ K_2$.  Then
\begin{align*}
 |V(L_r)|&=3r+6,&
 |E(L_r)|&=\tbinom r2+5r+6,\\
 \gamma_2(L_r)&=2r+4,&\alpha(L_r)&=r+2,\\
 |A_\delta(L_r)|&=2r+4,& |B(L_r)|&=4,\\
 \overline{\operatorname{ecc}}_{L_r}(A_\delta)
 &=\frac{3r+8}{r+2}.&&
\end{align*}
The proposed right side is $(3r+12)/2$, so its deficit is $(r-4)/2$.
\end{corollary}
\begin{proof}
The $r$ clique vertices of $H_r$ are universal and have eccentricity
one.  Its other two vertices have eccentricity two and form the
periphery.  Apply Proposition~\ref{early:prop:triangle-corona}; the
edge count is $\binom r2+2r+3(r+2)$.
\end{proof}

Thus $L_5$ has order $21$, size $41$, and
$\gamma_2=14>27/2$.  The graph $L_6$ has order $24$, size $51$, and
$\gamma_2=16>15$, an integer-valued violation.  The deficit is
unbounded with $r$.

\section{Triangle-free obstructions and the bipartite boundary of 394}
\label{r11td:extension}
The clique module in Corollary~\ref{early:cor:394} already refutes 394.
We now show that failure persists without triangles, while bipartiteness
suffices for validity. These are extensions of the same refutation, not
newly resolved numbered entries.
Write $D_2=\{v:d_G(v)=2\}$, $\lambda_G(v)=\alpha(G[N_G(v)])$, and
$M=\{v:\lambda_G(v)=\min_x\lambda_G(x)\}$. The proposed inequality is
\[
 \gamma_2(G)\le \mu(G)+|E(G[D_2])|+|M|.
\]
Here $\mu$ is matching number and every vertex outside a 2-dominating
set must have at least two neighbors in that set.

\subsection{Arbitrary internal graphs in the spoke construction}
Let $H$ be a connected graph on roots $c_1,\ldots,c_r$, where $r\ge1$.
For each $i$, choose a graph $F_i$ on
$Q_i=\{q_{i1},\ldots,q_{ik_i}\}$ with minimum degree at least two.
Add vertices $t_{i1},\ldots,t_{ik_i}$ and the edges
$c_it_{ij},t_{ij}q_{ij}$. Apart from the edges of $F_i$ and $H$, add
only one further vertex $p$, adjacent to a designated root $c_1$.
Call the resulting graph $G$, and let $H^+$ be $H$ with this pendant
vertex $p$ attached. Put $K=\sum_i k_i$.

\begin{theorem}\label{r11td:spokes}
The construction satisfies
\[
 |V(G)|=2K+r+1,\qquad
 \gamma_2(G)=K+r+1,\qquad
 \mu(G)=K+\mu(H^+).
\]
Moreover, $M=\{p\}$ and $D_2=\bigcup_i\{t_{i1},\ldots,t_{ik_i}\}$
is independent. Thus the excess of the left side of 394 over its
right side is exactly $r-\mu(H^+)$.
\end{theorem}
\begin{proof}
Consider a 2-dominating set $D$. If $c_i\in D$, every pair
$\{t_{ij},q_{ij}\}$ meets $D$, since otherwise $t_{ij}$ has only one
selected neighbor. This gives at least $k_i+1$ selected vertices in
that module. If $c_i\notin D$, all $t_{ij}$ must belong to $D$.
At least one $q_{ij}$ must also be selected: otherwise each omitted
$q_{ij}$ has only its one selected neighbor $t_{ij}$.
Again the module contributes at least $k_i+1$. The pendant vertex
$p$ must be selected, proving the lower bound. Equality follows by
selecting $p$, every root, and every vertex in every $Q_i$.

A matching contains at most $k_i$ internal edges in the module of
order $2k_i+1$. Its remaining edges form a matching in $H^+$.
Conversely, all $K$ edges $t_{ij}q_{ij}$ are disjoint from a maximum
matching of $H^+$. This proves the matching identity.

Each root has at least three pairwise nonadjacent neighbors among
its $t$-vertices. Each $t_{ij}$ has the two nonadjacent neighbors
$c_i,q_{ij}$. Each $q_{ij}$ has its neighbor $t_{ij}$ and at least one
neighbor in $F_i$, which are nonadjacent. Thus every vertex other than
$p$ has local independence at least two, whereas $\lambda(p)=1$.
Finally $d(q_{ij})\ge3$ and $d(c_i)\ge3$, so exactly the $t$-vertices
have degree two and no two of them are adjacent.
\end{proof}
This proof uses no clique assumption on $F_i$. Its counting and
matching arguments are elementary; no claim of historical priority is
made for the general construction.

\begin{corollary}\label{r11td:c4}
Take every $F_i=C_4$, let $H$ be a star on $r$ roots, and attach $p$
to its center. Then $G$ is triangle-free and
\[
 n=9r+1,\qquad \gamma_2(G)=5r+1,\qquad \mu(G)=4r+1.
\]
Its deficit in 394 is $r-1$, so every $r\ge2$ is a counterexample.
In particular $r=2$ gives a
19-vertex triangle-free graph with $\gamma_2=11$ and proposed bound $10$.
\end{corollary}
\begin{proof}
The graph $H^+$ is a star and has matching number one. No triangle
uses a spoke, and both $H$ and every $F_i$ are triangle-free.
Apply Theorem~\ref{r11td:spokes}.
\end{proof}

\subsection{A smaller triangle-free family}
Let $J$ have vertex set $\{0,1,\ldots,6\}$, root $0$, and edge set
\[
 E(J)=\{01,03,06,12,14,23,35,45,46\}.
\]
This is a triangle-free graph with independent degree-two set
$A=\{2,5,6\}$.

\begin{lemma}\label{r11td:rigid}
A set that 2-dominates every omitted vertex of $J$ other than its root
has at least four vertices. The set $\{0,1,3,4\}$ 2-dominates $J$.
\end{lemma}
\begin{proof}
Suppose at most three vertices are selected and let $a$ be the number
selected from $A$. An omitted vertex of $A$ forces both its neighbors
to be selected. If $a=0$, their combined neighborhoods
$\{0,1,3,4\}$ force four selections. If $a=1$, the two omitted
vertices of $A$ have at least three distinct neighbors, giving at least
four selections together. If $a=2$, the remaining omitted vertex
forces two more selections. If $a=3$, the selected set is exactly
$A$, but the nonroot vertex $1$ has only its neighbor $2$ selected.
These exhaust the possibilities. The stated four-set is checked
directly: each omitted vertex $2,5,6$ has both its neighbors in it.
\end{proof}

Take $r$ disjoint copies of $J$, connect their roots as a star centered
at the first root, and attach one pendant vertex $p$ to that root.
Denote the resulting connected graph by $G_r$.

\begin{theorem}\label{r11td:small}
For every $r\ge1$, $G_r$ is triangle-free and
\[
 n=7r+1,\quad \gamma_2(G_r)=4r+1,\quad
 \mu(G_r)=3r+1,\quad |M|=1,\quad |E(G_r[D_2])|=0.
\]
Thus the deficit is $r-1$, so 394 fails for every $r\ge2$. At $r=2$ this is a 15-vertex witness with
$\gamma_2=9$ and proposed bound $8$; no minimum-order claim is made.
\end{theorem}
\begin{proof}
Only roots receive external neighbors, so Lemma~\ref{r11td:rigid}
forces four selections in every copy. The pendant $p$ is mandatory.
Select $p$ and the displayed four-set in each copy to attain this
bound. Every matching has at most three internal edges per copy and
at most one external edge, since all external edges share the central
root. Attain $3r+1$ by matching that root to $p$, using
$\{12,35,46\}$ in its copy, and using the same three edges in all
other copies. The original degree-two vertices remain independent;
all roots have degree at least four, all other nonpendant vertices
have degree two or three, and the graph is triangle-free. Its local
independence equals its degree, so $M=\{p\}$.
\end{proof}

\subsection{Bipartite validity and an exact deletion bound}
Let $P$ be the pendant set of a connected graph of order greater than
two. A set $X$ is an odd-cycle transversal if $G-X$ is bipartite.

\begin{proposition}\label{r11td:oct}
For every such $X$,
\[
 \gamma_2(G)\le |P\cup X|+\mu\bigl(G-(P\cup X)\bigr).
\]
In particular 394 holds for every connected bipartite graph of order
greater than two.
\end{proposition}
\begin{proof}
Put $I=P\cup X$. By K\"onig's theorem, the bipartite graph $G-I$
has a vertex cover $C$ of size $\mu(G-I)$. Every vertex outside
$I\cup C$ has all its neighbors in $I\cup C$; its original degree
is at least two because it is not pendant. Hence $I\cup C$ is
2-dominating. If $G$ is bipartite and has leaves, take $X=\varnothing$;
then $\gamma_2\le|P|+\mu(G-P)\le|P|+\mu(G)$ and $M=P$, since
local independence equals degree. If there are no leaves, a minimum
vertex cover itself is 2-dominating, giving $\gamma_2\le\mu(G)$.
Both cases imply the proposed bound in 394.
\end{proof}
The proposition is an elementary application of the established
bipartite matching--vertex-cover theorem, not a claimed new matching
principle. Its value here is the exact comparison with the examples.

\begin{proposition}\label{r11td:octsharp}
Both triangle-free families above have odd-cycle transversal number
exactly $r$, and the bound in Proposition~\ref{r11td:oct} is attained.
\end{proposition}
\begin{proof}
For the $C_4$ spoke family, every module contains a five-cycle formed
by an edge $q_{ij}q_{i\ell}$ and its two spokes. These $r$ odd cycles
are vertex-disjoint. Deleting all roots makes the graph bipartite,
so the transversal number is $r$. Take those roots as $X$. After also
deleting $p$, the edges $t_{ij}q_{ij}$ form a perfect matching of size
$4r$. The deletion bound is therefore $(r+1)+4r=5r+1$.

For $G_r$, every copy contains the five-cycle $1,2,3,5,4,1$, giving
the same lower bound $r$. Delete vertex $4$ from each copy.
The remaining copy is a four-cycle $0,1,2,3,0$ with leaves $5$ and $6$
attached to $3$ and $0$; it is bipartite. Its bipartition may be reversed
independently in each copy so that all edges of the root star cross.
The pendant $p$ also receives the opposite color to its root.
Thus these $r$ deletions suffice. After also deleting $p$, the edges
$06,12,35$ in every copy form a perfect matching of size $3r$.
The bound is $(r+1)+3r=4r+1$, as asserted.
\end{proof}
The standard-library checker \nolinkurl{verification/extensions/module394/verify_extension394.py} enumerates every
subset of the 15- and 19-vertex examples to compute their exact
2-domination numbers, and independently computes maximum matchings,
local independence, and minimum odd-cycle transversals. These finite
checks corroborate the proofs for the infinite families.


\section{Total domination and distances between centers: statement 269}
\label{early:sec:269}

Graphs in this section are connected and have order at least two.
A total dominating set $D$ requires every vertex, including each vertex
of $D$, to have a neighbor in $D$.  Its minimum cardinality is
$\gamma_t(G)$.  For a center containing at least two vertices, put
\[
 \overline d_C(G)=\frac{1}{\binom{|C(G)|}{2}}
       \sum_{\{x,y\}\subseteq C(G)}d_G(x,y).
\]
Distances are measured in $G$, and only pairs of distinct centers are
averaged.  For the reductions below we set this average to zero when
the center is a singleton; all counterexamples have every vertex
central and do not require this convention.  Statement 269 proposes
\begin{equation}
 \gamma_t(G)\geq\ceil{1+\overline d_C(G)}
 \quad\text{for connected $C_4$-free graphs.}
 \label{early:eq:269-original}
\end{equation}
Here $C_4$-free excludes an ordinary four-cycle, whether induced or not,
as in the original definition~\cite{src:catalogue,src:definitions}.

We first consolidate the structural reduction from the earlier baseline
and Round~2.  It shows precisely why the counterexamples subsequently
found in Round~5 occur at total domination number and radius four.

\subsection{Equality in the radius bound}

\begin{lemma}[Spanning trees preserving a total dominating set]
\label{early:lem:td-tree}
Let $D$ be a total dominating set of a connected graph $G$.  There is a
spanning tree $T$ in which $D$ remains totally dominating.  Every such
tree satisfies $\diam(T)\leq2|D|-1$.  In particular,
$\rad(G)\leq\gamma_t(G)$.
\end{lemma}
\begin{proof}
Choose a spanning tree of each component of $G[D]$, and attach each
vertex outside $D$ to one chosen neighbor in $D$.  This gives a
spanning forest.  Add edges corresponding to a spanning tree of its
connected quotient in $G$.  Every member of $D$ has a neighbor in its
chosen component tree, since $G[D]$ has no isolated vertices, and every
other vertex has its chosen neighbor.  Thus $D$ remains totally
dominating.

For any path $P$ in a tree, a vertex has at most two neighbors on $P$;
if it is outside $P$, it has at most one.  Count incidences from $D$
to vertices of $P$.  Every vertex of $P$ has such an incidence, giving
$|V(P)|\leq2|D|$.  Hence $\diam(T)\leq2|D|-1$, and the formula for
the radius of a tree gives $\rad(G)\leq\rad(T)\leq|D|$.
\end{proof}

\begin{proposition}[The equality quotient]
\label{early:prop:equality-quotient}
Suppose $G$ is connected and
$\gamma_t(G)=\rad(G)=t$.  For any minimum total dominating set $D$:
\begin{enumerate}
\item $t=2k$ is even, and $G[D]$ consists of $k$ disjoint edges;
\item every vertex outside $D$ has neighbors in exactly one of these
      edges, possibly in both endpoints of that edge;
\item partition $V(G)$ into these $k$ edges and their respective
      outside neighbors.  The simple quotient $Q$ of the partition
      is a path, a cycle, or a single vertex.
\end{enumerate}
Every spanning tree preserving total domination by $D$ has diameter
$2t-1$.  Both endpoints of its central edge are centers of $G$.
\end{proposition}
\begin{proof}
For any such tree, Lemma~\ref{early:lem:td-tree} and
$t=\rad(G)\leq\rad(T)$ force $\diam(T)=2t-1$.  On a diametral
path $P$ there are exactly $2t$ vertices.  Equality in the incidence
count shows that every member of $D$ is internal to $P$, has both
path neighbors on $P$, and every vertex of $P$ has exactly one
neighbor in $D$.  Starting at an endpoint, membership in $D$ along
$P$ follows the repeating pattern $0,1,1,0$.  The other endpoint has
value zero and its predecessor has value one, forcing $2t$ to be a
multiple of four.  Thus $t$ is even.

If a component of $G[D]$ had at least three vertices, its chosen
component tree could be taken to contain a three-vertex path in $D$.
All its vertices lie on $P$, and uniqueness of tree paths puts that
three-vertex path on $P$.  Its middle vertex would have two neighbors
in $D$, a contradiction.  Hence $G[D]$ consists of disjoint edges.

If $x\notin D$ had neighbors $d_i,d_j$ in distinct such edges, attach
$x$ to $d_i$ and include the edge $xd_j$ in the quotient spanning
tree.  This is possible because that quotient edge is not a loop and
can be extended to a spanning tree.  The resulting tree path between
$d_i$ and $d_j$ passes through $x$.  Since both endpoints lie on $P$,
so does this path, and $x$ has two $D$-neighbors on $P$, again a
contradiction.  This proves the second assertion and makes the
partition in the third assertion unambiguous.

Every spanning tree of $Q$ can now be realized by the construction of
Lemma~\ref{early:lem:td-tree}.  Its diametral path contains all of $D$
and therefore meets every part.  Each part induces a connected
subtree, so the path cannot leave a part and subsequently reenter it.
Its projection is a Hamilton path in the chosen quotient tree, which
must itself be a path.  If $Q$ had a vertex of degree at least three,
the three incident edges could be extended to a spanning tree that
was not a path.  Thus the connected graph $Q$ has maximum degree at
most two, proving the classification.

Finally, the endpoints of the central edge of a tree of diameter
$2t-1$ have tree eccentricity $t$.  Their eccentricities in $G$ are
at most $t$ and at least $\rad(G)=t$, so they are centers of $G$.
\end{proof}

\begin{lemma}
\label{early:lem:radius-two}
If $G$ is connected and $C_4$-free, and
$\gamma_t(G)=\rad(G)=2$, then its center is a clique and
$\overline d_C(G)\leq1$.
\end{lemma}
\begin{proof}
Let $D=\{u,v\}$, where $uv$ is an edge.  Divide the outside vertices
into $A$, adjacent to $u$ only; $B$, adjacent to $v$ only; and $W$,
adjacent to both.  The radius is not one, so both $A$ and $B$ are
nonempty.  A four-cycle forbids edges from $A$ to $B$, forbids edges
from $W$ to $A\cup B$, and implies $|W|\leq1$.  Therefore the
distance from any member of $A$ to any member of $B$ is at least
three.  No member of $A\cup B$ is central.  The vertices in
$\{u,v\}\cup W$ all have eccentricity two and form a clique, proving
the assertion.
\end{proof}

\subsection{The radius-four reduction}

\begin{theorem}[The reduction recorded in Round 2]
\label{early:thm:269-reduction}
If a connected $C_4$-free graph violates
\eqref{early:eq:269-original}, then
$\gamma_t(G)=\rad(G)=4$.
\end{theorem}
\begin{proof}
Write $t=\gamma_t(G)$ and $r=\rad(G)$.  Distances between centers
are at most $r$, so if $t>r$, then
$\ceil{1+\overline d_C(G)}\leq r+1\leq t$.  Thus a counterexample
must have $t=r$.  Proposition~\ref{early:prop:equality-quotient} makes
$t$ even, and Lemma~\ref{early:lem:radius-two} excludes $t=2$.
It remains to prove the inequality when $t=r=2k$ with $k\geq3$.
Fix a minimum total dominating set with edges $a_ib_i$, and its
quotient $Q$ from Proposition~\ref{early:prop:equality-quotient}.

\emph{Contacts at an internal part.}
Suppose a part has two different neighbors in $Q$.  Choose one
interpart edge to each of them, and choose a quotient spanning tree
containing both.  In a spanning tree preserving $D$, its diametral
path contains both $D$-vertices of the part and traverses both chosen
interpart edges.  The two contacts must therefore attach to opposite
endpoints of its $D$-edge.  The choices of edges and attachments were
arbitrary.  Consequently all contacts toward one neighboring part
are adjacent to exactly one fixed endpoint of the $D$-edge, and all
contacts toward the other neighbor are adjacent to the opposite
endpoint.  No common neighbor of both endpoints can be a contact.

\emph{Case 1: $Q$ is a cycle.}
Orient its parts cyclically, with incoming contacts at $a_i$ and
outgoing contacts at $b_i$.  Choose one interpart edge at each
adjacency.  The $D$-edges, their contact attachment edges, and the
chosen interpart edges form a simple cycle $Z$ of length $4k$
containing all of $D$.  Delete one interpart edge of $Z$, and attach
each remaining vertex as a leaf to a $D$-neighbor.  This gives a
spanning tree preserving $D$.  Its central edge lies on the remaining
$4k$-vertex path.  If $k$ is odd, that edge is a $D$-edge; varying
the deleted interpart edge proves that all $2k$ members of $D$ are
central.  If $k$ is even, the central edge is the opposite interpart
edge.  Because the chosen interpart edges were arbitrary, every
endpoint of every interpart edge is central.  In particular there
are at least $2k$ centers in this case as well.

We use the following elementary metric fact.  For $m$ points, with
repetitions allowed, on a circle of circumference $L$, the sum of
their ordered pairwise circular distances is at most $Lm^2/4$.
Indeed, parameterize antipodal cuts by a half-circle of length $L/2$.
If a cut separates $s$ points from the other $m-s$, the number of
separated ordered pairs is $2s(m-s)\leq m^2/2$.  For a fixed pair,
the measure of separating cuts is its circular distance.  Integration
proves the stated bound.

Project each center to a neighbor in $D$, leaving centers already in
$D$ fixed, and view the images on $Z$.  Set $m=|C(G)|$ and
$s=|C(G)\cap D|$.  Travel to the images, along the shorter arc of
$Z$, and back.  The circular estimate and the attachment costs give
\begin{equation}
 \overline d_C(G)\leq
 \frac{km}{m-1}+2-\frac{2s}{m}.
 \label{early:eq:269-circle}
\end{equation}
For odd $k\geq3$, we have $s=2k$ and $m\geq2k$, whence
\[
 \overline d_C(G)\leq k+2+\frac{k}{m-1}-\frac{4k}{m}
 \leq k+2\leq2k-1.
\]
For even $k\geq4$, discard the last term of
\eqref{early:eq:269-circle} and use $m\geq2k$ to obtain
\[
 \overline d_C(G)\leq k+2+\frac{k}{2k-1}\leq2k-1.
\]

\emph{Case 2: $Q$ is a path on parts $1,\ldots,k$.}
At internal parts, orient $a_i$ toward the left and $b_i$ toward the
right as above.  We first show that all contacts at the first part
are adjacent to one fixed endpoint of its $D$-edge, and to no other.
Otherwise the first-part contacts collectively meet both endpoints,
possibly by a contact adjacent to both.  The left contacts at part
two are adjacent to $a_2$.  Thus $a_2$ is at distance at most three
from both first-part $D$-vertices, and at most four from every vertex
in the first part.  Choose a path through parts $2,\ldots,k$ using
one interpart edge at each step.  Assign positions
$a_i=4i-3$, $b_i=4i-2$, with the contacts at the intervening integer
positions.  Its vertex at position $2k$ exists for $k\geq3$.
This vertex is within $2k-5+4=2k-1$ of every first-part vertex.
For every other vertex, travel along the chosen path to a
$D$-neighbor; the distance is at most
\[
 \max\{2k-5+1,(4k-2)-2k+1\}=2k-1.
\]
This contradicts radius $2k$.  Interchanging the endpoints of the
first $D$-edge if necessary, all first-part contacts are consequently
private neighbors of $b_1$.  The symmetric argument makes all
last-part contacts private neighbors of $a_k$.

Assign real coordinates to every vertex in part $i$ as follows:
\[
\begin{array}{c|ccccc}
 \text{vertex type}&a_i&b_i&
 \text{private to }a_i&
 \text{private to }b_i&
 \text{common}\\\hline
 \text{coordinate}&4i-3&4i-2&4i-4&4i-1&4i-\tfrac52.
\end{array}
\]
The $C_4$-free assumption forbids an edge between the two private
classes, forbids an edge from a common neighbor to a private neighbor,
and permits at most one common neighbor.  All remaining within-part
edges have coordinate difference at most one.  Every interpart edge
joins coordinates $4i-1$ and $4i$, by the contact orientations.
Thus the coordinate function changes by at most one along every edge.

There must be a private neighbor $z$ of $a_1$, of coordinate zero.
If there were none, choose one edge between each consecutive pair
of parts and use the resulting path through all $D$-edges.  Its
vertex at coordinate $2k$ is at distance at most $2k-1$ from $a_1$
and all other $D$-vertices.  All outside vertices in the first part
can be reached through $b_1$, while all other outside vertices can
be reached through one of their $D$-neighbors, also within $2k-1$.
This again contradicts the radius.  Symmetrically there is a private
neighbor $w$ of $b_k$, of coordinate $4k-1$.

Every center has distance at most $2k$ from both $z$ and $w$.
The coordinate bound forces its coordinate into $[2k-1,2k]$.
If $k$ is odd, this interval contains only the middle $D$-edge and
its possible common neighbor, a clique.  Thus
$\overline d_C(G)\leq1$.  If $k$ is even, it contains only the
outgoing private vertices of part $k/2$ and the incoming private
vertices of part $k/2+1$.  A center on the left must have an
interpart neighbor on the right: a shortest path to $w$ has the
maximum allowed coordinate increase at every step, including its
first step.  The same holds with left and right reversed.  Moreover
both endpoints of every edge across this middle adjacency are
central, by the central-edge assertion of
Proposition~\ref{early:prop:equality-quotient}, realized with that
chosen edge.  Centers on the same side have distance at most two
through their common $D$-neighbor.  If $x$ and $y$ are on opposite
sides, take a cross-neighbor $x'$ of $x$; the path from $x$ to $x'$
and then through the right-hand $D$-vertex to $y$ has length at
most three.  Hence $\overline d_C(G)\leq3\leq2k-1$.

Both quotient cases give $\overline d_C(G)\leq t-1$ for $t\geq6$.
Together with the initial reductions, this leaves only $t=r=4$.
\end{proof}

\subsection{An explicit family in the remaining case}

\begin{definition}
\label{early:def:269-family}
For positive integers $Y,Z$, start with four vertices $a,b,c,d$ and
edges $ab,cd$.  For $1\leq i\leq Y$, add a path
$a,u_i,v_i,d$.  For $1\leq j\leq Z$, add a path
$b,p_j,q_j,c$, a vertex $x_j$ adjacent to $b,p_j$, and a vertex
$y_j$ adjacent to $c,q_j$.  There are no further edges.  Denote the
graph by $R(Y,Z)$.  Its order and size are
\[
 n=4+2Y+4Z,\qquad m=2+3Y+7Z.
\]
\end{definition}

\begin{lemma}
\label{early:lem:269-c4td}
The graph $R(Y,Z)$ is connected and $C_4$-free, and
$\gamma_t(R(Y,Z))=4$.
\end{lemma}
\begin{proof}
Connectivity is immediate.  Temporarily omit the $x_j,y_j$.  Every
cycle in the remaining graph either uses two length-three paths
with the same endpoints and has length six, or uses one path of
each kind together with $ab,cd$ and has length eight.  This base
graph has neither triangles nor four-cycles.  Reinsert the omitted
vertices one at a time.  Each is adjacent to the endpoints of a
base edge, and no base edge receives two such vertices.  A new
four-cycle would require a previously existing length-two path
between those endpoints.  Its middle vertex cannot be a base
vertex, as that would make a base triangle, or a previously added
vertex, as that would repeat the same base edge.  Therefore no
four-cycle is created.

The set $\{a,b,c,d\}$ totally dominates $R(Y,Z)$.  Conversely the
four open neighborhoods
\[
 \{a,v_1\},\quad\{d,u_1\},\quad
 \{b,q_1,x_1\},\quad\{c,p_1,y_1\}
\]
of $u_1,v_1,p_1,q_1$, respectively, are pairwise disjoint.
Every total dominating set meets all four, so its size is at least four.
\end{proof}

\begin{lemma}[Distances after adding a vertex on an edge]
\label{early:lem:edge-triangle-distance}
Adding vertices each adjacent to the endpoints of an existing base
edge does not change distances between base vertices.  A new vertex
with endpoint set $A$ has distance $1+\min_{a\in A}d_H(a,z)$ from
a base vertex $z$, and two distinct new vertices with endpoint sets
$A,A'$ have distance
$2+\min_{a\in A,a'\in A'}d_H(a,a')$, where $H$ is the base graph.
\end{lemma}
\begin{proof}
A passage between base vertices through a new vertex can be replaced
by the existing edge between its two neighbors.  A shortest path
therefore never needs a new internal vertex.  Minimize over the first
and last neighbor choices for paths with new endpoints.
\end{proof}

For completeness, the following tables determine every distance in
$R(Y,Z)$.  Among the hubs,
\[
 d(a,b)=d(c,d)=1,\qquad d(a,d)=d(b,c)=3,\qquad
 d(a,c)=d(b,d)=4.
\]
Distances to hubs, obtained from the length-three base paths and
Lemma~\ref{early:lem:edge-triangle-distance}, are
\[
\begin{array}{c|rrrr}
 &a&b&c&d\\\hline
 u_i&1&2&3&2\\
 v_i&2&3&2&1\\
 p_j&2&1&2&3\\
 q_j&3&2&1&2\\
 x_j&2&1&3&4\\
 y_j&4&3&1&2
\end{array}
\qquad
\begin{array}{c|cc}
 \text{type pair}&\text{equal indices}&\text{unequal indices}\\\hline
 p,q&1&3\\
 p,x&1&2\\
 p,y&2&3\\
 q,x&2&3\\
 q,y&1&2\\
 x,y&3&4
\end{array}
\]
Distinct vertices of any one of the six non-hub types are at distance
two.  Also $d(u_i,v_h)$ is one for $i=h$ and three otherwise, and
\begin{align*}
 (d(u_i,p_j),d(u_i,q_j),d(u_i,x_j),d(u_i,y_j))&=(3,4,3,4),\\
 (d(v_i,p_j),d(v_i,q_j),d(v_i,x_j),d(v_i,y_j))&=(4,3,4,3).
\end{align*}
These formulas follow by choosing the closest hub endpoints in the
base graph; Lemma~\ref{early:lem:edge-triangle-distance} then handles
every pair involving an added triangle vertex.  They exhaust all pairs.

\begin{proposition}[Center and exact distance distribution]
\label{early:prop:269-count}
Every vertex of $R(Y,Z)$ has eccentricity four.  If $N_j$ is the
number of unordered vertex pairs at distance $j$, then
\begin{align*}
 N_1&=2+3Y+7Z,\\
 N_2&=Y^2+3Y+4Z^2+4Z,\\
 N_4&=4YZ+Z^2+Z+2,\\
 N_3&=\binom n2-N_1-N_2-N_4.
\end{align*}
The sum of unordered distances is
\begin{equation}
 W=5Y^2+28YZ+21Z^2+12Y+25Z+16,
 \label{early:eq:269-wiener}
\end{equation}
and its excess over $3\binom n2$ is
\begin{equation}
 F(Y,Z)=-Y^2+4YZ-3Z^2-9Y-17Z-2.
 \label{early:eq:269-excess}
\end{equation}
\end{proposition}
\begin{proof}
The tables give an upper bound of four on all distances.  Partners
at distance four are $a,c$; $b,d$; $u_i,q_1$; $v_i,p_1$;
$p_j,v_1$; $q_j,u_1$; $x_j,d$; and $y_j,a$.
Thus the radius and diameter are four and the center is all of $V$.

The edge count gives $N_1$.  Distance-two pairs among the $u,v$ types
contribute $Y(Y-1)$, and between these types and hubs contribute
$4Y$.  Equal decorated types contribute $2Z(Z-1)$; unequal-index
pairs of types $p,x$ and $q,y$ contribute another $2Z(Z-1)$;
equal-index pairs of types $p,y$ and $q,x$ contribute $2Z$; and
decorated-to-hub pairs contribute $6Z$.  These sum to $N_2$.
At distance four there are two hub pairs, $2Z$ decorated-to-hub
pairs, $4YZ$ plain-to-decorated pairs, and $Z(Z-1)$ unequal-index
pairs of type $x,y$.  This gives $N_4$.  All other pairs have
distance three.  Finally
\[
 W-3\binom n2=N_4-2N_1-N_2,
\]
which expands to \eqref{early:eq:269-excess}; adding $3\binom n2$
gives \eqref{early:eq:269-wiener}.
\end{proof}

\begin{theorem}[Refutation of 269]
\label{early:thm:269-counterexamples}
There are infinitely many connected $C_4$-free graphs with
$\gamma_t=\rad=\diam=4$ for which
$\ceil{1+\overline d_C}=5$.  In particular $R(62,35)$ has order
$268$, size $433$, and
\[
 \overline d_C=\frac{17890}{5963}=3+\frac1{5963}.
\]
No member of the family $R(Y,Z)$ of smaller order violates
statement 269.
\end{theorem}
\begin{proof}
Take $Y=7s$, $Z=4s$.  Then $n=30s+4$ and
\[
 F(7s,4s)=15s^2-131s-2.
\]
This equals $34>0$ at $s=9$, and its increment when $s$ increases
by one is $30s-116>0$ for $s\geq9$.  Hence the average distance
exceeds three for all $s\geq9$.  It is strictly below four because
there are edges, while all distances are at most four.  Therefore
the indicated ceiling is five, contradicting total domination
number four.  The average tends to $3+1/30$ along this subfamily.

For $Y=62,Z=35$, the exact distance counts are
\[
 (N_1,N_2,N_3,N_4)=(433,9070,16333,9942).
\]
There are $35778$ unordered pairs and $W=107340$, giving the
displayed reduced fraction.

To prove minimality within this family, put $N=Y+2Z=(n-4)/2$.
Substitution in \eqref{early:eq:269-excess} and completion of a square
give
\[
 F(Y,Z)=-15\left(Z-\frac{8N+1}{30}\right)^2
       +\frac{4N^2-524N-119}{60}.
\]
The numerator in the last term is negative throughout $0\leq N\leq131$:
it is a convex quadratic with value $-119$ at both endpoints.
Consequently no such parameter pair has positive excess.  All
orders in the family are even; order below $268$ has $N\leq131$.
The displayed order-$268$ member proves the claimed family minimum.
This argument makes no claim about smaller counterexamples outside
the family.
\end{proof}

\section{Even-distance counts and complement neighborhoods: statement 302}
\label{early:sec:302}

For a vertex $z$ of a connected graph, define
\[
 e_G(z)=|\{u:d_G(z,u)\text{ is even}\}|,
 \qquad e_{\min}(G)=\min_z e_G(z).
\]
The vertex $z$ itself is included, since its distance is zero.
For a graph $H$, put
\[
 f_\lambda(H)=
 |\{z:\lambda_H(z)=\max_u\lambda_H(u)\}|,
 \qquad \lambda_H(z)=\alpha(H[N_H(z)]).
\]
This frequency concerns local independence numbers, not a spectral
multiplicity.  Statement 302 proposes
\begin{equation}
 \gamma_t(G)\leq e_{\min}(G)+f_\lambda(\overline G).
 \label{early:eq:302-original}
\end{equation}
These are the original invariant definitions~\cite{src:catalogue,early:glossary}.

\begin{definition}
\label{early:def:302-family}
Let $r\geq2$ and $q\geq\max\{r,3\}$.  Start with a clique
$Q=\{a,b_1,\ldots,b_q\}$ and a vertex $v$ anticomplete to $Q$.
For each $1\leq i\leq r$, add vertices $w_i,t_i$ and edges
$vw_i,aw_i,w_it_i,t_ib_i$.  For every $1\leq i\leq r$ and
$1\leq j\leq q$, add a vertex $\ell_{ij}$ with exactly the two
neighbors $t_i,b_j$.  Denote this graph by $S_{r,q}$.  It has
\[
 n=2+q+2r+rq,\qquad m=\binom{q+1}{2}+4r+2rq.
\]
\end{definition}

\begin{proposition}
\label{early:prop:302-invariants}
For $S=S_{r,q}$,
\[
 \gamma_t(S)=\min\{2r,q+2\},\qquad
 e_{\min}(S)=r+2,\qquad f_\lambda(\overline S)=1.
\]
The minimum even-distance count is attained uniquely at $v$.
Moreover $\rad(S)=2$ and $\diam(S)=3$.
\end{proposition}
\begin{proof}
Let $D$ be a total dominating set.  If $t_i\notin D$ for some $i$,
each $\ell_{ij}$ forces $b_j\in D$, so all $q$ vertices $b_j$ are
selected.  The vertex $v$ also forces some $w_h\in D$.  To dominate
that selected vertex, at least one of $v,a,t_h$ must be selected,
giving $|D|\geq q+2$.  If every $t_i$ is selected, their neighborhoods
\[
 N(t_i)=\{w_i,b_i\}\cup\{\ell_{ij}:1\leq j\leq q\}
\]
are pairwise disjoint and avoid all the $t_i$.  Each must meet $D$,
so $|D|\geq2r$.  The two sets
\[
 \{t_1,\ldots,t_r,b_1,w_2,\ldots,w_r\},\qquad
 \{b_1,\ldots,b_q,a,w_1\}
\]
are total dominating sets of sizes $2r$ and $q+2$, respectively.
Here $r\geq2$ ensures that the first set dominates $v$.
This proves the first formula.

All pairs are at distance at most three.  For example, clique
vertices connect all the $t_i$ and $\ell_{ij}$ within three steps,
$w_i$ reaches such vertices through $a$, and the distance layers from
$v$ are
\[
 \{v\},\quad \{w_i\},\quad\{a\}\cup\{t_i\},
 \quad\{b_j\}\cup\{\ell_{ij}\}.
\]
The vertex $a$ is within two of all vertices, whereas $v$ is at
distance three from every $b_j$.  No vertex is universal.  This
proves the assertions about radius and diameter.

Thus $e_S(z)$ equals one plus the number of distance-two vertices.
Enumerating them from the defining neighborhoods gives
\begin{equation}
\begin{array}{c|l}
 z&e_S(z)\\\hline
 v&r+2\\
 a&rq+r+2\\
 w_i&2q+r\\
 t_i&q+r+1\\
 b_j&rq+r+1-\mathbf1_{\{j\leq r\}}\\
 \ell_{ij}&2q+r+\mathbf1_{\{j\leq r,\ j\ne i\}}.
\end{array}
\label{early:eq:302-even-table}
\end{equation}
For clarity, the distance-two vertices from $w_i$ are the other
$r-1$ vertices $w_h$, all $q$ vertices $b_j$, and its row of $q$
vertices $\ell_{ij}$.  From $t_i$ they are $v,a$, the $q-1$ vertices
$b_j$ with $j\ne i$, and the $r-1$ vertices $\ell_{hi}$ with
$h\ne i$.  From $b_j$ they are all $w_i$, all leaf-type vertices
outside column $j$, and all $t_i$ except $t_j$ when $j\leq r$.
From $\ell_{ij}$ they are $Q\setminus\{b_j\}$, $w_i$, its other
row and column vertices, and $t_j$ precisely when $j\leq r$ and
$j\ne i$.  The rows for $v,a$ follow from their displayed layers
and neighborhoods.  These lists prove
\eqref{early:eq:302-even-table}.  Since $q\geq3$, every entry other
than the one for $v$ is strictly greater than $r+2$.

Finally,
\[
 \lambda_{\overline S}(z)
   =\omega\bigl(S[V(S)\setminus N_S[z]]\bigr).
\]
The clique $Q$ has order $q+1\geq4$ and is the unique clique of
that order.  Indeed $N(v)$ and $N(w_i)$ are independent; the only
edge in $S[N(t_i)]$ is $b_i\ell_{ii}$; and every $\ell_{ij}$ has
degree two.  Any clique containing a vertex outside $Q$ therefore
has order at most three.  The vertex $v$ is the only vertex outside
$Q$ anticomplete to $Q$: every $w_i$ sees $a$, every $t_i$ sees
$b_i$, and every $\ell_{ij}$ sees $b_j$.  Consequently the maximum
local independence number in the complement is $q+1$, attained
only at $v$, proving $f_\lambda(\overline S)=1$.
\end{proof}

\begin{corollary}[Refutation of 302]
\label{early:cor:302}
For $r\geq4$, take $q=2r-2$.  Then $S_{r,2r-2}$ has
\[
 n=2r(r+1),\quad m=6r^2-3r+1,\quad
 \gamma_t=2r,
\]
whereas the right side of \eqref{early:eq:302-original} is $r+3$.
The deficit $r-3$ is unbounded.  The first member $S_{4,6}$ has
order $40$, size $85$, and $\gamma_t=8>7$.
\end{corollary}
\begin{proof}
Substitute the parameter choice into
Proposition~\ref{early:prop:302-invariants} and the order and size formulas.
\end{proof}

The related valid bound
$\gamma_t(G)\leq\omega(G)+f_\lambda(\overline G)$ is an earlier
theorem of DeLaVi\~na, Liu, Pepper, Waller, and West
\cite[Theorem~10]{early:dlpww-total}.  The present examples attain it:
$\omega(S_{r,2r-2})=2r-1$ and $f_\lambda(\overline S)=1$.
This comparison identifies the loss caused by replacing the clique
number with the minimum even-distance count; it does not count the
earlier theorem as a new result.

\section{Annihilation and cycles of length divisible by four: statement 298}
\label{early:sec:298}

Let $G$ have $n$ vertices, $m$ edges, and degree sequence
$d_1\leq\cdots\leq d_n$.  Its annihilation number is
\[
 a(G)=\max\left\{k:\sum_{i=1}^k d_i\leq m\right\}.
\]
Equivalently, it is the maximum cardinality of a vertex set whose
degree sum in $G$ is at most $m$.  For a connected nontrivial graph
$a(G)<n$, and we write $d_{\rm cut}(G)=d_{a(G)+1}$ for its cutoff
degree.  The original statement 298, for connected graphs of order
greater than two, is
\begin{equation}
 \gamma_t(G)\leq a(G)+
 \mathbf1_{\{G\text{ contains no ordinary }C_4\}}.
 \label{early:eq:298-original}
\end{equation}
We prove the stronger assertion that the indicator can be omitted
whenever $G$ contains a cycle of any length divisible by four.

\subsection{Credited antecedents and the classical input}

The global bound $\gamma_t\leq a+1$ and the stronger bound
$\gamma_t\leq a$ when $d_{\rm cut}\geq3$ are prior results credited
in the Round~7 source to J.~J.~Gebendorfer's manuscript of
11 October 2026~\cite{early:gebendorfer}.  We do not count those two
prior results as new contributions of this consolidation.  To expose
the dependency of the present result, we reconstruct their required
proof chain below, following the detailed proof audit included in the
archive~\cite{early:298-audit}.  The prior PDF and text snapshots named in that audit were absent from
the supplied ZIPs. During preparation of this revised edition, the primary
manuscript archive was independently retrieved from the attachment to
issue 7242. Its SHA-256 digest matched the recorded value, and its six-page
written proof was read against the required dependency chain. The
reconstruction below was checked against that primary text; every needed
step is also proved here. No Lean rebuild is claimed.

The one classical theorem used without reproving it is the
Chv\'atal--McDiarmid transversal bound: if a finite 3-uniform
hypergraph has $n_H$ vertices and $m_H$ edges, then
\begin{equation}
 4\tau(H)\leq n_H+m_H.
 \label{early:eq:cm}
\end{equation}
Here $\tau$ is the minimum size of a set meeting every edge.
Isolated vertices and repeated edges are allowed.  The original
theorem is in~\cite{early:chvatal-mcdiarmid}; the exact formulation
and multiset convention also appear explicitly in
\cite{early:henning-lowenstein}.  Alternatively, repeated edges
can be deleted before applying the simple-hypergraph version, since
this does not change $\tau$ and only decreases the right side.

\begin{lemma}[Mixed transversal bound, credited antecedent]
\label{early:lem:mixed-transversal}
Let $H$ be a finite hypergraph whose nonempty edges have sizes one,
two, or three.  If $e_j$ counts the edges of size $j$, with
multiplicity, and $n_H=|V(H)|$, then
\[
 6\tau(H)\leq2n_H+4e_1+2e_2+e_3.
\]
\end{lemma}
\begin{proof}
Suppose a counterexample exists and choose one minimizing
$\Phi=2n_H+4e_1+2e_2+e_3$.  It has no isolated vertex, since such
a vertex can be deleted.  Write $d_j(v)$ for the number of incident
$j$-edges.  If $4d_1(v)+2d_2(v)+d_3(v)\geq4$, select $v$ and
delete it together with all incident edges.  The potential drops by
at least six, and a transversal of the smaller hypergraph together
with $v$ contradicts minimality.  Hence this weighted incidence is
at most three everywhere.  There are no singleton edges; each
vertex is in at most one 2-edge; and a vertex in a 2-edge is in at
most one 3-edge.

Suppose $uv$ is a 2-edge.  If neither endpoint lies in a 3-edge,
select one endpoint and delete both vertices and the 2-edge.  The
potential drops by six, again a contradiction.  Otherwise assume
$u$ lies in a 3-edge $f$.  Select $u$ and delete $u,uv,f$, lowering
the potential by five.  If $v$ is now isolated, delete it as well,
giving a drop of seven.  Otherwise $v$ lies in exactly one remaining
3-edge $g\ne f$.  Delete $v$ and shrink $g$ to the 2-edge
$g\setminus\{v\}$.  This lowers the potential by one more: deleting
the vertex saves two while the new edge weight increases from one
to two.  The shrinking preserves the transversal number of the
remaining hypergraph, since a transversal using $v$ can replace it
by any other member of $g$, and every transversal of the shrunken
edge also meets the old edge.  Adding the selected $u$ therefore
lifts a transversal with an additional cost of at most one, whereas
the total potential drop is at least six.  This is again impossible.

There are consequently no 2-edges.  The hypergraph is 3-uniform
with maximum degree at most three, so $m_H\leq n_H$.  The classical
bound \eqref{early:eq:cm} now gives
\[
 6\tau(H)\leq\frac32(n_H+m_H)\leq2n_H+m_H=\Phi,
\]
the final contradiction.
\end{proof}

\begin{lemma}[Annihilation budget, credited antecedent]
\label{early:lem:annihilation-budget}
Let $G$ be connected of order $n\geq3$, with $\ell$ degree-one
vertices and $b$ degree-two vertices.  If $d_{\rm cut}(G)\geq3$,
then
\[
 a(G)\geq\floor{\frac n2+\frac\ell3+\frac b6}.
\]
\end{lemma}
\begin{proof}
All the vertices of degrees one and two belong to the affordable
initial degree segment.  Put $L=\ell+2b$, let
$r=n-\ell-b$ be the number of remaining vertices, and let $E$ be
their degree sum.  Then $L\leq m$, $L+E=2m$, and
$E\geq L$, $E\geq3r>0$.  Set
\[
 t=\floor{\frac r2-\frac{rL}{2E}}.
\]
This is an integer between zero and $r$.  The sum of the $t$
smallest degrees among the $r$ high-degree vertices is at most
$tE/r\leq(E-L)/2=m-L$.  Adjoining them to all the low-degree
vertices gives an annihilation set.  Therefore
\begin{align*}
 a(G)&\geq\ell+b+t\\
 &=\floor{\frac n2+\frac\ell2+\frac b2
                  -\frac{r(\ell+2b)}{2E}}\\
 &\geq\floor{\frac n2+\frac\ell3+\frac b6},
\end{align*}
where the last step uses $r/E\leq1/3$.
\end{proof}

\begin{theorem}[A degree-count bound, credited antecedent]
\label{early:thm:td-degree-count}
Every connected graph $G$ of order $n\geq3$, with $\ell$ leaves
and $b$ degree-two vertices, satisfies
\[
 6\gamma_t(G)\leq3n+2\ell+b.
\]
\end{theorem}
\begin{proof}
We induct on $n$.  For a star the total domination number is two;
the inequality is immediate, including $P_3$, where the right side
equals $14$.  Suppose $G$ is not a star and has a strong support
vertex $v$, meaning one with at least two leaf neighbors.  Such a
vertex has a nonleaf neighbor, by connectivity, so its degree is
at least three.  Delete one of its leaf neighbors to obtain $H$.
The graph remains connected, its leaf count decreases by one, and
its number of degree-two vertices increases by at most one.  Every
total dominating set of $H$ contains $v$, forced by the remaining
leaf, and is consequently also a total dominating set of $G$.
The induction hypothesis gives
\[
 6\gamma_t(G)\leq3(n-1)+2(\ell-1)+(b+1)
 \leq3n+2\ell+b.
\]

It remains to consider graphs with no strong support.  Let $S$ be
the support set and put $s=|S|=\ell$.  Delete all leaves to form
the connected core $K$, which has at least two vertices.  Let
$S_2,S_3$ consist of supports whose degrees in $G$ are two and
three, with sizes $s_2,s_3$.  Put
\[
 T=N_K(S_2)\setminus S,\quad c=|T|,\quad D_0=S\cup T,
 \quad t_0=|N_K(D_0)|.
\]
Each member of $S_2$ has just one core neighbor, so $c\leq s_2$.
The disjoint sets $T,S_2$ both lie in $N_K(D_0)$, including when
two supports are adjacent.  Hence $t_0\geq c+s_2$.

Let $U=V(K)\setminus N_K(D_0)$ be the vertices not yet dominated,
and let $Q_0=V(K)\setminus D_0$.  For $u\in U$, all its core
neighbors lie in $Q_0$, and it has at least two: a core vertex of
degree one would be a member of $S_2$, already dominated.  Build a
hypergraph on $Q_0$ with one edge for each $u\in U$, taking its
whole core neighborhood if that neighborhood has size two and
any three of its neighbors otherwise.  There are
\[
 n_H=n-2s-c,\qquad h=|U|=n-s-t_0
\]
vertices and edges respectively.  The number $h_2$ of 2-edges
satisfies $h_2\leq b-s_2+s_3$: a core-degree-two vertex is either
a nonsupport of original degree two or a support of original
degree three.

If $X$ is a transversal of this hypergraph, $D_0\cup X$ totally
dominates $G$.  Vertices in $N_K(D_0)$ see $D_0$, all remaining
core vertices see $X$ by their neighborhood rows, and leaves see
their supports.  This includes vertices selected into $D_0\cup X$
themselves.  The mixed transversal bound gives
\begin{align*}
 6\gamma_t(G)
 &\leq6(s+c)+2(n-2s-c)+h+h_2\\
 &=3n+s+4c-t_0+h_2\\
 &\leq3n+s+b+3c-2s_2+s_3\\
 &\leq3n+s+b+s_2+s_3\\
 &\leq3n+2s+b=3n+2\ell+b.
\end{align*}
The argument also covers an empty row set or variable set; the
empty transversal then has size zero.
\end{proof}

\begin{corollary}[High-cutoff bound, credited antecedent]
\label{early:cor:high-cutoff}
For a connected graph of order at least three with
$d_{\rm cut}(G)\geq3$, one has $\gamma_t(G)\leq a(G)$.
\end{corollary}
\begin{proof}
The integer $\gamma_t(G)$ is at most
$\floor{n/2+\ell/3+b/6}$ by
Theorem~\ref{early:thm:td-degree-count}; this is at most $a(G)$ by
Lemma~\ref{early:lem:annihilation-budget}.
\end{proof}

\subsection{Low-cutoff transfer and incidence coloring}

\begin{lemma}[The cutoff-two formula]
\label{early:lem:cutoff-two}
If $G$ is connected of order at least three, then
$d_{\rm cut}(G)\geq2$.  If equality holds and $\ell$ is its
number of leaves, then
\[
 a(G)=\floor{\frac{m+\ell}{2}}.
\]
\end{lemma}
\begin{proof}
Distinct leaves have distinct incident edges when $n\geq3$.
Thus $\ell\leq m$ and all leaves form an annihilation set, which
excludes cutoff one.  At cutoff two, the first $a(G)+1$ degrees
consist of all $\ell$ ones followed by twos.  The defining
inequalities are
$2a(G)-\ell\leq m<2(a(G)+1)-\ell$, giving the formula.
\end{proof}

\begin{lemma}[Annihilation under a suitable edge deletion]
\label{early:lem:annihilation-transfer}
Suppose $d_{\rm cut}(G)=2$, $uv\in E(G)$, and $d_G(v)\geq3$.
For $H=G-uv$ one has $a(H)\leq a(G)$.
\end{lemma}
\begin{proof}
Write $m=|E(G)|$ and choose an annihilation set $A$ of $H$ with
$|A|=a(H)$.  Its degree sum in $H$ is at most $m-1$, so restoring
$uv$ gives a sum in $G$ of at most $m+1$.  If it is at most $m$,
there is nothing to prove.  Otherwise it is exactly $m+1$, and
both endpoints of the restored edge belong to $A$.

Let $L_2=\{x:d_G(x)\leq2\}$.  The first $a(G)+1$ degrees lie in
$L_2$ and sum to more than $m$, so the degree sum of $L_2$ is at
least $m+1$.  The set $A$ cannot contain all of $L_2$ as well as
$v$: its degree sum would then be at least $m+4$.  Choose
$w\in L_2\setminus A$ and replace $v$ by $w$.  This preserves
cardinality and decreases the degree sum by at least one, producing
an annihilation set of $G$ of size $a(H)$.
\end{proof}

\begin{lemma}[A two-coloring of neighborhood incidences]
\label{early:lem:incidence-coloring}
Let $G$ be a tree of order at least three, or a connected unicyclic
graph whose cycle has length $4k$ for some $k\geq1$.  If $S$ is
its support set and $\ell$ its number of leaves, then
\[
 \gamma_t(G)\leq\floor{\frac{n+|S|}{2}}
                 \leq\floor{\frac{n+\ell}{2}}.
\]
\end{lemma}
\begin{proof}
Put $Q=V(G)\setminus S$ and $U=V(G)\setminus N_G(S)$, where
$N_G(S)$ is an open neighborhood union.  For each $u\in U$, all
its neighbors are in $Q$, and $d_G(u)\geq2$: a leaf has a support
neighbor and would lie in $N_G(S)$.

Build the bipartite incidence graph with row vertices indexed by
$U$, variable vertices indexed by $Q$, and an edge for each
adjacency in $G$.  The graph $G$ is bipartite.  Its full
row-variable neighborhood graph splits into two disjoint copies
of $G$, one in each bipartition orientation.  The restricted
incidence graph therefore has only tree components or unicyclic
components whose cycle has length $4k$.  Every row has degree at
least two.

Color variables with two colors so every row sees both colors.
In a tree component, root at a variable.  Each row has a parent
variable and at least one child variable; give one child the color
opposite its parent and color the other children arbitrarily,
continuing outward.  Isolated variables cause no difficulty.
In a cyclic component, there are $2k$ variables around the cycle.
Color them alternately, which is consistent because $2k$ is even.
Every cycle row now sees both colors.  Extend to branches by the
same tree rule; branches issuing from a cycle row may start with
arbitrary variable colors because that row is already satisfied.

Let $Q_0,Q_1$ be the color classes.  Each set $S\cup Q_i$ totally
dominates $G$: every vertex in $N_G(S)$ sees $S$, and every vertex
of its complement $U$ sees both variable colors.  This reasoning
also dominates selected vertices themselves.  The two sizes sum
to $n+|S|$, so the smaller is at most its half rounded down.
Finally, distinct supports have disjoint nonempty sets of leaf
neighbors, so $|S|\leq\ell$.
\end{proof}

The tree case of Lemma~\ref{early:lem:incidence-coloring} is part of
the credited antecedent argument.  Its extension to a unicyclic
$4k$-cycle is the additional incidence-coloring step used by the
Round~7 result.  The proof explains the divisibility condition:
the variable vertices on each incidence cycle must have even number.

\begin{theorem}[Global $a+1$ bound, credited antecedent]
\label{early:thm:global-a-plus-one}
Every connected graph of order at least two satisfies
$\gamma_t(G)\leq a(G)+1$.
\end{theorem}
\begin{proof}
For $K_2$ this is $2\leq1+1$.  For order at least three,
Corollary~\ref{early:cor:high-cutoff} handles cutoff at least
three, so only cutoff two need be considered.  If $G$ is a tree,
$m=n-1$ and Lemmas~\ref{early:lem:incidence-coloring} and
\ref{early:lem:cutoff-two} give
\[
 \gamma_t(G)\leq\floor{\frac{n+\ell}{2}}
 =\floor{\frac{m+\ell+1}{2}}\leq a(G)+1.
\]
For a cycle $C_n$, $a(C_n)=\floor{n/2}$.  Label its vertices
$0,\ldots,n-1$ cyclically and write $n=4q+s$ with $0\leq s\leq3$.
Select $4j,4j+1$ for $0\leq j<q$.  For $s=1$, also select $4q$;
for $s=2$ or $3$, also select $4q,4q+1$.  These are total
dominating sets of size at most $\floor{n/2}+1$, including $C_3$.

If a counterexample remained, choose one with the fewest edges.
It is neither a tree nor a cycle, and has cutoff two.  Some cycle
has a vertex $v$ of degree at least three: otherwise every vertex
of any cycle would have degree two and connectivity would make
$G$ that cycle.  Delete a cycle edge incident with $v$, obtaining
a connected graph $H$.  Minimality, monotonicity of total domination
under edge addition, and
Lemma~\ref{early:lem:annihilation-transfer} give
\[
 \gamma_t(G)\leq\gamma_t(H)\leq a(H)+1\leq a(G)+1,
\]
a contradiction.
\end{proof}

\subsection{The prescribed-cycle theorem and the original inequality}

\begin{corollary}
\label{early:cor:unicyclic-4k}
If $G$ is connected and unicyclic with cycle length divisible by
four, then $\gamma_t(G)\leq a(G)$.
\end{corollary}
\begin{proof}
At high cutoff use Corollary~\ref{early:cor:high-cutoff}.  At
cutoff two, $m=n$, and
Lemmas~\ref{early:lem:cutoff-two} and
\ref{early:lem:incidence-coloring} both give the threshold
$\floor{(n+\ell)/2}$.
\end{proof}

\begin{lemma}
\label{early:lem:removable-cycle-edge}
Let $G$ be connected, let $C$ be a cycle, and suppose $G$ is not
unicyclic.  There is a nonbridge edge outside $E(C)$ with an
endpoint of degree at least three.
\end{lemma}
\begin{proof}
Take a different cycle $C'$, which has an edge outside $E(C)$.
If $C'$ meets $C$, start at an intersection and traverse $C'$ up
to its first edge outside $E(C)$.  Its initial vertex belongs to
$C$ and has that edge in addition to its two incident $C$-edges.
If the cycles are vertex-disjoint, take a shortest path between
them.  Its endpoint on $C'$ has a neighbor outside $C'$, so has
degree at least three.  Choose either incident $C'$-edge there.
In either case the chosen edge is outside $C$ and lies on $C'$,
so is not a bridge.
\end{proof}

\begin{theorem}[Prescribed cycles of length divisible by four]
\label{early:thm:4k-cycle}
Every connected graph containing a cycle of length divisible by
four satisfies $\gamma_t(G)\leq a(G)$.
\end{theorem}
\begin{proof}
Choose a counterexample with the fewest edges and fix a cycle
$C$ of length $4k$ in it.  Lemma~\ref{early:lem:cutoff-two} and
Corollary~\ref{early:cor:high-cutoff} force cutoff two.
Corollary~\ref{early:cor:unicyclic-4k} excludes the unicyclic case.
By Lemma~\ref{early:lem:removable-cycle-edge}, delete a nonbridge
edge outside $C$ with a high-degree endpoint.  The resulting
graph $H$ is connected, still contains $C$, and has fewer edges.
No condition on its cutoff is needed.  Minimality and the
annihilation-transfer lemma yield
\[
 \gamma_t(G)\leq\gamma_t(H)\leq a(H)\leq a(G),
\]
a contradiction.
\end{proof}

\begin{corollary}[Statement 298]
\label{early:cor:298}
Every connected graph $G$ of order greater than two satisfies
\eqref{early:eq:298-original}.
\end{corollary}
\begin{proof}
If $G$ contains an ordinary four-cycle, apply
Theorem~\ref{early:thm:4k-cycle} with $k=1$.  Otherwise apply the
credited global bound, Theorem~\ref{early:thm:global-a-plus-one}.
The four-cycle need not be induced.
\end{proof}

\begin{example}[Sharpness of both branches]
For $k\geq1$, the cycle $C_{4k}$ has
$\gamma_t=a=2k$.  The annihilation value follows from its constant
degree two.  The lower bound for total domination follows by
counting at most two dominated vertices per selected vertex,
and the repeated pattern $1100$ attains it.

For $C_{4k+2}$, the annihilation number is $2k+1$ and the total
domination number is $2k+2$.  To see the strict lower bound, suppose
a total dominating set had $2k+1$ vertices.  Every nonempty run
of selected vertices around the cycle has length at least two,
and every intervening run of unselected vertices has length at
most two.  There would be at most $k$ selected runs and thus at
most $2k$ unselected vertices, whereas there are $2k+1$.
This is impossible.  The construction in
Theorem~\ref{early:thm:global-a-plus-one} attains $2k+2$.
These cycles are $C_4$-free, so the additive one in the
four-cycle-free branch cannot be removed in general.
\end{example}



% Integrated source fragment: closed_late_structural.tex
\section{Counterexample families from Rounds 11--14}\label{late:counterexamples}
The four source statements treated in this section are mathematically
refuted by explicit infinite families. Their exact definitions are retained,
including the distinction between the periphery and arbitrary boundary
vertices, the largest-output interpretation of Maxine, and the
half-coefficient open-neighbourhood version of 354. These are source-qualified
conclusions; archived open flags are not assertions of current curator
status or worldwide publication priority.


\subsection{Conjecture 171: peripheral averages and even-distance counts}\label{late:171}

\subsubsection{Statement and scope}\label{late:171:statement-and-scope}

The archived original row 171 states, for a simple connected graph
\(G\),

\[L_s(G)\geq \frac{\max_{v\in V(G)}\operatorname{dist}_{\rm even}(v)-1}{\operatorname{dist}_{\rm avg}(B)}.\]

Here \(L_s\) is the maximum number of leaves in a spanning tree. The set
\(B\) is the periphery: vertices of maximum eccentricity. The
denominator is the average of all nonzero distances \textbf{between
pairs of vertices of \(B\)}, not the average between \(B\) and all
vertices and not the average distance to the set \(B\). The
even-distance count includes distance zero. Our principal counterexample
remains a strict counterexample if distance zero is omitted.

The source glossary defines the historical word \emph{boundary} as
the maximum-eccentricity periphery, not all leaves of a tree; the archived
row and definitions are recorded in~\cite{src:catalogue}.
The construction below refutes that precise formulation.

\subsubsection{The 407-vertex witness}\label{late:171:the-407-vertex-witness}

Construct a tree \(T\) as follows. All paths mentioned have no vertices
in common except their prescribed common center \(c\).

\begin{enumerate}
\def\labelenumi{\arabic{enumi}.}

\item
  From \(c\), draw a path of length 14 ending at a vertex \(a\); attach
  27 new pendant vertices \(p_1,\ldots,p_{27}\) to \(a\).
\item
  From \(c\), draw a path of length 15 ending at a vertex \(q\).
\item
  From \(c\), draw 25 further paths, each of length 14, ending at
  vertices \(s_1,\ldots,s_{25}\).
\end{enumerate}

These operations attach trees at a single existing vertex, so \(T\) is a
finite simple connected tree. Its order is

\[n=1+14+27+15+25\cdot14=407.\]

\begin{figure}[ht]
\centering
\setlength{\unitlength}{1mm}
\begin{picture}(145,54)
\put(70,30){\circle*{1.5}} \put(72,32){$c$}
\put(30,30){\circle*{1.5}} \put(28,33){$a$}
\put(130,30){\circle*{1.5}} \put(131,32){$q$}
\put(30,30){\line(1,0){40}} \put(42,33){\small 14 edges}
\put(70,30){\line(1,0){60}} \put(91,33){\small 15 edges}
\put(30,30){\line(-2,1){20}} \put(10,40){\circle*{1.2}}
\put(30,30){\line(-1,0){20}} \put(10,30){\circle*{1.2}}
\put(30,30){\line(-2,-1){20}} \put(10,20){\circle*{1.2}}
\put(0,44){\small 27 leaves $p_i$}
\put(2,25){$\vdots$}
\put(70,30){\line(-1,-1){20}} \put(50,10){\circle*{1.2}}
\put(70,30){\line(0,-1){20}} \put(70,10){\circle*{1.2}}
\put(70,30){\line(1,-1){20}} \put(90,10){\circle*{1.2}}
\put(77,15){$\cdots$}
\put(43,3){\small 25 side arms, each 14 edges}
\end{picture}
\caption{Compressed construction: degree-two vertices are suppressed.
Three representative leaves or arms stand for the indicated multiplicities.
Only the 27 left leaves and $q$ are peripheral.}
\end{figure}


\paragraph{Maximum spanning-tree
leaves}\label{late:171:maximum-spanning-tree-leaves}

A tree has itself as its unique spanning tree. Its leaves are precisely
the 27 vertices \(p_i\), the vertex \(q\), and the 25 vertices \(s_j\).
Hence

\[L_s(T)=27+1+25=53.\]

\paragraph{Periphery and its average
distance}\label{late:171:periphery-and-its-average-distance}

For every \(i\),

\[d(p_i,q)=15+15=30.\]

Every vertex has distance at most 15 from \(c\), so the diameter is 30.
Moreover, \(p_i\) and \(q\) have eccentricity 30. Every other vertex
\(x\) has \(d(c,x)\leq14\), and consequently

\[d(x,y)\leq d(x,c)+d(c,y)\leq14+15=29\qquad\text{for all }y.\]

Thus no other vertex is peripheral, and

\[B=\{p_1,\ldots,p_{27},q\},\qquad |B|=28.\]

Two distinct \(p_i,p_j\) have distance 2, while \(d(p_i,q)=30\).
Therefore the average over unordered distinct pairs (equivalently over
ordered nonzero pairs) is

\[\operatorname{dist}_{\rm avg}(B)
=\frac{2\binom{27}{2}+30\cdot27}{\binom{28}{2}}
=\frac{702+810}{378}=4.\]

\paragraph{Maximum even-distance
count}\label{late:171:maximum-even-distance-count}

In a tree, the vertices at even distance from a given root are precisely
the bipartition class containing that root. Count the two bipartition
classes using \(c\) as the initial root. The class containing \(c\) has

\[1+7+7+25\cdot7=190\]

vertices: \(c\), seven vertices on the length-14 left stem, seven
vertices on the length-15 right path, and seven vertices on each further
length-14 path. None of the \(p_i\) belong to this class, since their
distance from \(c\) is 15. The other class therefore has \(407-190=217\)
vertices. Both classes are nonempty, so

\[\max_v\operatorname{dist}_{\rm even}(v)=217.\]

For example, any neighbor of \(c\) attains this maximum.

\paragraph{Strict contradiction}\label{late:171:strict-contradiction}

The right side of Conjecture 171 is

\[\frac{217-1}{4}=54,\]

whereas \(L_s(T)=53\). Hence the conjectured inequality is false:

\[\boxed{53<54.}\]

If distance zero were excluded from the even-distance count, the right
side would instead be \(215/4=53.75\), still strictly larger than 53.

\subsubsection{An infinite counterexample
family}\label{late:171:an-infinite-counterexample-family}

Keep the first two parts of the construction fixed and use \(k\)
additional length-14 paths instead of 25. For every integer \(k\geq24\),
the same argument gives

\[n=57+14k,\qquad L_s=k+28,\qquad |B|=28,\qquad \operatorname{dist}_{\rm avg}(B)=4.\]

The bipartition classes have orders \(15+7k\) and \(42+7k\),
respectively. Thus the proposed right side exceeds the left side by

\[\frac{41+7k}{4}-(k+28)=\frac{3k-71}{4}>0.\]

In particular, \(k=24\) gives a 393-vertex counterexample; the \(k=25\)
witness was selected above because its claimed lower bound is the
integer 54 and the violation is independent of the distance-zero
convention.


\subsubsection{No positive constant-factor repair}
The same construction gives a stronger consequence. For every odd integer
$R\geq19$, use a length-$(R-1)$ left stem with $m=2R-3$ pendant leaves,
a length-$R$ right arm, and $k=R$ further length-$(R-1)$ arms.
Exactly the $m$ left leaves and the right endpoint are peripheral.
The same distance and parity counts give
\[
n=R^2+3R-3,\quad L_s=3R-2,\quad
\operatorname{dist}_{\rm avg}(B)=4,\quad
\max_v\operatorname{dist}_{\rm even}(v)=\frac{R^2+5R-6}{2}.
\]
The proposed right side exceeds $L_s$ by
\[
\frac{R^2-19R+8}{8}>0.
\]
For an integer-valued parametrization set $R=8t+19$, $t\geq0$. Then
\[
n=64t^2+328t+415,\qquad L_s=24t+55,
\]
\[
\text{proposed right side}=8t^2+43t+56,\qquad
\text{excess}=8t^2+19t+1.
\]
Its ratio to $L_s$ tends to infinity. Therefore there is no fixed
constant $c>0$ for which
\[
L_s(G)\geq c\,
\frac{\max_v\operatorname{dist}_{\rm even}(v)-1}
{\operatorname{dist}_{\rm avg}(B)}
\]
holds for all trees. This is a consequence of the same refutation, not
an additional numbered conjecture closure. No global minimum-order
claim is made for any displayed witness.


\subsection{Conjecture 382e: Maxine and two-domination}\label{late:382e}

\subsubsection{Statement, definitions, and prior construction}
All graphs are finite and simple. A set $D$ is 2-dominating if every
vertex outside $D$ has at least two neighbors in $D$. Its minimum size
is $\gamma_2(G)$. Ordinary domination requires one such neighbor and has
minimum size $\gamma(G)$. A \emph{Maxine run} repeatedly deletes a vertex
of maximum degree in the current induced graph until no edges remain.
We use $\operatorname{Maxine}(G)$ for the largest possible number of remaining vertices
over all tie-breaking choices. Original WOWII row 382e asserts
\begin{equation}\label{late:382e:eq:target}
 \gamma_2(G)\leq\operatorname{Maxine}(G)+\gamma(G)
\end{equation}
for every connected graph of order greater than two \cite{src:catalogue}.
Our examples have equal output size for \emph{every} run, so no
choice-dependent interpretation can remove the violation.

DeLaVi\~na, Larson, Pepper, and Waller prove the related valid inequality
$\gamma_2\leq\gamma+\alpha$ as a special case of their Theorem 6 and
Corollary 7 \cite{early:dpw-two}; this is not the Maxine inequality.
The forced-deletion mechanism has prior art: Cao and Eppstein's 2018
comments describe retaining the edges of a triangle-free cubic graph
and adding a degree-two vertex on each edge. Minimum-degree deletion
removes the new vertices first \cite{late:382e:ce}. We credit that mechanism.
Here the private vertices instead have \emph{three} prescribed neighbors,
including a common distinguished vertex. The changed neighborhoods and
a universal-vertex construction give the domination obstruction below.
We do not claim a smallest counterexample or first-ever worldwide discovery;
source-priority conclusions must be limited to the accompanying bounded audit.

\begin{example}[The attributed peeling precursor]\label{late:382e:precursor}
Let $k\ge3$, retain the edges of $K_{k,k}$, and for each core edge $ab$
add a private vertex adjacent exactly to $a,b$. Call the result $Q_k$
and put $P_k=\overline{Q_k}$. If $i$ denotes the minimum order of a
maximal independent set, then
\[
 i(P_k)=\alpha(P_k)=3,\qquad \operatorname{Maxine}(P_k)=2.
\]
For $J=K_1\vee P_k$ with new vertex $u$,
\[
 \gamma(J)=1,\quad \{u\}\text{ is its unique minimum dominating set},
 \quad i(J-u)=3>\operatorname{Maxine}(J)=2,
\]
but $\gamma_2(J)=3=\operatorname{Maxine}(J)+\gamma(J)$.
Thus these examples invalidate the corresponding independent-domination
bridges while still satisfying 382e. The construction and its peeling
behavior are attributed to Cao and Eppstein~\cite{late:382e:ce}.
\end{example}
\begin{proof}
Every maximal clique of $Q_k$ is a core edge together with its private
vertex. Consequently every maximal independent set of $P_k$ has order
three. In $Q_k$, private vertices have degree two, while each core vertex
retains at least its $k$ original core neighbors. Every minimum-degree
deletion therefore removes all private vertices first. From $K_{k,k}$,
minimum-degree deletion alternates between the two sides whenever their
sizes differ, eventually reaching $K_{1,1}$. Complementation exchanges
minimum-degree and maximum-degree deletion, and a clique in $Q_k$ is an
independent set in $P_k$. Every Maxine run thus ends with two vertices.

No vertex of $P_k$ is universal, since every vertex of $Q_k$ has at least
two neighbors. The new $u$ is therefore both the unique dominating
singleton of $J$ and its unique maximum-degree vertex. It is deleted
first by every Maxine run, proving the asserted residual values.
For two-domination, choose $u$ and the private vertices belonging to two
disjoint core edges. Every remaining core vertex is adjacent in $P_k$
to at least one of the two private vertices, and every other private
vertex is adjacent to both. This is a 2-dominating set of order three.
Every old vertex has at least two old nonneighbors, so a set of two
vertices containing an old vertex $x$ leaves some old nonneighbor of $x$
outside the set; that vertex has at most one selected neighbor. Hence
no two-element 2-dominating set exists.
\end{proof}

\subsubsection{The base graph}
Fix $k\geq4$. Let
\[
 A=\{a_0,a_1,\ldots,a_{k-1}\},\qquad
 B=\{b_0,b_1,\ldots,b_{k-1}\},\qquad
 X=\{x_{ij}:1\leq i\leq k-1,\ 0\leq j\leq k-1\}.
\]
Define $H_k$ as follows: retain every edge between $A$ and $B$;
for each $x_{ij}$ add exactly the three edges to $a_0,a_i,b_j$;
add no other edges. Put $F_k=\overline{H_k}$. Thus
$|F_k|=2k+k(k-1)=k^2+k$.
Equivalently, $A$, $B$, and $X$ each induce a clique in $F_k$;
there are no $A$--$B$ edges, and $x_{ij}$ is adjacent to every core
vertex except $a_0,a_i,b_j$. In particular,
\begin{equation}\label{late:382e:eq:special}
 N_{F_k}(a_0)=A\setminus\{a_0\}.
\end{equation}

\begin{lemma}\label{late:382e:lem:maxine}
Every Maxine run on $F_k$ leaves exactly two vertices, one in $A$ and one in $B$.
\end{lemma}
\begin{proof}
On any fixed surviving vertex set, maximum-degree deletion in a graph
is minimum-degree deletion in its complement. As long as a vertex of
$X$ remains in $H_k$, it has degree three, while every core vertex
has degree at least $k\geq4$, because the retained $K_{k,k}$ is still
present. Consequently every run deletes all of $X$ before any core
vertex. The remaining graph in $F_k$ is $K_k\sqcup K_k$.
Deleting positive-degree vertices cannot delete the last, isolated
vertex of either clique. The terminal set therefore has one vertex
from each clique.
\end{proof}

\begin{lemma}\label{late:382e:lem:domination}
For every $k\geq4$, $\gamma(F_k)=3$ and $\gamma_2(F_k)=4$.
\end{lemma}
\begin{proof}
First, $\{a_0,x_{10},x_{11}\}$ dominates $F_k$: $a_0$ dominates $A$,
each vertex of $B$ is adjacent to at least one of the two selected private
vertices, and every other private vertex is adjacent to both.
No pair dominates. A pair in $X$ misses $a_0$. A pair in the same core
part misses the other core part. A core pair $a_i,b_j$ misses $x_{ij}$
if $i\ne0$, and misses $x_{1j}$ if $i=0$. A pair $a_i,x_{pq}$ misses
$b_q$, while a pair $b_j,x_{pq}$ misses $a_0$. These cases exhaust all
pairs. Since a dominating singleton could be enlarged to a dominating
pair, $\gamma(F_k)=3$.

The four vertices $\{a_1,a_2,b_0,b_1\}$ form a 2-dominating set.
An unselected core vertex has two selected neighbors in its own part.
A private vertex misses at most one selected vertex in $A$ and at most
one in $B$, so it has at least two selected neighbors.

Suppose, for a contradiction, that a three-element set $D$ is
2-dominating. If $a_0\notin D$, equation~\eqref{late:382e:eq:special} forces
at least two vertices of $D$ into $A\setminus\{a_0\}$. The one possible
remaining vertex cannot supply two selected neighbors to any unselected
vertex of $B$. Such a vertex exists, so this is impossible.
Thus $a_0\in D$. The possibilities for the other two vertices are listed
below; in each row the indicated unselected vertex has fewer than two
neighbors in $D$.
\[
\begin{array}{c|l}
\text{Locations of }D\setminus\{a_0\}&\text{Witness to failure}\\\hline
 A,A & \text{any vertex of }B\\
 A,B & \text{a vertex of }B\setminus D\\
 B,B & \text{a vertex of }A\setminus\{a_0\}\\
 A,\ x_{ij} & b_j\\
 B,\ x_{ij} & a_i\\
 x_{ij},\ x_{pq} & b_j
\end{array}
\]
The needed unselected core vertices exist because $k\geq4$.
No smaller 2-dominating set exists either, since adjoining vertices
preserves 2-domination. Hence $\gamma_2(F_k)=4$.
\end{proof}

\subsubsection{Connected counterexamples and amplification}
Let $rF_k$ denote the disjoint union of $r$ copies of $F_k$.
Form $J_{k,r}=K_1\vee(rF_k)$ by adding a vertex $u$ adjacent to every
vertex in those copies.

\begin{theorem}\label{late:382e:thm:main}
For every $k\geq4$ and $r\geq1$, $J_{k,r}$ is connected and satisfies
\[
 |J_{k,r}|=1+r(k^2+k),\qquad
 \gamma(J_{k,r})=1,\qquad
 \operatorname{Maxine}(J_{k,r})=2r,\qquad
 \gamma_2(J_{k,r})=3r+1.
\]
Every Maxine run has the stated output size. Consequently,
\[
 \gamma_2(J_{k,r})-\bigl(\operatorname{Maxine}(J_{k,r})+\gamma(J_{k,r})\bigr)=r>0,
\]
so WOWII 382e is false, with unbounded additive defect.
\end{theorem}
\begin{proof}
The vertex $u$ is universal. No vertex in a copy of $F_k$ is universal,
since it has a nonneighbor already in that copy. Hence $u$ is the unique
maximum-degree vertex and is deleted first by every Maxine run.
The remaining graph is $rF_k$. Whenever a vertex in one component is
deleted, it has maximum degree within that component; ignoring the
interleaving therefore gives a valid local Maxine run. Lemma~\ref{late:382e:lem:maxine}
shows that each component contributes exactly two terminal vertices.
Thus $\operatorname{Maxine}(J_{k,r})=2r$, independently of ties.

Clearly $\gamma(J_{k,r})=1$. If a 2-dominating set contains $u$, its
intersection with each copy of $F_k$ must dominate that copy, and conversely
this condition suffices. Its minimum possible size is therefore $1+3r$.
If the set omits $u$, it must 2-dominate each copy separately, requiring
at least $4r$ vertices. Conversely, the union of minimum 2-dominating
sets in the copies also gives at least two neighbors to $u$.
It follows that
\[
 \gamma_2(J_{k,r})=\min\{1+3r,4r\}=3r+1,
\]
using $r\geq1$ and Lemma~\ref{late:382e:lem:domination}.
\end{proof}


\subsection{Conjecture 154: the ratio of radial-circle orders}\label{late:154}

\subsubsection{Exact statement and scope}\label{late:154:exact-statement-and-scope}

For a simple connected graph \(G\), let \(L_s(G)\) be the maximum number
of leaves in a spanning tree. Let \(r\) be its radius. A radial circle
is the set of vertices at distance \(r\) from a central vertex. Write
\(R_{\min}\) and \(R_{\max}\) for the smallest and largest orders of
these radial circles. Original WOWII 154 asserts

\[L_s(G)\geq\frac{1+R_{\max}}{R_{\min}}.\]

The construction below has \(L_s=13\), \(R_{\min}=1\) and
\(R_{\max}=13\). Thus it refutes the original assertion by the strict
inequality \(13<14\). It does \textbf{not} refute WOWII 155, whose
different right side counts distinct radial-circle orders.

The distinct original formulations of 154 and 155 are retained from the
archived rows~\cite{src:catalogue}. The following proof applies to the stated
nontrivial connected-graph formulation of 154.

\subsubsection{Construction}\label{late:154:construction}

In fact there is an infinite family, indexed by an integer \(t\geq2\).

Start with seven distinct vertices

\[v,w,a,u,b_1,b_2,c\]

and the six edges

\[vw,\ va,\ au,\ wb_1,\ wb_2,\ ac.\]

Thus these seven vertices induce a tree. For each stem
\(x\in\{b_1,b_2,c\}\) and each index \(j\in\{1,\ldots,t\}\), introduce
two new vertices \(p_{x,j},z_{x,j}\) and the four edges

\[xp_{x,j},\quad xz_{x,j},\quad p_{x,j}z_{x,j},\quad uz_{x,j}.\]

Finally introduce a vertex \(q\) with the single edge \(uq\). There are
no other edges. Denote the resulting graph by \(G_t\). It is visibly
simple and connected, with

\[n=8+6t.\]

The principal witness is \(G_2\), of order 20 and size 31.

\subsubsection{Connected domination number is exactly
seven}\label{late:154:connected-domination-number-is-exactly-seven}

Let \(\gamma_c\) denote minimum connected dominating set cardinality.
The original seven core vertices form a connected dominating set, so
\(\gamma_c(G_t)\leq7\).

We prove the reverse bound. Let \(D\) be any connected dominating set.
Necessarily \(u\in D\): otherwise the pendant vertex \(q\) must belong
to \(D\) to be dominated, but then it has no neighbor in \(D\) and
cannot be part of a connected dominating set of this graph.

For a stem \(x\), let its branch be

\[X_x=\{x\}\cup\{p_{x,j},z_{x,j}:1\leq j\leq t\}.\]

Its parent is \(w\) when \(x=b_1,b_2\), and \(a\) when \(x=c\). Two
observations give a branch-cost lower bound.

\begin{itemize}

\item
  If \(x\notin D\), then domination of each \(p_{x,j}\) requires \(D\)
  to meet \(\{p_{x,j},z_{x,j}\}\). These pairs are disjoint and
  \(t\geq2\), so \(|D\cap X_x|\geq2\).
\item
  If \(x\in D\) but its parent is absent from \(D\), then selecting no
  further vertex of \(X_x\) would leave \(x\) isolated in \(G[D]\),
  because its only neighbors outside \(X_x\) are its parent. Since
  \(u\in D\) and \(u\ne x\), this is impossible. Again
  \(|D\cap X_x|\geq2\).
\end{itemize}

Consequently a branch costs at least one selected vertex if its parent
is in \(D\), and at least two otherwise.

Write \(W=\mathbf1_{w\in D}\), \(A=\mathbf1_{a\in D}\) and
\(V=\mathbf1_{v\in D}\). Summing the three disjoint branch costs and the
required core selections gives

\[|D|\geq1+(4-2W)+(2-A)+W+A+V=7-W+V.\]

If \(|D|\leq6\), the only possibility is \(W=1\), \(V=0\), \(|D|=6\),
with equality in every cost estimate. In particular, each of the
branches \(X_{b_1},X_{b_2}\) contains exactly one selected vertex. It
must be its stem \(b_i\), since excluding that stem requires at least
two selected vertices. Thus neither branch contains any selected \(p\)-
or \(z\)-vertex. Since \(v\notin D\), the set \(\{w,b_1,b_2\}\) is then
a connected component of \(G[D]\) not containing the required vertex
\(u\). This contradicts connectivity.

Therefore

\[\boxed{\gamma_c(G_t)=7.}\]

For a connected graph of order at least three, \(L_s=n-\gamma_c\).
Indeed the internal vertices of any spanning tree form a connected
dominating set; conversely a connected dominating set can be spanned by
a tree and all remaining vertices attached as leaves. Hence

\[\boxed{L_s(G_t)=6t+1.}\]

\subsubsection{Radius and a radial circle of order
one}\label{late:154:radius-and-a-radial-circle-of-order-one}

The distance layers from \(v\) are exactly

\[\begin{aligned}
V_0&=\{v\},&V_1&=\{w,a\},\quad V_2=\{u,b_1,b_2,c\},\\
V_3&=\{q\}\cup\{p_{x,j},z_{x,j}:x\in\{b_1,b_2,c\},1\le j\le t\}.
\end{aligned}\]

Thus \(\operatorname{ecc}(v)=3\) and its third sphere has order
\(6t+1\).

There is no vertex of eccentricity at most two. Any such vertex would
have to be within distance two of the pendant vertex \(q\), and hence
belong to

\[\{q,u,a\}\cup\{z_{x,j}:x\in\{b_1,b_2,c\},1\leq j\leq t\}.\]

Each candidate fails:

\begin{itemize}

\item
  \(d(u,w)=3\) and \(d(q,w)=4\);
\item
  \(d(a,p_{b_1,1})=3\);
\item
  for any \(z_{x,j}\), choose a different stem \(y\ne x\). The vertices
  \(z_{x,j}\) and \(p_{y,1}\) are neither adjacent nor have a common
  neighbor, so their distance is at least three.
\end{itemize}

Since \(v\) has eccentricity three, the radius is exactly three.

The vertex \(u\) also has eccentricity three, and its unique vertex at
distance three is \(w\). Indeed \(u\) is adjacent to \(a\), every
\(z_{x,j}\) and \(q\); it reaches \(v\) through \(a\), every stem
through either \(a\) or a \(z\)-vertex, and every \(p\)-vertex through
its \(z\)-vertex, all within two steps. Its distance to \(w\) is three.
Therefore

\[\boxed{R_{\min}=1.}\]

\subsubsection{Maximum radial-circle order and
contradiction}\label{late:154:maximum-radial-circle-order-and-contradiction}

The radial circle at \(v\) has order \(6t+1=L_s(G_t)\). Every radial
circle has order at most \(L_s\): in a BFS spanning tree rooted at its
center, all vertices of the final distance layer are leaves. Thus

\[\boxed{R_{\max}=6t+1.}\]

Conjecture 154 would require

\[6t+1\geq\frac{1+(6t+1)}1=6t+2,\]

which is false for every integer \(t\geq2\). At \(t=2\), the explicit
contradiction is

\[\boxed{13<14.}\]

This is a full counterexample family to the original numbered
statement 154, not merely to an auxiliary strengthening.


\subsection{Conjecture 354 in the historical half-coefficient, open-neighbourhood form}\label{late:354}

\subsubsection{Exact statement and source qualification}
Write $\gamma_t(T)$ for the total domination number, $C(T)$ for the full center, and $D_2(T)=\{x:\deg_T(x)=2\}$. For a set $A$, let $N(A)=\bigcup_{a\in A}N(a)$ be the union of open neighborhoods. The notation $c(H)$ counts all connected components of $H$, including isolated vertices.

The pinned original WOWII row \#354, dated 18 February 2009, states that every tree $T$ of order greater than two satisfies
\begin{equation}
\gamma_t(T)\ge \frac{|C(T)|}{2}\,c\bigl(T[N(D_2(T))]\bigr).
\label{late:354:eq:target}
\end{equation}
This neighborhood is neither the external neighborhood $N(A)\setminus A$ nor the closed neighborhood $N[A]$.

The literal half coefficient and the union of open neighborhoods are
retained from the archived HTML formula~\cite{src:catalogue}. The source
audit found that an ASCII export dropped the half glyph; this is not
evidence for changing the target formula. In particular, a theorem about
a coefficient-one or closed-neighborhood variant does not decide
\eqref{late:354:eq:target}. The bounded historical audit did not complete
the comparison with the full Desormeaux--Henning 2016 article, so this
section makes no publication-priority claim.

\subsubsection{Feit's earlier construction and the modification}
The eleven-vertex branch used below is Feit's earlier construction.
His 34-vertex example for WoWII 240/346 consists of a root with three
copies of this branch. Feit's example and an unbounded-gap extension
are credited in DeLaVi\~na, Larson, Pepper and Waller
\cite[p.~7 of the author manuscript]{late:totaldom}. The branch is
therefore an attributed input to the bicentral modification below.

Feit's tree has total domination number 14 and 15 open-neighborhood components, but only one center: its radius is five and diameter ten. Its right-hand side in \eqref{late:354:eq:target} is therefore $15/2$, so that older example does not refute the half-coefficient formula.

Our principal 36-vertex tree is a \emph{bicentral modification of Feit's construction}: split the common root into adjacent vertices, distribute the three branches one on one side and two on the other, and add a leaf to the one-branch side. The last operation prevents either central vertex from having degree two. The contribution asserted here is this source-qualified modification/application, not the original branch or the earlier component-bound counterexample.

\subsubsection{The branch and the family}
Let $B$ have eleven distinct vertices
\[
a,b,c,d,e,f,g,h,i,j,\ell
\]
and edges
\[
ab,\ ac,\ bd,\ de,\ df,\ fg,\ ch,\ hi,\ hj,\ j\ell.
\]
The vertex $a$ is the attachment vertex. When $a$ is joined to a new vertex $r$, its branch vertices have the following distances from $r$:
\begin{center}
\begin{tabular}{cl}
\toprule
Distance & Vertices \\
\midrule
1 & $a$ \\
2 & $b,c$ \\
3 & $d,h$ \\
4 & $e,f,i,j$ \\
5 & $g,\ell$ \\
\bottomrule
\end{tabular}
\end{center}
The branch vertices of degree two are exactly $b,c,f,j$. Vertices $a,d,h$ have degree three; vertices $e,g,i,\ell$ are leaves. Consequently the open union neighborhood contributed by the four degree-two vertices is
\begin{equation}
\{a,d,h,g,\ell\}.
\label{late:354:eq:five}
\end{equation}
These five vertices are pairwise nonadjacent.

Fix $k\ge2$. Begin with an edge $uv$. Attach one disjoint copy of $B$ to $u$ and $k$ disjoint copies to $v$, joining each copy's vertex $a$ to its designated endpoint. Add a new leaf $z$ adjacent to $u$. There are no other edges. Denote the resulting tree by $T_k$.

Each branch is attached by one edge, so $T_k$ is a tree with
\[
|V(T_k)|=2+1+11(k+1)=11k+14,\qquad |E(T_k)|=11k+13.
\]
The principal example $T_2$ has 36 vertices and 35 edges.

\subsubsection{Degree-two neighborhood components}
The degrees of $u$ and $v$ are respectively 3 and $k+1$, neither of which is two. The new vertex $z$ has degree one. Therefore each branch contributes exactly $b,c,f,j$ to $D_2(T_k)$, and there are no other degree-two vertices.

By \eqref{late:354:eq:five}, the set $N(D_2(T_k))$ consists of five vertices in each branch. It contains none of $u,v,z$. Within each branch the five vertices are independent, and different branches have no edges between them. Hence
\begin{equation}
c\bigl(T_k[N(D_2(T_k))]\bigr)=5(k+1).
\label{late:354:eq:components}
\end{equation}
This calculation uses the open-union neighborhood directly, without replacing it by an external or closed neighborhood.

\subsubsection{The center}
Each branch has maximum depth five from its designated endpoint, and has vertices at that depth. Depth-five vertices on opposite sides of the edge $uv$ are at distance 11. Vertices on one side are at distance at most 10, and vertices on opposite sides are at distance at most $5+1+5=11$. Thus the diameter is 11.

The eccentricities of $u$ and $v$ are both six. Any other vertex $x$ has positive depth $q$ from the endpoint on its side, so its distance to a depth-five vertex on the opposite side is $q+1+5\ge7$. This includes the extra leaf $z$. Consequently,
\begin{equation}
C(T_k)=\{u,v\},\qquad \operatorname{rad}(T_k)=6.
\label{late:354:eq:center}
\end{equation}

\subsubsection{Exact total domination number}
Let $D$ contain $u,v$ and the four vertices $d,f,h,j$ from each branch. Its order is
\[
|D|=2+4(k+1)=4k+6.
\]
The centers dominate each other, and $u$ dominates $z$. Within a branch, $a$ is dominated by its center, $b$ by $d$, and $c$ by $h$. The adjacent pair $d,f$ dominates itself and the leaves $e,g$; the adjacent pair $h,j$ dominates itself and the leaves $i,\ell$. Thus $D$ is a total dominating set, proving $\gamma_t(T_k)\le4k+6$.

Conversely, the four leaves $e,g,i,\ell$ of each branch force their distinct neighbors $d,f,h,j$ into every total dominating set. The leaf $z$ forces $u$. These give $4(k+1)+1$ forced vertices. None of them is a neighbor of $u$: the neighbors of $u$ are $v$, the vertex $a$ of its single branch, and $z$. Total domination of $u$ therefore requires at least one additional vertex. It follows that
\begin{equation}
\gamma_t(T_k)=4k+6.
\label{late:354:eq:gamma}
\end{equation}
Only the upper certificate is necessary to refute \eqref{late:354:eq:target}; the lower argument establishes the exact value without an exhaustive search.

\subsubsection{The contradiction}
By \eqref{late:354:eq:components}, \eqref{late:354:eq:center}, and \eqref{late:354:eq:gamma}, the right-hand side of \eqref{late:354:eq:target} is $5(k+1)$, whereas its left-hand side is $4k+6$. Thus
\[
\frac{|C(T_k)|}{2}c\bigl(T_k[N(D_2(T_k))]\bigr)-\gamma_t(T_k)
=5k+5-(4k+6)=k-1>0.
\]
Every $T_k$ with $k\ge2$ is therefore a counterexample to the pinned historical formulation. The additive error is unbounded. In particular, $T_2$ gives
\[
\boxed{\gamma_t(T_2)=14<15=\tfrac12|C(T_2)|c\bigl(T_2[N(D_2(T_2))]\bigr).}
\]
No minimum-order claim is made.


\section{Leaf inequalities 177 and 180 on claw-free graphs}\label{late:claw}

\subsection{Scope}\label{late:claw:scope}

\textbf{Two subclass theorems, not full resolutions of the original
conjectures.} Both original inequalities allow arbitrary connected
graphs. This section proves them for connected claw-free graphs. No
numbered full closure or auxiliary-lemma priority claim is made.

Graphs are finite and simple. Write \(n=|V(G)|\), \(L\) for the maximum
number of leaves in a spanning tree, \(b\) for the largest order of an
induced bipartite subgraph, \(a=\alpha(G)\), and \(\gamma_c\) for
connected domination number. Let \(\delta\) be minimum degree and
\(\sigma=d_2\) the second-smallest degree, with multiplicity. For a root
\(v\), let \(O(v)\) be its odd-distance vertices and \(E(v)\) its
even-distance vertices; distance zero counts in \(E(v)\).

The following theorem proves the two source inequalities on the specified
class.
\begin{theorem}\label{late:claw:main}
Every nontrivial connected claw-free graph satisfies both
\[
 L+b\ge2a+\sigma\qquad\text{and}\qquad
 L+b\ge1+a+\max_{v\in V(G)}|E(v)|.
\]
These are, respectively, the inequalities in 177 and 180.
\end{theorem}

For \(n\ge3\), \(L=n-\gamma_c\). Put \(\eta=b-\gamma_c-1\), so
\(L+b=n+1+\eta\). The published DeLaVi\~na--Waller construction gives
\(\eta\ge0\). The two inequalities are therefore equivalent to

\[n+1+\eta\ge2a+\sigma,\qquad \eta\ge a-\min_v|O(v)|.\]

The graph \(K_2\) satisfies both directly. We henceforth assume
\(n\ge3\).

\subsection{Prior construction and the exact refinements
used}\label{late:claw:prior-construction-and-the-exact-refinements-used}

The construction below is the published method of DeLaVi\~na and Waller
\cite{late:DW}, with its prescribed-root retention and the odd-path
modification made explicit. It is not claimed as a new foundational
inequality.

\subsubsection{Rooted connected
domination}\label{late:claw:rooted-connected-domination}

For every prescribed vertex \(r\) of a nontrivial connected graph, there
is a connected dominating set containing \(r\) of size at most \(b-1\).

Start with a maximal connected induced bipartite set \(B_0\) containing
\(r\), grown from \(r\) and one neighbor. As long as the union of the
selected pieces does not dominate, choose a vertex \(x_j\) at distance
two from that union and a connector \(c_j\) adjacent to \(x_j\) and to
an earlier piece. In the graph remaining after deleting all earlier
pieces and their neighborhoods, grow a maximal connected induced
bipartite set \(B_j\) containing \(x_j\). Let the final pieces be
\(B_0,\ldots,B_k\). They are pairwise anticomplete, so their union \(B\)
is induced bipartite and \(|B|\le b\); there are at most \(k\) distinct
connectors.

Every uncolored vertex has neighbors of both colors in the first piece
it meets. Indeed, at the construction stage of that piece it had no
neighbor in an earlier piece and was eligible to be added; if it met
only one color, it could have been added without destroying
bipartiteness. In particular, a singleton piece is never the first piece
met by an uncolored vertex. Connectors stay uncolored because they
neighbor an earlier piece.

Choose a spanning tree of each piece, rooted at \(r\) in \(B_0\) and at
\(x_j\) in \(B_j\). Delete one leaf different from its root from every
nonsingleton piece; delete the sole vertex from any singleton later
piece. Retain all connectors. The retained initial piece is nonempty and
contains \(r\). Every nonempty later piece retains its root and connects
through \(c_j\) to an earlier retained piece: the connector retains a
neighbor in its first piece because it had neighbors in both colors.
Singleton pieces cannot interrupt this attachment, since none is a
connector's first piece.

The retained set dominates: each uncolored vertex keeps a neighbor in
its first piece; deleted tree leaves have retained neighbors; a deleted
singleton root is adjacent to its connector. Its order is at most
\(|B|+k-(k+1)\le b-1\).

\subsubsection{Zero slack excludes simplicial vertices except in
complete
graphs}\label{late:claw:zero-slack-excludes-simplicial-vertices-except-in-complete-graphs}

When \(\eta=0\), the rooted bound shows every vertex lies in a minimum
connected dominating set. Suppose a simplicial vertex \(v\) lies in such
a set \(D\). If \(D=\{v\}\), then \(v\) is universal and simplicial, so
\(G\) is complete. Otherwise choose a neighbor \(u\in D\) of \(v\).
Deleting \(v\) preserves connectedness of \(G[D]\), because its
neighborhood is a clique. It also preserves domination: \(u\) dominates
\(v\) and every vertex whose only neighbor in \(D\) was \(v\). This
contradicts minimality. Thus a noncomplete zero-slack graph has no
simplicial vertex and, in particular, \(\delta\ge2\).

\subsubsection{A maximal odd bipartite path forces positive
slack}\label{late:claw:a-maximal-odd-bipartite-path-forces-positive-slack}

If an inclusion-maximal connected induced bipartite subgraph is an
odd-order path of order at least three, then \(\gamma_c\le b-2\).

Use that path as \(B_0\) in the preceding construction. Instead of
deleting one initial tree leaf, delete both path endpoints. They have
the same color, so the retained nonempty path contains the entire
opposite color class. Every uncolored vertex meeting the initial path
retains a neighbor there. All connectivity and domination arguments
above remain valid, while one additional vertex was deleted. The
resulting connected dominating set has order at most
\(|B|+k-(k+2)\le b-2\).

\subsubsection{The published weaker form of
177}\label{late:claw:the-published-weaker-form-of-177}

For completeness, \(L+b\ge2a+1\) follows by the following short
argument, also published in \cite{late:DW}. Fix a maximum independent set
\(I\) and let \(C=V\setminus I\) have \(m\) induced connected
components. Selecting one vertex from each component together with \(I\)
gives an induced bipartite set of order \(a+m\), so \(b\ge a+m\). The
set \(C\) dominates \(G\), and at most \(m-1\) vertices of \(I\) suffice
to connect its components. Hence \(\gamma_c\le |C|+m-1\) and
\(L\ge a-m+1\). Adding gives the assertion.

\subsection{Independence core and odd layers in claw-free
graphs}\label{late:claw:independence-core-and-odd-layers-in-claw-free-graphs}

\subsubsection{Odd-layer injection}\label{late:claw:odd-layer-injection}

For every independent set \(I\) in a connected claw-free graph and every
root \(v\),

\[|I|\le |O(v)|+\mathbf 1_{v\in I}.\]

For each member \(x\in I\) at positive even distance \(2k\), choose a
neighbor \(p(x)\) in layer \(2k-1\). Two independent vertices cannot
choose the same predecessor \(p\): a preceding neighbor of \(p\) in
layer \(2k-2\) is nonadjacent to both of them, producing a claw. Map
each odd-distance member of \(I\) to itself. Its image cannot coincide
with a predecessor image, because that predecessor is adjacent to a
member of \(I\). This injects \(I\setminus\{v\}\) into \(O(v)\).

\subsubsection{An essential independent vertex forces an odd maximal
bipartite
path}\label{late:claw:an-essential-independent-vertex-forces-an-odd-maximal-bipartite-path}

The exchanges used here are part of the classical alternating-path approach
to independent sets in claw-free graphs; see Minty~\cite{late:Minty}.
We prove the precise assertion required below.


Suppose \(v\) belongs to every maximum independent set of a connected
claw-free graph. Fix one maximum independent set \(I\), and choose a
longest induced path \(P\) containing \(v\), alternating between \(I\)
and its complement, with both endpoints in \(I\). Initially the
singleton \(\{v\}\) is allowed.

Every internal non-\(I\) vertex of \(P\) already has two neighbors in
\(I\) on the path. Claw-freeness therefore prohibits further neighbors
in \(I\) outside the path.

The path has order at least three. A neighbor \(w\) of \(v\) lies
outside \(I\). If \(v\) were its only neighbor in \(I\), replacing \(v\)
by \(w\) would yield a maximum independent set avoiding \(v\). Thus
\(w\) has another neighbor \(z\in I\), giving the three-vertex
alternating induced path.

Suppose an outside vertex \(w\) could extend \(P\) while preserving
connectedness and induced bipartiteness. Every connected bipartite
claw-free graph has maximum degree at most two. Thus \(w\) either meets
one endpoint alone, or meets exactly the two endpoints, closing an even
cycle.

If it meets one endpoint and has a further neighbor
\(z\in I\setminus P\), the vertex \(z\) has no neighbor in \(P\): not
among its independent vertices, nor among its internal non-\(I\)
vertices, which already have two independent neighbors. Appending
\(w,z\) extends the alternating odd path, a contradiction. If it has no
such further neighbor, replace all \(I\) vertices of \(P\) by its
non-\(I\) vertices together with \(w\). The two sets have equal size;
the replacement is independent and has no neighbors in \(I\setminus P\).
It yields a maximum independent set avoiding \(v\), again a
contradiction.

If \(w\) meets both endpoints, it already has two neighbors in \(I\) and
therefore no neighbor in \(I\setminus P\). The identical exchange around
the even cycle yields the same contradiction.

No vertex can extend \(P\) bipartitely. Any larger connected induced
bipartite supergraph would supply a first outside neighbor of \(P\) that
could be added. Thus \(P\) is inclusion-maximal.

\subsubsection{Consequence at zero
slack}\label{late:claw:consequence-at-zero-slack}

In a claw-free graph with \(\eta=0\), the preceding lemma and Section
1.3 rule out every essential independent vertex. Consequently, each
vertex can be avoided by a maximum independent set. The odd-layer
injection then gives

\[a\le |O(v)|\quad\hbox{for every }v.\]

Equivalently, the intersection of all maximum independent sets---the
independence core---is empty.

\subsection{WOWII 180 on the claw-free
class}\label{late:claw:wowii-180-on-the-claw-free-class}

If \(\eta\ge1\), the injection gives \(a\le |O(v)|+1\le |O(v)|+\eta\).
If \(\eta=0\), Section~\ref{late:claw:consequence-at-zero-slack} gives \(a\le |O(v)|\). Taking the minimum
odd-distance count proves 180 for every connected claw-free graph.

This does not treat graphs containing induced claws.

\subsection{WOWII 177 on the claw-free
class}\label{late:claw:wowii-177-on-the-claw-free-class}

Complete graphs satisfy the inequality with equality. For \(\sigma=1\),
Section~\ref{late:claw:the-published-weaker-form-of-177} proves it. Assume henceforth \(a\ge2\) and \(\sigma\ge2\).

For a maximum independent set \(I\), every outside vertex has at most
two neighbors in \(I\). At most one vertex of \(I\) has degree below
\(\sigma\), so

\begin{equation}\label{late:claw:incidence}
 (a-1)\sigma+\delta\le e(I,V\setminus I)\le2(n-a).
\end{equation}

\subsubsection{A general degree bound}\label{late:claw:a-general-degree-bound}

Every noncomplete connected claw-free graph with \(\sigma\ge2\)
satisfies

\begin{equation}\label{late:claw:degree-bound}
 2a+\sigma\le n+2.
\end{equation}

If \(\sigma=2\), \eqref{late:claw:incidence} gives \(n\ge2a-1+\delta/2\), hence \(n\ge2a\) by
integrality. If \(a=2\), a noncomplete graph has two nonuniversal
vertices, so \(\sigma\le n-2\). If \(a=3\) and \(\sigma\ge3\), \eqref{late:claw:incidence} gives
\(n\ge\sigma+3+\delta/2\), hence \(n\ge\sigma+4\). Finally, if \(a\ge4\)
and \(\sigma\ge3\), \eqref{late:claw:incidence} gives

\[n-(2a+\sigma-2)\ge\frac{(a-3)\sigma+\delta-2a+4}{2}\ge\frac{a-4}{2}\ge0.\]

Thus \(\eta\ge1\) already proves 177.

\subsubsection{The zero-slack equality
profiles}\label{late:claw:the-zero-slack-equality-profiles}

Assume \(\eta=0\). Section~\ref{late:claw:zero-slack-excludes-simplicial-vertices-except-in-complete-graphs} gives \(\delta\ge2\) and no simplicial
vertex. When \(a\ge4\) and \(\sigma\ge3\), the final lower bound above
improves to \((a-3)/2>0\), so integrality yields \(2a+\sigma\le n+1\).
It remains to exclude equality in \eqref{late:claw:degree-bound} in the following three profiles.

\textbf{Profile A: \(\sigma=2\), \(n=2a\).} Equality in \eqref{late:claw:incidence} forces every
vertex of \(I\) to have degree two and every vertex of
\(J=V\setminus I\) to have exactly two neighbors in \(I\). If two
vertices of \(J\) are adjacent, they must have a common neighbor in
\(I\); otherwise one endpoint and its two independent neighbors form a
claw. The common neighbor then has degree two with adjacent neighbors
and is simplicial. Hence there are no edges inside \(J\). The graph is
connected and 2-regular bipartite, so is an even cycle. But an even
cycle has \(b=n\), \(\gamma_c=n-2\), and \(\eta=1\), contradiction.

\textbf{Profile B: \(a=3\), \(\sigma\ge3\), \(n=\sigma+4\).} Equality in
\eqref{late:claw:incidence} forces \(\delta=2\), one vertex \(x\in I\) of degree two, the other
two vertices \(y,z\in I\) of degree \(\sigma\), and exactly two
\(I\)-neighbors at every vertex of \(J\). Write \(N(x)=\{u,v\}\) and
\(W=J\setminus\{u,v\}\), where \(|W|=\sigma-1\). Every member of \(W\)
meets both \(y,z\). The degree counts force \(u\) to meet \(y\) but not
\(z\), and \(v\) to meet \(z\) but not \(y\), after relabeling.
Nonsimpliciality of \(x\) gives \(uv\notin E(G)\). Since \(u\) has
degree at least \(\sigma\ge3\), it has a neighbor \(w\in W\). The set
\(\{u,w,z\}\) is connected and dominating: it covers \(x,y\) through
\(u\), \(v\) and all of \(W\) through \(z\). Meanwhile \(y,u,x,v,z\)
induce a five-vertex path. Thus \(\gamma_c\le3\) and \(b\ge5\), giving
\(\eta\ge1\), contradiction.

\textbf{Profile C: \(a=2\), \(\sigma=n-2\).} A universal vertex gives
\(\gamma_c=1\) and \(b\ge a+1=3\), again positive slack. Without a
universal vertex, the complement has minimum degree at least one and at
most one vertex of degree greater than one. It is a disjoint union of a
star and matching edges. It cannot be a single spanning star, since that
would disconnect \(G\). Two complement edges from different components
induce a four-cycle in \(G\), so \(b\ge4\). Choose a vertex \(u\) of
degree \(n-2\), its unique nonneighbor \(v\), and a neighbor \(w\) of
\(v\). Then \(uw\) is an edge and \(\{u,w\}\) is connected dominating:
\(u\) dominates every vertex except \(v\), which is dominated by \(w\).
Hence \(\gamma_c\le2\), again giving positive slack.

All equality profiles are impossible at \(\eta=0\). Therefore
\(2a+\sigma\le n+1=L+b\), completing the claw-free proof of 177.


\section{Induced forests, degree sequences, and matching structure}
\label{late:induced-forests}
The results of Rounds 10, 12, 16, 17, 19, and 20 are presented here with
one common residue framework. The induced-matching candidate is left
open in general. Forest equality and a smallest-counterexample
restriction come first, followed by split-composition inheritance and
unigraphs. The first triangle-free residue equality then gives the
diameter-at-most-six case of 61. The final two subclasses use different
matching hypotheses and yield the stronger half-diameter term: graphs
with one feedback vertex use acyclic matchings, while the
matching-geodesic case uses an ordinary maximum matching on a
diametral path. These hypotheses are retained separately.

\subsection{Residue, induced matchings, and equality for forests}\label{late:forest}

\subsubsection{Scope, notation, and the unresolved candidate}

All graphs are finite, simple, and undirected. We allow empty and disconnected
graphs except where connectedness is explicitly assumed. Write $n(G)$ for
order, $\alpha(G)$ for independence number, $\nu(G)$ for ordinary matching
number, and $f(G)$ for the largest order of an induced forest. An
\emph{induced matching} is a matching whose endpoints induce exactly its
matching edges; its maximum cardinality is $\im(G)$. A feedback vertex is a
vertex whose deletion leaves a forest. A nontrivial star will mean $K_{1,d}$
with $d\ge2$, keeping isolated vertices and edges separate.

For a nonincreasing graphical degree list $D=(d_1,\ldots,d_n)$ with $d_1>0$,
the Havel--Hakimi step removes $d_1$, subtracts one from the next $d_1$
entries, and sorts again. The Havel--Hakimi theorem guarantees that the result
is graphical \cite{late:forest:Hakimi}. Repeating until only zeros remain defines the
\emph{residue} $\res(D)$ as the number of those zeros; set $\res(\varnothing)=0$ and
$\res(G)=\res(D(G))$. The elimination sequence lists the removed positive heads,
followed by the final zeros. It has length $n$ and total sum $|E(G)|$.

\begin{conjecture}\label{late:forest:cand}
For every finite simple graph $G$,
\begin{equation}\label{late:forest:eq:candidate}
f(G)\ge \res(G)+\im(G).
\end{equation}
\end{conjecture}

\noindent\textbf{Status distinction.}
Candidate~\ref{late:forest:cand} is left unresolved in this paper. Our proved results are the
forest equality characterization in Theorem~\ref{late:forest:thm:forest} and the
conditional smallest-counterexample restriction in
Theorem~\ref{late:forest:thm:noapex}. No worldwide novelty claim for these auxiliary
statements is made on the basis of a limited literature search.

WOWII Conjecture 61 asks, for a connected nontrivial graph of diameter $D$,
whether
\begin{equation}\label{late:forest:eq:wow61}
f(G)\ge \res(G)+\lceil D/3\rceil;
\end{equation}
see the recorded formulation \cite{late:forest:WOW61}. A diametral geodesic is induced.
Taking its edges with indices $0,3,6,\ldots$ gives an induced matching of size
$\lceil D/3\rceil$. Thus a general proof of \eqref{late:forest:eq:candidate} would prove
\eqref{late:forest:eq:wow61} directly. The unproved part is
\eqref{late:forest:eq:candidate} itself, not this implication.

\subsubsection{Classical residue dependencies}

For lists of equal length and equal sum, sorted nonincreasingly, write
$A\succeq B$ when every prefix sum of $A$ is at least the corresponding prefix
sum of $B$. We use the following established result, explicitly credited
rather than reproved here.

\begin{theorem}[Graphical elimination majorization; Triesch]\label{late:forest:thm:triesch}
If graphical degree lists $A,B$ have equal length and equal total degree, and
$A\succeq B$, then their Havel--Hakimi elimination sequences satisfy the same
majorization relation.
\end{theorem}

This is Triesch's theorem \cite{late:forest:Triesch}, stated with attribution as
Theorem 1.6 in Hiller's research manuscript \cite{late:forest:Hiller}. The equal-sum and
graphical hypotheses are essential to the way it is used here. The proof
does not assume that individual Havel--Hakimi successor lists preserve an
arbitrary majorization relation.

\begin{lemma}[Residue order]\label{late:forest:lem:resorder}
Under the hypotheses of Theorem~\ref{late:forest:thm:triesch}, $\res(A)\ge \res(B)$.
\end{lemma}
\begin{proof}
The two elimination sequences have the same total sum. If the second has $t$
positive entries, its first $t$ entries already sum to that total.
Majorization forces the first sequence's first $t$ entries to equal the total
as well, so it has at most $t$ positive entries. Its zero count is therefore
at least that of the second sequence.
\end{proof}

\begin{lemma}[Largest-entry decrements]\label{late:forest:lem:decrement}
Fix a nonnegative list and an integer $d$ no larger than its number of positive
entries. Among the lists obtained by decrementing $d$ distinct positive
entries by one, decrementing the $d$ largest entries produces a
majorization-minimal sorted list.
\end{lemma}
\begin{proof}
If a smaller entry is decremented and a larger entry is not, exchange these
choices. For original entries $a\ge b$, this changes the surviving pair
$(a,b-1)$ to $(a-1,b)$. When $a=b$ the change is only a permutation; otherwise
it is a unit balancing transfer and cannot increase sorted-prefix sums.
Repeated exchanges reach the largest-entry choice.
\end{proof}

\begin{lemma}[Positive maximum-degree deletion]\label{late:forest:lem:deletion}
If $\Delta(G)>0$ and $d_G(v)=\Delta(G)$, then $\res(G-v)\ge \res(G)$.
\end{lemma}
\begin{proof}
Remove the common maximum degree head $d=d_G(v)$. The canonical
Havel--Hakimi step decrements the $d$ largest tail entries; actual vertex
deletion decrements the $d$ neighbor entries. By
Lemma~\ref{late:forest:lem:decrement}, the latter list majorizes the former. Both are
graphical lists of the same order and sum. Lemma~\ref{late:forest:lem:resorder} gives
$\res(G-v)\ge \res(\operatorname{HH}(D(G)))=\res(G)$.
\end{proof}

The positive-degree hypothesis cannot be omitted: deleting a vertex from an
edgeless graph decreases residue by one.

\begin{corollary}[Classical residue bound]\label{late:forest:cor:alpha}
For every graph, $\res(G)\le\alpha(G)$.
\end{corollary}
\begin{proof}
Induct on order. For an edgeless graph both sides equal its order. Otherwise
choose a maximum-degree vertex $v$. By Lemma~\ref{late:forest:lem:deletion} and induction,
$\res(G)\le \res(G-v)\le\alpha(G-v)\le\alpha(G)$.
\end{proof}

This is the established Havel--Hakimi residue bound, proved in the classical
literature, including Griggs and Kleitman \cite{late:forest:GK}; the derivation above
records our exact dependency chain.

\begin{lemma}[Classical disjoint-union subadditivity]\label{late:forest:lem:union}
For any two disjoint graphs $G,H$,
\[
\res(G\sqcup H)\le \res(G)+\res(H).
\]
\end{lemma}
\begin{proof}
Induct on total order, writing $A,B$ for the component degree lists. If both
are zero lists, equality is immediate. Otherwise choose a largest pooled
positive head, say in $A$, of value $d$. Let $C$ be the global canonical
reduction of $A\cup B$. Instead decrementing the $d$ largest remaining
entries of $A$ alone gives $S=\operatorname{HH}(A)\cup B$. This is feasible because $A$ is
graphical. Both $C$ and $S$ are graphical; they have the same length and sum.
By Lemmas~\ref{late:forest:lem:decrement} and \ref{late:forest:lem:resorder}, $S\succeq C$ and
$\res(C)\le \res(S)$. Induction on the shorter total length now gives
\[
\res(A\cup B)=\res(C)\le \res(S)
\le \res(\operatorname{HH}(A))+\res(B)=\res(A)+\res(B).
\]
\end{proof}

Lemma~\ref{late:forest:lem:union} is the $k=1$ case of the established Disjoint Union
Lemma, Theorem 11 of Amos, Davila and Pepper \cite{late:forest:ADP}. Their Lemma 8 gives
the maximum-degree residual comparison, and their Theorem 16 combines the two
facts at a maximum-degree cut vertex. The proof here makes the dependencies
transparent; it is not a claim to a new residue theorem.

\subsubsection{Weak majorization by matching padding}
The preceding majorization theorem requires equal degree sums. The
following extension, extracted from the Round 10 auxiliary note, removes
that requirement on graphical inputs. It remains a sequence theorem;
it does not supply the graph-level forest augmentation needed for a
general proof of 61.

For a nonincreasing list $A$, let $P_k(A)$ be its first $k$ entries'
sum, with zero padding when necessary. Write $E(A)$ for its canonical
Havel--Hakimi elimination sequence, including terminal zeros. This
sequence is nonincreasing, and its total is one half the degree sum.

\begin{lemma}[Matching-padding identity]\label{late:forest:padding-identity}
For a graphical list $B$ and an integer $m\ge0$,
\begin{equation}\label{late:forest:padding-eq}
 E(B\cup1^{2m})
 =\operatorname{sort}\bigl(E(B)\cup1^m\cup0^m\bigr).
\end{equation}
Adding isolated vertices likewise appends only zeros to the elimination
sequence.
\end{lemma}
\begin{proof}
Run the canonical reductions on the original $B$ portion first. At a
positive head $h$ of that portion, graphicality guarantees at least $h$
positive entries after the head: a realization has $h$ distinct neighbors
of the head vertex. Each entry selected by the $B$-only Havel--Hakimi
step is therefore at least one. The appended ones can tie with these
entries but cannot strictly displace them. Resolve ties in favor of the
$B$ entries, leaving every appended one unchanged. This is a valid
canonical reduction of the combined multiset.

When the $B$ portion reaches zeros, the $2m$ appended ones produce
exactly $m$ further eliminated heads equal to one and $m$ additional
terminal zeros. A tie-breaking choice cannot change the reduced degree
multiset or the later head values. This proves
\eqref{late:forest:padding-eq}. Appended zeros cannot affect a positive
head or a decrement, proving the final assertion.
\end{proof}

\begin{theorem}[Weak prefix order for graphical elimination]
\label{late:forest:weak-majorization}
Let $A,B$ be graphical degree lists, possibly of different lengths.
If $P_k(A)\ge P_k(B)$ for every $k$, then
\[
 P_k(E(A))\ge P_k(E(B))\qquad\text{for every }k.
\]
Equality of the total degree sums is not required.
\end{theorem}
\begin{proof}
Pad $A,B$ with zeros to a common length $n$. Their nonnegative degree-sum
difference is even; write it as $2m$. Form sorted graphical lists
\[
 A'=A\cup0^{2m},\qquad B'=B\cup1^{2m}.
\]
They have equal lengths and sums: the first is realized by adjoining
$2m$ isolates to a realization of $A$, and the second by adjoining $m$
disjoint edges to a realization of $B$.

Let $p$ be the number of positive entries in $A$. If $k\le p$, a
selection of $k$ entries from $B'$ uses some $q\le k$ entries from
$B$ and $k-q$ appended ones. Its sum is at most
\[
 P_q(B)+(k-q)\le P_q(A)+(k-q)\le P_k(A),
\]
because the first $k$ entries of $A$ are positive integers. In
particular, this bounds the maximum-sum selection and gives
$P_k(B')\le P_k(A')$. For $k>p$, the latter prefix is already the
entire sum of $A'$, equal to the sum of $B'$. Thus $A'\succeq B'$.

Triesch's theorem, Theorem~\ref{late:forest:thm:triesch}, now gives
$P_k(E(A'))\ge P_k(E(B'))$. Adding isolates does not change the left
prefix sum. By Lemma~\ref{late:forest:padding-identity}, the multiset
defining $E(B')$ contains every entry of $E(B)$ and extra nonnegative
entries. Therefore its largest $k$ entries sum to at least $P_k(E(B))$.
Together these facts give the asserted weak prefix order.
\end{proof}

\begin{corollary}[A residual degree-sum bound]\label{late:forest:gap}
Under the hypotheses of Theorem~\ref{late:forest:weak-majorization}, let
$A_k,B_k$ be the residual lists after $k$ canonical steps, padding a
completed history by zero heads. Then
\[
 \sum A_k-\sum B_k\le\sum A-\sum B.
\]
\end{corollary}
\begin{proof}
Each eliminated head $h$ removes $2h$ degree-sum units. Hence
\[
 \sum A_k-\sum B_k
 =\sum A-\sum B
  -2\bigl(P_k(E(A))-P_k(E(B))\bigr),
\]
and the theorem bounds the parenthesized expression below by zero.
\end{proof}
This is not a claim that the residual lists themselves continue to
majorize each other at every step. Nor is a statement about arbitrary
nongraphical input lists being made. The proof depends on the credited
equal-sum graphical theorem and the exact padding identity.

\subsubsection{Equality for forests}

We use the forest identity $\alpha(F)=|V(F)|-\nu(F)$. For completeness, a
forest's minimum vertex-cover number equals its matching number: at a leaf
$x$ with neighbor $y$, some minimum vertex cover contains $y$ rather than
$x$, and some maximum matching contains $xy$. Both parameters therefore obey
the recurrence $1$ plus the corresponding parameter on $F-\{x,y\}$, with
edgeless base zero. Independence number is order minus vertex-cover number.

\begin{theorem}\label{late:forest:thm:forest}
For every finite forest $F$,
\begin{equation}\label{late:forest:eq:forest}
|V(F)|\ge \res(F)+\im(F).
\end{equation}
Equality holds if and only if every component is an isolated vertex or an
edge, except possibly one component that is a star $K_{1,d}$ with $d\ge2$.
\end{theorem}

\begin{proof}
The inequality follows immediately from
\begin{equation}\label{late:forest:eq:chain}
\res(F)\le\alpha(F)=|V(F)|-\nu(F)
\le |V(F)|-\im(F).
\end{equation}
We determine when it can be tight.

\medskip\noindent\emph{Step 1: connected nonstar trees are strict.}
If $\nu(T)>\im(T)$, strictness follows already from \eqref{late:forest:eq:chain}.
Otherwise let $a=\nu(T)=\im(T)$, and fix an induced matching $M$ of size $a$.
It is also a maximum ordinary matching.

Every edge $xy\in M$ has a leaf endpoint. If both endpoints had other
neighbors $u,w$, inducedness of $M$ would put $u,w$ outside its endpoint set.
They are distinct because $T$ has no triangle. Replacing $xy$ by $ux,yw$
would augment $M$, a contradiction.

Choose the nonleaf endpoint of each matching edge as its support, choosing
either endpoint if the edge is a component on its own. The set $A$ of
supports is independent, has cardinality $a$, and each support has a private
leaf partner. It is a vertex cover: unmatched vertices are independent, and
a leaf partner is incident only with its matching edge. Thus every edge meets
$A$.

Since $T$ is connected and not a star, $a\ge2$, and two distinct supports
$x,y$ have a common neighbor $z$ on a path between supports. Let $x',y'$ be
their private leaf partners. Perform the degree-preserving switch
\[
xx',\ yy'\quad\longmapsto\quad xy,\ x'y'.
\]
The new edges were absent, because $A$ is independent and $x',y'$ are private
leaves. The resulting simple graph $T'$ has the same degree sequence and
residue as $T$. It contains, on mutually disjoint vertex sets, the triangle
$xyz$, the edge $x'y'$, and the other $a-2$ edges of $M$. Every vertex cover of
$T'$ therefore has size at least $2+1+(a-2)=a+1$, regardless of additional
cross edges. Corollary~\ref{late:forest:cor:alpha} gives
\[
\res(T)=\res(T')\le\alpha(T')\le |V(T)|-a-1.
\]
Hence $\res(T)+\im(T)\le |V(T)|-1$, as required.

\medskip\noindent\emph{Step 2: any nonstar component forces strictness.}
Induced matching number is additive across components. If a forest has a
nonstar component, Lemma~\ref{late:forest:lem:union} and Step 1 give
\[
\res(F)+\im(F)
\le\sum_{T\text{ component}}\bigl(\res(T)+\im(T)\bigr)
<\sum_T |V(T)|=|V(F)|.
\]
We do not assert $\res(T)<\alpha(T)$ for every nonstar tree: for example $P_4$
has $\res(P_4)=\alpha(P_4)=2$, and its strictness instead comes from $\nu>\im$.

\medskip\noindent\emph{Step 3: two nontrivial star components force strictness.}
Suppose all components are stars, edges, or isolates, and two star centers
$x,y$ each have at least two leaves. Choose distinct leaves $x_1,x_2$ at $x$
and $y_1,y_2$ at $y$, and switch
\[
xx_1,\ yy_1\quad\longmapsto\quad xy,\ x_1y_1.
\]
In this degree-equivalent realization $F'$, the three edges
$xx_2,yy_2,x_1y_1$, together with one edge from every other nontrivial
component, form a matching of size $\nu(F)+1$. In every graph, a vertex cover
has size at least that of any matching. Consequently
\[
\res(F)=\res(F')\le\alpha(F')\le |V(F)|-\nu(F)-1.
\]
For a union of stars, $\nu(F)=\im(F)$, so \eqref{late:forest:eq:forest} is strict.

\medskip\noindent\emph{Step 4: the remaining forests attain equality.}
For $K_{1,d}\sqcup eK_2\sqcup iK_1$ with $d\ge2$, the degree list is
$(d,1^{d+2e},0^i)$. Its first Havel--Hakimi step leaves $2e$ ones and $d+i$
zeros. Thus
\[
\res=d+e+i,\qquad \im=1+e,\qquad n=d+1+2e+i,
\]
giving equality. If there is no nontrivial star, $eK_2\sqcup iK_1$ has
$\res=e+i$, $\im=e$, and $n=2e+i$. This includes the empty forest.
\end{proof}

\subsubsection{A smallest-counterexample restriction}

The assertions in this section are conditional on a counterexample to
Candidate~\ref{late:forest:cand} existing. They do not assume its truth on graphs of the
same or larger order.

\begin{lemma}\label{late:forest:lem:tightdelete}
Let $G$ have minimum order among all counterexamples to
Candidate~\ref{late:forest:cand}. Then $G$ is connected and is not a forest. For every
maximum-degree vertex $v$, with $H=G-v$, one has
\begin{align}
\res(H)&=\res(G),\label{late:forest:eq:r-equal}\\
\im(H)&=\im(G)-1,\label{late:forest:eq:im-equal}\\
f(G)=f(H)&=\res(G)+\im(G)-1.\label{late:forest:eq:f-equal}
\end{align}
Moreover, $H$ is candidate-tight, residue is additive across its components,
and every component is candidate-tight.
\end{lemma}
\begin{proof}
If $G$ were disconnected, minimality on its components, additivity of $f$ and
$\im$, and Lemma~\ref{late:forest:lem:union} would prove the candidate for $G$.
Theorem~\ref{late:forest:thm:forest} excludes forests. In particular $\Delta(G)>0$.

Deleting one vertex destroys at most one edge of any induced matching.
Minimality and Lemma~\ref{late:forest:lem:deletion} therefore give
\[
f(G)\ge f(H)\ge \res(H)+\im(H)
\ge \res(G)+\im(G)-1.
\]
Since $G$ fails the integer-valued inequality \eqref{late:forest:eq:candidate}, every
inequality in this chain is equality. This proves
\eqref{late:forest:eq:r-equal}--\eqref{late:forest:eq:f-equal}. Applying the same component sum to
$H$ shows that neither a strict residue-sum inequality nor a surplus in any
component is possible.
\end{proof}

\begin{theorem}\label{late:forest:thm:noapex}
If $G$ is a smallest counterexample to Candidate~\ref{late:forest:cand}, then $G-v$
contains a cycle for every maximum-degree vertex $v$.
\end{theorem}
\begin{proof}
Suppose $H=G-v$ were a forest. By Lemma~\ref{late:forest:lem:tightdelete}, it is tight in
\eqref{late:forest:eq:forest}, so Theorem~\ref{late:forest:thm:forest} applies.
Equation~\eqref{late:forest:eq:im-equal} says that a maximum induced matching of $G$
consists of an edge $vu$ and $\im(H)$ edges in $H$.

The vertex $u$ must be isolated in $H$. If it belonged to an edge component,
that component's one matching edge could not be retained. If it belonged to
the possible star component, every star edge would meet or have an endpoint
adjacent to $u$, so the star's one matching edge could not be retained either.
There is no way to compensate elsewhere, because the other retained edges
already attain the sum of component induced matching numbers.

There is also no edge component in $H$: retaining its sole edge would force
$v$ to have no neighbor in that component, contradicting connectedness of
$G$. If there were no nontrivial star component either, $G$ would be a star,
contrary to Lemma~\ref{late:forest:lem:tightdelete}. Hence $H$ consists of a star $K_{1,d}$
with center $w$, together with $p\ge1$ isolates.

To retain a star edge in the induced matching, $v$ is nonadjacent to $w$ and
misses at least one star leaf. Connectedness forces it to meet some star leaf.
Let it meet exactly $a$ of them. All $p$ isolates of $H$ are neighbors of $v$.
Since $G$ is not a forest, $a\ge2$; also $a\le d-1$, so $d\ge3$.
Maximum degree of $v$ implies $a+p\ge d$.

The degree list of $G$ is thus
\[
(a+p,\ d,\ 2^a,\ 1^{d-a+p}),
\qquad a+p\ge d\ge3.
\]
Choose the maximum head $a+p$ in the first Havel--Hakimi step. It decrements
the degree $d$, all $a$ entries equal to two, and $p-1$ ones. The residual
list is
\[
(d-1,\ 1^{d+1},\ 0^{p-1}).
\]
The next positive head $d-1$ leaves exactly two ones; the following head
removes their edge. There are exactly three positive heads in total. Hence
\[
\res(G)=n(G)-3=d+p-1,
\qquad \res(H)=d+p,
\]
contradicting \eqref{late:forest:eq:r-equal}.
\end{proof}

\begin{remark}[Exact scope]
Theorem~\ref{late:forest:thm:noapex} is a restriction on a \emph{smallest counterexample}
to the unproved candidate. It is not a proof of the candidate for every graph
of feedback vertex number one: a feedback vertex in such a graph need not
have maximum degree. It does not resolve original WOWII 61.
\end{remark}

\subsubsection{Why a simpler maximum-degree induction does not suffice}

A tempting replacement for the unresolved step would assert that some
maximum-degree vertex can be avoided by a maximum induced matching whenever
the graph has no leaves. This is false even at arbitrarily large minimum
degree.

\begin{proposition}\label{late:forest:prop:retired}
For integers $m\ge2$ and $k\ge2m$, let
$G=I_k\vee(mK_2)$, the join of an independent set of size $k$ and $m$ disjoint
edges. Every maximum-degree vertex belongs to every maximum induced matching,
and $\delta(G)=2m$.
\end{proposition}
\begin{proof}
The matching endpoints have degree $k+1>2m$, while independent-set vertices
have degree $2m$, so the endpoints are exactly the maximum-degree vertices.
The original $m$ edges form an induced matching. Any induced matching using
a join edge has size only one: its independent-set endpoint is adjacent to
every matching endpoint, and there are no edges inside the independent set.
Since $m\ge2$, the unique maximum induced matching is therefore the set of all
$m$ original edges.
\end{proof}

For clarity, this family does satisfy Candidate~\ref{late:forest:cand}. Label its matching
edges $a_jb_j$ for $1\le j\le m$. For each $j<m$, choose two new, distinct
vertices $x_j,y_j$ of $I_k$, and switch
$b_jx_j,a_{j+1}y_j$ to $b_ja_{j+1},x_jy_j$. This is possible since
$k\ge2m$. In the degree-equivalent realization, the matching endpoints induce
a path on $2m$ vertices, and $I_k$ induces $m-1$ disjoint edges and isolates.
Each vertex has lost at most one join incidence, with distinct missing
partners. An independent set meeting both sides has size at most two; one
within the path has size at most $m$, and one within the other side has size
at most $k-m+1$. Hence $\res(G)\le k-m+1$. The original graph contains a star
on $k+1$ vertices, so $\res(G)+\im(G)\le k+1\le f(G)$.

The family shows that a minimum-degree threshold cannot justify the proposed
deletion shortcut, even though the target inequality itself holds on the
family. The remaining mathematical gap is to rule out a smallest
counterexample whose every maximum-degree deletion leaves tight cyclic
components as in Lemma~\ref{late:forest:lem:tightdelete}.


\subsection{Inheritance under split composition and the unigraph class}\label{late:unigraph}

\subsubsection{Definitions and scope}

A \emph{unigraph} is a graph determined up to isomorphism by its degree
sequence. We retain the residue, forest number, and induced-matching
notation from Section~\ref{late:forest}, including the unresolved
candidate~\eqref{late:forest:eq:candidate}. Recall the classical bound
\begin{equation}\label{late:unigraph:eq:classical}
 \res(G)\le\alpha(G)
\end{equation}
from Corollary~\ref{late:forest:cor:alpha}; Favaron, Mah\'eo and
Sacl\'e~\cite{late:FMS} are an earlier primary source.
The subclass result is the following.

\begin{theorem}\label{late:unigraph:thm:main}
Inequality~\eqref{late:forest:eq:candidate} holds for every unigraph. It also holds
for every graph obtained from a unigraph by any finite sequence of split
compositions; the outer split factors need not be unigraphs.
\end{theorem}

A \emph{split graph} has a specified partition $A\mathbin{\dot\cup} B$ into an
independent set $A$ and a clique $B$. For a nonempty graph $J$, its split
composition $S(A,B)\circ J$ is formed from disjoint copies of $S$ and $J$
by adding every edge between $B$ and $V(J)$, and no edge between $A$ and
$V(J)$. Either $A$ or $B$ may be empty.

We explicitly credit two inputs from Barrus~\cite{late:unigraph:Barrus}:
\begin{align}
\res(S)&=\alpha(S) &&\text{for split graphs $S$}
&&\text{\cite[Proposition 2.2]{late:unigraph:Barrus}},\label{late:unigraph:eq:split}\\
\res(S(A,B)\circ J)&=|A|+\res(J) &&\text{for nonempty $J$}
&&\text{\cite[Corollary 3.3]{late:unigraph:Barrus}}.\label{late:unigraph:eq:compose}
\end{align}
These residue identities, and the canonical classification used below,
are existing results, not contributions claimed by this section.

\subsubsection{Inheritance under split composition}

For any graph with an edge,
\begin{equation}\label{late:unigraph:eq:forest-independence}
f(G)\ge \alpha(G)+1.
\end{equation}
Indeed, adjoining any vertex outside a maximum independent set produces
an induced star together with isolated vertices. A split graph with an
edge has $\im(G)=1$, since every edge meets its clique part. Equations
\eqref{late:unigraph:eq:split} and \eqref{late:unigraph:eq:forest-independence} therefore prove
\eqref{late:forest:eq:candidate} for split graphs with an edge. For edgeless graphs,
$f=\res=|V|$ and $\im=0$.

\begin{proposition}\label{late:unigraph:prop:inherit}
If a nonempty graph $J$ satisfies~\eqref{late:forest:eq:candidate}, then so does
$S(A,B)\circ J$ for every split graph with specified partition $A\mathbin{\dot\cup} B$.
\end{proposition}
\begin{proof}
If $J$ is edgeless, the composition is split and has already been handled.
Suppose that $J$ has an edge, and put $G=S(A,B)\circ J$. Deleting $B$
leaves $J$ and $|A|$ isolates, so
\[
f(G)\ge |A|+f(J).
\]
Every induced matching of size at least two lies entirely in $J$.
An edge not lying in $J$ has an endpoint in $B$, and that endpoint is
adjacent to every vertex of $J$ and every other vertex of $B$; it therefore
conflicts with every other possible matching edge. Since $J$ has an edge,
$\im(G)=\im(J)$. By~\eqref{late:unigraph:eq:compose},
\[
f(G)\ge |A|+\res(J)+\im(J)=\res(G)+\im(G).
\]
\end{proof}

\begin{corollary}\label{late:unigraph:cor:indecomp}
A smallest counterexample to~\eqref{late:forest:eq:candidate}, if one exists, is
indecomposable under nontrivial split composition.
\end{corollary}
\begin{proof}
In a nontrivial representation $S(A,B)\circ J$, the inner graph $J$ is
smaller. Minimality and Proposition~\ref{late:unigraph:prop:inherit} would close the
inequality. The edgeless-inner case is split and is also covered.
\end{proof}

\subsubsection{The established unigraph classification}

Tyshkevich's canonical decomposition~\cite{late:unigraph:Tyshkevich}, in the form used
by Barrus~\cite[Theorems 4.1--4.2]{late:unigraph:Barrus}, expresses every unigraph as a
sequence of split compositions over an indecomposable unigraph core.
If that core is nonsplit, then, up to complementation, it is one of
\begin{equation}\label{late:unigraph:eq:cores}
C_5,\qquad mK_2\ (m\ge2),\qquad
mK_2\mathbin{\dot\cup} K_{1,n}\ (m\ge1,\ n\ge2),\qquad U_m\ (m\ge1).
\end{equation}
Here $U_m$ consists of a chordless four-cycle and $m$ triangles sharing
one common vertex and otherwise vertex-disjoint. The classification is
an external theorem; we do not purport to reprove it. By
Proposition~\ref{late:unigraph:prop:inherit}, it remains to check these cores and their
complements. Split cores were handled above.

\begin{lemma}\label{late:unigraph:lem:one}
A nonempty graph has residue one if and only if it is complete.
\end{lemma}
\begin{proof}
The complete-graph direction is immediate. For the converse, induct on
the order $N$. The one-vertex case is immediate. If $N>1$ and the residue
is one, the first HH head is positive and its residual list has residue
one, hence is the complete-graph list of length $N-1$ by induction.
If the first head $d$ were less than $N-1$, a decremented tail entry
would originally have degree $N-1>d$, contradicting maximality of the
head. Thus $d=N-1$ and every original entry is $N-1$.
\end{proof}

\paragraph{The five-cycle and the star forests}

For $C_5$, the values are $\res=2$, $\im=1$ and $f=4$.
For the two forest families in~\eqref{late:unigraph:eq:cores}, direct computation gives
\begin{align*}
(f,\res,\im)(mK_2)&=(2m,m,m),\\
(f,\res,\im)(mK_2\mathbin{\dot\cup} K_{1,n})&=(2m+n+1,m+n,m+1).
\end{align*}
The residue calculation removes the star head, when present, and then
pairs off the remaining ones. Both forest families satisfy
\eqref{late:forest:eq:candidate} with equality.

The complement of either listed forest has independence two. It has no
induced $2K_2$, since that would correspond to an induced four-cycle in
the forest. Thus its induced matching number is one. By
\eqref{late:unigraph:eq:classical} and Lemma~\ref{late:unigraph:lem:one}, its residue is two, and
\eqref{late:unigraph:eq:forest-independence} gives $f\ge3$. This proves the required
inequality for both complementary families as well.

\paragraph{The shared-cycle graph and its complement}

The degree list of $U_m$ is
\[
(2m+2,2^{\,2m+3}),
\]
where exponents indicate multiplicities. One HH step leaves
$(2,1^{2m+2})$, the degree list of $mK_2\mathbin{\dot\cup} P_3$. Hence
$\res(U_m)=m+2$. Deleting the common vertex leaves a forest on $2m+3$
vertices; since the original graph contains a cycle, this is maximum.
An induced matching avoiding that vertex takes one edge from each of the
$m$ triangles and one edge from the residual $P_3$, for $m+1$ edges.
If the common vertex is used, no second edge can coexist: its only
nonneighbor is the opposite four-cycle vertex. Therefore
\[
f(U_m)=2m+3=\res(U_m)+\im(U_m),\qquad \im(U_m)=m+1.
\]

The complementary degree list is
\[
((2m+1)^{\,2m+3},1).
\]
For $m\ge1$, two HH steps produce
\[
((2m-1)^{\,2m+1},1),
\]
the same list with parameter $m-1$. At $m=0$ this is $1^4$, of residue
two. Consequently $\res(\overline{U_m})=2$.
The original four-cycle induces $2K_2$ in the complement, giving both
an induced matching of size two and an induced forest of order four.
There cannot be an induced matching of size three: its six endpoints
would induce minimum degree four in $U_m$, although every vertex of
$U_m$ except the common vertex has total degree two. It follows that
\[
\im(\overline{U_m})=2,\qquad
f(\overline{U_m})\ge4=\res(\overline{U_m})+\im(\overline{U_m}).
\]
All cores in the credited classification are now covered. Repeated
application of Proposition~\ref{late:unigraph:prop:inherit} proves
Theorem~\ref{late:unigraph:thm:main}.

\subsubsection{Relation to the original diameter conjecture}

As explained in Section~\ref{late:forest}, a general proof of
\eqref{late:forest:eq:candidate} would imply original~61.
Theorem~\ref{late:unigraph:thm:main} establishes only the specified subclass.

More importantly, connected unigraphs already have diameter at most
three. Each connected nonsplit core in~\eqref{late:unigraph:eq:cores} or its complement
has diameter at most three. In a connected nontrivial split composition,
the outer clique is nonempty, every outer independent vertex has a
neighbor in it, and every inner vertex is adjacent to it. Two outer
independent vertices are therefore at distance at most three, and all
other pairs at distance at most two. Connected split graphs have the same
diameter bound. Thus~\eqref{late:forest:eq:wow61} on unigraphs is already the elementary
inequality~\eqref{late:unigraph:eq:forest-independence} combined with
\eqref{late:unigraph:eq:classical}. It is not a new original-conjecture advance.

The actual auxiliary strengthening here is the induced-matching term,
which is unbounded on the family $U_m$. This is an auxiliary strengthening, not a full resolution of a numbered
WoWII conjecture.


\subsection{The first residue equality in triangle-free graphs}\label{late:equality}

\subsubsection{Definitions and scope}
Retain the notation of Section~\ref{late:forest}. The residue depends
only on the degree sequence, and the classical bound is
\begin{equation}\label{late:equality:eq:res-alpha}
 \res(H)\le\alpha(H).
\end{equation}
No residue result beyond Corollary~\ref{late:forest:cor:alpha} is needed here.

For every connected graph with an edge,
\begin{equation}\label{late:equality:eq:basic}
 f(G)\ge\alpha(G)+1\ge\res(G)+1.
\end{equation}
Indeed, a maximum independent set together with one vertex outside it
induces a star and isolated vertices. Our main result identifies equality
in the last bound for bipartite graphs.

\begin{theorem}\label{late:equality:thm:main}
Let $G$ be a connected bipartite graph with at least one edge. Then
\[
 f(G)=\res(G)+1
 \quad\Longleftrightarrow\quad
 G\text{ is a star or }C_4.
\]
Thus every other connected bipartite graph satisfies
$f(G)\ge\res(G)+2$.
\end{theorem}

Here a star includes $K_{1,1}$. The single-vertex graph is excluded because
its forest order and residue are both one. Our proof is a classification
of this first equality case, not a proof of the unrestricted candidate
$f\ge\res+\im$.

\paragraph{An attributed preliminary observation.}
Gunderson, Meagher, Morris and Pantangi already observed that when
\(f(G)=\alpha(G)+1\), relative to any maximum independent set \(I\),
adjacent vertices outside \(I\) share a neighbor in \(I\), and nonadjacent
ones share at least two~\cite[Section 5]{late:Gunderson}. In particular,
the triangle-free implication \(f=\alpha+1\Rightarrow G\) bipartite follows
from that prior observation. Its short proof is retained below.
The subsequent residue-equality argument is a separate step.

\subsubsection{Maximum independent sets}
\begin{lemma}\label{late:equality:lem:sides}
Let $G$ be connected and bipartite, with an edge and bipartition $(A,B)$.
If $f(G)=\alpha(G)+1$, every maximum independent set is either $A$ or $B$.
\end{lemma}
\begin{proof}
Let $I$ be a maximum independent set. Suppose it omits both a vertex
$a\in A$ and a vertex $b\in B$. Their neighborhoods within $I$ are
disjoint, since $N(a)\cap I\subseteq B$ and $N(b)\cap I\subseteq A$.
The induced graph on $I\cup\{a,b\}$ is therefore two stars, possibly
joined by the edge $ab$, together with isolated vertices. It is a forest
of order $\alpha(G)+2$, a contradiction. Hence $I$ contains an entire
bipartition class. Every vertex has a neighbor, so independence prevents
adding any vertex from the other class. Thus $I=A$ or $I=B$.
\end{proof}

We next prove a residue bound under this independent-set restriction,
without assuming anything about $f$.
\begin{theorem}\label{late:equality:thm:strict}
Let $G$ be connected and bipartite, with an edge. Suppose every maximum
independent set is a bipartition class. If $G$ is neither a star nor $C_4$,
then
\[
 \res(G)\le\alpha(G)-1.
\]
\end{theorem}

\subsubsection{The unequal-side case}
\begin{proof}[Proof of Theorem~\ref{late:equality:thm:strict} when the sides are unequal]
Write $|A|=p>|B|=q$. The hypothesis implies that $A$ is the unique maximum
independent set and $\alpha(G)=p$. For every nonempty $X\subseteq B$,
the independent set $X\cup(A\setminus N(X))$ differs from $A$, so has
order at most $p-1$. Thus
\begin{equation}\label{late:equality:eq:strict-Hall-unequal}
 |N(X)|\ge |X|+1\qquad(\varnothing\ne X\subseteq B).
\end{equation}
If $q=1$, connectivity makes $G$ a star, contrary to the assumption.
Hence $q\ge2$. Connectivity gives a vertex $w\in A$ with two distinct
neighbors $u,v\in B$; otherwise no path could join two vertices of $B$.
By~\eqref{late:equality:eq:strict-Hall-unequal} and Hall's theorem, $G-w$ has a matching
$M$ saturating $B$. Let $au,bv\in M$. The vertices $a,b,w$ are distinct.

Delete the edges $au,bv$ and insert $ab,uv$, obtaining $G'$. Both new
edges were absent, because they join vertices in the same bipartition
class. Thus this is a valid degree-preserving switch. The graph $G'$
contains the following vertex-disjoint subgraphs: the triangle on
$\{w,u,v\}$, the edge $ab$, and the $q-2$ other matching edges of $M$.
Every vertex cover of $G'$ has at least
\[
 2+1+(q-2)=q+1
\]
vertices. Since a set is a vertex cover exactly when its complement is
independent, $\alpha(G')\le p-1$. Equality of degree sequences and
\eqref{late:equality:eq:res-alpha} now give
\[
 \res(G)=\res(G')\le\alpha(G')\le p-1,
\]
as required.
\end{proof}

\subsubsection{The balanced case and conformal cycles}
An even cycle is called \emph{conformal} here if deleting its vertices
leaves a graph with a perfect matching. The cycle need not be induced. A balanced bipartite graph whose only
maximum independent sets are its two sides is classically called
\emph{bistable}. The equivalence of this condition with strict Hall
inequalities (and with connectedness and every edge belonging to a
perfect matching) is established in Levit--Mandrescu
\cite[Theorem~2.5]{late:equality:LM}. We include the short direction needed here and
then give the explicit conformal-cycle and degree-switch constructions.

\begin{lemma}\label{late:equality:lem:conformal}
Let $G$ be bipartite with $|A|=|B|=q\ge3$. Suppose
\begin{equation}\label{late:equality:eq:strict-Hall-balanced}
 |N(X)|\ge |X|+1\qquad
 (\varnothing\ne X\subsetneq A).
\end{equation}
Then $G$ contains a conformal even cycle of length at least six.
\end{lemma}
\begin{proof}
The displayed inequalities also ensure $N(A)=B$: applying one of them
to a set of size $q-1$ already gives all $q$ vertices of $B$.
Hall's theorem therefore supplies a perfect matching
$M=\{a_i b_i:1\le i\le q\}$. Define a directed graph $D$ on the matching
indices, with an arc $i\to j$ for $i\ne j$ whenever $a_i b_j$ is an edge
of $G$. Every nonempty proper index set has an outgoing arc by
\eqref{late:equality:eq:strict-Hall-balanced}. Hence $D$ is strongly connected.

If $D$ has a directed cycle on $k\ge3$ indices, its arcs together with
the corresponding $k$ matching edges form an even cycle of length $2k$.
The other matching edges show that this cycle is conformal.

Otherwise $D$ has no directed cycle of length at least three. For every
arc $i\to j$, a shortest directed return path from $j$ to $i$ has length
one: a longer return path would give such a forbidden directed cycle.
Thus every arc has its reverse. The underlying undirected graph is
connected and acyclic, because an undirected cycle could be oriented to
give a directed cycle of length at least three. It is a tree. Since
$q\ge3$, it contains a three-vertex path $i-j-k$. The six distinct
vertices of those matching pairs form the cycle
\[
 a_i,b_i,a_j,b_k,a_k,b_j,a_i.
\]
Again the other edges of $M$ certify conformality.
\end{proof}

\begin{proof}[Proof of Theorem~\ref{late:equality:thm:strict} when the sides are equal]
Write $|A|=|B|=q=\alpha(G)$. For a nonempty proper subset $X$ of $A$,
the independent set $X\cup(B\setminus N(X))$ equals neither $A$ nor $B$.
It therefore has order at most $q-1$, proving
\eqref{late:equality:eq:strict-Hall-balanced}. If $q=1$, $G=K_2$ is a star. If $q=2$,
the strict Hall inequalities force all four possible edges, so $G=C_4$.
These cases are excluded, hence $q\ge3$.

By Lemma~\ref{late:equality:lem:conformal}, choose a conformal cycle with consecutive
vertices $c_0,c_1,\ldots,c_{2k-1}$, where $k\ge3$. Delete $c_0c_1,c_3c_4$
and insert $c_0c_4,c_1c_3$. The inserted edges were absent because their
endpoints lie in the same bipartition class. Degrees are preserved.
The resulting graph $G'$ contains two vertex-disjoint odd cycles:
\[
 c_1,c_2,c_3,c_1,
 \qquad
 c_4,c_5,\ldots,c_{2k-1},c_0,c_4.
\]
Their orders are $3$ and $2k-3$, respectively, so their vertex covers
require at least $2+(k-1)=k+1$ vertices in total. A perfect matching
outside the original cycle supplies $q-k$ further disjoint edges.
Every vertex cover of $G'$ consequently has size at least $q+1$.
Therefore
\[
 \res(G)=\res(G')\le\alpha(G')\le q-1,
\]
which completes the proof of Theorem~\ref{late:equality:thm:strict}.
\end{proof}

\subsubsection{Equality and carefully limited consequences}
\begin{proof}[Proof of Theorem~\ref{late:equality:thm:main}]
If $f(G)=\res(G)+1$, then~\eqref{late:equality:eq:basic} forces
$f(G)=\alpha(G)+1$ and $\alpha(G)=\res(G)$. Lemma~\ref{late:equality:lem:sides} and
Theorem~\ref{late:equality:thm:strict} exclude every graph except a star or $C_4$.
Conversely, a star with $q$ leaves has forest order $q+1$ and residue $q$;
$C_4$ has forest order three and residue two.
\end{proof}

\begin{corollary}\label{late:equality:cor:triangle-free}
Let $G$ be a connected triangle-free graph with an edge. Then
$f(G)=\res(G)+1$ if and only if $G$ is a star or $C_4$.
Every other such graph satisfies $f(G)\ge\res(G)+2$.
\end{corollary}
\begin{proof}
Equality in~\eqref{late:equality:eq:basic} forces $f(G)=\alpha(G)+1$.
Let $I$ be a maximum independent set. If $V(G)\setminus I$ contains an
edge $uv$, triangle-freeness gives $N(u)\cap N(v)\cap I=\varnothing$.
Then $G[I\cup\{u,v\}]$ is two stars joined by $uv$, together with
isolates, and is a forest of order $\alpha(G)+2$, a contradiction.
Thus $V(G)\setminus I$ is independent, so $G$ is bipartite and
Theorem~\ref{late:equality:thm:main} applies. The converse is already proved.
\end{proof}

Original WOWII Conjecture~61 asks for
\begin{equation}\label{late:equality:eq:61}
 f(G)\ge\res(G)+\lceil\operatorname{diam}(G)/3\rceil
\end{equation}
in connected graphs. For diameter at most three this already follows
from~\eqref{late:equality:eq:basic}, with the one-vertex case checked directly.
Corollary~\ref{late:equality:cor:triangle-free} additionally proves~\eqref{late:equality:eq:61} for connected
triangle-free graphs of diameter four, five, or six: none is a star or $C_4$,
and the required ceiling is two. Thus the whole triangle-free subclass with
diameter at most six is covered, but diameter at least seven is not.

The same candidate~\eqref{late:forest:eq:candidate}
holds for connected triangle-free graphs with $\im(G)\le2$. The case
$\im=1$ follows from~\eqref{late:equality:eq:basic}; when $\im=2$, the exceptions in
Corollary~\ref{late:equality:cor:triangle-free} are impossible. The one-vertex case is immediate.
No assertion is made here for larger induced matching number or for
general graphs containing triangles. In particular, no full original WOWII conjecture is
counted as closed by these subclass results. These are mathematical subclass
consequences, not a claim of established publication novelty.


\subsection{Acyclic matchings and triangle-free graphs with one feedback vertex}\label{late:apex}

\subsubsection{Definitions, credited input, and exact scope}
All graphs are finite and simple. Write $n(G)=|V(G)|$, $\alpha(G)$ for the
independence number, $\nu(G)$ for the ordinary matching number, and $f(G)$
for the maximum order of an \emph{induced} forest. The feedback-vertex
number is $n(G)-f(G)$. Thus a graph has feedback-vertex number at most one
if it is a forest or becomes a forest after deleting one vertex.
No maximum-degree condition on that vertex is imposed.

A matching is \emph{acyclic} if its endpoints induce a forest. Its maximum
possible size is denoted by $\nu_{\mathrm{ac}}(G)$. This is the established acyclic
matching notion introduced in the subgraph-restricted matching framework
\cite{late:apex:GHHL}; see also \cite{late:apex:BR}. It is stronger than requiring that the
matching be the unique perfect matching on its endpoints. Always
\begin{equation}\label{late:apex:eq:acnu}
 \nu_{\mathrm{ac}}(G)\le\nu(G).
\end{equation}

We again use the classical residue bound
(Corollary~\ref{late:forest:cor:alpha}),
\begin{equation}\label{late:apex:eq:classical}
 \res(H)\le\alpha(H)
\end{equation}
for every graph $H$. Also,
\begin{equation}\label{late:apex:eq:matching}
 \alpha(H)\le n(H)-\nu(H),
\end{equation}
since every vertex cover meets each edge of a maximum matching and the
complement of an independent set is a vertex cover.

The use of degree-equivalent realizations to upper-bound residue through
independence, and the two-edge-switch method itself, are classical;
Fajtlowicz discusses both in the July 2004 \emph{Written on the Wall}
notes \cite[p.~37]{late:apex:Fajtlowicz}. They are not new ingredients claimed here.

\subsubsection{An unmatched nonleaf forces strictness}
\begin{lemma}\label{late:apex:lem:unmatched}
Let $G$ be triangle-free, with $n$ vertices. If some maximum matching
leaves a vertex of degree at least two unmatched, then
\[
 \res(G)\le n-\nu(G)-1.
\]
\end{lemma}
\begin{proof}
Let $M$ be such a maximum matching, let $x$ be the unmatched vertex, and
choose distinct neighbors $u,v$ of $x$. Two unmatched vertices cannot be
adjacent, since their edge could be added to $M$. Thus $u$ and $v$ have
partners $a$ and $b$ in $M$, respectively. Triangle-freeness implies that
$uv$ is absent, so $ua$ and $vb$ are distinct matching edges. In particular,
$x,u,v,a,b$ are five distinct vertices.

First suppose $ab\in E(G)$. Then
\[
 x,u,a,b,v,x
\]
is a five-cycle. It is vertex-disjoint from the $\nu(G)-2$ remaining
edges of $M$. Every vertex cover of $G$ therefore has at least
\[
 3+(\nu(G)-2)=\nu(G)+1
\]
vertices. Hence $\alpha(G)\le n-\nu(G)-1$, and
\eqref{late:apex:eq:classical} proves the assertion.

Now suppose $ab\notin E(G)$. Delete $ua,vb$ and insert $uv,ab$, producing
$G'$. The added edges were absent and their four endpoints are distinct,
so this is a valid degree-preserving switch. The graph $G'$ contains the
mutually vertex-disjoint triangle on $\{x,u,v\}$, edge $ab$, and the
$\nu(G)-2$ other edges of $M$. Every vertex cover of $G'$ has at least
\[
 2+1+(\nu(G)-2)=\nu(G)+1
\]
vertices. Consequently
\[
 \res(G)=\res(G')\le\alpha(G')\le n-\nu(G)-1.
\]
The two cases exhaust all possibilities.
\end{proof}

\begin{corollary}\label{late:apex:cor:tight-unmatched}
If a triangle-free graph satisfies $\res(G)=n(G)-\nu(G)$, every vertex
left unmatched by any maximum matching has degree at most one.
\end{corollary}

The unmatched vertices of an acyclic ordinary maximum matching form an
independent feedback vertex set, as explicitly observed by
F\"urst--Rautenbach \cite[Conclusion]{late:apex:FR}. We recall that standard
observation in the proof below; the strict residue conclusion additionally
uses Lemma~\ref{late:apex:lem:unmatched}.

\begin{corollary}\label{late:apex:cor:cyclic}
If $G$ is triangle-free and contains a cycle, and
$\nu_{\mathrm{ac}}(G)=\nu(G)$, then
\[
 \res(G)\le n(G)-\nu(G)-1.
\]
\end{corollary}
\begin{proof}
Choose an acyclic matching $M$ of size $\nu(G)$. Its endpoint-induced
graph is a forest. The unmatched set is independent, because $M$ is a
maximum matching. A cycle of $G$ cannot be wholly contained in the
endpoint-induced forest, so it contains an unmatched vertex $x$.
This vertex has at least two neighbors, and Lemma~\ref{late:apex:lem:unmatched}
applies.
\end{proof}

The hypothesis in Corollary~\ref{late:apex:cor:cyclic} is that an \emph{ordinary
maximum} matching can be chosen acyclic. It does not assume that all
maximum matchings are acyclic.

\subsubsection{Triangle-free apex forests}
\begin{theorem}\label{late:apex:thm:apex}
If $G$ is triangle-free and has feedback-vertex number at most one, then
\begin{equation}\label{late:apex:eq:main}
 f(G)\ge\res(G)+\nu_{\mathrm{ac}}(G).
\end{equation}
\end{theorem}
\begin{proof}
If $G$ is a forest, every matching is acyclic and $f(G)=n(G)$.
Combining~\eqref{late:apex:eq:classical} and~\eqref{late:apex:eq:matching} gives
\[
 f(G)=n(G)\ge\res(G)+\nu(G)=\res(G)+\nu_{\mathrm{ac}}(G).
\]

Otherwise $G$ is cyclic and its feedback-vertex number is exactly one,
so $f(G)=n(G)-1$. If $\nu_{\mathrm{ac}}(G)\le\nu(G)-1$, then
\[
 \res(G)+\nu_{\mathrm{ac}}(G)
 \le(n(G)-\nu(G))+(\nu(G)-1)
 =n(G)-1=f(G).
\]
If $\nu_{\mathrm{ac}}(G)=\nu(G)$, Corollary~\ref{late:apex:cor:cyclic} instead yields
\[
 \res(G)+\nu_{\mathrm{ac}}(G)
 \le(n(G)-\nu(G)-1)+\nu(G)
 =n(G)-1=f(G).
\]
These cases are exhaustive by~\eqref{late:apex:eq:acnu}.
\end{proof}

The theorem allows arbitrarily many cycles through the feedback vertex;
it is not limited to unicyclic graphs. Its forest case is simply the
classical independence bound and is not counted as a separate new result.

\begin{corollary}\label{late:apex:cor:diameter}
Let $G$ be a connected triangle-free graph with feedback-vertex number
at most one, and let $D=\operatorname{diam}(G)$. Then
\begin{equation}\label{late:apex:eq:half}
 f(G)\ge\res(G)+\lceil D/2\rceil.
\end{equation}
In particular, original WOWII Conjecture~61 holds for every graph in
this subclass, without a diameter restriction.
\end{corollary}
\begin{proof}
A diametral path is induced and has $D+1$ vertices. A maximum matching
of that path has size
\[
 \left\lfloor\frac{D+1}{2}\right\rfloor
 =\left\lceil\frac D2\right\rceil.
\]
The path matching size is standard; see also \cite[Theorem~3]{late:apex:GHHL}.
Its endpoints induce a subgraph of the induced path, hence a forest.
Thus $\nu_{\mathrm{ac}}(G)\ge\lceil D/2\rceil$, and Theorem~\ref{late:apex:thm:apex} gives
\eqref{late:apex:eq:half}. Original~61 asks for
$f(G)\ge\res(G)+\lceil D/3\rceil$, which is weaker. The one-vertex
case has $D=0$ and $f=\res=1$.
\end{proof}

\subsubsection{Limitations of the stronger inequality}
No proof is claimed here for arbitrary triangle-free graphs of
feedback-vertex number at least two. Neither the general half-diameter
bound nor the general acyclic-matching bound is established by this
paper. In particular, this is not a full closure of original~61.

The triangle-free hypothesis cannot simply be dropped. A triangle with
a pendant leaf at each of two distinct triangle vertices has
$f=4$, $\res=3$, $\nu_{\mathrm{ac}}=2$, and diameter three. It violates both
\eqref{late:apex:eq:main} and~\eqref{late:apex:eq:half}. Ordinary matching number also cannot
replace acyclic matching number, even within the feedback-vertex-one
subclass: $C_4$ has $f=3$, $\res=2$, $\nu=2$, and $\nu_{\mathrm{ac}}=1$.


\subsection{When a diametral path attains the matching number}\label{late:geodesic}

\subsubsection{Definitions, credited input, and precise scope}
All graphs are finite and simple. Write $n(G)=|V(G)|$, $\alpha(G)$ for the
independence number, $\nu(G)$ for the \emph{ordinary} matching number,
and $f(G)$ for the largest number of vertices inducing a forest. A
feedback-vertex set is a set whose deletion leaves a forest; its minimum
size is
\[
 \operatorname{fvs}(G)=n(G)-f(G).
\]
A vertex cover meets every edge, so its minimum size is $n(G)-\alpha(G)$.
The feedback-vertex and vertex-cover parameters are different.

The residue depends only on the degree sequence. Our residue input is
again Corollary~\ref{late:forest:cor:alpha}:
\begin{equation}\label{late:geodesic:eq:resalpha}
 \res(H)\le\alpha(H)\qquad\text{for every finite simple graph }H.
\end{equation}
The degree-equivalent-realization viewpoint and two-edge-switch method
are classical, and are explicitly discussed in Fajtlowicz's notes
\cite[p.~37]{late:apex:Fajtlowicz}. They are not claimed as new methods here.

A diametral path of length $D$ contains a matching of size
$\lfloor(D+1)/2\rfloor=\lceil D/2\rceil$, giving the elementary lower
bound $\nu(G)\ge\lceil D/2\rceil$ for a connected graph. We treat precisely
the equality case.
\begin{theorem}\label{late:geodesic:thm:main}
Let $G$ be a finite simple connected triangle-free graph, and put
$D=\diam(G)$ and $k=\lceil D/2\rceil$. If $\nu(G)=k$, then
\begin{equation}\label{late:geodesic:eq:main}
 f(G)\ge\res(G)+k.
\end{equation}
\end{theorem}
Original WOWII Conjecture~61 asks for
$f(G)\ge\res(G)+\lceil\diam(G)/3\rceil$ for every simple connected graph.
Consequently Theorem~\ref{late:geodesic:thm:main} establishes that inequality only on its
stated subclass, with a stronger diameter term. No conclusion for
triangle-free graphs with $\nu(G)>\lceil D/2\rceil$ is obtained below.
The proof is combinatorial, with \eqref{late:geodesic:eq:resalpha} as its credited input;
finite computations are supporting controls rather than its basis.

\subsubsection{The structure forced by a diametral path}
For the remainder of the proof assume $D\ge1$ and fix a diametral path
\[
 P=p_0p_1\cdots p_D.
\]
Every subpath is a shortest path between its ends; in particular, $P$ is
induced. Set $R=V(G)\setminus V(P)$.

\begin{lemma}\label{late:geodesic:lem:outside}
The set $R$ is independent. Every $x\in R$ has either one neighbor on
$P$, or exactly two neighbors $p_i,p_{i+2}$. It has no other neighbors.
\end{lemma}
\begin{proof}
A matching of $k$ edges in $P$ is maximum in $G$. Any edge inside $R$
could be added to it, which is impossible. Thus $R$ is independent.
Connectedness ensures that each $x\in R$ has a neighbor on $P$.
If $p_i,p_j\in N(x)$ with $i<j$, the two-edge path $p_ixp_j$ shows
$j-i\le2$, since otherwise it would shorten the corresponding geodesic
subpath. Consecutive neighbor indices would form a triangle with $x$.
Thus the neighborhood is a singleton or a pair at index difference two.
\end{proof}

We now choose $k$ distinct path vertices called \emph{supports} and
pair each with a distinct \emph{partner} in a fixed $k$-edge path matching
$M$. The supports and partners are disjoint sets. In both parity cases
all partners will form an independent set.

\paragraph{Even diameter: one support chain}
Suppose $D=2k$. For each $j\in\{0,\ldots,k\}$ there is a $k$-edge
matching in $P$ exposing precisely $p_{2j}$: match the two even-order
subpaths on its sides. If $p_{2j}$ had a neighbor in $R$, that edge would
extend the matching. Therefore outside vertices have only odd-indexed
neighbors on $P$.

For $1\le i\le k$ take
\[
 u_i=p_{2i-1},\qquad a_i=p_{2i-2},\qquad
 M=\{u_ia_i:1\le i\le k\}.
\]
The supports are $u_1,\ldots,u_k$, their partners $a_1,\ldots,a_k$ are
independent, and $p_{2k}$ is unmatched. Consecutive supports $u_i,u_{i+1}$
are joined by one or more internally disjoint two-edge paths: the path
through $p_{2i}$, and any paths through outside common neighbors.
All other vertices are pendant at a support. This includes the two path
endpoints. There are no other edges, by Lemma~\ref{late:geodesic:lem:outside} and the
inducedness of $P$.

\paragraph{Odd diameter: two support chains and an optional bridge}
Suppose $D=2k-1$, and fix
\[
 M=\{p_{2j}p_{2j+1}:0\le j<k\}.
\]
Let $E$ and $O$ be respectively the even and odd path indices having a
neighbor in $R$. There are no $e\in E$, $o\in O$ with $e<o$. Indeed,
choose outside neighbors $x$ of $p_e$ and $y$ of $p_o$. They are distinct
by Lemma~\ref{late:geodesic:lem:outside}, since an outside vertex has neighbors of only
one parity. The path
\[
 x,p_e,p_{e+1},\ldots,p_o,y
\]
is an $M$-augmenting path: its end edges are outside $M$, and the path
segment starts and ends with edges of $M$. Exchanging alternating edges
would increase the matching size, a contradiction.

It follows that all odd outside attachment indices precede all even
ones. Choose $h\in\{0,\ldots,k\}$ so that every odd attachment index is
less than $2h$, and every even one is at least $2h$. If $O\ne\varnothing$,
one may take $h=(\max O+1)/2$; if $O=\varnothing$, take $h=0$.
For each $0\le j<k$ designate the support and its partner by
\[
 (u_j,a_j)=
 \begin{cases}
 (p_{2j+1},p_{2j}),&j<h,\\
 (p_{2j},p_{2j+1}),&j\ge h.
 \end{cases}
\]
Every outside vertex is adjacent only to supports. The odd supports on
the left form one chain, and the even supports on the right form another.
Within either chain, consecutive supports are joined by one or more
internally disjoint two-edge paths. All remaining outside vertices are
pendant at a support.

If both chains are nonempty, their sole connection is the path edge
$p_{2h-1}p_{2h}$. This edge is a bridge: there are no path chords, no
edges within $R$, and no outside vertex adjacent to both parities of
supports. If $h=0$ or $h=k$, one chain is empty and this bridge is absent.
All partners are independent. Partners within one chain have the same
parity; across the boundary the nearest two have indices $2h-2$ and
$2h+1$, whose difference is three. The inducedness of $P$ rules out all
such partner edges.

\begin{remark}\label{late:geodesic:rem:structure}
This is a necessary description relative to the fixed diametral path,
not a converse assertion that arbitrary pendant additions preserve its
diameter. Empty attachment sets, a chain with one support, and an empty
chain are allowed. If $R$ is empty, the description reduces to $P$ itself.
\end{remark}

\subsubsection{The exact feedback-vertex number}
Call the collection of two-edge paths between two consecutive supports
in the same chain a \emph{link}. Its internal vertices are precisely the
common neighbors used by these paths, and it induces $K_{2,c}$ for some
$c\ge1$. Exactly one of its internal vertices belongs to $P$; the others
belong to $R$. Each internal vertex belongs to exactly one link.
A link is \emph{bad} if $c\ge2$.

Let $Q$ be the graph on all $k$ supports whose edges are the bad links.
The optional bridge between the two odd-case chains is not an edge of
$Q$. Thus $Q$ is a disjoint union of paths, including isolated vertices.
Put $t=\nu(Q)$.

\begin{lemma}\label{late:geodesic:lem:fvs}
With the notation above, $\operatorname{fvs}(G)=t$ and $f(G)=n(G)-t$.
\end{lemma}
\begin{proof}
Every cycle of $G$ lies in one link. To see this, the support-chain
skeleton, including the optional bridge, is a path; a support separating
two links is a cut vertex separating their internal vertices. A simple
cycle cannot cross such a support and return through it, nor can it use
a bridge. Pendant vertices do not lie on cycles. A link $K_{2,c}$ is
cyclic exactly when $c\ge2$.

A vertex cover of $Q$ is consequently a feedback-vertex set of $G$.
For a path with $\ell$ edges, a matching and a vertex cover of size
$\lceil\ell/2\rceil$ are obtained by alternating edges and alternating
vertices, respectively. No vertex cover can be smaller than a matching.
Summing over path components, $Q$ therefore has a vertex cover of size
$t$, giving $\operatorname{fvs}(G)\le t$.

Conversely, choose a matching of $t$ bad links. In each link choose two
of its internal vertices. They and its two supports form a four-cycle.
The resulting $t$ four-cycles are vertex-disjoint: their support pairs
are disjoint, and internal vertices belong to distinct links. Every
feedback-vertex set meets all these cycles, so $\operatorname{fvs}(G)\ge t$.
\end{proof}

\subsubsection{Simultaneous degree-preserving switches}
\begin{proof}[Proof of Theorem~\ref{late:geodesic:thm:main}]
When $D=0$, connectedness gives $G=K_1$, so $f(G)=\res(G)=1$ and $k=0$.
Assume $D\ge1$ and use the structure and matching $M$ above. Write
$n=n(G)$, and choose a matching $L$ of $t$ bad links in $Q$.

For each $uv\in L$, choose an extra common neighbor $x\notin V(P)$.
Such a vertex exists because the link has at least two internal vertices
and exactly one of them is on $P$. Let $a,b$ be the partners of $u,v$ in
$M$. The five vertices $x,u,v,a,b$ are distinct. Replace the two edges
$ua,vb$ by the two edges $uv,ab$.

Both new edges were absent: $u,v$ have path indices differing by two,
and all partners are independent. The selected support pairs are
disjoint, their partners are distinct, and the outside witnesses are
distinct because each belongs to only one link. Supports, partners and
outside witnesses are disjoint classes. Hence these switches have
pairwise disjoint five-vertex domains and may be performed simultaneously,
producing a finite simple graph $G'$ with exactly the same degree at
every vertex as $G$.

In $G'$ the following subgraphs are pairwise vertex-disjoint:
\begin{itemize}
\item $t$ triangles $xuv$, one for each selected link;
\item the $t$ new partner edges $ab$;
\item the $k-2t$ unchanged edges of $M$ on the remaining supports.
\end{itemize}
Their edges survive all the switches. Additional edges between these
subgraphs are harmless: they are not asserted to be induced components.
Since $L$ is a matching on $k$ supports, $2t\le k$. Every vertex cover
of $G'$ must therefore contain at least
\[
 2t+t+(k-2t)=k+t
\]
vertices. By \eqref{late:geodesic:eq:resalpha} and Lemma~\ref{late:geodesic:lem:fvs},
\[
 \res(G)=\res(G')\le\alpha(G')\le n-k-t=f(G)-k.
\]
This proves \eqref{late:geodesic:eq:main}. When $t=0$, no switch is made and the unchanged
matching alone is the required cover certificate. The switched graph
need not remain triangle-free, connected, or have the same diameter;
none of these extra conditions is needed for \eqref{late:geodesic:eq:resalpha}.
\end{proof}

\subsubsection{Boundary cases and an unbounded family}
The small diameter cases are included: $D=1$ gives $K_2$, and $D=2$ under
the matching hypothesis gives a star. Neither case has a bad link.
Both odd split endpoints $h=0,k$, missing attachment sets, and $2t=k$
are included in the argument. Link multiplicities may exceed two; a
selected bad link always has an outside witness, and only two of its
internal vertices are needed for the disjoint-cycle certificate.

\begin{proposition}\label{late:geodesic:prop:family}
For each $k\ge1$ there is a graph satisfying the theorem with
$n=3k$, $D=2k$, $\nu=k$, and $\operatorname{fvs}=\lfloor k/2\rfloor$.
\end{proposition}
\begin{proof}
Take $k$ supports in path order, replace each consecutive-support link
by two internally disjoint two-edge paths, and attach one pendant vertex
to each end support. For $k=1$, attach two distinct pendant vertices to
the single support. The graph is triangle-free. It has
$k+2(k-1)+2=3k$ vertices. The two pendant vertices have distance $2k$,
and every other vertex lies between their levels in the chain, so the
diameter is $2k$. The supports form a vertex cover of size $k$, while
the diametral path has a matching of size $k$, proving $\nu=k$.
Every link is bad, so $Q$ is the path on $k$ vertices and
Lemma~\ref{late:geodesic:lem:fvs} gives $\operatorname{fvs}=\lfloor k/2\rfloor$.
\end{proof}
In particular, the theorem is not restricted to feedback-vertex number
at most one. The proposition states structural parameters, not equality
in \eqref{late:geodesic:eq:main} for every member.

\paragraph{Scope and priority.}
The unrestricted half-diameter bound for triangle-free graphs remains
unproved here, as does original~61 for arbitrary connected graphs.
The classical residue bound, path matching observation, vertex-cover
certificate, and degree-switch methodology are credited or elementary
inputs. The source checkpoint is a bounded literature search, not a
certification of worldwide novelty or publication priority.


\section{Conjecture 155 when spanning trees have at most three leaves}\label{late:leaves}

For a nontrivial connected graph $G$, let $L(G)$ be the maximum number of
leaves of a spanning tree, let $r=\rad(G)$, and let $C(G)$ be the full
center. Define
\[
 R(v)=|\{x:d(v,x)=r\}|\quad(v\in C(G)),\qquad
 q(G)=|\{R(v):v\in C(G)\}|.
\]
Thus $q$ counts the \emph{distinct orders} of radial circles. Original
statement 155 is
\begin{equation}\label{late:leaves:target}
 L(G)\ge1+q(G).
\end{equation}
It contains no ratio of radial-circle orders; that ratio belongs to 154.

\begin{theorem}\label{late:leaves:main}
For every nontrivial connected graph with $L(G)\le3$, inequality
\eqref{late:leaves:target} holds. More precisely, $L=2$ implies $q=1$,
and $L=3$ implies $q\le2$.
\end{theorem}

\subsection{A full-star criterion and the resulting families}
\begin{lemma}\label{late:leaves:stars}
For a connected graph of maximum degree at most three, the following
are equivalent:
\begin{enumerate}
\item $L(G)\le3$;
\item for every two distinct degree-three vertices, the union of their
full incident-edge stars contains a cycle;
\item every adjacent pair of degree-three vertices has a common neighbor,
and every nonadjacent pair has at least two common neighbors.
\end{enumerate}
\end{lemma}
\begin{proof}
Every forest in a connected graph extends to a spanning tree. If two full
stars form a forest, an extension retains degree three at both centers.
A subcubic tree has two more leaves than degree-three vertices, so the
extension has at least four leaves. Conversely, a spanning tree with at
least four leaves has two degree-three vertices, and its edge set contains
their full stars, which form a forest. Two adjacent stars contain a cycle
exactly when their centers have a common neighbor; two nonadjacent stars
contain a cycle exactly when their centers have at least two common
neighbors. These observations prove all equivalences.
\end{proof}

If $L\le3$, then $\Delta(G)\le3$: the star at any vertex extends to a
spanning tree, and a tree has at least as many leaves as the degree of
each of its vertices. Let $B$ be the set of degree-three vertices.

If $B$ is independent, fix $x\in B$. Its three neighbors have degree at
most two, so together they have at most three edges to $B\setminus\{x\}$.
Every further vertex of $B$ requires at least two such edges by
Lemma~\ref{late:leaves:stars}. Consequently $|B|\le2$.

Suppose $|B|\ge3$. There is an edge $ab$ with both endpoints in $B$.
They have a common neighbor $c$. If $c\in B$, there is a triangle in $B$.
Otherwise $c$ has degree two and is adjacent only to $a,b$. A third
vertex $z\in B$ cannot miss both $a$ and $b$: to share two neighbors with
$a$, it would have to be adjacent to $b$ or $c$, which is then impossible.
Suppose, by symmetry, that $za$ is an edge. The adjacent pair $a,z$ needs
a common neighbor; of the other neighbors $b,c$ of $a$, only $b$ can meet
$z$, because $c$ is saturated. Hence $a,b,z$ form a triangle in $B$.

Fix a triangle $abc$ in $B$. Every further $z\in B$ must meet the
triangle, since otherwise it could share at most the one remaining
neighbor of $a$ with $a$. If $za$ is an edge, the adjacent-star condition
forces $z$ to meet $b$ or $c$ as well. Each additional vertex therefore
uses at least two of the triangle's three free incidences. At most one
is possible, proving $|B|\le4$.

Every component outside $B$ is a path, unless $B$ is empty and the whole
graph is a cycle. Its ends attach to free incidences of vertices in $B$
or terminate as leaves. A \emph{tail} means a positive-length pendant
path. A \emph{theta graph} has two endpoints joined by three internally
disjoint paths of the specified lengths. The exhaustive list is:
\begin{enumerate}
\item $|B|=0$: a path or a cycle;
\item $|B|=1$: a three-leaf tree, or a cycle with one tail;
\item $|B|=2$ with the vertices adjacent: a triangle with tails at two
vertices, or a theta graph of lengths $1,2,\ell$, with $\ell\ge2$;
\item $|B|=2$ with the vertices nonadjacent: a four-cycle with tails at
opposite vertices, or a theta graph of lengths $2,2,\ell$, with
$\ell\ge2$;
\item $|B|=3$: a triangle with three tails, or a theta graph of lengths
$1,2,\ell$, with a tail at the internal vertex of the length-two path;
\item $|B|=4$: $K_4$, a graph obtained by subdividing exactly one edge
of $K_4$, or a diamond $K_4-e$ with tails at its two degree-two vertices.
\end{enumerate}
For the last case, the fourth vertex meets two or three vertices of the
triangle. Three gives $K_4$. Two gives a diamond; its remaining two
incidences either join by a path or terminate in separate tails.
For $|B|=3$, the triangle's three free incidences are either all tails,
or two join by a path and the third is a tail. For $|B|=2$, the one or
two required common neighbors give exactly the listed alternatives;
three common neighbors is the theta graph with lengths $2,2,2$.
This proves the classification directly from the full-star criterion.

\subsection{Paths, cycles, trees, and a cycle with one tail}
If $L=2$, then $\Delta\le2$, so $G$ is a path or cycle. A cycle is
vertex-transitive, and all its radial circles have the same order.
A path has one or two centers; the radial-circle order is two at a
unique center, or one at each of two centers. Thus $q=1$.
A tree has at most two centers, the middle vertex or vertices of a
diametral path, so a three-leaf tree has $q\le2$.

Now attach a tail of length $t\ge1$ at a vertex $c$ of a cycle of length
$m$, and put $s=\lfloor m/2\rfloor$. At a cycle vertex $x$,
\[
 \operatorname{ecc}(x)=\max\{s,d_{C_m}(x,c)+t\}.
\]
At depth $j$ on the tail, eccentricity is $\max\{t-j,s+j\}$.
If $t>s$, the center lies on the tail, and minimizing the maximum of
these two opposite affine functions gives at most two centers.
If $t\le s$, the radius is $s$, no positive-depth tail vertex is
central, and the central cycle vertices satisfy $d(x,c)\le s-t$.
Each cycle vertex has $\epsilon$ farthest cycle vertices, where
$\epsilon=1$ for even $m$ and $\epsilon=2$ for odd $m$. At a central
cycle vertex the tail endpoint is the only possible additional farthest
vertex, occurring exactly when $d(x,c)+t=s$. Thus all radial orders are
in $\{\epsilon,\epsilon+1\}$, and $q\le2$.

\subsection{A triangle with tails}
Let $A\ge B\ge C\ge0$ be tail lengths at the triangle vertices $a,b,c$;
zero means no tail, and at least two lengths are positive. The
eccentricity of $a$ is $\max\{A,B+1\}$, and those of $b,c$ are $A+1$.
If $A=B$, all three triangle vertices are central with radius $A+1$,
and no positive-depth tail vertex is central. A radial circle consists
of the endpoints of maximum-length tails at the other triangle vertices,
so has order one or two.
If $A>B$, the eccentricity of $a$ is at most $A$, whereas every vertex
outside its tail has eccentricity at least $A+1$. The center lies on
that tail, including $a$, and at depth $j$ eccentricity is
$\max\{A-j,B+1+j\}$. There are at most two minimizing vertices.
Therefore $q\le2$ in either case.

\subsection{A four-cycle or diamond with opposite tails}
Let the nonadjacent attachment vertices $c,d$ have tail lengths
$A\ge B\ge1$, and let $a,b$ be the other two vertices of the four-cycle.
The edge $ab$ is optional, producing the diamond when present. In both
cases,
\[
 \operatorname{ecc}(c)=\max\{A,B+2\},\quad
 \operatorname{ecc}(d)=A+2,\quad
 \operatorname{ecc}(a)=\operatorname{ecc}(b)=A+1.
\]
If $A=B$, the center is exactly $\{a,b\}$. If $A=B+1$, it is exactly
$\{a,b,c\}$, with radius $A+1\ge3$. The vertices $a,b$ have the endpoint
of the $A$-tail as their unique farthest vertex; $c$ has the endpoint of
the $B$-tail as its unique farthest vertex. All radial orders are one.
If $A\ge B+2$, the center lies on the $c$-tail, including $c$, where the
eccentricity at depth $j$ is $\max\{A-j,B+2+j\}$. There are at most two
centers. The exclusion of other centers follows because a positive-depth
tail vertex reaches the opposite tail endpoint through a core path of
length two. Thus $q\le2$ throughout.

\subsection{A theta graph of lengths 1, 2, and an arbitrary third length}
Let $a,b$ be the theta endpoints, let $a-c-b$ be the length-two path,
and let the third path have length $\ell\ge2$. Allow a tail of length
$t\ge0$ at $c$. The edge $ab$ and the long path form a cycle $Q$ of
length $m=\ell+1$, and distances within $Q$ are unchanged by $c$ and
the tail. Put $s=\lfloor m/2\rfloor$ and
\[
 h(x)=d(x,c)=1+\min\{d_Q(x,a),d_Q(x,b)\}\quad(x\in V(Q)).
\]
Then
\begin{equation}\label{late:leaves:theta-cycle}
 \operatorname{ecc}(x)=\max\{s,h(x)+t\}.
\end{equation}
At depth $j$ on the tail, including $j=0$ at $c$, eccentricity is
\begin{equation}\label{late:leaves:theta-tail}
 \max\{t-j,H+j\},\qquad H=\max_{x\in V(Q)}h(x).
\end{equation}

If $m=2s$, each value $1,\ldots,s$ of $h$ occurs twice and $H=s$.
For $t<s$, the radius is $s$. The centers are $c$ and the cycle vertices
with $h(x)\le s-t$. Each such cycle vertex has one cycle antipode and
gains the tail endpoint exactly when $h(x)+t=s$; the vertex $c$ has its
two farthest cycle vertices. All radial orders are one or two.
For $t=s$, the sole center is $c$, with radial order three.
For $t>s$, every cycle vertex has eccentricity at least $t+1$, whereas
$c$ has eccentricity $t$. The center is on the tail, and
\eqref{late:leaves:theta-tail} gives at most two centers.

If $m=2s+1$, each value $1,\ldots,s$ occurs twice and $s+1$ occurs once,
so $H=s+1$. For $t<s$, the radius is $s$, and the centers are exactly
the cycle vertices with $h(x)\le s-t$. Each has two cycle antipodes and
gains the tail endpoint precisely when $h(x)+t=s$. All radial orders
are two or three. For $t=s$, the center is $\{a,b,c\}$, with radius
$s+1$: $a,b$ have only the tail endpoint at that distance, and $c$ has
only the unique cycle vertex with $h=s+1$. Every radial order is one.
For $t\ge s+1$, every cycle vertex has eccentricity at least $t+1$ and
$c$ has eccentricity $t$; again the center consists of at most two tail
vertices. These cases include $t=0$, and always give $q\le2$.

\subsection{The remaining theta graphs and subdivisions of one edge of a clique}
For a theta graph of lengths $2,2,\ell$, $\ell\ge2$, form a cycle $Q$
from the long path and one length-two path. Its length is $m=\ell+2\ge4$;
let $c$ be the internal vertex of its length-two path. The other
length-two path adds a vertex $d$ with the same neighbors as $c$.
Optionally add $cd$. With that edge, the graph is obtained by subdividing
one edge of $K_4$; without it, it is the given theta graph.

Distances within $Q$ are unchanged. For $x\in V(Q)\setminus\{c\}$,
$d(x,d)=d_Q(x,c)$ whether or not $cd$ is present. The distance between
$c,d$ is two without that edge and one with it. Put
$s=\lfloor m/2\rfloor\ge2$. Every vertex has eccentricity $s$.
Each vertex of $Q$ has $\epsilon$ cycle antipodes, with $\epsilon=1$
for even $m$ and $\epsilon=2$ for odd $m$. A vertex other than $c$ gains
$d$ as a farthest vertex only when its distance to $c$ is $s$.
The vertices $c,d$ likewise have either $\epsilon$ or $\epsilon+1$
farthest vertices, according as their mutual distance equals $s$.
Hence all radial orders lie in $\{\epsilon,\epsilon+1\}$.
The unsubdivided $K_4$ has every radial-circle order equal to three.

\begin{proof}[Conclusion of the proof of Theorem~\ref{late:leaves:main}]
The classification above is exhaustive for $L\le3$. The case $L=2$
has $q=1$, and all families with $L=3$ have $q\le2$. Therefore
$L\ge1+q$ throughout the claimed class.
\end{proof}
This proves an arbitrary-order subclass theorem. It does not establish
155 for $L\ge4$ or the separate, unproved radial-order lemma discussed
in the source notes.


\section{Degree frequencies and quartile sufficient conditions for 389a}\label{late:389a}

\subsection{The source convention and the scope of the reduction}
The original statement 389a asks whether every connected graph of order
$n>2$ satisfies
\begin{equation}\label{late:389a:target}
 \gamma_2(G)\le \mu(q+a)-1,
\end{equation}
where $a$ is the number of vertices of degree \emph{at most} two, $\mu$
is the size of the mode-min degree class, and $q$ is the first-quartile
degree. All modal classes have the same size, so $\mu$ is simply the
largest frequency of any degree. We use $\mu$ here to distinguish this
frequency from the induced-forest number $f(G)$ used in the preceding
sections. The condition on $a$ is different from the exactly-degree-two
set in the formulation of 354.

The archived definition orders the degrees as $d_1\le\cdots\le d_n$ and
specifies the entry at rank $n/4$, without a rounding convention when
$4\nmid n$~\cite{src:catalogue}. Accordingly, throughout this section,
for $n\ge4$ we assume that $q$ is a positive integer satisfying
\begin{equation}\label{late:389a:quartile}
 \bigl|\{v:d(v)\le q\}\bigr|\ge\lfloor n/4\rfloor.
\end{equation}
This includes both the lower-rank and upper-rank choices, and the literal
source definition when $4\mid n$. No interpolation convention is assumed.
For $n=3$, the connected graphs are $P_3$ and $K_3$; each has
$\gamma_2=2$ and $a=3$, with $\mu=2$ and $3$, respectively, so
\eqref{late:389a:target} holds for every positive integer $q$.

Connectedness puts all degrees in $\{1,\ldots,n-1\}$, and hence
\begin{equation}\label{late:389a:elementary}
 \mu\ge2,\qquad a\le2\mu.
\end{equation}
If $q\ge3$, counting the degree classes in \eqref{late:389a:quartile}
gives
\begin{equation}\label{late:389a:rank}
 \lfloor n/4\rfloor\le a+\mu(q-2),\qquad
 n\le4a+4\mu(q-2)+3.
\end{equation}

\subsection{Two upper bounds for two-domination}
\begin{lemma}[A cut bound]\label{late:389a:cut}
For any graph with $n$ vertices, of which $a$ have degree at most two,
\[
 \gamma_2(G)\le\left\lfloor\frac{n+a}{2}\right\rfloor.
\]
\end{lemma}
\begin{proof}
Take a cut locally maximal under moving a single vertex to the other side.
Every vertex of degree at least three has at least two neighbors across
the cut, since otherwise moving it would increase the cut. One side
contains at least $\lceil(n-a)/2\rceil$ vertices of degree at least three.
Omit exactly those vertices and select all remaining vertices. Each
omitted vertex has at least two selected neighbors across the cut, so the
selected set is 2-dominating and has the required size.
\end{proof}

\begin{lemma}[A frequency-sensitive alteration bound]\label{late:389a:random}
Let $G$ have no isolated vertices, let $a$ count its vertices of degree at
most two, and let $\mu$ be its largest degree frequency. For $0<p<1$,
\begin{equation}\label{late:389a:probability}
 \gamma_2(G)\le pn+(1-p)a+
 \mu(1-p)^3\left(1+\frac2p\right).
\end{equation}
In particular,
\begin{equation}\label{late:389a:sixth}
 \gamma_2(G)\le\frac n6+\frac{5a}{6}+\frac{1625\mu}{216}.
\end{equation}
\end{lemma}
\begin{proof}
Select each vertex independently with probability $p$, producing a set
$S$. Add every unselected vertex having fewer than two neighbors in $S$.
The resulting set is 2-dominating. Put $z=1-p$. A vertex of degree $d$
is added in the second step with probability
\[
 z^{d+1}+dpz^d.
\]
For $d=1$ this equals $z$, and for $d=2$ it equals
$z^2(1+p)\le z$. At most $\mu$ vertices have any given degree $d\ge3$.
The expected size of the final set is therefore at most
\[
 pn+za+\mu\sum_{d=3}^{\infty}(z^{d+1}+dpz^d).
\]
Using the geometric series and its derivative gives
\[
 \sum_{d=3}^{\infty}(z^{d+1}+dpz^d)
 =3z^3+\frac{2z^4}{p}
 =z^3\left(1+\frac2p\right).
\]
Some outcome has size at most its expectation, proving
\eqref{late:389a:probability}. Substituting $p=1/6$ proves
\eqref{late:389a:sixth}.
\end{proof}
These are elementary applications of the usual cut and alteration
methods; the methods themselves are not claimed as new.

\subsection{Explicit sufficient conditions}
\begin{theorem}\label{late:389a:criteria}
Under \eqref{late:389a:quartile}, suppose $q\ge3$.
Inequality~\eqref{late:389a:target} holds if either of the following is true:
\begin{align}
 \mu(a+4-q)&>\frac{5a+3}{2},\label{late:389a:C}\\
 \mu\left(\frac q3-\frac{1337}{216}\right)
       +a\left(\mu-\frac32\right)&>\frac12.
       \label{late:389a:P}
\end{align}
\end{theorem}
\begin{proof}
Lemma~\ref{late:389a:cut} and \eqref{late:389a:rank} imply
\[
 \gamma_2(G)\le\frac{5a+4\mu(q-2)+3}{2}.
\]
Condition~\eqref{late:389a:C} says precisely that this upper bound is
strictly less than the integer $\mu(q+a)$. Integrality then yields
\eqref{late:389a:target}. Similarly, \eqref{late:389a:sixth} and
\eqref{late:389a:rank} give
\begin{equation}\label{late:389a:rankprob}
 \gamma_2(G)\le\frac{2\mu q}{3}+\frac{3a}{2}
                     +\frac{1337\mu}{216}+\frac12.
\end{equation}
Condition~\eqref{late:389a:P} makes this bound strictly smaller than
$\mu(q+a)$, giving the same conclusion.
\end{proof}

\begin{corollary}\label{late:389a:q20}
Under the stated quartile convention, 389a holds whenever $q\ge20$.
\end{corollary}
\begin{proof}
The left side of \eqref{late:389a:P} minus $1/2$, using
$\mu\ge2$ and $a\ge0$, is at least
\[
 2\left(\frac{20}{3}-\frac{1337}{216}\right)-\frac12
 =\frac{49}{108}>0.
\]
\end{proof}

\begin{corollary}\label{late:389a:smallq}
Inequality~\eqref{late:389a:target} holds for $q=2$, and for $q=3$ with
$\mu\ge3$. It also holds for $q=3$, $\mu=2$, $a=0$.
\end{corollary}
\begin{proof}
If $q=2$, then \eqref{late:389a:quartile} gives $n\le4a+3$ and $a\ge1$.
The cut bound becomes $\gamma_2\le\lfloor(5a+3)/2\rfloor$.
For $a=1,2,3,4$, these upper bounds are $4,6,9,11$, while
$\mu(a+2)-1\ge2(a+2)-1$ gives $5,7,9,11$. For $a\ge5$, use
$\mu\ge a/2$ and
\[
 \frac{a(a+2)}2-1-\frac{5a+3}2
 =\frac{a^2-3a-5}{2}>0.
\]
For $q=3$, condition~\eqref{late:389a:C} reduces to
$\mu(a+1)>(5a+3)/2$, which holds when $\mu\ge3$, and also when
$\mu=2$, $a=0$.
\end{proof}

\begin{corollary}\label{late:389a:a7}
Under the stated quartile convention, 389a holds whenever $a\ge7$.
\end{corollary}
\begin{proof}
First let $q\ge3$. Since $a\le2\mu$, we have $\mu\ge4$. Define
\[
 U=a+4-\frac{5a+3}{2\mu},\qquad
 L=\frac{1337}{72}-3a+\frac{9a+3}{2\mu}.
\]
Conditions~\eqref{late:389a:C} and \eqref{late:389a:P} are respectively
$q<U$ and $q>L$. Their intervals overlap, since
\[
 U-L=4a-\frac{1049}{72}-\frac{7a+3}{\mu}
 \ge28-\frac{1049}{72}-13=\frac{31}{72}>0.
\]
For the inequality, the expression increases with $\mu$, and for
$\mu\ge4$ its coefficient of $a$ is $4-7/\mu>0$. Thus every real $q$
satisfies at least one sufficient condition. The case $q=2$ is
Corollary~\ref{late:389a:smallq}.

For $q=1$, let $\ell$ be the number of leaves. Then
$n\le4\ell+3$, $\ell\ge1$, and $\mu\ge\ell$. A connected graph of
order greater than two has a vertex of degree at least two. Omitting that
one vertex gives a 2-dominating set, so
$\gamma_2\le n-1\le4\ell+2$. Since $a\ge7$,
\[
 \mu(a+1)-1\ge\ell(a+1)-1\ge4\ell+2,
\]
as required.
\end{proof}

\subsection{Residual parameter profiles}
Any counterexample under \eqref{late:389a:quartile} must have $a\le6$
and $q\le19$, avoid the cases in Corollary~\ref{late:389a:smallq}, and,
when $q\ge3$, satisfy both
\begin{equation}\label{late:389a:residual}
 \begin{split}
 \mu(a+4-q)&\le\frac{5a+3}{2},\\
 \mu\left(a+\frac q3-\frac{1337}{216}\right)&\le\frac{3a+1}{2}.
 \end{split}
\end{equation}
If either coefficient of $\mu$ is positive, these inequalities give a
finite bound on $\mu$, after which \eqref{late:389a:rank} bounds $n$.
For $q\ge3$ both coefficients can be nonpositive only in the following
ranges:
\[
\begin{array}{c|c}
a&\text{possibly unbounded degree frequency}\\\hline
0&4\le q\le18\\
1&5\le q\le15\\
2&6\le q\le12\\
3&7\le q\le9.
\end{array}
\]
These are necessary parameter restrictions, not a resolution of 389a.
They do not assert that all graphs realizing the remaining profiles
have been tested. The accompanying exact-rational checker verifies the
displayed algebra on its declared parameter grid; it does not enumerate
graphs. The case $q=1$ outside the proved $a\ge7$ range remains subject
to its separate rank condition and is not covered by
\eqref{late:389a:residual}.


% Integrated source fragment: local_133.tex
\section{Local independence and induced-path density: Conjecture 133}
\label{sec:local-density}
For this part, let $p(G)$ be the maximum \emph{vertex order} of an induced
path, $r(G)=\rad(G)$, and
\[
 \lambda_G(v)=\alpha(G[N_G(v)]),\qquad
 \Lambda(G)=\sum_v\lambda_G(v),\qquad \ell(G)=\Lambda(G)/|V(G)|.
\]
The average $\ell$ is never rounded. A graph is called $C_4$-free here
only when it has no four-cycle as an ordinary subgraph, including one
with chords; $P_j$-free means no induced path on $j$ vertices.
The archived formulation of Conjecture 133 is
\begin{equation}\label{local:target}
 p(G)\ge r(G)+\ell(G)^{\chi_{C_4}(G)},
\end{equation}
where $\chi_{C_4}=1$ exactly in the $C_4$-free case
\cite{src:catalogue,src:definitions}. If a four-cycle occurs, a geodesic
gives $p\ge r+1$; the substantive branch is $p\ge r+\ell$.
In a $C_4$-free graph every neighborhood is a matching plus isolates.
Writing $t(v)$ for the number of triangles through $v$, and $t$ for
the number of triangles, therefore gives
\begin{equation}\label{local:potential}
 \lambda(v)=d(v)-t(v),\quad \Lambda=2m-3t,\quad
 \Lambda(G)-\Lambda(G-v)=2\lambda(v)-t(v).
\end{equation}
The next results retain the sharper deficit and equality information
from the earlier rounds before proving the stronger global $P_7$-free
bound. A clique cutset is a clique whose removal disconnects a connected
graph, and an atom is a connected graph with no clique cutset.


\subsection{\texorpdfstring{An elementary $P_6$-free bound}{An elementary P6-free bound}}
\begin{theorem}\label{local9:thm:p6}
If $G$ is connected, $C_4$-free, and $P_6$-free, then $\ell(G)\le3$, with equality
exactly for the Petersen graph $P$. More precisely, for
$D(G)=3n-2m+3t$ one has
\[
 D(K_1)=3,\quad D(P)=0,\quad D(P-v)=3,
\]
and $D(G)\ge4$ for every other such graph. If $G$ contains a triangle,
then $D(G)\ge6$.
\end{theorem}

The general atom structure studied by Brandst\"adt--Ho\`ang
\cite{src:bh}, and restated in \cite[Theorem 9]{src:ghp}, motivated the following
specialization. Under our stronger four-cycle exclusion the needed
contact and domination arguments are elementary; those external structure
theorems are not dependencies of Theorem~\ref{local9:thm:p6}.

\begin{lemma}\label{local9:lem:atoms}
A $C_4$-free, $P_6$-free atom is one of the graphs in
Table~\ref{local9:tab:atoms}.
\end{lemma}
For the table, label a five-cycle $c_0,\ldots,c_4$ modulo five. Given an
index set $I$, add $a_i$ adjacent to $c_i$ for each $i\in I$, and join
$a_i,a_j$ exactly when $i-j=\pm2\pmod5$. The nonempty permitted sets $I$
induce an edge, $P_3$, $P_4$, or $C_5$ in the pentagram on the indices.
Write the resulting graph as $Q_{|I|}$. Thus $Q_4=P-v$ and $Q_5=P$.
The graph $T$ consists of the five-cycle, a vertex $a$ adjacent to $c_3$,
and a vertex $b$ adjacent to $c_0,c_1,a$.
\begin{table}[ht]
\centering
\caption{The atoms needed for the $P_6$-free deficit bound.}\label{local9:tab:atoms}
\begin{tabular}{lrrrr}\toprule
Graph&$n$&$m$&$t$&$D=3n-2m+3t$\\\midrule
$K_1$&1&0&0&3\\
$K_2$&2&1&0&4\\
$K_3$&3&3&1&6\\
$C_6$&6&6&0&6\\
$C_5$&5&5&0&5\\
$Q_2$&7&8&0&5\\
$Q_3$&8&10&0&4\\
$Q_4=P-v$&9&12&0&3\\
$Q_5=P$&10&15&0&0\\
$T$&7&9&1&6\\\bottomrule
\end{tabular}
\end{table}

\begin{proof}[Proof of Lemma~\ref{local9:lem:atoms}]
Suppose first that the atom $H$ contains an induced five-cycle as above.
An exterior vertex cannot contact two cycle vertices at cycle-distance
two, since that creates a $C_4$. Thus its nonempty contact set is either
$\{c_i\}$ or $\{c_i,c_{i+1}\}$. The cycle is dominating. Otherwise choose
adjacent exterior vertices $y,x$ with $y$ having no cycle contact and $x$
having a contact. Append to $y,x$ four consecutive cycle vertices, starting
at a contact of $x$ and moving away from its other contact if present.
This is an induced $P_6$.

Let $A_i$ and $B_i$ be the singleton-contact and adjacent-pair-contact
classes, respectively. All indices in this paragraph are modulo five.
The following restrictions account for every possible adjacency.
\begin{enumerate}
\item $|B_i|\le1$, since two members would have two common neighbors.
Each $A_i$ induces a matching, being contained in $N(c_i)$.
\item $A_i$ is anticomplete to $A_{i\pm1}$, since an edge would form a
$C_4$ with the corresponding cycle edge. It is complete to $A_{i\pm2}$:
a nonedge $xy$, with $x\in A_i,y\in A_{i+2}$, gives the induced path
$x,c_i,c_{i-1},c_{i-2},c_{i+2},y$.
\item $B_i$ and $B_{i+1}$ cannot both be present. Their vertices, if
adjacent, form a $C_4$ with $c_i,c_{i+1}$; if nonadjacent they are the ends
of the induced path
$b_i,c_i,c_{i-1},c_{i-2},c_{i+2},b_{i+1}$.
Edges between $B_i$ and $B_{i\pm2}$ are also forbidden by a $C_4$ using
the cycle edge between their closest contacts. Thus there are no edges
between distinct $B$-classes.
\item A vertex of $A_i$ can contact a $B$-class only at $B_{i+2}$; every
other such edge creates a $C_4$ using a cycle edge or two-edge segment.
At most one member of $A_i$ can contact the vertex in $B_{i+2}$, since
otherwise that vertex and $c_i$ have two common neighbors.
\end{enumerate}

Every $x\in A_i$ must have a neighbor outside $A_i\cup\{c_i\}$.
Otherwise $N(x)$ consists of $c_i$ and possibly its matching partner in
$A_i$, a clique whose deletion separates $x$ from the remaining cycle.
If $|A_i|\ge2$, restriction 2 and $C_4$-freeness force
$A_{i-2}=A_{i+2}=\varnothing$. The only possible exterior neighbor is
then the unique vertex of $B_{i+2}$, and at most one member of $A_i$ can
contact it. This contradicts the required exterior neighbor for every
member. Hence $|A_i|\le1$.

If $B_i$ is present, its vertex must contact $A_{i-2}$; otherwise its
neighborhood is the clique $\{c_i,c_{i+1}\}$, a clique cutset. This edge
excludes $A_i,A_{i+1}$ by $C_4$-freeness. Applying the same observation to
other $B$-classes excludes $B_{i\pm2}$, while restriction 3 excludes
$B_{i\pm1}$. Thus $B_i$ is the only $B$-class present. The remaining
possible classes $A_{i-1},A_{i+2}$ must either both be absent or both be
present, because each needs an exterior neighbor. If both are present,
\[
 b_i,c_i,c_{i-1},a_{i-1},a_{i+2},c_{i+2}
\]
is an induced $P_6$. Hence both are absent and $H=T$.

If no $B$-class is present, every present $A_i$ must have a neighbor in
$A_{i-2}\cup A_{i+2}$. The selected indices therefore induce a subgraph
of the pentagram without isolated vertices. The possibilities are the
empty set, an edge, $P_3$, $P_4$, and $C_5$, giving the stated five
triangle-free graphs.

Finally suppose $H$ contains no induced $C_5$. If $H$ is chordal, it is a
clique: a noncomplete connected chordal graph has a clique cutset. Its
order is at most three. Otherwise an induced cycle must be $C_6$, since
$C_4,C_5$ are excluded and any longer induced cycle contains a $P_6$.
If there is an exterior vertex, connectedness supplies one with a cycle
contact. It has at most two contacts, which if two are adjacent or
opposite. One contact or two adjacent contacts produce an induced $P_6$
by appending five successive cycle vertices in a suitable direction.
Opposite contacts produce an induced $C_5$. Thus $H=C_6$. The counts in
the table are immediate.
\end{proof}

\begin{lemma}\label{local9:lem:noextension}
Neither $P$ nor $P-v$ admits a one-vertex extension with a nonempty clique
as its contact set while remaining $P_6$-free.
\end{lemma}
\begin{proof}
Represent $P-v$ by the six-cycle $0,1,\ldots,5$ and three vertices
$a_0,a_1,a_2$, where $a_i$ contacts $i,i+3$. Adding a vertex $z$ adjacent
to the three $a_i$ gives $P$. Rotation of the six-cycle covers all contact
types in the following table. Each sequence is an induced $P_5$ beginning
at one contact and avoiding the other when there are two.
\begin{center}
\begin{tabular}{ll}\toprule
Contact set, up to rotation&Induced $P_5$\\\midrule
$\{0\}$&$0,1,a_1,4,3$\\
$\{a_0\}$&$a_0,0,1,a_1,4$\\
$\{z\}$, in $P$ only&$z,a_0,0,1,2$\\
$\{0,1\}$&$0,a_0,3,2,a_2$\\
$\{0,a_0\}$&$0,1,a_1,4,3$\\
$\{z,a_0\}$, in $P$ only&$z,a_1,1,0,5$\\\bottomrule
\end{tabular}
\end{center}
The graphs are triangle-free, so these are all nonempty clique-contact
types. Prepending the added vertex gives an induced $P_6$.
\end{proof}

\begin{proof}[Proof of Theorem~\ref{local9:thm:p6}]
Prove the deficit assertions simultaneously by induction on $n$. Atoms
are covered by Table~\ref{local9:tab:atoms}. Otherwise let $K$ be a clique
cutset, partition the components of $G-K$ into two nonempty unions
$X_1,X_2$, and put $G_i=G[K\cup X_i]$. Both sides are connected proper
induced subgraphs with at least two vertices. Neither can be $P$ or $P-v$:
an outside vertex adjacent to the separator would contradict
Lemma~\ref{local9:lem:noextension}. Thus $D(G_i)\ge4$ by induction.
Inclusion--exclusion for vertices, edges, and triangles gives
\begin{equation}\label{local9:eq:glue3}
 D(G)=D(G_1)+D(G_2)-D(K).
\end{equation}
For $|K|=1,2$, the subtracted values are $3,4$, giving $D(G)\ge5,4$.
For $|K|=3$, both sides have a triangle, so their deficits are at least
six, and subtraction of $D(K_3)=6$ leaves at least six. If $G$ has a
triangle and $|K|\le2$, one side has deficit at least six and the other
at least four; this gives at least seven or six. The induction is complete.
Equation~\eqref{local:potential} gives $D=n(3-\ell)$.

For sharpness, realize $P$ as $KG(5,2)$: its vertices are the two-element
subsets of a five-element set, adjacent when disjoint. It is cubic and
triangle-free; adjacent vertices have no common neighbor, and nonadjacent
vertices have exactly one. In particular it is $C_4$-free. In a hypothetical
induced $P_6$, its six pairs of vertices at path-distance at least three
would require distinct common neighbors outside the path. An exterior
vertex contacts at most two path vertices, whose indices must differ by
at least three, so it can witness at most one of those pairs. Only four
exterior vertices are available. This contradiction proves $P_6$-freeness,
and $\ell(P)=3$.
\end{proof}



\subsection{\texorpdfstring{A $P_7$-free consequence of a published structure theorem}{A P7-free consequence of a published structure theorem}}
\begin{theorem}\label{local9:thm:p7}
If $G$ is connected, $C_4$-free, and $P_7$-free, then $\ell(G)<4$. More precisely,
$F(G)=4n-\Lambda(G)$ is at least eight when $G$ is noncomplete and at least
nine when $G$ contains a triangle. The complete cases $K_1,K_2,K_3$ have
$F=4,6,9$, respectively.
\end{theorem}
The required prior result is Chen--Wu--Xu \cite[Theorem 1.3]{src:cwxpre}:
a connected graph with no induced $P_7,C_4$, or gem, and no clique cutset,
is a clique blow-up of $P$ or has a \emph{bisimplicial} vertex, whose
neighborhood is covered by two cliques. The journal version appeared in
2024 \cite{src:cwx}. A gem contains a four-cycle as a non-induced subgraph,
so this theorem applies under our hypotheses.

\begin{proof}
Induct simultaneously on the two deficit assertions. Complete graphs are
immediate. First let $G$ be a noncomplete atom. A clique blow-up of $P$
can be $C_4$-free only when every bag is a singleton: two vertices in a
bag and one vertex from each of two neighboring bags would form a $C_4$.
Thus this case is $P$, with $F(P)=10$ and no triangle.

Otherwise take a bisimplicial vertex $v$, so $\lambda(v)\le2$. Using
\eqref{local:potential}, with $q=t(v)$, gives
\begin{equation}\label{local9:eq:delete}
 F(G)-F(G-v)=4-2d(v)+3q=4-2\lambda(v)+q\ge0.
\end{equation}
An atom has no cutvertex, so $H=G-v$ is connected. A noncomplete atom
has at least four vertices; hence $H$ has at least three and $F(H)\ge8$
by induction, including its possible complete case $K_3$. This proves
$F(G)\ge8$. If a triangle survives in $H$, then $F(H)\ge9$. If $G$
has a triangle but $H$ does not, then $q\ge1$, making the increment in
\eqref{local9:eq:delete} at least one. Either way $F(G)\ge9$.

For a clique cutset decomposition as before,
\[
 F(G)=F(G_1)+F(G_2)-F(K).
\]
If $|K|=1$, each side has at least two vertices and deficit at least six,
so subtraction of four leaves at least eight. If $|K|=2$, each side has
at least three vertices and deficit at least eight, so subtraction of six
leaves at least ten. If $|K|=3$, both sides have a triangle and deficit at
least nine, leaving at least nine after subtraction. When a triangle occurs
and $|K|\le2$, at least one side has deficit at least nine, giving at
least eleven in either case. This proves the induction and $\Lambda(G)<4n$.
\end{proof}



\subsection{The stronger global bound without an induced seven-vertex path}
\begin{theorem}\label{local15:thm:general}
Every $C_4$-subgraph-free, $P_7$-free graph satisfies $\Lambda(G)\le3n$.
\end{theorem}
It suffices to treat connected graphs, since the two sides are additive
over components. The proof proceeds through dominating and
nondominating five-cycles, the branch with neither a five- nor a
seven-cycle, and finally seven-cycle contacts.


We begin with two five-cycle statements, whose complete proofs are
included below.
\begin{theorem}\label{local15:thm:extension}
Every connected $C_4$-subgraph-free, $P_7$-free graph containing an induced
$C_5$ and no induced $C_7$ satisfies $\Lambda(G)\le3n$.
\end{theorem}
\begin{theorem}\label{local15:lem:dom}
A $C_4$-subgraph-free, $P_7$-free graph with a dominating induced $C_5$
satisfies $\Lambda(G)\le3n$.
\end{theorem}
The stronger deficit estimate in the dominating case will also be
needed. The first five-cycle argument is independent of any structural
classification theorem.
\subsection{A dominating five-cycle: contact classes}
Fix a dominating induced cycle $c_0\dots c_4c_0$, with subscripts modulo
five. Every exterior contact set is a singleton or an adjacent pair:
a distance-two pair gives a $C_4$, and any three cycle vertices contain
such a pair. Let $A_i$ contain the vertices with contact set $\{c_i\}$,
and let $B_i$ contain those with contact set $\{c_i,c_{i+1}\}$.
Four-cycle exclusion gives the following properties.
\begin{enumerate}[label=\textup{(\roman*)}]
\item Each $G[A_i]$ is a matching plus isolated vertices. A three-vertex
path in $N(c_i)$ would complete a $C_4$.
\item $|B_i|\le1$, and there are no edges between different $B$ classes.
Two vertices with the same pair of contacts give a $C_4$; an edge between
different pair-contact classes completes a $C_4$ with suitable cycle
contacts.
\item For $i\ne j$, an $A_i$--$A_j$ edge is possible only if $j=i\pm2$.
Edges between any fixed two classes form a matching, since two neighbors
in $A_j$ give a $C_4$ with $c_j$.
\item An $A_i$--$B_j$ edge is possible only if $i=j+3$. Each $B_j$
vertex has at most one such neighbor, by the same common-neighbor test.
\end{enumerate}
Thus every $B$ vertex has degree two or three and belongs to exactly one
triangle, on its two cycle contacts.

In any $C_4$-subgraph-free graph each neighborhood is a matching plus
isolated vertices. Consequently
\begin{equation}\label{local15:eq:L}
 \Lambda(G)=2m-3t.
\end{equation}
Deleting a $B$ vertex decreases $\Lambda$ by either one or three. If
$G_0=G-B$ satisfies $\Lambda(G_0)\le3|V(G_0)|$, reinstating the $B$ vertices
gives $\Lambda(G)\le \Lambda(G_0)+3|B|\le3n$. We may therefore assume that all
exterior vertices are in $A=\bigcup_iA_i$.

\subsection{Diagonal components}
Let $a=|A|$, let $q$ count the edges within the individual $A_i$, and let
$D$ be the graph on $A$ containing only edges between different contact
classes. Write $d=|E(D)|$. The maximum degree of $D$ is two. Along a
nonbacktracking path of $D$, contact indices progress consistently by
$+2$ or by $-2$ modulo five: reversing direction would give a vertex
two neighbors in one contact class, and hence a $C_4$.

\begin{lemma}\label{local15:lem:blocks}
Every component of $D$ is $P_1$, $P_2$, $P_3$, $P_4$, or $C_5$.
\end{lemma}
\begin{proof}
Take five consecutive vertices $x_0,\dots,x_4$ of a diagonal path.
Their contacts, in order, are $c_i,c_{i+2},c_{i+4},c_{i+1},c_{i+3}$. If $x_0x_4$ is absent,
then
\[
 c_i,x_0,x_1,c_{i+2},c_{i+3},x_4,x_3
\]
is an induced $P_7$. No within-class matching edge can be a chord, since
the five exterior vertices have distinct contact classes. The only
possible additional diagonal chord is $x_0x_4$. If that edge is present,
the vertices form an entire $C_5$ component, since each already has
diagonal degree two. Every other component has at most four vertices,
which proves the assertion.
\end{proof}
For a component $Q$ of $D$, put $e(Q)=2|E(Q)|-|V(Q)|$. Its values on
$P_1,P_2,P_3,P_4,C_5$ are respectively $-1,0,1,2,5$.
The only triangles of $G$ arise from the $q$ within-class matching edges
and their cycle contacts. Hence \eqref{local15:eq:L} gives
\begin{equation}\label{local15:eq:deficit}
 3n-\Lambda(G)=5+a+q-2d=5+q-\sum_Q e(Q).
\end{equation}
It suffices to bound the sum of the positive excesses by $5+q$.

\subsection{Rooted paths force matching edges}
A component is \emph{positive} if it is $P_3$, $P_4$, or $C_5$.
For such a component $Q$, let $T(Q)$ be the set of contact indices of
endpoints of its induced three-vertex paths. Then
\[
 |T(P_3)|=2,\qquad |T(P_4)|=4,\qquad
 T(C_5)=\mathbb Z/5\mathbb Z.
\]

\begin{lemma}\label{local15:lem:force}
If two distinct positive components have intersecting $T$ sets, some
within-class matching edge joins them. Moreover a $C_5$ component
requires at least two such edges to each $P_3$ component, four to each
$P_4$ component, and five to each other $C_5$ component.
\end{lemma}
\begin{proof}
For the first assertion, choose a three-vertex path in each component
ending in the same class $A_i$, and join the paths through $c_i$.
The seven vertices are distinct. Each three-vertex path uses three
distinct contact classes, so $c_i$ has no other neighbor on either path.
No diagonal edge joins the two components. Thus a matching edge must
join the two paths, or their concatenation is an induced $P_7$.

For the stronger assertion, let $x\in A_i$ be an endpoint of a
three-vertex path in the other component. The $C_5$ component contains
one vertex $y\in A_i$, with rooted three-vertex paths in both circular
orientations. Choose the orientation opposite to that of the path rooted
at $x$. The two sets of non-root contact indices are disjoint. Therefore
the only possible matching chord in the joined seven-vertex path is
$xy$, and that edge must be present. A $P_3$ has two possible root
vertices, a $P_4$ has four, and another $C_5$ has five. This forces the
stated numbers of distinct edges. Edges required by distinct component
pairs are also distinct.
\end{proof}

\subsection{Counting the excess}
Let $h,k,s$ be the numbers of $C_5$, $P_4$, and $P_3$ components of $D$.
If $h\ge1$, Lemma~\ref{local15:lem:force} gives
\[
 q\ge5\binom h2+h(4k+2s).
\]
Consequently
\[
 5h+2k+s-q
 \le 5h-5\binom h2+(2-4h)k+(1-2h)s\le5.
\]
Here $h-\binom h2\le1$ for every positive integer $h$, and the remaining
coefficients are nonpositive.

Now suppose $h=0$. A four-element contact set intersects every other
positive contact set, which has size two or four. Thus
$q\ge\binom k2+ks$. If $k\ge1$, then
\[
 2k+s-q\le2k-\binom k2+(1-k)s\le3.
\]
If $k=0$, form the intersection graph of the $s$ two-element sets from
the $P_3$ components. Distinct connected components of this intersection
graph use disjoint collections of contact indices. Each uses at least
two of the five indices, so there are at most two components. The
intersection graph has at least $s-2$ edges, each requiring a distinct
matching edge by Lemma~\ref{local15:lem:force}. Therefore $s-q\le2$; the case
$s=0$ is immediate.

In all cases, the positive excess sum is at most $5+q$. The $P_1$ and
$P_2$ components have nonpositive excess, so \eqref{local15:eq:deficit} proves
$\Lambda(G)\le3n$. Reinstating the $B$ vertices proves
Theorem~\ref{local15:lem:dom}.

\subsection{The nondominating layers}
We prove Theorem~\ref{local15:thm:extension}. Fix an induced cycle
$c_0\dots c_4c_0$, and use the contact classes $A_i,B_i$ defined above,
now without assuming that this cycle dominates. Let $R$ be the set of
vertices with no cycle contact. Every first-layer vertex $x$ starts an
induced five-vertex path using $x$ and four consecutive cycle vertices;
for adjacent-pair contacts, omit the unused contact.

Every component $Q$ of $G[R]$ is complete to each first-layer vertex
that touches it. Otherwise an edge $ab$ in $Q$ with $xa$ present and $xb$
absent gives an induced $P_7$, starting $b,a,x$ and continuing along the
four cycle vertices. The neighborhood of $x$ is a matching plus isolated
vertices, so connectedness forces $Q$ to be $K_1$ or $K_2$. Every
component touches the first layer because $G$ is connected. In particular
there are no vertices at distance at least three from the cycle.

A $K_2$ component has exactly one first-layer neighbor: two would give its
two vertices two common neighbors and hence a $C_4$. It is a pendant
triangle. Removing it decreases $\Lambda$ by three and $n$ by two, so it is
safe for proving $\Lambda\le3n$.

\begin{lemma}\label{local15:lem:Rdegree}
Every singleton component of $G[R]$ has degree at most two. Its possible
pairs of neighbors have types
\[
 (A_i,A_{i\pm2}),\qquad(A_i,B_{i+2}),\qquad(B_i,B_{i\pm2}).
\]
In the first pair the neighbors are adjacent; in the last pair they are
nonadjacent.
\end{lemma}
\begin{proof}
Let $z$ be the singleton. Two neighbors $x,y$ cannot share a cycle contact,
or they give a $C_4$ with $z$. For $x\in A_0,y\in A_1$, the sequence
\[
 c_2,c_3,c_4,c_0,x,z,y
\]
is an induced $P_7$. For $x\in A_0,y\in A_2$ with $xy$ absent, the cycle
\[
 x,c_0,c_4,c_3,c_2,y,z,x
\]
is an induced $C_7$. The same cycle excludes an $A_0,B_1$ pair; in that
case $xy$ is already forbidden by the $C_4$ on $x,c_0,c_1,y$.
Reflection excludes the $A_0,B_3$ pair, while $B_0,B_4$ share a contact
with $A_0$. Two $B$ types must have disjoint cycle edges, so their indices
differ by $\pm2$. Their mutual edge is excluded by $C_4$-freeness.
These exhaust all pairs up to rotation and reflection.

No three types are pairwise compatible. The compatibility graphs among
$A$ types alone and among $B$ types alone are five-cycles. Moreover an
$A$ type is compatible with only one $B$ type, and conversely.
Repeating a type shares a cycle contact. Thus $d(z)\le2$.
\end{proof}

A singleton leaf has deletion cost $\Delta \Lambda=2$ for one vertex. A
degree-two singleton whose neighbors are adjacent has $\Delta \Lambda=1$.
Remove these safe vertices and all the $K_2$ components. Every remaining
$R$ vertex has degree two with nonadjacent neighbors and subdivides
$A_i$--$B_{i+2}$ or $B_i$--$B_{i\pm2}$. There is at most one subdivision
vertex for each fixed endpoint pair, by $C_4$ exclusion.

\subsection{Removing isolated marked vertices}
Call an $A$ vertex \emph{marked} if it is an endpoint of a remaining
$A$--$B$ subdivision. Each $A_i$ vertex supports at most one such
subdivision, since its only possible $B$ endpoint is the unique $B_{i+2}$,
and repeating the pair creates a $C_4$.

Let $D$ again denote the diagonal-edge graph on $A$. Suppose a marked
vertex $x$ has degree zero in $D$, with subdivision vertex $z$. Its
neighbors are its cycle contact, $z$, and possibly its within-class mate.
It cannot have a direct $B$ neighbor, since the only possible one is the
nonadjacent endpoint of the subdivision. Deleting $x,z$ therefore removes
three edges and no triangle, or four edges and one triangle. Its cost
$\Delta \Lambda$ is six or five for two vertices. Perform all such safe pair
deletions. Every remaining marked vertex now has a diagonal neighbor.
The fixed cycle and all remaining contact vertices stay connected.

\subsection{Restrictions imposed by a subdivision}
Let $z$ subdivide $B_j$--$x$, where $x\in A_{j+3}$, and call $B_j$
\emph{active}. All subscripts in the remainder are modulo five.

\paragraph{At most two marks at one active vertex.}
Suppose $x$ has a diagonal neighbor $y\in A_j$; the $A_{j+1}$ case is
symmetric. If another mark $x'$ at $B_j$ is not adjacent to $x$, with
subdivision vertex $z'$, then
\[
 c_{j+1},c_j,y,x,c_{j+3},x',z'
\]
is an induced $P_7$. In particular, $yx'$ is impossible because $y$
already has its unique neighbor $x$ in that contact class. Thus every
other mark must be the within-class mate of $x$. There are at most two
marks at $B_j$, and if there are two, they are matched.

\paragraph{Neighboring double-contact types are absent.}
If $B_{j+1}$ existed, then
\[
 c_j,B_j,z,x,c_{j+3},c_{j+2},B_{j+1}
\]
would be an induced $P_7$. Reflection excludes $B_{j-1}$. Consequently
active indices form an independent set in a five-cycle, so there are at
most two. One active $B$ leaves at most three $B$ vertices in total;
two active $B$ vertices leave exactly those two.

\paragraph{Other double-contact vertices have no direct single-contact neighbor.}
The only other possible types are $B_{j+2},B_{j+3}$. If $B_{j+2}$ has its
permitted direct neighbor $y\in A_j$, then
\[
 c_{j+1},c_j,y,B_{j+2},c_{j+3},x,z
\]
is an induced $P_7$. The potential chord $xy$ would give the $C_4$
$x,c_{j+3},B_{j+2},y,x$. Reflection handles $B_{j+3}$. Thus, with one
active $B$, at most that vertex has a direct $A$ neighbor. With two active
vertices, neither can have one.

\paragraph{A subdivision between double-contact vertices leaves only two such vertices.}
Suppose $z$ subdivides $B_0$--$B_2$. The other three types give these
induced $P_7$ sequences:
\begin{align*}
 B_1 &: \quad c_4,c_0,B_0,z,B_2,c_2,B_1,\\
 B_3 &: \quad c_2,B_2,z,B_0,c_0,c_4,B_3,\\
 B_4 &: \quad c_1,B_0,z,B_2,c_3,c_4,B_4.
\end{align*}
Rotation and reflection cover every possible pair. Hence a $B$--$B$
subdivision forces exactly two $B$ vertices, and there is at most one
such subdivision in total.

\subsection{The residual deficit}
Let $b$ be the number of $B$ vertices, $e$ the number of direct $A$--$B$
edges, $s$ the number of $A$--$B$ subdivisions, and $u$ the number of
$B$--$B$ subdivisions. Put
\[
 G_A=G[V(C)\cup A],\qquad D_A=3|V(G_A)|-\Lambda(G_A).
\]
Adding the $B$ vertices creates $2b+e$ edges and $b$ triangles. Each
remaining subdivision vertex creates two edges and no triangle. Hence
\begin{equation}\label{local15:eq:residual}
 3n-\Lambda(G)=D_A+2b-2e-s-u.
\end{equation}
The component proof above gives $D_A\ge2$ if $D$ has no $C_5$ component:
in that case the positive diagonal excess minus the matching count is
at most three, and \eqref{local15:eq:deficit} starts with five.

The restrictions just established give
\begin{equation}\label{local15:eq:cost}
 2b-2e-s-u\ge-2.
\end{equation}
Here is the complete case check. With no active vertex, $s=0$; use $e\le b$
if $u=0$, or $b=2,e\le2$ if $u=1$, obtaining a bound of zero or $-1$.
With one active vertex, $s\le2$ and $e\le1$. For $b=1$ we have $u=0$
and the bound is $-2$; for $b=2$, $u\le1$ gives $-1$; for $b=3$, $u=0$
gives two. With two active vertices, $b=2,e=0,s\le4,u\le1$, giving $-1$.
Thus \eqref{local15:eq:residual} and \eqref{local15:eq:cost} finish the case in which
$D$ has no five-cycle component.

\subsection{The Petersen alternative}
Suppose $D$ has a $C_5$ component, labelled $u_i\in A_i$ with edges
$u_iu_{i\pm2}$. Together with the original cycle these vertices induce
the Petersen graph. No other component of $D$ can have an edge: an edge
$xy$ with $x\in A_i,y\in A_{i+2}$ outside this component gives
\[
 c_i,x,y,c_{i+2},c_{i+3},u_{i+3},u_{i+1}
\]
as an induced $P_7$. The only possible within-class mates of $x,y$
have indices $i,i+2$, neither of which occurs among the two displayed
$u$ vertices.

Every remaining mark therefore belongs to the displayed component,
since marks with diagonal degree zero were removed. A subdivision $z$
between $B_j$ and $u_{j+3}$ gives the induced $P_7$
\[
 c_j,B_j,z,u_{j+3},c_{j+3},c_{j+2},u_{j+2}.
\]
Thus there is no $A$--$B$ subdivision. A subdivision $z$ between $B_0,B_2$
gives the induced $P_7$
\[
 c_0,c_1,u_1,u_3,c_3,B_2,z.
\]
There is no $B$--$B$ subdivision either. The fixed cycle now dominates
the reduced graph, and Theorem~\ref{local15:lem:dom} applies. Reversing all safe
deletions completes the proof of Theorem~\ref{local15:thm:extension}.


\subsection{The branch with neither a five-cycle nor a seven-cycle}
Assume that $G$ has no induced $C_5$ or $C_7$. If it also has no induced
$C_6$, it is chordal: a longer hole contains $P_7$, and a four-cycle is
excluded. A chordal graph has a simplicial vertex; here its degree is
at most two, since a clique of order four contains a four-cycle.
Deleting a simplicial vertex from a connected graph leaves a connected
graph, unless it was the only vertex. Reinstating a leaf adds two to
$\Lambda$, while reinstating a degree-two simplicial vertex adds one, by
\eqref{local15:eq:L}. Induction gives $\Lambda\le2(n-1)$.

We may therefore fix an induced cycle $C=c_0\dots c_5c_0$, with indices
modulo six. Each exterior cycle-contact set is a singleton, an adjacent
pair, or an opposite pair. Indeed, a distance-two pair gives a four-cycle,
and every three cycle vertices contain such a pair. Denote the singleton
classes by $A_i$, the adjacent-pair classes by $B_i$, and the opposite-pair
classes by $Z_i$, with contact sets respectively
\[
 \{c_i\},\qquad \{c_i,c_{i+1}\},\qquad \{c_i,c_{i+3}\}.
\]
Here $Z_i=Z_{i+3}$. Each $B_i$ and $Z_i$ has at most one vertex, by the
common-neighbor test. We also use $B_i,Z_i$ for those vertices when
present. There are at most three $Z$ vertices, and they are independent:
an edge between two different opposite-contact types completes a
four-cycle with adjacent cycle contacts.

Let $R$ consist of vertices with no cycle contact. Every $A$ or $B$
vertex starts an induced $P_6$ using five consecutive cycle vertices.
An edge from it to $R$ would therefore produce an induced $P_7$. Thus
only the $Z$ vertices can have neighbors in $R$.

\subsubsection{Removing first-layer vertices at nonnegative deficit cost}
Every $B$ vertex has degree two. To see this, four-cycle exclusion
permits a $B_0$--$A$ edge only to $A_3$ or $A_4$, and a $B_0$--$B$ edge
only to $B_3$. Such edges give an induced five-cycle; for example, with
$x\in A_3$ the cycle is $B_0,x,c_3,c_2,c_1,B_0$, and for the second
case it is $B_0,c_1,c_2,c_3,B_3,B_0$. All $B$--$Z$ edges give a
four-cycle. Edges to $R$ were excluded above. Therefore each $B$ is a
degree-two simplicial vertex, belonging to one triangle, and deleting it
decreases $\Lambda$ by one. Remove all $B$ vertices.

Edges among distinct $A$ classes can join only opposite classes:
adjacent classes give a four-cycle, and distance-two classes give an
induced five-cycle. These opposite edges form a matching, since every
vertex has at most one neighbor in a fixed other contact class.
Within each $A_i$ there is a matching plus isolated vertices.

An $A_i$ vertex can meet only $Z_i$ among the $Z$ vertices; another type
would complete a four-cycle. If $x\in A_i$ is adjacent to $Z_i$, it
cannot have a within-class mate $y$, because $x,y,c_i,Z_i,x$ would be a
four-cycle. Nor can it have an opposite-class neighbor $y$, because
$x,y,c_{i+3},Z_i,x$ would be a four-cycle. It has no neighbor in $R$.
Thus $x$ is another degree-two simplicial vertex. Delete all such
vertices, at cost one in $\Lambda$ per vertex.

The remaining $A$ vertices meet only the cycle and other $A$ vertices.
Let $a$ be their number, $q$ their within-class matching count, and $e$
the number of opposite-class edges. Reinstating them adds
\[
 \Delta \Lambda=2a+2e-q,
 \qquad \Delta(3n-\Lambda)=a-2e+q\ge0,
\]
since $2e\le a$. We may therefore remove them while retaining this
nonnegative contribution to the deficit.

\subsubsection{Components attached to opposite-contact vertices}
The reduced graph consists of $C$, at most three $Z$ vertices, and $R$.
Every component $Q$ of $G[R]$ touches at least one $Z$ vertex.
No vertex of $R$ can meet two distinct $Z$ vertices: their two opposite
pairs have adjacent cycle contacts, and these, together with the common
neighbor, give an induced $C_5$.

For each $Z$ vertex $z$, the graph $G[R\cup\{z\}]$ has no induced
four-vertex path with endpoint $z$. Otherwise append the three cycle
vertices of a rooted induced $P_4$ starting at $z$, giving $P_7$.
Fix a component $Q$ touching $z$, and put
\[
 X=N(z)\cap Q,\qquad Y=Q\setminus X.
\]
Every vertex of $Q$ is within distance two of $z$, by the forbidden
rooted path. Each $y\in Y$ has a unique neighbor in $X$; two would give
a four-cycle through $z$. Call that neighbor its parent. Both $G[X]$
and each set of children of one parent induce matchings plus isolated
vertices. No edge joins children with different parents, since then
$z,x,y,y'$ would be an induced four-vertex path. Connectedness of $Q$
now implies that $X$ is $K_1$ or $K_2$. Thus $Q$ consists of one or two
adjacent parent vertices and their private child sets.

If $Q$ touches only $z$, then $Q\cup\{z\}$ is a block graph with all
blocks edges or triangles. For $q=|Q|$ and $\tau$ its triangle count,
including triangles through $z$, the attachment has $q+\tau$ edges.
It adds $2q-\tau\le2q$ to $\Lambda$, and hence adds nonnegative deficit.

Suppose instead that $Q$ touches another $Z$ vertex $z'$. Its neighbor
$y$ in $Q$ is a child of some $x\in X$, since no vertex meets two $Z$
vertices. If $X$ contained another vertex $x'$, then $z',y,x,x'$ would
be an induced four-vertex path rooted at $z'$. Hence $X=\{x\}$.
There is at most one child adjacent to $z'$, by the common-neighbor
condition applied to $x,z'$. Every other child $w$ must be adjacent to
$y$, or $z',y,x,w$ would be the forbidden rooted path. Since the child
set induces a matching, there is at most one such $w$. Consequently
$Q$ is $K_2$ or $K_3$, and its different boundary $Z$ vertices attach
to different vertices of $Q$. The only contributions are
\[
\begin{array}{c|c}
\text{attachment}&\Delta(3n-\Lambda)\\\hline
K_2\text{ with two }Z\text{ neighbors}&6-2(1+2)=0\\
K_3\text{ with two }Z\text{ neighbors}&9-[2(3+2)-3]=2\\
K_3\text{ with three }Z\text{ neighbors}&9-[2(3+3)-3]=0.
\end{array}
\]
Thus every component adds nonnegative deficit. Finally the base $C_6$
with $k\le3$ opposite-contact vertices has
\[
 n=6+k,\qquad m=6+2k,\qquad t=0,\qquad 3n-\Lambda=6-k\ge3.
\]
Reinstating all deleted vertices proves $\Lambda\le3n$ in this branch.
Together with Theorem~\ref{local15:thm:extension}, this settles every graph
without an induced seven-cycle.

\subsection{Seven-cycle contacts and a connected deletion}
Fix an induced $C_7=c_0\dots c_6c_0$, with indices modulo seven.
Every nonempty exterior cycle-contact set is of one of the forms
\[
 X_i=\{c_i,c_{i+3}\},\qquad
 Y_i=\{c_i,c_{i+3},c_{i+4}\}.
\]
Indeed, distance-two contacts give a four-cycle. A singleton or adjacent
pair produces $P_7$ by appending six consecutive cycle vertices, omitting
the unused contact in the adjacent-pair case. Four contacts include a
distance-two pair, and the three-contact sets without such a pair are
exactly the displayed $Y_i$. No type is repeated, since that would give
two common cycle neighbors.

All exterior vertices having a cycle contact are independent. For an
edge $xy$ between two such vertices, no cycle edge may run from a
contact of $x$ to a distinct contact of $y$, since that edge and $xy$
complete a four-cycle. If $x$ has type $X_0$, the contacts of $y$ must
therefore lie in $\{c_0,c_3,c_5\}$; the only admissible type contained
there is $X_0$, already excluded by repetition. If $x$ has type $Y_0$,
the permitted set is just $\{c_0\}$, which contains no admissible type.

The following compatibility restrictions are necessary:
\begin{align}
 X_i,X_{i\pm2}&\text{ cannot coexist},\label{local15:eq:xx}\\
 X_i,Y_j&\text{ coexist only if }j-i\in\{4,5,6\},\label{local15:eq:xy}\\
 Y_i,Y_j&\text{ coexist only if }j-i\in\{2,5\}.
 \label{local15:eq:yy}
\end{align}
The common-two-contact cases give four-cycles. All remaining exclusions
follow by rotating or reflecting these induced $P_7$ witnesses, with
$x,y$ denoting the two indicated exterior types:
\[
\begin{array}{c|l}
X_0,X_2&c_1,c_0,x,c_3,c_4,c_5,y\\
X_0,Y_1&c_2,c_3,x,c_0,c_6,c_5,y\\
X_0,Y_2&c_1,c_0,x,c_3,c_4,c_5,y\\
Y_0,Y_1&c_2,c_3,x,c_0,c_6,c_5,y.
\end{array}
\]
The exact contact sets and the just-proved independence verify every
chord condition. The standard $X/Y$ partition and several of these
exclusions also occur in Chen--Wu--Xu \cite[proof of Theorem~1.1,
claims (1)--(6)]{src:cwx}; the exclusion of a four-cycle as a subgraph is
used here throughout.

If there is no $Y$, the $X$ indices form an independent set in the
seven-cycle with steps $\pm2$, and there are at most three. If $Y_0$
occurs, the only possible $X$ indices are $1,2,3$, with $X_1,X_3$
incompatible. The only other possible $Y$ indices are $2,5$, which are
incompatible with each other. If $Y_2$ also occurs, only $X_3$ remains
permitted, and the reflected case is the same. Thus the entire first
layer has at most three vertices.

\begin{proof}[Proof of Theorem~\ref{local15:thm:general}]
Induct on the order of a connected graph. All graphs without induced
$C_7$ have been settled. Otherwise use the contact description above.
If a $Y$ occurs, take it to be $Y_0$. The only other possible $Y$ types
are $Y_2,Y_5$, neither of which touches $c_3$ or $c_4$. The vertex
$c_3$ has degree three unless $X_3$ occurs, and $c_4$ has degree three
unless $X_1$ occurs. Since these two $X$ types cannot coexist, one of
$c_3,c_4$ has degree exactly three and lies in exactly one triangle.
Delete that vertex $v$. Its two cycle neighbors are connected by the
remaining cycle path, and its third neighbor $Y_0$ retains two cycle
contacts. Thus $G-v$ is connected. Since
\[
 \Lambda(G)-\Lambda(G-v)=2d(v)-3t(v)=3,
\]
induction proves $\Lambda(G)\le3(n-1)+3=3n$.

If there is no $Y$, there are two adjacent cycle vertices untouched by
all $X$ types. With at most one $X$ this is immediate. For two $X$
indices, up to rotation and reflection the pairs are $\{0,1\}$ or
$\{0,3\}$, leaving respectively $c_5,c_6$ or $c_1,c_2$ untouched.
Every permitted three-index set is a rotation of $\{0,1,4\}$, again
leaving $c_5,c_6$ untouched. Delete this adjacent degree-two pair. The
remaining cycle path is connected, and no exterior vertex touched the
deleted pair, so the remaining graph is connected. Exactly three edges
and no triangles are lost. Therefore induction gives
$\Lambda(G)\le3(n-2)+6=3n$. This completes the proof.
\end{proof}



\subsection{\texorpdfstring{A sharp elementary bipartite $P_8$-free bound}{A sharp elementary bipartite P8-free bound}}
\begin{theorem}\label{local9:thm:p8}
A connected bipartite $C_4$-free, $P_8$-free graph has average degree at
most three, with equality exactly for the Heawood graph.
\end{theorem}
\begin{proof}
Trees have average degree $2-2/n<2$. Otherwise the girth is six or eight:
a shortest cycle is induced and even, and a longer one contains $P_8$.
If the girth is eight, fix a shortest $C_8$. An exterior vertex has at
most one cycle contact: two contacts have even cycle-distance, and yield
a $C_4$ or $C_6$. A single contact followed by seven successive cycle
vertices gives an induced $P_8$. Connectedness therefore forces $G=C_8$.

Now fix an induced $C_6=c_0c_1\cdots c_5c_0$, with indices modulo six.
Let $A,B$ be the vertices at distance one and two from this cycle.
Every vertex in $A$ has exactly one cycle contact. Its contacts must have
the same parity on the cycle, and two would be at distance two, creating
a $C_4$. Put
\[
 A_i=\{x\in A:N(x)\cap V(C_6)=\{c_i\}\}.
\]
There are no vertices at distance at least three: a shortest path
$z,y,x,c_0$ to the cycle, followed by $c_1,c_2,c_3,c_4$, would be an
induced $P_8$. Thus $V(G)=V(C_6)\cup A\cup B$.

Every edge in $G[A]$ joins $A_i$ to $A_{i+3}$. Bipartiteness requires
opposite parity, and adjacent contact indices would yield a $C_4$ using
the cycle edge. Moreover $G[A]$ is a matching: two neighbors of
$x\in A_i$ would both lie in $A_{i+3}$, forming a $C_4$ with $c_{i+3}$.
Write $k=|E(G[A])|$. The set $B$ is independent. If $bb'$ were an edge
and $x\in A_i$ a neighbor of $b$, bipartiteness excludes $b'x$, and
\[
 b',b,x,c_i,c_{i+1},c_{i+2},c_{i+3},c_{i+4}
\]
would be an induced $P_8$.

A vertex $b\in B$ has at most one neighbor in any $A_i$, by $C_4$-freeness,
and can contact classes of only one parity. We claim that if it contacts
$A_i$, it is complete to $A_{i+2}$ and $A_{i+4}$. Indeed, for
$x\in A_i$ adjacent to $b$ and $y\in A_{i+4}$ nonadjacent to $b$,
\[
 b,x,c_i,c_{i+1},c_{i+2},c_{i+3},c_{i+4},y
\]
is an induced $P_8$. Here $xy$ is absent because exterior edges only join
opposite contact classes. Reverse the cycle direction to treat $A_{i+2}$.

Let $A^0=A_0\cup A_2\cup A_4$ and $A^1=A_1\cup A_3\cup A_5$,
and let $B^j$ consist of vertices whose neighbors lie in $A^j$. Write
\[
 a_j=|A^j|,\quad b_j=|B^j|,\quad e_j=|E(A^j,B^j)|,
\]
and let $q_j$ count the nonempty contact classes defining $A^j$.
Call parity $j$ special if $q_j\ge2$ and $b_j>0$. For a special parity,
the claim forces any $b\in B^j$ to contact every nonempty class. Applying
the claim from a second class also makes it complete to its original
class. Each class is therefore a singleton. Every member of $B^j$ contacts
all these vertices, so $C_4$-freeness gives $b_j=1$. Consequently
\begin{equation}\label{local9:eq:special}
 a_j=q_j\in\{2,3\},\quad b_j=1,\quad e_j=q_j,
 \quad R_j:=a_j+3b_j-2e_j=3-q_j\ge0.
\end{equation}
For a nonspecial parity, either $b_j=0$ or $q_j\le1$; in the latter case
every vertex of $B^j$ has degree one. Thus
\begin{equation}\label{local9:eq:nonspecial}
 R_j\ge a_j\ge0.
\end{equation}
Counting the cycle edges, cycle-to-$A$ edges, matching edges, and
$A$--$B$ edges gives
\begin{equation}\label{local9:eq:bipcount}
 3n-2m=6+R_0+R_1-2k.
\end{equation}
If neither parity is special, $R_0+R_1\ge a_0+a_1\ge2k$, giving a
positive deficit of at least six. If exactly one is special, its $A$-side
has at most three vertices, so $k\le3$, while the other side contributes
$R\ge k$. The deficit is then at least $6-k\ge3$. If both are special,
\[
 3n-2m=12-q_0-q_1-2k\ge0,
\]
since $q_0,q_1\le3$ and $k\le\min(q_0,q_1)$. Equality requires
$q_0=q_1=k=3$.

The equality graph has vertices $c_0,\ldots,c_5$, $a_0,\ldots,a_5$,
$z_0,z_1$, and exactly the edges
\[
 c_ic_{i+1},\quad c_ia_i\ (0\le i<6),\quad
 a_ia_{i+3}\ (0\le i<3),\quad
 z_0a_i\ (i=0,2,4),\quad z_1a_i\ (i=1,3,5).
\]
This is the Heawood graph, equivalently the incidence graph of the Fano
plane. The displayed graph is cubic, bipartite, and $C_4$-free. Every
pair in the same bipartition has exactly one common neighbor, as is seen
directly, or from $C_4$-freeness and
$7\binom32=\binom72$. In an induced path $v_0,\ldots,v_7$, the pairs
$v_6,v_0$ and $v_6,v_2$ would need distinct common neighbors outside the
path: they cannot coincide because $v_0,v_2$ already share $v_1$, and an
internal contact would be a chord. Together with $v_5,v_7$, these give
$v_6$ four neighbors, a contradiction. Hence the equality graph is indeed
$P_8$-free.
\end{proof}

\begin{remark}
The same layer argument gives a useful simpler refinement: a cyclic
connected bipartite $C_4$-free, $P_7$-free graph satisfies $3n-2m\ge6$.
A shortest cycle is $C_6$. A distance-two vertex would give $P_7$ by
appending five cycle vertices, so the cycle is dominating. Its exterior
vertices have unique contacts and induce a matching. With $k$ matching
edges and $s$ exterior isolated vertices, $n=6+2k+s$, $m=6+3k+s$, and
$3n-2m=6+s$.
\end{remark}



\section{Radius classifications for Conjecture 133}
\label{sec:local-radius}
The classical induced-path radius inequality used in this section is
\begin{equation}\label{local:chung}
 p(G)\ge2r(G)-1.
\end{equation}
It is recorded by Erd\H{o}s--Saks--S\'os with credit to Fan Chung
\cite[Section II, pp.~64--65]{src:ess}. We first prove an arbitrary-cycle
attachment theorem, then use connected domination to obtain complete
radius-four and radius-five classifications. These are statements about
all finite graphs with the stated hypotheses, not bounded-order searches.
\subsection{Dominating even cycles}


\begin{theorem}\label{local16:thm:main}
Fix $k\ge3$. Let $G$ be a connected $C_4$-subgraph-free graph with a
dominating induced cycle $C=c_0\dots c_{2k-1}c_0$. The graph has radius
$k+1$ and is $P_{2k+2}$-free if and only if it has the following form:
\begin{enumerate}[label=\textup{(\roman*)}]
\item For every $i$, add a nonempty set $A_i$, join all its vertices to
$c_i$, and place a matching plus isolated vertices on $A_i$.
\item For each cycle edge $c_ic_{i+1}$, optionally add one vertex $b_i$
adjacent to its endpoints.
\item Add no other vertices or edges.
\end{enumerate}
Every graph in the classification satisfies
\[
 p(G)=2k+1=2r(G)-1,\qquad \Lambda(G)=2n-t\le2n,
\]
where $n=|V(G)|$ and $t$ counts its triangles.
\end{theorem}

\subsection{Contact sets and radius witnesses}
Put $N=2k$, and take all indices modulo $N$. Suppose an exterior vertex
$x$ has two cycle contacts at circular distance $d\ge3$, choosing the
shorter arc so that $d\le k$. Every cycle vertex is within
\[
 1+\left\lfloor\frac{2k-d}{2}\right\rfloor\le k-1
\]
of $x$. Since $C$ dominates, $x$ has eccentricity at most $k$, contrary
to the assumed radius. A pair of contacts at distance two gives a
four-cycle. Thus every exterior contact set is a singleton or an
adjacent pair; three contacts would contain a prohibited pair.
Denote the singleton class by $A_i$ and the adjacent-pair vertex by
$b_i$, with contacts $c_i,c_{i+1}$. There is at most one vertex of the
latter type, by the common-neighbor condition for four-cycle exclusion.
Each $G[A_i]$ is a matching plus isolated vertices, since a three-vertex
path in $N(c_i)$ would complete a four-cycle.

The only possible cross-class edges are
\begin{align*}
 A_i\text{--}A_j&:\quad \text{circular distance at least two},\\
 A_i\text{--}b_j&:\quad j-i\in\{2,\ldots,N-3\},\\
 b_i\text{--}b_j&:\quad j-i\in\{3,\ldots,N-3\}.
\end{align*}
The excluded cases give a four-cycle using the two exterior endpoints
and an appropriate cycle edge, or the shared cycle contact and its
neighbor.

The cycle vertex $c_{i+k}$ is within distance $k$ of every vertex except
possibly those in $A_i$. Since its eccentricity is at least $k+1$,
there is a vertex $z_i\in A_i$ with
\begin{equation}\label{local16:eq:witness}
 \operatorname{dist}(c_{i+k},z_i)=k+1.
\end{equation}
In particular every $A_i$ is nonempty. The vertex $z_i$ has no exterior
neighbor outside $A_i$. Indeed, the cycle contact of any allowed $A_j$
neighbor is within cycle distance $k-2$ of $c_{i+k}$, giving a route of
length at most $k$ to $z_i$. The same argument applies to either contact
of an allowed $b_j$ neighbor. Both contradict \eqref{local16:eq:witness}.
Thus the only possible exterior neighbor of $z_i$ is one mate in $A_i$.

\subsection{Excluding every cross-class edge}
Every displayed sequence in this section has exactly $N+2$ distinct
vertices. Consecutive indices in an ellipsis follow the indicated
ordinary cyclic direction. The omitted cycle vertices separate the two
cycle segments. The lack of cross-class neighbors of the $z_i$ excludes
all further exterior chords.

\subsubsection{Two distant single-contact classes}
Suppose $x\in A_0$ and $y\in A_d$ are adjacent, where
$3\le d\le N-3$. Then
\[
\begin{split}
 z_{d-2},c_{d-2},c_{d-3},\ldots,c_0,
 x,y,\qquad\qquad\\
 c_d,c_{d+1},\ldots,c_{N-2},z_{N-2}
\end{split}
\]
is an induced $P_{N+2}$. The omitted cycle vertices are $c_{d-1}$ and
$c_{N-1}$, and the two witness classes differ from those of $x,y$.
This is impossible.

\subsubsection{Single-contact classes at distance two}
Let $x\in A_0$ and $y\in A_2$ be adjacent. The witness $z_0$ differs
from $x$, since $x$ has a cross-class neighbor. If $z_0x$ were absent,
then
\[
 z_{N-2},c_{N-2},c_{N-3},\ldots,c_2,y,x,c_0,z_0
\]
would be induced $P_{N+2}$. Hence $z_0x$ is an edge. Reflection gives
$z_2y$ as an edge. But now
\[
 z_3,c_3,c_4,\ldots,c_{N-1},c_0,x,y,z_2
\]
is induced $P_{N+2}$. The omitted cycle vertices are $c_1,c_2$;
$z_2$ has only its mate $y$ among the exterior path vertices.
Rotation and reflection cover every distance-two pair.

\subsubsection{Single-contact and double-contact classes}
Suppose $x\in A_0$ is adjacent to $b_d$, where $2\le d\le N-3$.
The sequence
\[
\begin{split}
 z_{d-1},c_{d-1},c_{d-2},\ldots,c_0,
 x,b_d,\qquad\qquad\\
 c_{d+1},c_{d+2},\ldots,c_{N-2},z_{N-2}
\end{split}
\]
is induced $P_{N+2}$. Here $c_d,c_{N-1}$ are omitted, so the unused
contact $c_d$ of $b_d$ is absent. Both witness classes differ from $A_0$.

\subsubsection{Two double-contact classes}
For an edge $b_0b_d$, with $3\le d\le N-3$, use
\[
\begin{split}
 z_{d-1},c_{d-1},c_{d-2},\ldots,c_1,
 b_0,b_d,\qquad\qquad\\
 c_{d+1},c_{d+2},\ldots,c_{N-1},z_{N-1}.
\end{split}
\]
The omitted cycle vertices $c_0,c_d$ are precisely the unused contacts
of the two double-contact vertices. This is again induced $P_{N+2}$.

The list exhausts all possible cross-class edges. None can therefore
occur, proving the asserted structure.

\subsection{Converse and the local-independence sum}
Consider a graph constructed as in Theorem~\ref{local16:thm:main}. An $A$
attachment lies in a pendant edge or triangle. A four-cycle containing
$b_i$ would require a second length-two path between $c_i$ and $c_{i+1}$,
and the construction has no such path. The remaining core cycle has
length at least six. Thus the graph is $C_4$-subgraph-free.

Every exterior vertex is simplicial: it has one neighbor or two adjacent
neighbors. It can therefore occur in an induced path only as an endpoint.
An induced path uses at most $N-1$ cycle vertices and two exterior
endpoints, giving $p\le N+1$. Since all $A_i$ are nonempty, an induced
path from $A_0$ through $c_0,c_1,\ldots,c_{N-2}$ to $A_{N-2}$ has
$N+1$ vertices. Hence $p=N+1$.

Every cycle vertex has eccentricity $k+1$, witnessed by the opposite
$A$ class. Every $A$ vertex has eccentricity $k+2$, also witnessed by an
opposite-class vertex. Each $b_i$ has eccentricity $k+1$, witnessed by
either farthest $A$ class. Thus the radius is exactly $k+1$.

Put $a=\sum_i|A_i|$, $q=\sum_i|E(G[A_i])|$, and let $b$ count the
vertices $b_i$. Then
\[
 n=N+a+b,\qquad m=N+a+q+2b,\qquad t=q+b.
\]
Every neighborhood in a $C_4$-subgraph-free graph is a matching plus
isolated vertices. If $t(v)$ counts triangles through $v$, then
$\lambda(v)=d(v)-t(v)$. Summing yields
\[
 \Lambda=2m-3t=2N+2a-q+b=2n-(q+b)=2n-t.
\]
This finishes the proof of Theorem~\ref{local16:thm:main}.



\begin{corollary}\label{local16:cor:wow}
For $k\ge3$, Conjecture 133 holds for every connected graph with a
dominating induced $C_{2k}$, radius $k+1$, and $\ell\le k+1$.
\end{corollary}
\begin{proof}
The four-cycle-present case follows from a geodesic. Otherwise
\eqref{local:chung} settles $\ell\le k$. If $k<\ell\le k+1$,
Theorem~\ref{local16:thm:main} excludes $P_{2k+2}$-freeness, since its
classified graphs have $\ell\le2<k$. Hence
$p\ge2k+2\ge(k+1)+\ell$.
\end{proof}


\subsection{Radius four: the classification used in the corollary}
Write $\operatorname{dist}$, $\operatorname{ecc}$, and $\operatorname{rad}$
for distance, eccentricity, and radius.
\begin{theorem}\label{local15:r4:thm:main}
A connected $C_4$-subgraph-free graph has radius four and is $P_8$-free if
and only if it is $C_8$, or has the following form:
\begin{enumerate}[label=\textup{(\roman*)}]
\item Start with an induced cycle $c_0c_1\dots c_5c_0$.
\item For each $i$, add a nonempty set $A_i$, join every vertex of $A_i$
to $c_i$, and put a matching plus isolated vertices on $A_i$.
\item On each cycle edge $c_ic_{i+1}$ optionally add one vertex $b_i$
adjacent to its endpoints.
\item Add no other edges.
\end{enumerate}
The six-cycle family satisfies $\Lambda(G)=2n-t$. In particular all graphs in
the classification have $\ell(G)\le2$.
\end{theorem}
\subsubsection{Reduction to a dominating cycle}
We use the classical induced-path bound
\begin{equation}\label{local15:r4:eq:radius}
 p(G)\ge2\operatorname{rad}(G)-1,
\end{equation}
recorded by Erd\H{o}s--Saks--S\'os with credit to Fan Chung
\cite[Section II, pp.~64--65]{src:ess}.

\begin{lemma}\label{local15:r4:lem:p6}
A connected $C_4$-subgraph-free $P_6$-free graph of radius three is $C_6$.
\end{lemma}
\begin{proof}
Theorem~6 of van 't Hof--Paulusma \cite{src:hof} supplies a dominating induced
six-cycle or a dominating complete bipartite subgraph. In a
$C_4$-subgraph-free graph the latter is a star, whose center has
 eccentricity at most two.

Fix a dominating induced cycle $c_0\dots c_5c_0$. Every exterior vertex
has at most two cycle neighbors, since three vertices on a six-cycle
include a pair at circular distance two, which would complete a $C_4$.
Two contacts at distance two are likewise forbidden. The remaining
contact sets are singletons, adjacent pairs, and opposite pairs.
A singleton or adjacent pair gives an induced $P_6$: take the exterior
vertex followed by five consecutive cycle vertices, leaving out the
unused contact in the adjacent-pair case. Thus every exterior vertex has
an opposite pair of contacts.

If an exterior vertex has contacts $c_0,c_3$, then $c_0$ is within distance
two of the whole cycle. Every opposite pair contains one of
$c_0,c_1,c_5$, so $c_0$ is also within distance two of every exterior
vertex. This contradicts radius three. Hence there is no exterior vertex.
\end{proof}

\begin{lemma}\label{local15:r4:lem:cycle}
Every connected $C_4$-subgraph-free $P_8$-free graph of radius four is
$C_8$ or has a dominating induced $C_6$.
\end{lemma}
\begin{proof}
Theorem~5 of Camby--Schaudt \cite{src:camby}, in the pinned arXiv version,
supplies a connected dominating induced subgraph $H$ that is $P_6$-free
or is $C_8$. Since $H$ dominates,
$\operatorname{rad}(G)\le\operatorname{rad}(H)+1$. In the $P_6$-free alternative,
\eqref{local15:r4:eq:radius} gives $\operatorname{rad}(H)\le3$, so $\operatorname{rad}(H)=3$ and
Lemma~\ref{local15:r4:lem:p6} yields $H=C_6$.

Suppose instead that $H=C_8$. The cycle-contact set of each exterior
vertex contains a pair at circular distance three or four. Indeed,
distance-two contacts create a $C_4$; without a distance-three/four pair,
the contacts are a singleton or adjacent pair, which yields an induced
$P_8$ by appending seven consecutive cycle vertices.
For each cycle vertex, its three cycle vertices at distance at least
three are consecutive. They cannot contain such a contact set. Thus
every cycle vertex is within distance three of every exterior vertex.
If an exterior vertex has two contacts at distance four, it shortens the
distance between those antipodal vertices to two. If its contacts have
distance three, it gives a path of length three from a contact to that
contact's antipode. In either case one cycle vertex is now within
distance three of the entire cycle and all exterior vertices. This
contradicts radius four. Therefore the $C_8$ alternative has no exterior
vertices.
\end{proof}


\begin{proof}[Completion of the radius-four classification]
Lemma~\ref{local15:r4:lem:cycle} leaves $C_8$ or a dominating induced
$C_6$. The latter is precisely the $k=3$ case of
Theorem~\ref{local16:thm:main}. That theorem supplies both the attachment
restrictions and their converse, including $\Lambda=2n-t$.
The conclusions for $C_8$ follow directly. This proves
Theorem~\ref{local15:r4:thm:main}.
\end{proof}


\subsection{The radius-five classification}
In this subsection, a cycle is \emph{distance-two dominating} if every
vertex of the graph has distance at most two from the cycle. This
distance condition is distinct from the two-neighbor condition in the
parameter $\gamma_2$.
For an even cycle $C=c_0\dots c_{N-1}c_0$, a \emph{one-level decoration}
has the following form. Each cycle vertex $c_i$ receives a nonempty
set $A_i$ of neighbors; the graph on $A_i$ is a matching plus isolates.
For each edge $c_ic_{i+1}$ there may also be one vertex $b_i$ adjacent
to its two endpoints. There are no other vertices or edges. All indices
are read modulo $N$.

A \emph{two-level decoration} is obtained in the same way, allowing some
$A_i$ initially to be empty, and then adding a set of children to any
first-layer vertex in an $A_i$ or of type $b_i$. Each child has exactly
one first-layer neighbor, its parent. The children of each parent induce
a matching plus isolates, and different child classes are anticomplete.
We require, for every $i$, that some vertex of $A_i$ has a child. No
other vertices or edges are permitted.

\begin{theorem}\label{local18r:thm:main}
A connected $C_4$-subgraph-free graph has radius five and is $P_{10}$-free
if and only if it is one of the following:
\begin{enumerate}[label=\textup{(\roman*)}]
\item $C_{10}$;
\item a one-level decoration of $C_8$;
\item a two-level decoration of $C_6$.
\end{enumerate}
Every graph in the list satisfies $p=9$ and $\Lambda=2n-t\le2n$.
\end{theorem}

\begin{corollary}\label{local18r:cor:wow}
Original WOWII133 holds for every connected graph with $r(G)=5$ and
$\ell(G)\le5$.
\end{corollary}

\subsection{Previously established inputs and the structural reduction}
We use the following two classification results, already proved in
Theorems~\ref{local15:r4:thm:main} and~\ref{local16:thm:main}. The remaining argument
classifies all possible second-layer attachments.
\begin{proposition}\label{local18r:prop:inputs}
\begin{enumerate}[label=\textup{(\alph*)}]
\item A connected $C_4$-subgraph-free, $P_8$-free graph of radius four
is either $C_8$ or a one-level decoration of $C_6$.
\item For $k\ge3$, a connected $C_4$-subgraph-free graph with a dominating
induced $C_{2k}$, radius $k+1$, and no induced $P_{2k+2}$ is a one-level
decoration of that cycle. Conversely every such decoration has those
properties, with $p=2k+1$ and $\Lambda=2n-t$.
\end{enumerate}
\end{proposition}

The induced-path radius bound, due to Fan Chung as credited in
Erd\H{o}s--Saks--S\'os \cite[Section II, pp.~64--65]{src:ess}, is
\begin{equation}\label{local18r:eq:chung}
 p(G)\ge2r(G)-1.
\end{equation}
Camby--Schaudt \cite[Theorem 5, arXiv version 1]{src:camby} shows that every
connected $P_{10}$-free graph has a connected dominating induced
subgraph $H$ which is $P_8$-free or is an induced $C_{10}$.
If $r(G)=5$ and $H$ is $P_8$-free, then
\[
 5\le r(H)+1,\qquad 2r(H)-1\le p(H)\le7,
\]
so $r(H)=4$. Proposition~\ref{local18r:prop:inputs}(a) gives a dominating $C_8$
in $G$, or a fixed induced $C_6$ whose distance from every vertex of $G$
is at most two.

If the dominating subgraph is $C_{10}$, there can be no exterior vertex.
Indeed, a singleton or adjacent-pair contact set yields an induced
$P_{10}$ by appending nine consecutive cycle vertices, and contacts at
cycle distance two yield a $C_4$. Hence an exterior vertex has two
contacts at cycle distance at least three. For any cycle vertex $c$,
the cycle vertices at distance at least four from $c$ form three
consecutive vertices. This set cannot contain an admissible contact
set: two contacts in it would be adjacent or at distance two.
Consequently $c$ is within distance four of every exterior vertex.
If an exterior vertex has contacts at cycle distance $d\in\{3,4,5\}$,
one contact reaches its antipode in
$2+(5-d)\le4$ edges through that exterior vertex. This contact is
therefore within distance four of the whole graph, a contradiction.
Thus $G=C_{10}$.

The dominating-$C_8$ branch follows from Proposition~\ref{local18r:prop:inputs}(b)
with $k=4$. It remains to prove the following lemma.
\begin{lemma}\label{local18r:lem:two}
Let $G$ be connected, $C_4$-subgraph-free, $P_{10}$-free, and of radius
five. If an induced $C_6$ is distance-two dominating, then $G$ is a two-level
decoration of that cycle.
\end{lemma}

\subsection{Contacts and clean radius witnesses}
Fix the cycle $C=c_0\dots c_5c_0$ in Lemma~\ref{local18r:lem:two}. Indices in
this proof are modulo six. Three cycle contacts include a pair at
distance two, which produces a $C_4$. Hence nonempty contact sets are
singletons, adjacent pairs, or opposite pairs. A vertex with opposite
contacts is within distance two of the cycle, and hence within distance
four of all of $G$, contrary to the radius. Thus the first layer has
classes $A_i$ with contact set $\{c_i\}$ and classes $B_i$ with contact
set $\{c_i,c_{i+1}\}$. Four-cycle exclusion gives $|B_i|\le1$; write its
vertex, when present, as $b_i$. Let $R$ be the set with no cycle contact.
Every vertex of $R$ has a first-layer neighbor.

Each $A_i$ induces a matching plus isolates. The only possible
first-layer cross edges have types
\begin{equation}\label{local18r:eq:cross}
 A_iA_j\ (j-i\in\{2,3,4\}),\qquad
 A_i b_j\ (j-i\in\{2,3\}),\qquad b_i b_{i+3}.
\end{equation}
All other types immediately give a four-cycle.

Every cycle vertex has eccentricity at most five by distance-two domination,
and at least five by the radius. A distance-five witness for $c_{i+3}$
must therefore belong to $R$, and all its first-layer neighbors must
belong to $A_i$: any other contact gives a route of length at most four.
Two neighbors in $A_i$ would give a four-cycle through $c_i$. We may
therefore choose six arms
\begin{equation}\label{local18r:eq:arms}
 c_i-a_i-z_i,\quad a_i\in A_i,\ z_i\in R,\quad
 N(z_i)\cap\text{(first layer)}=\{a_i\},\quad
 \operatorname{dist}(c_{i+3},z_i)=5.
\end{equation}
Every $a_i$ has no first-layer neighbor outside $A_i$. Indeed, any
permitted cross edge in \eqref{local18r:eq:cross} shortens the prescribed distance
in \eqref{local18r:eq:arms} to at most four. Within $A_i$, the vertex $a_i$ has
at most one mate.

We next prove that each $z_i$ is simplicial, of degree one or two.
Rotate to $i=0$, and suppose $z_0w$ is an edge with $w\in R$. A
first-layer neighbor $x$ of $w$ can only have type
$A_0,A_1,A_5,B_0,B_5$, since all other types put $z_0$ within distance
four of $c_3$. The following alternatives give induced ten-vertex paths:
\begin{align*}
 x\in A_0\setminus\{a_0\}:&\quad
 z_2,a_2,c_2,c_3,c_4,c_5,c_0,x,w,z_0;\\
 x\in A_1:&\quad
 z_0,w,x,c_1,c_0,c_5,c_4,c_3,a_3,z_3;\\
 x=b_0:&\quad
 z_0,w,b_0,c_0,c_5,c_4,c_3,c_2,a_2,z_2.
\end{align*}
Reflect the latter two for $A_5,B_5$. In the first row, the possible
chords $wa_2,wz_2,z_0z_2$ shorten one of the prescribed witness distances
to at most four. The same argument excludes $wa_3,wz_3,z_0z_3$ in the
second row and $wa_2,wz_2,z_0z_2$ in the third. Exact cycle contacts,
the absence of cross edges at clean parents $a_i$, and the unique
first-layer neighbor of $z_0$ exclude all other chords. The unused
contact $c_1$ of $b_0$ is omitted in the third row.

Thus every first-layer neighbor of $w$ is $a_0$. At least one exists,
so $wa_0$ is an edge. Two such vertices $w$ would have the two common
neighbors $a_0,z_0$, giving a four-cycle. This proves the assertion for
$z_0$, and by rotation for all $z_i$.

\subsection{Eliminating first-layer cross edges}
Every displayed path in this section has ten vertices. Clean parents
have no cross-class neighbors, clean terminals are simplicial, and the
cycle contacts are exact; these properties exclude the chords.
An edge $xy$ with $x\in A_0,y\in A_3$ gives
\[
 z_1,a_1,c_1,c_0,x,y,c_3,c_4,a_4,z_4.
\]
An edge from $x\in A_0$ to $b_2$ or $b_3$ gives, respectively,
\begin{align*}
 &z_1,a_1,c_1,c_0,x,b_2,c_3,c_4,a_4,z_4;\\
 &z_5,a_5,c_5,c_0,x,b_3,c_3,c_2,a_2,z_2.
\end{align*}
An edge $b_0b_3$ gives
\[
 z_2,a_2,c_2,c_1,b_0,b_3,c_4,c_5,a_5,z_5.
\]
Finally, suppose $x\in A_0,y\in A_2$ are adjacent. Necessarily
$x\ne a_0$, since $a_0$ has no cross edge. If $a_0x$ is absent, then
\[
 z_4,a_4,c_4,c_3,c_2,y,x,c_0,a_0,z_0
\]
is induced. Hence $a_0x$ is present, and reflection gives $a_2y$.
But now
\[
 z_3,a_3,c_3,c_4,c_5,c_0,x,y,a_2,z_2
\]
is induced. Rotations and reversals exhaust \eqref{local18r:eq:cross}.
Therefore the first layer contains only the individual $A_i$ matchings;
every $b_i$ has no first-layer neighbor.

\subsection{Unique parenthood in the second layer}
Suppose $w\in R$ has at least two first-layer neighbors. It has at most
one neighbor in each $A_i$, by four-cycle exclusion through $c_i$.
It is nonadjacent to every $z_i$, because a neighbor of $z_i$ in $R$
has only $a_i$ as its first-layer neighbor. We repeatedly use
\begin{equation}\label{local18r:eq:guard}
 x\in A_i,\ wx\in E(G)\quad\Longrightarrow\quad wa_{i+3}\notin E(G).
\end{equation}
Otherwise $c_i,x,w,a_{i+3},z_{i+3}$ contradicts the prescribed distance
five in \eqref{local18r:eq:arms}.

\subsubsection{Two neighbors of type $A$}
First suppose $x\in A_0,y\in A_2$ are neighbors of $w$. According to
the relation between $x$ and $a_0$, use one of the following paths:
\begin{align*}
 x=a_0:&\quad z_0,x,w,y,c_2,c_3,c_4,c_5,a_5,z_5;\\
 x\ne a_0,\ a_0x\in E:&\quad a_0,x,w,y,c_2,c_3,c_4,c_5,a_5,z_5;\\
 x\ne a_0,\ a_0x\notin E:&\quad
 z_0,a_0,c_0,x,w,y,c_2,c_3,a_3,z_3.
\end{align*}
In the latter two cases $wa_0$ would create a four-cycle through
$c_0,x$. The endpoint-parent chord $wa_5$ or $wa_3$ is excluded by
\eqref{local18r:eq:guard}. All first-layer cross edges have already been
eliminated, so these are induced $P_{10}$s. Thus two $A$ contacts cannot
have indices at distance two.

For $x\in A_0,y\in A_1$, if $wa_2$ is absent, use
\[
 z_2,a_2,c_2,c_3,c_4,c_5,c_0,x,w,y.
\]
If $wa_2$ is present, the distance-two case applies to $x,a_2$.
Thus adjacent $A$ indices are also impossible.
For $x\in A_0,y\in A_3$, the preceding cases exclude every $A$ neighbor
of $w$ in the other four classes, and
\[
 a_1,c_1,c_0,x,w,y,c_3,c_4,a_4,z_4
\]
is induced. Hence $w$ has at most one neighbor of type $A$.

\subsubsection{One neighbor of each type}
Let $x\in A_0$ be a neighbor of $w$. Types $b_0,b_5$ give a four-cycle
through $c_0$. Every clean parent outside $A_0$ is a nonneighbor of $w$
by the preceding exclusion of two $A$ neighbors.
For a second neighbor $b_2$, use
\[
 a_1,c_1,c_0,x,w,b_2,c_3,c_4,a_4,z_4;
\]
the $b_3$ case is its reflection. For $b_1$, use
\begin{align*}
 x=a_0:&\quad a_2,c_2,b_1,w,x,c_0,c_5,c_4,a_4,z_4;\\
 x\ne a_0,\ a_0x\notin E:&\quad
 a_0,c_0,x,w,b_1,c_2,c_3,c_4,a_4,z_4;\\
 x\ne a_0,\ a_0x\in E:&\quad
 a_0,x,w,b_1,c_2,c_3,c_4,c_5,a_5,z_5.
\end{align*}
Reflect for $b_4$. Whenever needed, $wa_0$ is excluded by a four-cycle.
The unused contact of the displayed $B$ vertex is omitted, and all
other chord exclusions were established above. Thus no mixed pair is
possible.

\subsubsection{Two neighbors of type $B$}
In this case $w$ has no $A$ neighbor, by the preceding subsection.
Two $B$ types sharing a cycle contact yield a four-cycle. Up to rotation
and reflection the remaining pairs are $b_0,b_2$ and $b_0,b_3$, giving
respectively
\begin{align*}
 &a_1,c_1,b_0,w,b_2,c_3,c_4,c_5,a_5,z_5;\\
 &a_2,c_2,c_1,b_0,w,b_3,c_4,c_5,a_5,z_5.
\end{align*}
Both are induced $P_{10}$s. Every vertex of $R$ therefore has exactly
one first-layer neighbor, called its parent.

\subsection{Child classes, converse, and counting}
Suppose $uv$ is an edge in $R$ with distinct parents $x,y$.
If one parent is of type $A$, orient so $x\in A_i$ and choose
$s\in\{i+2,i-2\}$ with $y\ne a_s$. Take the five-vertex cycle arc
from $c_s$ to $c_i$ that omits the other neighbor of $c_i$. Then
\[
 z_s,a_s,\ [\text{five cycle vertices from }c_s\text{ to }c_i],\ x,u,v
\]
is an induced $P_{10}$. Distinct parents exclude $xv$, unique parenthood
excludes the other first-layer chords, and the choice of $s$ together
with the simplicial terminal property excludes chords to the clean arm.
If both parents have type $B$, rotate so $x=b_0$ and use
\[
 z_2,a_2,c_2,c_3,c_4,c_5,c_0,b_0,u,v.
\]
The unused contact $c_1$ is omitted. Hence distinct child classes are
anticomplete. Each child class has maximum degree one by four-cycle
exclusion in its parent's neighborhood. The arms in \eqref{local18r:eq:arms}
ensure the required nonempty depth-two class at each cycle vertex.
This proves Lemma~\ref{local18r:lem:two}.

Conversely, consider any stated two-level decoration of $C_6$.
The $A$ attachments and child matchings introduce only pendant edges or
triangles. A four-cycle through $b_i$ would require a second length-two
path between $c_i,c_{i+1}$, which the construction does not contain.
Thus the graph is $C_4$-subgraph-free.
An induced path meeting the core meets it in one contiguous induced
cycle path, of at most five vertices, with at most two exterior vertices
at each endpoint. Traversing a matching edge and then the common parent
would produce a triangular chord. A path missing the core lies in one
attachment component and has at most four vertices. Therefore $p\le9$,
and depth-two vertices at the ends of a five-vertex cycle arc give
$p=9$.

Cycle, $A$, and $B$ vertices have eccentricities five, six, and five,
respectively. A child of an $A$ vertex has eccentricity seven, while a
child of a $B$ vertex has eccentricity six. Opposite or farthest
attachment classes supply these distances, and matching edges do not
shorten routes between classes. In particular the radius is five.
If $a,b,s$ count the $A$ vertices, the $B$ vertices, and their children,
and $q,h$ count the matching edges on the $A$ and child classes, then
\[
 n=6+a+b+s,\quad m=6+a+2b+s+q+h,\quad t=b+q+h.
\]
Equation~\eqref{local:potential} gives $\Lambda=2n-t$. The $C_{10}$ and decorated-$C_8$
cases have the same conclusions, by direct calculation and
Proposition~\ref{local18r:prop:inputs}(b). The structural reduction now proves
Theorem~\ref{local18r:thm:main}.

\begin{proof}[Proof of Corollary~\ref{local18r:cor:wow}]
If $G$ contains a four-cycle, a geodesic supplies $p\ge r+1=6$, as
required by the original exponent-zero case. Otherwise
\eqref{local18r:eq:chung} gives $p\ge9$, proving $p\ge5+\ell$ for $\ell\le4$.
If $4<\ell\le5$, Theorem~\ref{local18r:thm:main} rules out $P_{10}$-freeness,
since every classified graph has $\ell\le2$. Hence
$p\ge10\ge5+\ell$, proving the remaining range.
\end{proof}




\paragraph{Antecedents.}
Equal-length pendant-arm cores are classical ciliates. Relevant antecedents
include Fajtlowicz's characterization of radius-critical graphs and the
results on induced ciliates and radius equality discussed by DeLaVi\~na,
Larson, Pepper, and Waller \cite{src:fajtlowicz,src:dlpw}.
The attachment proofs above use the explicitly stated connected-domination
theorems.


\subsection{The unrounded average range up to four}
\begin{corollary}\label{local15:cor:four}
Conjecture 133 holds for every connected graph with $\ell(G)\le4$.
\end{corollary}
\begin{lemma}\label{local15:lem:small}
Every connected $C_4$-subgraph-free graph with $\ell\le3$ satisfies
$p\ge r+\ell$.
\end{lemma}
\begin{proof}
First, a connected $P_5$-free, $C_4$-subgraph-free graph has $\Lambda\le2n$.
If it has an induced $C_5$, any exterior neighbor of that cycle has a
singleton or adjacent-pair contact set, producing $P_5$ with four
consecutive cycle vertices. Connectedness therefore implies that the
whole graph is $C_5$. Without $C_5$, the graph is chordal, since every
longer hole contains $P_5$; the simplicial deletion argument above gives
$\Lambda\le2(n-1)$.

Also every connected $P_4$-free, $C_4$-subgraph-free graph has a universal
vertex. If it is not complete, choose an induced $u,w,v$ path.
Any vertex $x$ nonadjacent to $w$ cannot meet $u$ or $v$: such a contact
would give $P_4$ or $C_4$ on these four vertices. A connected $P_4$-free
graph has diameter at most two, so take a path $x,y,w$. The vertex $y$
cannot meet both $u,v$, by four-cycle exclusion. Choosing the one it
does not meet gives an induced $P_4$ starting $x,y,w$, a contradiction.
Thus $w$ is universal.

For radius one, a universal vertex has a matching plus isolates on
its neighborhood. If that matching has $q$ edges, then
$n=1+d$, $m=d+q$, $t=q$, and $\Lambda=2d-q<2n$.
A noncomplete connected graph has $p\ge3>1+\ell$ in this case.
A complete graph has $p=2$, $r=\ell=1$, except for the trivial
one-vertex case, which is immediate.

At radius two, $p\ge4$ by the universal-vertex observation. This is
sufficient for $\ell\le2$. If $2<\ell\le3$, the $P_5$-free density
bound forces $p\ge5\ge2+\ell$. At radius three, \eqref{local:chung}
gives $p\ge5$, sufficient for $\ell\le2$. For $2<\ell\le3$, Lemma~\ref{local15:r4:lem:p6} excludes $P_6$-freeness: its only radius-three possibility
is $C_6$, with $\ell=2$. Hence $p\ge6\ge3+\ell$. Finally, for
$r\ge4$, \eqref{local:chung} gives $p\ge2r-1\ge r+3\ge r+\ell$.
\end{proof}

\begin{proof}[Proof of Corollary~\ref{local15:cor:four}]
If a four-cycle occurs, \eqref{local:target} asks only for $p\ge r+1$, which
follows from a geodesic. Assume henceforth that $G$ is $C_4$-subgraph-free.
The range $\ell\le3$ is Lemma~\ref{local15:lem:small}. Suppose $3<\ell\le4$.
Radius one is impossible by the calculation in that lemma. For $r=2$
or $3$, Theorem~\ref{local15:thm:general} forces $p\ge7\ge r+\ell$.
For $r=4$, Theorem~\ref{local15:r4:thm:main} has $\ell\le2$,
so $p\ge8\ge4+\ell$. For $r\ge5$, \eqref{local:chung} gives
$p\ge2r-1\ge r+4\ge r+\ell$.
\end{proof}



\section{Closed core certificates and the final range for Conjecture 133}
\label{sec:local-certificates}
The finite objects in this section are proof certificates: every permitted
extension is accounted for, and the extension lemmas give them force at
arbitrary graph order. Explicit paths and adjacency-preserving permutations
are checked directly. The certificate files and independent verifiers are
included in the accompanying reproducibility package.
\subsection{Four cores without an induced eight-vertex path}


Let $H$ be the graph obtained from the Heawood graph by deleting one
vertex and joining its three neighbors pairwise. This is the classical
$Y$--$\Delta$ construction; neither the operation nor the graph is
claimed as new. It is the graph labeled $C_{13}$ in the Heawood family
(not the thirteen-cycle); see \cite[Figure 1.3]{src:heawoodfamily}, where its
$\Delta Y$ successor $C_{14}$ is identified as the Heawood graph.
We use the following adjacency table, in which the
unique triangle is $T=\{0,1,2\}$:
\begin{center}
\begin{tabular}{cl@{\qquad}cl}
\toprule
Vertex&Neighbors&Vertex&Neighbors\\\midrule
$0$&$1,2,3,9$&$7$&$3,4,5$\\
$1$&$0,2,4,12$&$8$&$5,9,12$\\
$2$&$0,1,5,10$&$9$&$0,8,11$\\
$3$&$0,6,7$&$10$&$2,6,11$\\
$4$&$1,7,11$&$11$&$4,9,10$\\
$5$&$2,7,8$&$12$&$1,6,8$\\
$6$&$3,10,12$&&\\\bottomrule
\end{tabular}
\end{center}
Its graph6 representation is \texttt{L\{O\_\_wC\_HC@BOg}.
It has $(n,m,t)=(13,21,1)$ and $\lambda(v)=3$ at every vertex.
The certificate verification also checks that it is $C_4$-subgraph-free
and $P_8$-free.

\begin{theorem}\label{local18p:thm:classification}
Let $G$ be connected, $C_4$-subgraph-free, and $P_8$-free. If
$\lambda_G(v)\ge3$ at every vertex, then $G$ is exactly one of:
\begin{enumerate}[label=\textup{(\roman*)}]
\item the Petersen graph;
\item the Heawood graph;
\item the Heawood graph with one edge contracted;
\item $H$.
\end{enumerate}
The triangle-free cases \textup{(i)--(iii)} are the theorem of
Chudnovsky--Stacho \cite[Theorem 1]{src:cs}; the proof below establishes
the triangle-containing case by a closed finite certificate and an
arbitrary-order extension argument.
\end{theorem}
\begin{corollary}\label{local18p:cor:density}
Every nonempty $C_4$-subgraph-free $P_8$-free graph satisfies $\Lambda(G)<4n$.
\end{corollary}

\subsection{A complete one-vertex extension rule}
\begin{lemma}\label{local18p:lem:extension}
Let $F$ be an induced subgraph of a $C_4$-subgraph-free graph $G$ with
$\lambda_G(v)\ge3$ at every vertex. If $u\in V(F)$ satisfies
$\lambda_F(u)<3$, then some $x\in V(G)\setminus V(F)$ is adjacent to $u$
and anticomplete to $N_F(u)$.
\end{lemma}
\begin{proof}
Both neighborhoods of $u$ induce matchings plus isolates. Suppose every
new neighbor of $u$ had a neighbor in $N_F(u)$. A new neighbor cannot
meet two old neighbors, and an old neighbor cannot meet two new ones.
Nor can any new neighbor paired with an old neighbor have an additional
neighbor inside $N_G(u)$. Thus every new neighbor merely pairs with an
old isolated neighbor. Replacing an isolate by a matched pair leaves
the independence number unchanged. This contradicts
$\lambda_G(u)\ge3>\lambda_F(u)$.
\end{proof}

The contact set $S=N_G(x)\cap V(F)$ of this new vertex satisfies
\begin{equation}\label{local18p:eq:contacts}
 u\in S,\qquad S\cap N_F(u)=\varnothing,\qquad
 N_F(v)\cap N_F(w)=\varnothing\quad(v,w\in S,\ v\ne w).
\end{equation}
Conversely, the last condition exactly excludes a four-cycle through
$x$ in the extension of $F$. Therefore the required contact sets are
precisely the cliques containing $u$ in the compatibility graph whose
edges join pairs with no common neighbor, subject to the middle
restriction in \eqref{local18p:eq:contacts}. Any deficient vertex $u$ may be
chosen; the choice need not be invariant under graph isomorphism.

\subsection{The closed finite certificate}
The supplemental file \texttt{p8\_core\_certificate.json.gz} is a finite
directed certificate with the labeled triangle as its root. Each state
specifies a graph by integer adjacency bitmasks. A nonterminal state
specifies a vertex $u$ with $\lambda_F(u)<3$ and contains one row for
\emph{every} contact set satisfying \eqref{local18p:eq:contacts}. Each extension
$F+x$ is accompanied by either:
\begin{enumerate}[label=\textup{(\alph*)}]
\item eight distinct vertices inducing $P_8$; or
\item a target state of order $|V(F)|+1$, with an explicit
adjacency-preserving permutation from $F+x$ to that state.
\end{enumerate}
Every terminal state has local independence at least three at every
vertex and comes with an explicit isomorphism to the displayed $H$.

\begin{center}
\begin{tabular}{lr}
\toprule
Certificate quantity&Count\\\midrule
States&192\\
Maximum state order&25\\
Required one-vertex extensions&93,895\\
Explicit induced-$P_8$ rejections&93,608\\
Explicit retained-extension isomorphisms&287\\
Terminal states&1\\
Unaccounted extensions or cutoff frontier&0\\\bottomrule
\end{tabular}
\end{center}

The maximum order 25 is an observed property of this \emph{closed}
certificate, not an assumption on the order of a target graph. The
builder uses no order cutoff. Canonical labeling is used only to build
a compact certificate; it is not trusted by the verifier.

The separate program \texttt{verify\_p8\_core\_certificate.py} uses only
the Python standard library. It independently enumerates all contact
sets in \eqref{local18p:eq:contacts}, requires exact equality with the listed
rows, checks every induced-path witness on all vertex pairs, and checks
every retained-extension permutation on all vertex pairs. It also
checks the state graphs, the root, reachability, increasing order along
retained arcs, and the terminal isomorphism. No graph-isomorphism or
canonical-labeling package is a verification dependency.
The compressed certificate has SHA256
\begin{center}\small\ttfamily
 d7be0e33ac427e6ed0aabeef4e69efd50be0ebd7a05063a30bee4aa42fe31982.
\end{center}

\begin{lemma}\label{local18p:lem:contains}
Every $C_4$-subgraph-free, $P_8$-free graph containing a triangle and
having $\lambda_G(v)\ge3$ everywhere contains an induced $H$.
\end{lemma}
\begin{proof}[Proof from the certificate]
Embed the root triangle in $G$. If an embedded state $F$ is nonterminal,
apply Lemma~\ref{local18p:lem:extension} at its designated vertex. The resulting
new vertex has a contact set present in the certificate. Its row cannot
be an induced-$P_8$ rejection, since that path would also be induced in
$G$. The retained-extension permutation therefore gives an induced
embedding of the specified next state in $G$. This increases the order
by one. A finite certificate with all nonterminal extensions accounted
for cannot support an infinite sequence of increasing orders. Thus a
terminal state is reached, and its isomorphism gives an induced $H$.
No bound on $|V(G)|$ is assumed.
\end{proof}

\subsection{Attachment restrictions at the terminal graph}
A second, small finite certificate is included in
\texttt{weighted\_core\_contact\_certificates.json}. The compatibility
graph on the 13 vertices of $H$ has 30 edges and no triangle; the latter
is checked on all $\binom{13}{3}=286$ triples. Consequently a nonempty
contact set of a new vertex in a $C_4$-subgraph-free extension has size
at most two. Of its $13+30=43$ possibilities, 40 have an explicit
induced-$P_8$ witness. The only survivors are
$\{0\},\{1\},\{2\}$; the three single-leaf extensions are also checked
to be $P_8$-free. The same standard-library verifier checks all these
facts and the rooted paths used below.

\begin{lemma}\label{local18p:lem:attachments}
Every connected $C_4$-subgraph-free $P_8$-free graph containing an induced
$H$ consists of $H$ together with, at each vertex of $T=\{0,1,2\}$,
a possibly empty matching plus isolates whose vertices are all adjacent
to that vertex of $T$. There are no other vertices or edges.
Conversely every graph of this form is $C_4$-subgraph-free and $P_8$-free.
\end{lemma}
\begin{proof}
The finite contact certificate shows that every exterior vertex with a
neighbor in $H$ has exactly one such neighbor, belonging to $T$.
If $y$ had distance two from $H$, choose a path $y-x-i$ with $i\in T$.
The following induced six-vertex paths start at the respective roots:
\[
\begin{array}{c|l}
i&\text{path in }H\\\hline
0&0,1,4,7,5,8\\
1&1,0,3,6,10,11\\
2&2,0,3,6,12,8
\end{array}
\]
Prepending $y,x$ gives an induced $P_8$, since $y$ has no $H$ contact
and $x$ has only $i$. Thus there are no distance-two vertices, and
connectedness excludes further layers.
Let $A_i$ be the exterior vertices with contact $i\in T$.
An edge between $A_i,A_j$ for $i\ne j$ forms a four-cycle through $ij$.
A three-vertex path $x-y-z$ within $A_i$ forms the four-cycle
$i-x-y-z-i$. Hence each $A_i$ is a matching plus isolates and the
classes are anticomplete.

Conversely, the attachments add only pendant edges or triangles, so
there is no four-cycle. An induced path meeting $H$ uses at most one
attachment vertex at either end: traversing a matching edge and then
the common root creates a chord. If both exterior endpoints have
distinct roots, those roots are adjacent, so their segment in $H$ has
only that edge. If exactly one endpoint is exterior, an induced $P_8$
would give a rooted seven-vertex path in $H$, contrary to the checked
single-leaf extension. Paths entirely outside $H$ have at most two
vertices. Thus the graph is $P_8$-free.
\end{proof}

Every added vertex in a proper extension from
Lemma~\ref{local18p:lem:attachments} has local independence one. Hence a graph
with local independence at least three everywhere cannot properly
contain $H$. This proves the triangle-containing case of
Theorem~\ref{local18p:thm:classification}, using Lemma~\ref{local18p:lem:contains}.
In the triangle-free case $\lambda(v)=d(v)$, and the remaining three
possibilities follow from Chudnovsky--Stacho \cite[Theorem 1]{src:cs}.



\begin{proof}[Proof of Corollary~\ref{local18p:cor:density}]
Induct on the order, allowing disconnected graphs. If some vertex has
$\lambda(v)\le2$, equation~\eqref{local:potential} gives
\[
 \Lambda(G)-\Lambda(G-v)=2\lambda(v)-t(v)\le4.
\]
The induction hypothesis proves the strict inequality; the one-vertex
base case is immediate. Otherwise every connected component is one of
the four graphs in Theorem~\ref{local18p:thm:classification}. Their pairs $(n,\Lambda)$
are, respectively,
\[
 (10,30),\qquad(14,42),\qquad(13,40),\qquad(13,39).
\]
Each satisfies $\Lambda<4n$, and summing over components completes the proof.
\end{proof}



\subsection{Eleven cores without an induced nine-vertex path}


Table~\ref{local19:tab:models} specifies eleven graphs. The final column counts
induced five-cycles. All tuples $(n,m,t,c_5)$ are distinct, and therefore
the models are pairwise nonisomorphic.
For completeness, their graph6 strings use the following convention:
the first character has code $n+63$; the remaining characters, after
subtracting 63 and writing each as six binary digits, encode the
adjacency bits in the order
$(0,1),(0,2),(1,2),(0,3),(1,3),(2,3),\ldots$, with trailing zero padding.
The certificate also stores explicit adjacency arrays.

\begin{table}[ht]
\centering\small
\begin{tabular}{clrrrrr}
\toprule
Model&graph6&$n$&$m$&$t$&$\Lambda$&$c_5$\\\midrule
$M_{0}$&\texttt{L\{OcA?\_F?b@S@o}&13&21&1&39&12\\
$M_{1}$&\texttt{M\{OcA?\_F?b@S@a@P?}&14&24&1&45&21\\
$M_{2}$&\texttt{O\{OcA?\_F?a@P@\_@OaDOPI}&16&30&4&48&24\\
$M_{3}$&\texttt{N\{OcA?\_E?g@\_@sAiAK\_}&15&27&1&51&30\\
$M_{4}$&\texttt{IheA@GUAo}&10&15&0&30&12\\
$M_{5}$&\texttt{MheA@GUA\_QCA?a@@\_}&14&23&0&46&17\\
$M_{6}$&\texttt{LheA@GUA\_QCA?J}&13&21&0&42&18\\
$M_{7}$&\texttt{KheA@GQA?K\_M}&12&18&0&36&9\\
$M_{8}$&\texttt{KheA@GQA?H\_Y}&12&18&0&36&8\\
$M_{9}$&\texttt{LheA@GQA?H?YOE}&13&20&0&40&8\\
$M_{10}$&\texttt{MhEKA?\_C\_Q?c?i?T?}&14&21&0&42&0\\
\bottomrule
\end{tabular}
\caption{The eleven explicitly specified models.}\label{local19:tab:models}
\end{table}

\begin{theorem}\label{local19:thm:main}
A connected $C_4$-subgraph-free, $P_9$-free graph has
$\lambda_G(v)\ge3$ at every vertex if and only if it is isomorphic to one
of the eleven models in Table~\ref{local19:tab:models}.
\end{theorem}
\begin{corollary}\label{local19:cor:density}
Every nonempty $C_4$-subgraph-free $P_9$-free graph satisfies $\Lambda(G)<4n$.
\end{corollary}
The classification proof uses the closed finite certificate described
below as a proof dependency, not merely as supplementary evidence.

\subsection{Two exhaustive extension rules}
The deficient-vertex extension rule is
Lemma~\ref{local18p:lem:extension}, whose proof uses only four-cycle exclusion.
The new vertex's contact set $S=N_G(x)\cap V(F)$ therefore satisfies
\begin{equation}\label{local19:eq:forcedcontacts}
 u\in S,\qquad S\cap N_F(u)=\varnothing,\qquad
 N_F(v)\cap N_F(w)=\varnothing\quad(v,w\in S,\ v\ne w).
\end{equation}
The last condition is necessary and sufficient to exclude a new
four-cycle. If the target $G$ is triangle-free, we additionally require
$S$ to be independent in $F$.

There is a separate rule when $F$ already has local independence at
least three. If $F$ is proper in a connected target $G$, connectedness
supplies an exterior vertex with a \emph{nonempty} contact set in $F$.
We must then consider every nonempty set satisfying the four-cycle
condition, and the independence condition in triangle-free mode.
There is no designated deficient vertex or neighborhood-anticompleteness
condition in this second rule. This is why an accepting finite core
cannot simply terminate the supergraph analysis.

\subsection{Initial states and contact partitions}
\subsubsection{Six roots}
A target containing a triangle starts with an induced $C_3$.
A triangle-free target has minimum degree at least three, so it contains
a cycle. A shortest cycle is induced, has length at least five, and
has length at most nine: an induced cycle on at least ten vertices
contains an induced $P_9$. Consequently the six roots are $C_3$ in
unrestricted mode, and $C_5,C_6,C_7,C_8,C_9$ in triangle-free mode.
The latter mode preserves triangle-freeness by requiring independent
contact sets.

\subsubsection{A region of possible contacts}
For a state $F$, define
\[
 \operatorname{Compat}(v)=\{w\ne v:N_F(v)\cap N_F(w)=\varnothing\}.
\]
In triangle-free mode, remove neighbors of $v$ from this set as well.
A node of a contact-partition tree represents disjoint sets $(S,R)$:
$S$ contains the contacts already selected, and $R$ contains undecided
vertices still available. Selected contacts are pairwise compatible,
and every vertex of $R$ is compatible with $S$.

At a deficient state, with chosen vertex $u$, start with
\[
 S=\{u\},\qquad R=\operatorname{Compat}(u)\setminus(N_F(u)\cup\{u\}).
\]
At a state without a deficient vertex, start with
$S=\varnothing$, $R=V(F)$; the empty final contact set will be discarded.

A binary split chooses any $v\in R$ and has exactly the children
\begin{align*}
 \text{include: }&(S\cup\{v\},\ (R\setminus\{v\})\cap\operatorname{Compat}(v)),\\
 \text{exclude: }&(S,\ R\setminus\{v\}).
\end{align*}
These two regions partition all admissible completions. Incompatible
contacts in the include branch are forced absent, not omitted without
justification.

\subsubsection{Three leaf types}
There are precisely three permitted leaf types.
\begin{enumerate}[label=\textup{(\roman*)}]
\item \textbf{Forbidden-path box.} An explicit induced $P_9$ is supplied
in the extension $F+x$ with current contact set $S$. If $x$ lies on the
path, no vertex of $R$ lies on it. Thus no further contact can add a
chord to that path: additional edges at $x$ only reach vertices outside
the displayed path. If $x$ is absent, the path is entirely in $F$ and
cannot change. The entire region is therefore forbidden.
\item \textbf{Retained extension.} Here $R=\varnothing$ and
$S\ne\varnothing$. An explicit adjacency-preserving permutation maps
$F+x$ to a target state of order $|V(F)|+1$ in the same mode.
\item \textbf{Empty contacts.} Here $S=R=\varnothing$, at a state with
no deficient vertex. This is not a connected extension and is discarded.
\end{enumerate}
Finding a path in a partial extension is not, by itself, enough to
reject that region. The explicit no-repair condition in leaf type
(i) is essential and is checked by the verifier.

\subsection{The finite certificate and its arbitrary-order force}
Each graph state either specifies a vertex of local independence at
most two or is an \emph{accepting state}. An accepting state has local
independence at least three everywhere and includes an explicit
isomorphism to one of the eleven models. It also includes a full
contact-partition tree covering every possible proper connected
extension. Thus accepting states are not unexamined endpoints.

The closed certificate has the following counts:
\begin{center}
\begin{tabular}{lr}
\toprule
Quantity&Count\\\midrule
Graph states&13,930\\
Maximum state order&45\\
Initial roots&6\\
Accepting model states&11\\
Contact-tree nodes&442,204\\
Binary split nodes&214,137\\
Forbidden-$P_9$ boxes&207,299\\
Retained-extension permutations&20,757\\
Empty-contact leaves&11\\
Unaccounted branches or cutoff frontier&0\\\bottomrule
\end{tabular}
\end{center}
The builder has no graph-order or state cutoff. The maximum order 45
is an observed feature of the closed certificate, not a bound assumed
for target graphs. Every canonical merge used in construction is
accompanied by an exact isomorphism; no canonical-labeling package is
trusted by verification.

The file \texttt{p9\_core\_certificate.json.gz} has SHA256
\begin{center}\small\ttfamily
9d76bd6ba46c60a43175f792e0e67517c4d19b22b068917eeee5e34327544e1b.
\end{center}
The standard-library-only verifier checks the state graphs, deficient
vertices, accepting-model maps, exact include/exclude regions,
forbidden paths and their no-repair conditions, retained-extension
permutations, roots, modes, and reachability. It directly verifies each
model's connectivity, four-cycle exclusion, $P_9$-freeness, local
independence, potential $\Lambda$, and induced-five-cycle count.
A separate table check verifies the graph6 transcriptions against the
certificate's adjacency arrays.

\begin{proof}[Proof of Theorem~\ref{local19:thm:main} from the certificate]
Choose a root cycle in $G$ as above and embed its state as an induced
subgraph. Suppose a state $F$ has been embedded. If $F=G$, then $F$
has no deficient vertex, so its accepting-model map identifies $G$
with a listed model.

Otherwise, if $F$ is deficient, Lemma~\ref{local18p:lem:extension} supplies a new
vertex with a legal contact set. If $F$ is accepting, connectedness
supplies an exterior neighbor of $F$ with a nonempty legal contact set.
Follow this exact set through the binary contact partition. A
forbidden-path box is impossible: its displayed path persists under
those contacts and would be induced in $G$. An empty-contact leaf is
also impossible. Therefore a retained leaf gives an induced embedding
of a state with one more vertex.

This process cannot continue indefinitely, because the certificate has
finitely many states and all retained arcs increase order. Since all
branches are accounted for, the process must reach $F=G$ at an accepting
state. Thus $G$ is isomorphic to a listed model, with no restriction on
its initial order. Conversely, all required properties of each model
are verified directly. Distinct $(n,m,t,c_5)$ tuples establish pairwise
nonisomorphism.
\end{proof}

\subsection{Density and the original conjecture}
\begin{proof}[Proof of Corollary~\ref{local19:cor:density}]
Induct on the order, allowing disconnected graphs. If some vertex has
$\lambda(v)\le2$, equation~\eqref{local:potential} gives
\[
 \Lambda(G)-\Lambda(G-v)=2\lambda(v)-t(v)\le4.
\]
For $n>1$ the induction hypothesis gives the strict bound $\Lambda(G)<4n$;
the one-vertex case is immediate. Otherwise every connected component
is one of the eleven models, each satisfying $\Lambda<4n$ by
Table~\ref{local19:tab:models}. Summing over components completes the proof.
\end{proof}



\subsection{Final consolidated consequence for the original conjecture}
\begin{theorem}\label{local:final}
The original, unrounded Conjecture 133 holds for every connected graph
with $\ell(G)\le5$.
It also holds in the additional ranges $r=2$, $\ell\le7$, and
$r=3$, $\ell\le6$.
\end{theorem}
\begin{proof}
A graph containing a four-cycle is covered by the geodesic bound.
Assume $C_4$-freeness. Corollary~\ref{local15:cor:four} settles
$\ell\le4$. For $4<\ell\le5$, Corollary~\ref{local19:cor:density}
forces $p\ge9$, sufficient when $r\le4$. At radius five,
Theorem~\ref{local18r:thm:main} excludes $P_{10}$-freeness because all
its graphs have $\ell\le2$. Thus $p\ge10\ge5+\ell$.
For $r\ge6$, inequality~\eqref{local:chung} gives
$p\ge2r-1\ge r+5\ge r+\ell$.
More generally, $\ell>4$ forces $p\ge9$, which also covers the stated
radius-two and radius-three extensions. Their lower-average endpoints
are already included in Corollary~\ref{local15:cor:four}.
\end{proof}
The unrestricted inequality remains outside the conclusions of this
paper. The eleven-core proof does not need an external graph
classification theorem. The preceding four-core result credits the
triangle-free theorem of Chudnovsky--Stacho, and its $Y$--$\Delta$--Heawood
graph is a classical member of the Heawood family
\cite{src:cs,src:heawoodfamily}. Individual known graphs and earlier
parameter ranges are not counted as additional resolutions of an
original catalogue item.


% Integrated source fragment: paired_384b.tex
\section{Triangle diversity and 2-domination: the remaining case of WoWII 384b}
\label{paired:section}

For a vertex $v$ of a finite simple graph, write $\theta_G(v)$ for the
number of triangles containing $v$, and put
\[
 T(G)=\{\theta_G(v):v\in V(G)\},\qquad t(G)=|T(G)|.
\]
Thus $t(G)$ counts distinct triangle-incidence values; it does not count
triangles. A set $D$ is \emph{2-dominating} if every vertex outside $D$
has at least two neighbors in $D$. Its minimum cardinality is
$\gamma_2(G)$. The archived statement of WoWII 384b is
\begin{equation}\label{paired:target}
 \gamma_2(G)\le\floor{n-\frac{t(G)}2}
 \qquad\text{for every connected graph of order }n>2.
\end{equation}
The results consolidated here leave this statement unresolved. They do,
however, reduce any possible failure to a sharply constrained finite
range. The distinction between a bound on a representative and a bound
on every counterexample is essential.

\begin{theorem}[Combined reduction for 384b]\label{paired:terminal}
Let $G$ be connected of order greater than two, and let $h$ be the
number of its vertices of degree at least three.
\begin{enumerate}
 \item Inequality~\eqref{paired:target} holds if $h\le10$ or $h\ge34$.
 \item If a counterexample exists, then a connected counterexample
 $G'$ exists with $|V(G')|\le37$ and $11\le h(G')\le33$.
 More specifically, every counterexample has such a representative
 with the same number of vertices of degree at least three.
 \item In every counterexample, putting $L=\{v:d_G(v)\le2\}$ and
 $H=G[V(G)\setminus L]$, one has
 \[
 n-\gamma_2(G)=\ceil{h/2},\qquad
 \gamma_2(G)=|L|+\floor{h/2},\qquad
 \mu(H)=\floor{h/2},
 \]
 where $\mu$ denotes the matching number. For odd $h$, all high-degree
 triangle counts are distinct and at least two, and $t=h+2$.
 For even $h$, at least $h-1$ high vertices have distinct counts at
 least two, and $t\in\{h+1,h+2\}$.
\end{enumerate}
\end{theorem}

The proof occupies the rest of this section. Universal structural
arguments are given explicitly. The finite parts are integer or rational
certificates with specified input coverage; they are not numerical
optimization conclusions. The verification files and their hashes are
specified in Section~\ref{paired:repro}. The terminal exclusion of $h=35$
uses a different sampler from the preceding exclusions and cannot be
obtained by treating its matching edges as independent.

\subsection{Omission sets, the critical profile, and matching}
\label{paired:critical-section}

Call $S\subseteq V(G)$ \emph{omittable} if
\begin{equation}\label{paired:omission}
 d_G(v)-|N_G(v)\cap S|\ge2\qquad(v\in S).
\end{equation}
Its complement is 2-dominating. We use
$k=n-\gamma_2(G)$, the maximum size of an omittable set, and consistently
use $H$ for the induced high-degree graph and $h=|V(H)|$.
Since $n$ and $\gamma_2(G)$ are integers,
\eqref{paired:target} is equivalent to $t\le2k$.

\begin{lemma}[Critical profile]\label{paired:critical}
A counterexample satisfies the assertions about $k$, $\gamma_2$, and
triangle counts in Theorem~\ref{paired:terminal}.
\end{lemma}
\begin{proof}
In a maximum cut, every vertex has at least half its neighbors across
the cut. Every high vertex consequently has at least two crossing
neighbors. The high vertices on either side are omittable, whence
$k\ge\ceil{h/2}$. A low vertex has triangle count zero or one, so
$t\le h+2$. In a counterexample $2k<t\le h+2$. For $h=2r$ this forces
$k=r$ and $t=h+1$ or $h+2$; at least $t-2\ge h-1$ distinct values
at least two are realized by high vertices. For $h=2r-1$ it forces
$k=r$ and $t=h+2$; every high vertex supplies a different value at
least two, and low vertices supply both zero and one. Subtracting $k$
from $n=|L|+h$ gives the claimed formula for $\gamma_2$.
\end{proof}

We use the following published bound of DeLaVi\~na, Goddard, Henning,
Pepper, and Vaughan~\cite[Theorem 3]{paired:DGH}. We give its specialized
proof to make clear that connectivity is not required when it is later
applied after deleting edges.

\begin{lemma}[Matching bound]\label{paired:matching-bound}
For any finite simple graph, with $H$ and $L$ defined by the same degree
threshold three,
\begin{equation}\label{paired:matching-ineq}
 \gamma_2(G)\le |L|+\mu(H).
\end{equation}
Consequently a critical counterexample has a perfect matching in $H$
when $h$ is even and a matching missing exactly one vertex when $h$ is odd.
\end{lemma}
\begin{proof}
Take a maximum cut of $H$ and let $B$ be its bipartite graph of crossing
edges. Every high vertex has at least $\ceil{d_H(v)/2}$ neighbors in $B$.
By the bipartite matching--vertex-cover theorem, $B$ has a vertex cover
$C$ with $|C|=\mu(B)\le\mu(H)$. Every $v\in V(H)\setminus C$ has all
its crossing neighbors in $C$, and therefore has at least
\[
 |N_G(v)\cap L|+\ceil{d_H(v)/2}\ge2
\]
neighbors in $L\cup C$; the inequality follows from $d_G(v)\ge3$.
Thus $L\cup C$ is 2-dominating. In a counterexample,
Lemma~\ref{paired:critical} and $\mu(H)\le\floor{h/2}$ force equality.
\end{proof}

For comparison, another established result is
\cite[Theorem 5]{paired:DGH}: if no two vertices of degree less than two
are adjacent, then $\gamma_2(G)\le(n+\alpha(G[L]))/2$. Its hypothesis
holds for connected $n>2$. In the present notation it gives
$k\ge\ceil{(h+|L|-\alpha(G[L]))/2}$. The sharper structural statements
below are proved directly; this published bound is not a substitute for
their proofs.

\subsection{Stability at equality in the maximum-cut bound}
\label{paired:stability-section}

For this subsection suppose $h>0$ and $k=\ceil{h/2}$. A degree-two
vertex is \emph{good} in a cut if both its neighbors are across the cut.
Define a multigraph $J$ on $V(H)$ by giving one edge $ab$ for each
degree-two vertex with high neighbors $a,b$. The pair $ab$ need not be
an edge of $H$. At this point parallel edges are allowed in the definition.

\begin{lemma}[Even stability]\label{paired:even-stability}
If $h=2k$, then $G[L]$ is independent, every degree-two vertex has two
high neighbors, and $J$ is a simple matching. Every high vertex has
strictly more crossing than noncrossing neighbors in every maximum cut.
\end{lemma}
\begin{proof}
Each high side of a maximum cut is omittable, so both sides have $k$
high vertices. A good degree-two vertex could be added to the high
vertices on its own side, contradicting maximality of $k$. Every
degree-two vertex thus has exactly one crossing neighbor; flipping it
preserves maximality of the cut.

A degree-two vertex cannot be adjacent to a leaf: the leaf edge crosses,
and the flip would give the leaf no crossing neighbor in a maximum cut.
If two degree-two vertices were adjacent, flipping one would give the
other zero or two crossing neighbors, both impossible. Two leaves
cannot be adjacent in a connected graph of order greater than two.
Hence $G[L]$ is independent. The high neighbors of a degree-two vertex
are on opposite sides; if they were on the same side, maximizing its
two incident edges would make it good.

If a high vertex $a$ supported two such vertices $x,y$, replace $a$
in its $k$-element high side by $x,y$. Both new vertices have their
neighbors selected, and no retained high vertex on that side is adjacent
to $x$ or $y$. This yields $k+1$ omissions, a contradiction. Thus $J$
is a simple matching. Finally, a high vertex with equally many neighbors
on each side could be flipped while preserving the maximum cut, giving
a high side of size $k+1$. This proves the strict majority assertion.
\end{proof}

\begin{lemma}[Odd low-degree structure]\label{paired:odd-stability}
If $h=2k-1$, then $G[L]$ has at most one edge. The multigraph $J$ is
simple and consists of a matching together with at most one component
isomorphic to $P_3$. If $G[L]$ has an edge, $J$ is a matching.
\end{lemma}
\begin{proof}
In a maximum cut let its high sides be $A,B$ with $|A|=k$ and
$|B|=k-1$. No good degree-two vertex lies on the side of $A$, and at
most one lies on the side of $B$: good vertices on one side are
independent and can be added to that high side.

Every low component is a path or isolated vertex. A low cycle would
be an entire component of $G$, contrary to $h>0$. A low path with an
edge is either a chain of at least two degree-two vertices between
high endpoints, possibly equal, or a chain of degree-two vertices
ending in a leaf and attached at its other end to a high vertex.
Keep the high endpoints fixed and maximize crossing edges along a
chain. If a chain between high endpoints has at least three internal
degree-two vertices, all its edges can cross or all but one can cross,
with the possible noncrossing edge placed at an end. In either case
two consecutive internal degree-two vertices are good on opposite
sides, which is impossible. A pendant chain with at least two internal
degree-two vertices has the same obstruction, since all its edges can
cross. Thus every nontrivial low component is a single edge.

For a path $a-x-y-b$ with $x,y$ low, opposite high colors would make
both internal vertices good. High endpoints both on the side of $B$
would permit a maximum two-crossing-edge orientation with a good vertex
on the side of $A$. Therefore both endpoints are on the side of $A$,
and the low edge can be oriented to give a good vertex on the side of
$B$. A leaf--degree-two edge similarly supplies a good degree-two
vertex on that side. Two low edges could be oriented independently
to supply two good vertices there, a contradiction. This proves the
assertion about $G[L]$.

There is no $J$ edge inside $B$, and at most one inside $A$. Among
crossing $J$ edges, an $A$ vertex has incidence at most one: replacing
it in $A$ by two incident low vertices would increase the omission
set. A $B$ vertex has incidence at most two, and two $B$ vertices
cannot both have incidence two: replacing the two by their four
incident low vertices gives $(k-1)-2+4=k+1$ omissions. Those four
vertices are distinct by the incidence restriction on $A$.

If there is no edge inside $A$, these conditions give a matching plus
at most one $P_3$. If there is one, its good low vertex can first be
added to $B$. Now any $B$ vertex with two crossing attachments would
enlarge this $k$-element omission set. Thus crossing $J$ edges form a
matching. Adding the single edge inside $A$ can create only a matching,
a $P_3$, or a $P_4$ with high-side pattern $B,A,A,B$. The last case
is impossible: replace the two middle $A$ vertices by the three low
vertices on its edges to get $k-2+3=k+1$ omissions.

Finally, if a low edge is present, use its good low vertex on the
side of $B$. This rules out an edge inside $A$ and forces crossing
incidence at most one also in $B$, giving a matching. All replacements
used here preserve domination at retained omitted high vertices:
newly omitted low vertices are adjacent, within the old high omission
set, only to vertices explicitly removed from that set.
\end{proof}

\begin{proposition}[Cubic vertices in the even case]\label{paired:cubic-filter}
An even-$h$ counterexample has between one and three vertices of full
degree three. In every maximum cut, at most six high vertices have
at least two neighbors outside the opposite high side. More precisely,
the number of such vertices on one side is at most twice the number
of cubic vertices on the other side.
\end{proposition}
\begin{proof}
Among $h-1$ high vertices with distinct triangle counts at least two,
at most two can be cubic, since a cubic count is at most three. There
is at most one exceptional high vertex, proving the upper bound three.

Suppose all high degrees were at least four. In a maximum cut the
strict-majority assertion of Lemma~\ref{paired:even-stability} gives
at least three crossing neighbors at every high vertex. Write $A,B$
for the two $k$-element high sides. If $v\in B$ had at least two
neighbors outside $A$, then $A\cup\{v\}$ would be omittable: $v$ is
supported and each retained vertex of $A$ loses at most one of its
at least three selected neighbors. Therefore each high vertex has
at most one neighbor that is low or on its own high side. In
particular $H[A]$ and $H[B]$ are matchings.

For a high vertex with a same-side mate, triangles using that mate
number at most $k$, and the remaining triangles use an edge of the
opposite-side matching, numbering at most $\floor{k/2}$. Such a vertex
has no low neighbor. A high vertex without a same-side mate has at
most $\floor{k/2}$ triangles using two high neighbors and at most
one further triangle using a low neighbor. Thus all triangle counts
are at most $k+\floor{k/2}$, and
$t\le k+\floor{k/2}+1\le2k$ for $k\ge2$. When $k=1$, all counts
are at most one, with the same conclusion. This contradiction proves
existence of a cubic vertex.

For the final assertion, let $v\in B$ have at least two neighbors
outside $A$. The attempted omission set $A\cup\{v\}$ must fail at
some $u\in A$ adjacent to $v$ with exactly two neighbors outside $A$.
Strict crossing majority forces $u$ to be cubic. Each such $u$ can
block at most two candidates $v$. Apply this argument on both sides
and the already proved bound of three cubic vertices. Only one new
high vertex is added at a time; no simultaneous two-vertex augmentation
is asserted.
\end{proof}

\subsection{Triangle kernels and the small high-degree orders}
\label{paired:kernel-section}

In a connected graph with $h>0$, every triangle using low vertices
has one of two forms. A \emph{pendant triangle} has a high root and
two adjacent degree-two vertices. An \emph{edge attachment} has one
degree-two vertex attached to both ends of an edge of $H$. Let $P,Q$
count these two kinds of triangles.

\begin{lemma}[Finite kernel reduction]\label{paired:kernel}
In a critical counterexample $P+Q\le k$. Let $K$ be the graph induced
by the high vertices and all low vertices belonging to triangles.
Then $|V(K)|=h+2P+Q\le h+2k$. Set
\[
 b(v)=\begin{cases}\max(3,d_K(v)),&v\in V(H),\\2,&v\in V(K)\setminus V(H).
 \end{cases}
\]
Every $S\subseteq V(K)$ satisfying
$b(v)-|N_K(v)\cap S|\ge2$ for all $v\in S$ is omittable in $G$.
For even $h$ the edge attachments form a matching and $P=0$.
For odd $h$ they form a matching with at most one $P_3$ component;
if $P=1$ they form a matching, and $P\le1$.
\end{lemma}
\begin{proof}
Omit every edge-attachment vertex and one private vertex from every
pendant triangle. These choices are independent and retain both
neighbors selected, so their number is at most $k$. Every high
vertex has full degree at least $b(v)$, every omitted vertex lies
in $K$, and any additional vertex is kept selected. Thus the displayed
inequalities imply~\eqref{paired:omission}. The attachment restrictions
are immediate from Lemmas~\ref{paired:even-stability} and
\ref{paired:odd-stability}. This is a one-way safe reduction; we do
not assert that every enumerated kernel is realizable as a counterexample.
\end{proof}

For an odd critical core, write $a_v=\theta_H(v)$, let $p_v$ be the
number of pendant triangles at $v$, and let $Q$ also denote the graph
of edge attachments. Then
\begin{equation}\label{paired:offset}
 \theta_G(v)=a_v+p_v+d_Q(v).
\end{equation}
Each offset is zero or one except that at most one can be two. For
the elementary $h\le7$ checks it is harmless to enumerate the larger
class of all attachment multisets with $P+Q\le k$.

\begin{lemma}[Safe sorted-count filters]\label{paired:sorted-filter}
For odd $h$, sort the core counts $a_1\le\cdots\le a_h$ and set
$b_0=1$, $b_i=\max(a_i,b_{i-1}+1)$. If offsets of the form just
described can make distinct final counts at least two, then
$b_i-a_i\le2$ for every $i$ and at most one difference equals two.
For even $h$, after deleting a suitable one of the core counts, the
same greedy construction on $h-1$ counts must have all differences
at most one.
\end{lemma}
\begin{proof}
Consider a feasible assignment of final counts and an inversion
$a_i\le a_j$ but $x_i>x_j$. If $a_i=a_j$, interchange $x_i,x_j$,
which only interchanges offsets. If $a_i<a_j$, after interchanging
both offsets are nonnegative and at most one: $x_j-a_i\le x_i-a_i-1$
and $x_i-a_j\le x_i-a_i-1$. Thus removing inversions does not create
an extra offset two. The resulting increasing assignment dominates
the greedy sequence termwise, proving the odd assertion. In the
even case choose the $h-1$ distinct high counts at least two from
Lemma~\ref{paired:critical}; all their attachment offsets are at most
one, and the same interchange argument applies.
\end{proof}

Here is the explicit witness used in the larger small-core checks.
For a permitted attachment pattern set
$d_v^*=\max(3,d_H(v)+2p_v+d_Q(v))$. Test every $R\subseteq V(H)$
with $d_v^*-|N_H(v)\cap R|\ge2$ for $v\in R$. With $R$, omit one
private vertex of each pendant triangle whose root is outside $R$,
and every edge-attachment vertex whose two roots are outside $R$.
The resulting omission set has size
\begin{equation}\label{paired:kernel-witness}
 |R|+\sum_{v\notin R}p_v+
 |\{uv\in E(Q):u,v\notin R\}|.
\end{equation}
It is safe because the new omitted low vertices are adjacent to no
retained vertex of $R$. For a high-only witness, use
$d_v^*=\max(3,d_H(v))$, which is still a lower bound on full degree.

\begin{proposition}[Verified small-core theorem]\label{paired:small}
Every connected graph with $h\le10$ satisfies~\eqref{paired:target}.
\end{proposition}
\begin{proof}
For $h=0$, the graph is a path or cycle. All triangle counts are equal,
so $t=1$, and a degree-two vertex can be omitted. Suppose $h>0$ and
that a counterexample exists. Lemmas~\ref{paired:critical}--
\ref{paired:kernel} give the finite verification described next.

Encode a labeled core on $r$ vertices by the bits of its unordered
pairs in lexicographic order. Triangle counts are calculated by
counting edges inside each neighborhood. For $1\le h\le6$, enumerate
all $2^{\binom h2}$ masks. For odd $h$, enumerate attachment multisets
with $P+Q\le\ceil{h/2}$; the necessary inequality
\[
 \sum_{v\in H}\theta_H(v)+2k\ge h(h+3)/2
\]
can first discard a core, because each attachment raises the high
count sum by at most two. For even $h$, enumerate only matching
edge attachments. Retain only patterns for which the set of high
counts, together with zero and with one if a triangle attachment is
present, has more than $2k$ values. Adding a value not realized by
the graph enlarges this set and is a safe relaxation. In each retained
kernel enumerate $(k+1)$-subsets and test the inequalities in
Lemma~\ref{paired:kernel}.

For $h=7$ use the 1,044 supplied seven-vertex masks and the same
attachment procedure. Coverage of these masks is separately verified:
all $7!$ relabelings of all masks mark every one of the $2^{21}$
labeled masks. For $h=8$, extend each of these seven-vertex masks
by each of the $2^7$ possible neighborhoods of an eighth vertex,
then enumerate matching attachments. Every eight-vertex graph is
isomorphic to one of these extensions; duplicates cause no difficulty.

For $h=9$, use the 12,346 supplied eight-vertex masks. Their coverage
is certified by an explicit permutation for each of the preceding
133,632 extensions, mapping it to a listed mask. Extend by all $2^8$
neighborhoods. Apply Lemma~\ref{paired:sorted-filter}, search first
for a high-only omission set of size six, and, if needed, enumerate
the two attachment families in Lemma~\ref{paired:kernel} and all
$2^9$ choices in~\eqref{paired:kernel-witness}.

For $h=10$, Proposition~\ref{paired:cubic-filter} supplies a full
cubic vertex. Relabel it as the tenth vertex. Its core neighborhood
has at most three vertices, so only
$\sum_{j=0}^3\binom9j=130$ extensions per nine-vertex mask are needed.
The supplied list of 274,668 nine-vertex masks is covered by explicit
permutations for all $12,346\cdot2^8$ extensions of the eight-vertex
list. Check the even filter of Lemma~\ref{paired:sorted-filter} for
each possible exceptional vertex. Every surviving extension already
has six high-only omissions using the lower degree bound
$\max(3,d_H(v))$; attachments do not need to be enumerated at this step.

The complete enumerations give the following counts. ``Critical''
means that the triangle-count filter allows failure, not that a
counterexample exists.
\begin{center}\small
\begin{tabular}{rrrrr}
\toprule
$h$&Core masks or extensions&Attachment patterns&Critical patterns&Failures\\
\midrule
1&1&2&0&0\\
2&2&3&0&0\\
3&8&0&0&0\\
4&64&304&0&0\\
5&1,024&24,416&1,800&0\\
6&32,768&708,608&25,920&0\\
7&1,044&2,446,446&61,253&0\\
8&133,632&17,150,528&857,204&0\\
9&3,160,576&20,064,253&296,450&0\\
10&35,706,840&0&0&0\\
\bottomrule
\end{tabular}
\end{center}
For $h=7$, 170 cores pass the triangle-sum filter. For $h=9$,
342,889 extensions pass the sorted-offset filter and 340,059 already
have a high-only witness. For $h=10$, all 3,971,627 offset-compatible
extensions have a high-only witness. The last row's zero critical
attachment patterns records that no attachment search was needed.
All coverage and witness checks are exact and were rerun for this
consolidation. Their algorithms and hashes are recorded in
Section~\ref{paired:repro}. A failure would require an unresolved
kernel, and none occurs.
\end{proof}

\subsection{A component restriction at the lower boundary}
\label{paired:components-section}

\begin{lemma}\label{paired:small-components}
In an odd-$h$ counterexample, every component of $H$ has at least five
vertices. A component with only matching edge attachments, and no
pendant triangle or fork, has at least seven vertices. In particular,
if $h=11$, then $H$ is connected.
\end{lemma}
\begin{proof}
All high counts are distinct and at least two. A one-vertex core
component has at most one pendant increment and cannot reach two.
A two-vertex component has at most one edge attachment; even a
pendant increment leaves one endpoint with count at most one. A
three-vertex component has all core counts zero unless it is a
triangle, in which case all are one. Offsets at most one, with at
most one two, cannot give three distinct counts at least two. In
order four, a core count three forces $K_4$, where only the three
values $3,4,5$ are possible after offsets. Otherwise core counts are
at most two and there are only the values $2,3$ and possibly one
four. Thus every component has order at least five.

Consider now a component $K$ of order five with only matching
increments. Distinct corrected counts at least two have sum at least
20. If $z$ is the number of triangles of $K$ and $M$ its attachment
matching, then $3z+2|M|\ge20$. Since $|M|\le2$, $z\ge6$.
Deleting two edges from $K_5$ destroys at least five triangles, so
$K$ is $K_5$ or $K_5-e$. Their count multisets are respectively
$6^5$ and $3^2,5^3$; adding zero or one cannot separate a triple
of equal counts.

For order six, suppose first that the largest core count is at most
six. The six distinct corrected counts must then be $2,3,4,5,6,7$,
of sum 27. The equality $3z+2|M|=27$ gives $|M|=0$ or $3$. The
first cannot produce seven, so $M$ is perfect and the base counts
are exactly $1,2,3,4,5,6$. A vertex with base count six and degree
four belongs to a $K_5$, forcing repeated counts at least six. Thus
it is universal. Its five-vertex neighborhood $W$ has six edges and
exactly one triangle, since the total graph has seven triangles.
The other five counts are $d_W(w)+\theta_W(w)=1,2,3,4,5$.
Count five forces degree four there (degree
at most three plus at most one triangle is at most four), hence a
universal vertex of that neighborhood. A universal vertex in a
five-vertex six-edge graph belongs to two triangles, a contradiction.

If the largest count is at least seven, its vertex is universal and
the six-vertex component has at least twelve edges. Its complement
has at most three edges. Exhausting the possible complements gives
\begin{center}\small
\begin{tabular}{ll}
\toprule
Complement&Triangle-count multiset in $K$\\\midrule
$6K_1$&$10^6$\\
$K_2+4K_1$&$6^2,9^4$\\
$2K_2+2K_1$&$5^4,8^2$\\
$P_3+3K_1$&$3,6^2,8^3$\\
$3K_2$&$4^6$\\
$P_3+K_2+K_1$&$2,4^2,5^2,7$\\
$P_4+2K_1$&$3^2,5^2,7^2$\\
$K_{1,3}+2K_1$&$1,6^3,7^2$\\
$K_3+3K_1$&$3^3,7^3$\\\bottomrule
\end{tabular}
\end{center}
A triple of equal values cannot be separated by binary increments.
In the $P_3+K_2+K_1$ row, four base values in $\{4,5\}$ have only
three possible corrected values. In the $P_4+2K_1$ row, separating
all three pairs requires exactly three matched vertices, impossible.
This excludes all components of orders five and six without an
exceptional attachment.

There is at most one exceptional attachment component globally:
either one pendant triangle, with the other attachments a matching,
or one fork and no pendant triangle. An exceptional component needs
at least five vertices and every other component at least seven.
Two components would therefore require at least twelve vertices.
This proves connectivity when $h=11$; it does not exclude that order.
\end{proof}

\subsection{Triangle rank and basic probabilistic alteration}
\label{paired:rank-section}

We next turn triangle diversity into degree lower bounds. Throughout,
an out-of-range binomial coefficient is zero.

\begin{lemma}[Untouched-neighbor rank bound]\label{paired:rank}
Let $v$ have degree $d$ and triangle count $u$ in a graph. At least
$2u-d(d-2)$ neighbors of $v$ have triangle count at least $u$.
Consequently, in a critical counterexample, the $i$th of the increasing
distinct high counts at least two has full degree at least
\begin{equation}\label{paired:profile-definition}
 a_i(h)=\min\{d\ge3:\ \tbinom d2\ge i+1,
                    \ d(d-2)\ge3i+2-h\}.
\end{equation}
For odd $h$ these indices are $1,\ldots,h$; for even $h$ they are
$1,\ldots,h-1$, with one additional high vertex of degree at least three.
\end{lemma}
\begin{proof}
There are $r=\binom d2-u$ missing edges in the graph induced by
$N(v)$. At most $2r$ neighbors are endpoints of missing edges.
Every untouched neighbor $w$ is adjacent to the other $d-1$
neighbors of $v$. The $d-1$ triangles using $vw$, together with the
neighborhood edges not incident with $w$, give at least
$(d-1)+(u-(d-1))=u$ triangles at $w$. There are at least
$d-2r=2u-d(d-2)$ such neighbors.

For the vertex of rank $i$, $u\ge i+1$. Every untouched neighbor
has count at least two and is therefore high. Among the other high
vertices at most $h-i$ can have count at least $u$: these are the
later ranked vertices and, in the even case, possibly the one
unranked vertex. Hence $2u-d(d-2)\le h-i$, which gives
\eqref{paired:profile-definition}. The binomial inequality is immediate.
\end{proof}

The simpler profile $\rho(j)=\min\{d\ge3:\binom d2\ge j\}$ suffices
for the first range of exclusions. Select all low vertices initially,
select each high vertex independently with probability $p$, and add
each unselected high vertex having fewer than two selected neighbors.
For $z=1-p$, the probability of adding a high vertex of full degree
$d$ is at most
\begin{equation}\label{paired:g}
 g_d=z^{d+1}+dpz^d,
 \qquad g_d-g_{d+1}=dp^2z^d>0.
\end{equation}
Indeed, replacing a preselected low neighbor by an independently
selected artificial neighbor can only increase the bad event; with
$d$ independent neighbors the two terms are the probabilities of
zero or one selected neighbor while the vertex itself is unselected.
The resulting 2-dominating set proves
\begin{equation}\label{paired:independent-bound}
 k\ge (1-p)h-\sum_{v\in H}g_{d_G(v)}.
\end{equation}

\begin{proposition}\label{paired:large44}
Every connected graph with $h\ge44$ satisfies~\eqref{paired:target}.
\end{proposition}
\begin{proof}
First put $p=3/8$ and define
$A_h=\sum_{j=2}^{h+1}g_{\rho(j)}$ and $B_h=g_3+A_{h-1}$.
The critical odd and even profiles bound the repair sums by $A_h$
and $B_h$, respectively. To contradict the critical value of $k$,
it suffices that $h/8-A_h-1/2>0$ in odd order or $h/8-B_h>0$ in
even order. Exact evaluation gives
\[
 \frac{52}{8}-B_{52}=\frac{3607919267}{34359738368},\qquad
 \frac{53}{8}-A_{53}-\frac12=\frac{3547612629}{34359738368}.
\]
When either order is increased by two, the margin increases by at
least
\[
 \frac14-2g_{11}=\frac{3367232921}{17179869184}>0,
\]
because the two newly required triangle counts have degree lower
bound at least eleven. This is a parity induction proving the result
for every $h\ge52$, not a finite parameter extrapolation.

For $44\le h\le51$, select a uniformly random $s$-subset of the high
vertices. If a high vertex has full degree at least $d$, its repair
probability is at most
\begin{equation}\label{paired:R0}
 R_0(h,s,d)=\frac{\binom{h-1-e}{s}+e\binom{h-1-e}{s-1}}
                       {\binom hs},\qquad e=\min(d,h-1).
\end{equation}
To see this when some neighbors are low, fill the shortfall in actual
high neighbors by distinct other high vertices. The preselected low
neighbors provide at least as much support as the chosen filler
vertices. If no filler position remains, the preselected support
already makes the event impossible, so the displayed upper bound
remains safe. Counting samples that omit the vertex and select zero
or one of its $e$ artificial neighbors gives~\eqref{paired:R0}.

The expected selected-high count after alteration is bounded by
$s+\sum R_0(h,s,\rho(j))$, with the odd or even critical profile.
The exact margins below the forbidden threshold $\floor{h/2}$ are
\begin{center}\small
\begin{tabular}{rrl@{\qquad}rrl}
\toprule
$h$&$s$&Margin&$h$&$s$&Margin\\\midrule
44&16&$1745160/12198197$&45&16&$1763081/12198197$\\
46&17&$36171337/109383285$&47&17&$479277388/1339779509$\\
48&17&$16780745963/32154708216$&49&18&$337323149/617148728$\\
50&18&$75018443/101504725$&51&18&$78679171/101504725$\\\bottomrule
\end{tabular}
\end{center}
Every margin is positive, completing the finite remaining range.
\end{proof}

\subsection{Preselected low neighbors and the orders 41--43}
\label{paired:boundary-section}

If at least one low neighbor is already selected, the corresponding
one-neighbor repair bound is
\begin{equation}\label{paired:R1}
 R_1(h,s,d)=\frac{\binom{h-1-e}{s}}{\binom hs},
 \qquad e=\min(d-1,h-1),
\end{equation}
and with two preselected low neighbors repair is unnecessary. These
bounds follow from the same filling argument, reserving the credited
low neighbors before filling. For any independent set $I$ of degree-two
vertices, omit $x\in I$ if all its high neighbors are initially selected;
any low neighbors are kept. These omissions are simultaneous, and
each costs at least
\begin{equation}\label{paired:low-credit}
 q(h,s)=\frac{s(s-1)}{h(h-1)}
\end{equation}
in expected selected size. A pendant triangle permits one private
vertex to be omitted when its root is sampled, with probability
$p(h,s)=s/h$, and its root needs no repair while both private
vertices remain selected. When the private vertex is omitted its
root was initially selected, so there is no conflict between these
two savings.

\begin{lemma}\label{paired:4342-degree2}
There is no counterexample with $h=43$, or with $h=42$ and a
degree-two vertex. For $h=41$, a counterexample would have exactly
one degree-two vertex.
\end{lemma}
\begin{proof}
For a fixed critical $\rho$-profile put $u_i=R_0(h,s,d_i)-R_1(h,s,d_i)$.
An ordinary degree-two vertex, with two high neighbors, saves at least
the sum of the two smallest $u_i$, in addition to its omission
credit $q(h,s)$. A pendant triangle saves at least the smallest
$R_0(h,s,d_i)$, together with $p(h,s)$.

Odd critical order requires a low vertex of triangle count one, hence
an ordinary triangle attachment or a pendant triangle. For $h=43$,
take $s=16$. The exact margins below 21 are respectively
$1057200/7762489$ and $2799123/7762489$. For $h=42$, every degree-two
vertex is ordinary by even stability; the same sample size gives
margin $75065/541569$ below 21.

For $h=41$, take $s=15$. Two nonadjacent degree-two vertices give
two simultaneous credits~\eqref{paired:low-credit}; even without
crediting repair reductions, the exact margin below 20 is
$162283/3197836$. If there are at least two degree-two vertices but
none are nonadjacent, there are exactly two: three pairwise adjacent
degree-two vertices would be a separate triangle component. The
required low triangle count one then makes the two vertices a
pendant triangle, whose margin is $28995/155992$. There must be a
degree-two vertex to realize the low count one. Thus only the
exactly-one case remains. All four evaluations are substitutions
into~\eqref{paired:R0}--\eqref{paired:low-credit} with the specified
profiles and were reproduced exactly.
\end{proof}

\begin{lemma}[The leaf-only case at $h=42$]\label{paired:h42}
There is no counterexample with $h=42$ and no degree-two vertex.
\end{lemma}
\begin{proof}
All low vertices are leaves. Criticality forces them to be present,
and forces all 42 high triangle counts to be distinct and positive.
These are also their core counts. In increasing order, $u_i\ge i$.
Applying the untouched-neighbor argument inside $H$ gives
$d_H(v_i)(d_H(v_i)-2)\ge3i-42$. Apart from a possible core-degree-two
vertex at rank one, a lower profile is
\[
 (3^3,4^3,5^4,6^5,7^6,8^7,9^7,10^5,11^2).
\]
The exception can only have count one and must have a leaf neighbor
in order to be high in $G$.

For a uniform 16-subset of the core define
\[
 F_0(d)=\frac{\binom{41-d}{16}+d\binom{41-d}{15}}{\binom{42}{16}},
 \qquad F_1(d)=\frac{\binom{41-d}{16}}{\binom{42}{16}}.
\]
For the displayed profile the baseline is
$16+\sum F_0=705171683/33577278$. The least saving from one leaf
neighbor is
$c=\min(F_0(d)-F_1(d))=10120/987567$. If all core degrees are at
least three, at least one leaf supplies this saving. If the rank-one
exception occurs, its actual contribution is at most $F_1(2)$,
and its saving against the nominal $F_0(3)$ is
$F_0(3)-F_1(2)=160/861>c$. In either case the expected high cost is
at most $21-295235/33577278<21$, contradicting criticality.
\end{proof}

\begin{lemma}[The one-attachment case at $h=41$]\label{paired:h41}
There is no counterexample with $h=41$ and exactly one degree-two vertex.
\end{lemma}
\begin{proof}
Write $x$ for that vertex. It supplies the required low triangle count
one, so its high neighbors $a,b$ are adjacent. The low count zero
forces a leaf $z$ at a high root $c$. Ignore any other leaves in the
repair test. Order the 41 distinct high counts increasingly as $t_i$,
so $t_i\ge i+1$, and put $\epsilon_i=1$ at the two ranks $a,b$ and
zero elsewhere. The core count at rank $i$ is $t_i-\epsilon_i$.
The untouched-neighbor argument inside $H$ now gives
\[
 \binom{d_i}{2}\ge i+1-\epsilon_i,\quad
 d_i(d_i-2)\ge3i+2-41-3\epsilon_i,\quad d_i\ge3-\epsilon_i.
\]
For completeness, when $\epsilon_i=1$, a neighbor with core count
at least $t_i-1$ has full count at least that value, and at most
$41-i+1$ other ranks can qualify; this is the extra one that produces
the coefficient $-3\epsilon_i$. Let $\ell_i^{\epsilon}$ be the
least integer satisfying these three inequalities.

Choose a uniform 15-subset $S$ of the core and put
\[
 F_r(d)=\frac{\binom{40-d}{15}+\mathbf1_{r=0}d\binom{40-d}{14}}
                    {\binom{41}{15}}\quad(r=0,1),
 \qquad F_r(d)=0\quad(r\ge2).
\]
For each placement of the ranks $a<b$ and $c$, the expected high
cost, after omitting $x$ when both its roots are sampled, is at most
\begin{equation}\label{paired:U41}
 U(a,b,c)=15+\sum_{i=1}^{41}
 F_{\epsilon_i+\mathbf1_{i=c}}(\ell_i^{\epsilon_i})
 -\frac{15\cdot14}{41\cdot40}.
\end{equation}
Looping over the $\binom{41}{2}\cdot41=33,620$ placements with integer
binomial coefficients gives
\begin{align*}
 \max U&=\frac{127990267}{6395672}=20+\eta,
       &\eta&=\frac{76827}{6395672},\\
 \max_{\min(a,b,c)\le34}U&=\frac{127854567}{6395672}<20,\\
 F_0(3)-F_0(4)&=\frac{3150}{28823}>\eta,
 &F_0(9)-F_0(10)&=\frac{108675}{6395672}>\eta.
\end{align*}
The first two maxima are attained at $(40,41,41)$ and $(28,40,40)$.
These formulas and the finite loop completely specify the arithmetic
certificate; there are only 22 placements with $U\ge20$.

Thus a counterexample requires $a,b,c\ge35$. The ranks at most 34
have no distinguished low support. Since $\ell_1^0=\ell_2^0=3$,
the first saving inequality forces their core degrees to be exactly
three. Their full counts are distinct and at least two, so $t_1=2$
and $t_2=3$. Since $\ell_i^0\le9$ for $i\le34$ and
$\ell_{34}^0=9$, the second saving inequality forces all these
degrees to be at most nine and the degree at rank 34 to equal nine.
The exact untouched-neighbor inequality at that vertex gives
$7\ge2t_{34}-9(9-2)$, so $t_{34}=35$. Its neighborhood is $K_9$
minus one edge. The seven untouched neighbors are precisely ranks
35 through 41, and they induce a clique.

The neighborhood of rank one, whose degree is three and count is
two, is a $P_3$. At least one of its two nonadjacent endpoints $w$
is outside the seven-vertex clique. Then $3\le d_H(w)\le9$,
$w$ and $v_1$ have exactly one common core neighbor, and neither
is a distinguished boundary root. On the event that neither lies
in $S$ and each has exactly one neighbor in $S$, both are initially
bad. Add $w$ but do not add $v_1$. The latter now has its original
selected neighbor and the added neighbor $w$. All other bad vertices
are still added; all other vertices left outside already had two
initially selected neighbors. Thus this modification is valid and
saves one. Optional omission of $x$ affects only roots in $S$.

For $d=d_H(w)$ the event has probability
\[
 \frac{\binom{39-d}{14}+(d-2)\binom{39-d}{13}}{\binom{41}{15}}
 \ge\frac{198375}{12791344}\qquad(3\le d\le9).
\]
The minimum is at $d=9$; all seven evaluations are in the exact
certificate. The improved expectation is at most
$20+(153654-198375)/12791344<20$, a contradiction.
\end{proof}

Together, Proposition~\ref{paired:large44} and
Lemmas~\ref{paired:4342-degree2}--\ref{paired:h41} exclude every $h\ge41$.

\subsection{Matching-coupled repair with a fixed-size sample}
\label{paired:uniform-pairs-section}

The next improvement uses the actual perfect or near-perfect matching
forced by Lemma~\ref{paired:matching-bound}. We first give the common
support and repair argument, which is used with several sampling laws.

\begin{lemma}[Artificial support and pair repair]\label{paired:artificial}
Suppose a high vertex $v$ has a degree lower bound $a(v)$, and that
$r(v)$ distinct low neighbors are credited as initially selected.
If $1\le a(v)-r(v)\le h-1$, an artificial set
$A_v\subseteq V(H)\setminus\{v\}$ of size $a(v)-r(v)$ can be chosen
such that
\begin{equation}\label{paired:artificial-domination}
 |N_G(v)\cap(S\cup L)|\ge r(v)+|A_v\cap S|
 \qquad(S\subseteq V(H)).
\end{equation}
When $v$ has a partner in a fixed matching, $A_v$ can include that
actual partner. The sets need not be symmetric.

Call an endpoint outside $S$ bad if its artificial support is at most
one. Add bad endpoints, except that when both endpoints of a matching
edge are bad and at least one has support one, add only one: add the
zero-support endpoint if there is one, and otherwise either endpoint.
Repair an unmatched bad vertex individually. The resulting set is
2-dominating.
\end{lemma}
\begin{proof}
Take actual high neighbors first, including the partner if present,
and fill any shortfall with distinct other high vertices. Since
$d_G(v)\ge a(v)$, the number of fillers is at most the number of
uncredited low neighbors. For every $S$, those permanently selected
neighbors provide at least as much support as the selected fillers.
This proves~\eqref{paired:artificial-domination}. If actual high
neighbors exceed the required size, choose a suitable subset instead.

When one bad endpoint is left outside, it had one initial selected
neighbor and its actual partner is now added. Thus it has two
selected neighbors. All other vertices left outside were already
supported, and additions on other matching edges only help. Since
matching edges are disjoint, all repairs are simultaneously valid.
\end{proof}

For a uniform $s$-subset of the high vertices, first take $r(v)=0$.
Let $a,b$ be the artificial degrees of a matched pair and let $c$ be
the intersection size of their artificial neighborhoods outside the
pair. Define the integer numerators
\begin{align}
 F_{h,s}(a)&=\binom{h-1-a}{s}+a\binom{h-1-a}{s-1},
       \label{paired:F-uniform}\\
 Q_{h,s}(a,b,c)&=(a+b-2-c)\binom{h-a-b+c}{s-1}\notag\\
 &\hspace{17mm}+(a-1-c)(b-1-c)\binom{h-a-b+c}{s-2}.
       \label{paired:Q-uniform}
\end{align}
The probability that $v$ is bad is $F_{h,s}(a)/\binom hs$. The
numerator $Q_{h,s}$ counts the event that both endpoints are outside
$S$, each has support at most one, and at least one has support one.
One selected vertex in the union of the outside neighborhoods gives
the first term. Two selected vertices, one in each exclusive part,
give the second. Consequently the expected pair repair cost is at
most $(F_{h,s}(a)+F_{h,s}(b)-Q_{h,s}(a,b,c))/\binom hs$.

Choose rational weights $\beta_a$ satisfying
\begin{equation}\label{paired:uniform-dual}
 \beta_a+\beta_b\le Q_{h,s}(a,b,c)
 \quad\text{for }\max(0,a+b-h)\le c<\min(a,b).
\end{equation}
For a perfect matching put $\zeta=0$; for a near-perfect matching
choose $\zeta\ge\max_a\beta_a$. Summing over the actual matching,
including the unmatched correction, yields the universal bound
\begin{equation}\label{paired:uniform-bound}
 \mathbb E(|D|-|L|)\le s+
 \frac{\sum_a n_a(F_{h,s}(a)-\beta_a)+\zeta}{\binom hs}.
\end{equation}
Weights are allowed to be negative. Their role is solely the displayed
inequality, verified by exact substitution for every allowable overlap.

\begin{proposition}[The order $h=40$]\label{paired:h40}
There is no counterexample with exactly forty high vertices.
\end{proposition}
\begin{proof}
The rank profile~\eqref{paired:profile-definition}, including the one
exceptional type three, has multiplicities
\[
 (n_3,\ldots,n_{10})=(3,3,4,5,6,7,6,6).
\]
Take $s=14$ and $D=\binom{40}{14}=23206929840$. The following integer
weights satisfy~\eqref{paired:uniform-dual}:
\begin{center}\small
\begin{tabular}{cr@{\qquad}cr}
\toprule
$a$&$\beta_a$&$a$&$\beta_a$\\\midrule
3&923276120&7&236242200\\
4&775290440&8&49774300\\
5&571290100&9&$-137436500$\\
6&413052400&10&$-389725340$\\\bottomrule
\end{tabular}
\end{center}
There are 192 overlap inequalities over 36 unordered type pairs.
Their exact binomial evaluation is the finite certificate. The checker
also compares the formula with 1,297 direct small-set controls.
Equation~\eqref{paired:uniform-bound} gives
\[
 \mathbb E(|D|-|L|)\le\frac{92365610053}{4641385968}
 =20-\frac{462109307}{4641385968}<20,
\]
contrary to the critical value $|L|+20$.
\end{proof}

\subsection{Credited boundary support at \texorpdfstring{$h=39$}{h=39}}
\label{paired:h39-section}

At odd critical order the mandatory low count one supplies either
a pendant triangle (case P) or an edge attachment (case Q). The
mandatory low count zero supplies extra boundary support in case Q.

\begin{proposition}[The order $h=39$]\label{paired:h39}
There is no counterexample with exactly thirty-nine high vertices.
\end{proposition}
\begin{proof}
The rank-profile multiplicities are
$(2,3,4,5,6,7,6,6)$ for degrees $3,\ldots,10$, and $H$ has a
near-perfect matching. If a pendant triangle exists, credit its two
private selected neighbors at its root. Otherwise choose a private
degree-two vertex $x$ with adjacent high roots $u,v$. A low vertex
of count zero exists; a path from its low component to $H$ supplies
a low vertex $z\ne x$ adjacent to a high root $w$. Credit $x$ at
$u,v$ and $z$ at $w$. Thus the credited support is either one at
three distinct roots, or two at one of $u,v$ and one at the other.

Choose a uniform 14-subset and use
Lemma~\ref{paired:artificial} with types $(a,r)$, artificial-set size
$a-r$. Write $D=\binom{39}{14}=15084504396$ and
\[
 F_r(a)=\begin{cases}
 \binom{38-a}{14}+a\binom{38-a}{13},&r=0,\\
 \binom{39-a}{14},&r=1,\\
 0,&r=2.
 \end{cases}
\]
For a pair $(a,r),(b,t)$, set $d=a-r$, $e=b-t$. If their outside
artificial sets have intersection $c$, the pair saving has numerator
\begin{equation}\label{paired:Q-credit}
 Q(a,r;b,t;c)=
 \sum_{\substack{i,j,\ell\in\{0,1\}\\r+i+j\le1,\ t+j+\ell\le1\\
                         r+t+i+\ell+2j\ge1}}
 \binom{d-1-c}{i}\binom cj\binom{e-1-c}{\ell}
 \binom{39-d-e+c}{14-i-j-\ell}.
\end{equation}
Here $i,j,\ell$ count selected vertices in the first exclusive part,
the common part, and the second exclusive part. The constraints are
exactly that both endpoints are bad and at least one has support
one. This derives the formula, including its use with credited support.

In P omit one private vertex if its root is initially sampled, saving
$14/39=\binom{38}{13}/D$. In Q omit $x$ if both roots are initially
sampled, saving $\binom{37}{12}/D$. There is one additional possible
exchange in Q. Suppose $w=u$, so $u$ has support credit two, and
its matching partner $q$ is an ordinary type $(a,0)$. On the event
\[
 u,q\notin S,\quad v\in S,\quad |A_q\cap S|=1,
\]
the ordinary repair would add $q$. Add $u$ instead. Now $q$ has
its one initial neighbor and $u$, while both roots of $x$ are
selected; omit $x$ to save one extra vertex. If $v\in A_q$ this
event has numerator $\binom{38-a}{13}$, and otherwise numerator
$(a-1)\binom{37-a}{12}$. Thus its numerator is at least
\begin{equation}\label{paired:exchange39}
 E(a)=\min\{\binom{38-a}{13},(a-1)\binom{37-a}{12}\}.
\end{equation}
The event is disjoint from the preceding ordinary omission event.
There is at most one support-two root, so this exchange is never
counted on more than one matching edge. The symmetric case $w=v$
is identical; no extra credit is used if the support-two root is
unmatched or matched to the other credited root.

Use the following integer weights, all scaled by $D$:
\begin{center}\small
\begin{tabular}{rrr}
\toprule
$a$&$\beta_{a,0}$&$\beta_{a,1}$\\\midrule
3&383804427&0\\
4&364489027&0\\
5&291476815&0\\
6&203279727&0\\
7&104890051&$-101591575$\\
8&41134373&$-212045947$\\
9&$-62663615$&$-284255827$\\
10&$-161677327$&$-323629527$\\\bottomrule
\end{tabular}
\end{center}
Set $\zeta=383804427$. In Q take $\beta_{a,2}=0$; in P take
$\beta_{a,2}=-\zeta$. All weights are at most $\zeta$. Add $E(b)$
to~\eqref{paired:Q-credit} when $(r,t)=(2,0)$ in Q, and add $E(a)$
when $(r,t)=(0,2)$; call the resulting numerator $\widehat Q$.
For every applicable pair and overlap the exact inequalities are
\[
 \beta_{a,r}+\beta_{b,t}\le\widehat Q(a,r;b,t;c).
\]
There are 1,249 Q inequalities and 565 P inequalities. They are all
checked by directly evaluating the finite sums of binomial coefficients.

Summing over the near-perfect matching, and writing $n_{a,r}$ for
the credited-type multiplicities, yields the upper bound
\begin{equation}\label{paired:bound39}
 14+\frac{\sum_{a,r}n_{a,r}(F_r(a)-\beta_{a,r})+\zeta-O}{D},
 \qquad O=\begin{cases}\binom{37}{12},&\mathrm Q,\\
                         \binom{38}{13},&\mathrm P.
          \end{cases}
\end{equation}
All allocations of the credited roots to the profile require 351 Q
type placements and eight P type placements. Repeated types at
distinct roots and coincidence of $w$ with a Q root are included.
The exact maxima of~\eqref{paired:bound39} are
\[
 \frac{5617583945}{295774596}\quad(\mathrm Q),\qquad
 \frac{16662763795}{887323788}\quad(\mathrm P),
\]
attained at three distinct type-ten roots and at a type-ten pendant
root, respectively. Both are below 19; the smaller strict margin is
$2133379/295774596$. All other low vertices remain selected, and
each low omission affects only roots that are selected after the
specified repair. Hence this is a valid 2-dominating construction,
contradicting the critical value $|L|+19$.
\end{proof}

\subsection{Independent choices on matching edges}
\label{paired:independent-matching-section}

A fixed-size vertex sample has negative dependence across every pair
of vertices. A different useful distribution makes choices independent
on different matching edges, while permitting at most one selected
endpoint on an edge. The support indicators within an artificial
neighborhood are still not independent, and their joint distribution
must be evaluated.

Use Lemma~\ref{paired:artificial} with no explicitly credited support.
For a focal matched pair let $X,Y$ be its outside artificial support
counts, capped at two, and define
\begin{equation}\label{paired:gpair}
 g(x,y)=\begin{cases}2,&x=y=0,\\0,&x=y=2,\\1,&\text{otherwise}.
 \end{cases}
\end{equation}
If the focal pair is inactive, this is its pair-repair cost. If just
its first endpoint is selected, the second already has its partner
selected and needs repair exactly when $Y=0$; the other orientation
is symmetric.

We give an explicit finite classification that will also be used for
the dependent sampler in Section~\ref{paired:fixed-active-section}.
Color every outside vertex according to membership in the two outside
artificial neighborhoods by
\[
 0=(0,0),\quad A=(1,0),\quad B=(0,1),\quad C=(1,1).
\]
If the artificial types are $a,b$ and their outside intersection is
$c$, the nonzero color populations are
\begin{equation}\label{paired:populations}
 n_A=a-1-c,\quad n_B=b-1-c,\quad n_C=c,
 \qquad 0\le c<\min(a,b).
\end{equation}
If there is an unmatched singleton, assign it any available color
and remove it from that color population. The external matched edges
between nonzero colors are specified by six nonnegative integers
\begin{equation}\label{paired:pair-patterns}
 m_{AA},m_{AB},m_{AC},m_{BB},m_{BC},m_{CC},
\end{equation}
subject to
\begin{align*}
 2m_{AA}+m_{AB}+m_{AC}&\le n_A,\\
 2m_{BB}+m_{AB}+m_{BC}&\le n_B,\\
 2m_{CC}+m_{AC}+m_{BC}&\le n_C.
\end{align*}
Every remaining nonzero-colored vertex is paired with a zero. These
six counts, the common count, and the singleton color specify all
information needed for the capped distribution of $(X,Y)$. Allowing
unlimited zeros is a safe enlargement when edges are sampled
independently, since zero--zero edges have no effect on that
distribution. Later, when the number of active edges is fixed, this
relaxation will no longer be permissible.

\begin{proposition}[The even orders 34, 36, and 38]\label{paired:even-sampler}
There is no counterexample with $h\in\{34,36,38\}$.
\end{proposition}
\begin{proof}
The critical core has a perfect matching. Independently on each edge,
choose neither endpoint with probability $1/5$ and either specified
endpoint with probability $2/5$. Then $\mathbb E|S|=2h/5$, and the
focal pair's expected repair cost is
\begin{equation}\label{paired:even-W}
 W=\frac15\mathbb E\left[g(X,Y)+2\mathbf1_{X=0}+2\mathbf1_{Y=0}\right].
\end{equation}
For an external edge with endpoint color vectors $(i,j),(i',j')$,
its probability-generating factor is
\begin{equation}\label{paired:even-factor}
 \frac{1+2x^iy^j+2x^{i'}y^{j'}}5.
\end{equation}
Multiply these factors and cap each exponent at two. This gives an
exact nine-state distribution, without making independent-neighbor
assumptions.

For $3\le a\le b\le10$, enumeration of all common counts and all
six-tuples~\eqref{paired:pair-patterns} gives 12,556 patterns over
36 unordered type pairs. Evaluating~\eqref{paired:even-factor} and
\eqref{paired:even-W} on every pattern proves
$W\le\lambda_a+\lambda_b$ for the rational potentials
\begin{center}\small
\begin{tabular}{rr@{\qquad}rr}
\toprule
$a$&$10^9\lambda_a$&$a$&$10^9\lambda_a$\\\midrule
3&295986688&7&64396800\\
4&206041600&8&43481600\\
5&140654080&9&30569984\\
6&95296000&10&25344512\\\bottomrule
\end{tabular}
\end{center}
The finite procedure uses integer factor coefficients and exact
fractions. In addition, 234 small feasible pairings and their 27,324
sampling outcomes were directly checked against the generating
functions. These controls test implementation; the six-count
classification is what proves coverage.

The rank-profile multiplicities for $a=3,\ldots,10$ and the resulting
strict bounds are
\begin{center}\small
\begin{tabular}{rll}
\toprule
$h$&Profile multiplicities&Expected high cost\\\midrule
34&$(3,3,4,5,6,6,5,2)$&$33195294/1953125$\\
36&$(3,3,4,5,6,7,5,3)$&$6978444/390625$\\
38&$(3,3,4,5,6,7,6,4)$&$36563928/1953125$\\\bottomrule
\end{tabular}
\end{center}
Each bound is $2h/5+\sum n_a\lambda_a$. Their margins below $h/2$
are respectively $7831/1953125$, $52806/390625$, and
$545447/1953125$, all positive. Criticality is contradicted.
\end{proof}

\begin{lemma}[Odd independent-edge certificate]\label{paired:odd-table}
For an odd critical core whose rank-profile types lie in
$\{3,\ldots,10\}$, independently choose neither endpoint of a
near-perfect matching edge with probability $1/4$ or either specified
endpoint with probability $q=3/8$, and choose the unmatched vertex
independently with probability $q$. Pair repair has expected total
high cost at most
\begin{equation}\label{paired:odd-B}
 B_h=\frac{3h}{8}+\sum_a n_a\lambda_a+\zeta,
\end{equation}
where
\begin{center}\small
\begin{tabular}{rr@{\qquad}rr}
\toprule
$a$&$10^9\lambda_a$&$a$&$10^9\lambda_a$\\\midrule
3&341385154&7&88287832\\
4&248550156&8&61517419\\
5&177279951&9&43816963\\
6&125316621&10&32792315\\\bottomrule
\end{tabular}
\end{center}
and $\zeta=107833599/10^9$.
\end{lemma}
\begin{proof}
For a focal pair,
\begin{equation}\label{paired:odd-W}
 W=\frac18\mathbb E[2g(X,Y)+3\mathbf1_{X=0}+3\mathbf1_{Y=0}].
\end{equation}
The external-edge factor is
$(2+3x^iy^j+3x^{i'}y^{j'})/8$, and the singleton factor for color
$(i,j)$ is $(5+3x^iy^j)/8$. Enumerate all singleton colors and the
patterns~\eqref{paired:pair-patterns}, allowing unlimited zero fillers.
There are 37,022 patterns, again covering all 36 unordered type
pairs. Exact multiplication and~\eqref{paired:odd-W} verify
$W\le\lambda_a+\lambda_b$ on every pattern.

For the unmatched vertex of type $a$, suppose $j$ matching edges
contain two artificial neighbors. The probability-generating function
of its selected support is
\[
 \left(\frac{2+6x}{8}\right)^j
 \left(\frac{5+3x}{8}\right)^{a-2j},
 \qquad 0\le j\le\floor{a/2}.
\]
The repair probability is $5/8$ times the sum of coefficients of
degrees zero and one. Exact evaluation verifies it is at most
$\lambda_a+\zeta$ for every $a,j$. Summing the pair inequalities
and this one unmatched inequality proves~\eqref{paired:odd-B}.
As controls, 333 small pairings and all their 34,938 sampling outcomes
were compared with the generating-function values.
\end{proof}

The following special configuration supplies the extra saving needed
at odd orders. It also explains why an essential edge cannot simply
be treated as if its endpoints were sampled on distinct matching edges.

\begin{lemma}[Essential triangle edge and cubic repair]\label{paired:essential}
Suppose a critical odd-$h$ counterexample has a degree-two vertex
$x$ with adjacent high roots $u,v$. If every maximum matching of
$H$ uses $uv$, then some root, say $u$, has full degree three. Its
third neighbor $w$ is high, both $uw,vw$ are edges, and
$\theta_G(u)=2$. In particular $u$ has rank-profile type three.

For a sampler that never selects both endpoints of a matching edge,
the usual repair on the matched pair $uv$ can be replaced by a
repair of net cost one exactly when neither root nor $w$ is initially
selected, and net cost zero otherwise. Here net cost counts added
high vertices minus the possible omission of $x$.
\end{lemma}
\begin{proof}
If both roots had full degree at least four, deleting $uv$ would
preserve the high set. It would reduce $\mu(H)$ because every
maximum matching uses that edge. Since deleting edges cannot reduce
$\gamma_2$, Lemma~\ref{paired:matching-bound} would imply
\[
 \gamma_2(G)\le\gamma_2(G-uv)
 \le |L|+\mu(H-uv)<|L|+\mu(H)=\gamma_2(G),
\]
a contradiction. Thus one root is cubic. Its other neighbors are
$x,v,w$; $x$ is adjacent only to $u,v$. Having at least two triangles
forces $vw\in E(G)$ and gives exactly two triangles at $u$.

If $w$ were low, then $x,w$ would be two degree-two vertices with
the same high roots, contrary to Lemma~\ref{paired:odd-stability}.
There is also a direct cut argument: $u$ has no neighbor outside
this diamond, so with $v$ fixed a maximum cut puts $u,v$ together
and $x,w$ opposite, with four local crossing edges instead of three.
The side of $x,w$ has at least $k-1$ high vertices and can omit
both $x,w$, exceeding $k$ (or already exceeds $k$ using one if its
high side has size $k$). Hence $w$ is high.

If neither root is selected and $w$ is selected, keep $x$ and add
neither root; both roots are supported by $x,w$. If neither root
nor $w$ is selected, add both roots and omit $x$, with net cost one.
If exactly one root is selected, add the other and omit $x$, with
net cost zero. In every case an omitted $x$ has both roots selected;
other pair repairs remain valid. Distinct high counts make $u$ the
rank-one vertex, and~\eqref{paired:profile-definition} gives type three.
\end{proof}

\begin{proposition}[The order $h=37$]\label{paired:h37}
There is no counterexample with exactly thirty-seven high vertices.
\end{proposition}
\begin{proof}
The rank profile is $(2,3,4,5,6,7,5,5)$, so
Lemma~\ref{paired:odd-table} gives $B_{37}=18.090352599$; the decimal
is exact, with denominator $10^9$. A mandatory low triangle exists.
If it is pendant, omit one private vertex when its root is sampled,
giving $B_{37}-3/8=17.715352599<18$.

Otherwise take an edge attachment $x$ on $uv$. If a maximum matching
avoids $uv$, use it: its two roots lie on distinct matching edges
or one is unmatched, and their selections are independent with
probability $q=3/8$. Omit $x$ when both are selected, giving
$B_{37}-9/64=17.949727599<18$.

If $uv$ belongs to every maximum matching, use
Lemma~\ref{paired:essential}. Focal inactivity has probability $1/4$,
and $w$, lying on another edge or unmatched, is independently
unselected with probability $5/8$. The special net repair costs
$5/32$ in expectation. Replacing the generic pair potential gives
\[
 B_{37}-\lambda_3-\min_a\lambda_a+5/32=17.87242513<18.
\]
Each case contradicts the critical value $|L|+18$.
\end{proof}

\subsection{Exactly thirteen active edges: the exclusion of \texorpdfstring{$h=35$}{h=35}}
\label{paired:fixed-active-section}

The independent sampler does not exclude $h=35$ without an order
assumption. The last round replaces it by a finite-capacity sampler
with an exact number of active matching edges. We spell out its
dependence rather than importing the previous probability calculation.

\begin{proposition}[The order $h=35$]\label{paired:h35}
Every connected graph with exactly thirty-five high vertices satisfies
\eqref{paired:target}.
\end{proposition}
\begin{proof}
Assume a counterexample. Its core has a near-perfect matching of
17 edges and one unmatched vertex, and its rank profile is
\[
 (n_3,\ldots,n_{10})=(2,3,4,5,6,7,5,3).
\]
Choose exactly 13 of the 17 matching edges uniformly, and on each
active edge choose one endpoint independently and uniformly. Select
the unmatched vertex independently with probability $\vartheta=13/32$.
Then matched vertices have marginal probability $13/34$, the
singleton has marginal $13/32$, and
$\mathbb E|S|=13+13/32$.

Use the artificial sets and pair repair of
Lemma~\ref{paired:artificial}. For a focal matching edge, inactivity
has probability $4/17$ and leaves exactly 13 active edges among the
16 external edges. Activity has probability $13/17$ and leaves
exactly 12 active external edges. Its expected repair cost is
\begin{equation}\label{paired:fixed-W}
 W=\frac4{17}\mathbb E_{13}g(X,Y)
   +\frac{13}{34}\bigl(\Pr_{12}(X=0)+\Pr_{12}(Y=0)\bigr).
\end{equation}
The subscripts are conditioning on the number of active external
edges. These two distributions are different.

Use colors~\eqref{paired:populations}, remove the chosen singleton
color, and enumerate the six counts~\eqref{paired:pair-patterns}.
Pair unused nonzero colors with zeros. If these nonzero external
edges number $E$, impose $E\le16$ and set $Z=16-E$. The remaining
$Z$ edges are zero--zero edges. Unlimited extra zeros would change
the conditional sampling law and are not allowed here.

For the $E$ nonzero edges form the integer polynomial
\begin{equation}\label{paired:fixed-poly}
 P(z,x,y)=\prod_{\{u,v\}}
      \left[2+z(x^{i_u}y^{j_u}+x^{i_v}y^{j_v})\right],
\end{equation}
capping the $x,y$ exponents at two. Write $P_{j,x,y}$ for its
coefficients. Conditional on $r$ active external edges, the capped
support distribution before the singleton is
\begin{equation}\label{paired:fixed-coef}
 \Pr_r(X=x,Y=y)=
 \frac{\sum_j P_{j,x,y}\binom Z{r-j}}{2^E\binom{16}{r}}.
\end{equation}
To derive this, assign two possible orientations also to inactive
edges, so all activation subsets and orientations have the common
denominator $2^{16}\binom{16}{r}$. Each zero--zero edge contributes
$2(1+z)$; cancelling their factor $2^Z$ gives precisely
\eqref{paired:fixed-coef}. Finally mix the singleton color with
weights $19/32$ and $13/32$, capping again after its possible shift.

For all $3\le a\le b\le10$, all common counts, all available
singleton colors, and all feasible six-count patterns, exact
evaluation of~\eqref{paired:fixed-W}--\eqref{paired:fixed-coef} gives
$W\le\lambda_a+\lambda_b$, with
\begin{center}\small
\begin{tabular}{rr@{\qquad}rr}
\toprule
$a$&$10^6\lambda_a$&$a$&$10^6\lambda_a$\\\midrule
3&333547&7&71002\\
4&234704&8&45945\\
5&161174&9&30585\\
6&108489&10&22040\\\bottomrule
\end{tabular}
\end{center}
There are 37,014 patterns and 36 type pairs. Exactly eight patterns
from the unlimited-zero classification exceed the finite capacity
and are excluded. This is an exact coverage restriction, not a
probabilistic approximation.

For an unmatched vertex of type $a$, let $j$ matching edges have
two marked artificial neighbors. The 17 edges comprise $j$ doubly
marked, $a-2j$ singly marked, and $17-a+j$ zero--zero edges. Apply
the same polynomial method conditional on 13 active edges among
all 17. Multiplying its support-zero-or-one probability by $19/32$
gives a repair cost at most $\lambda_a+\zeta$, where
$\zeta=89937/10^6$. All $0\le j\le\floor{a/2}$ are checked.
Summing over the near-perfect matching gives
\begin{equation}\label{paired:B35}
 B=13+\frac{13}{32}+\sum_a n_a\lambda_a+\zeta
 =\frac{8510603}{500000}=17.021206.
\end{equation}
The finite certificate also compares 333 small feasible patterns
with direct enumeration of 17,664 fixed-active outcomes.

It remains to use the mandatory low triangle. In case P, omit one
private vertex when the root is initially selected. Its selection
probability is at least $13/34$, including an unmatched root, and
\[
 B-13/34=141430251/8500000<17.
\]
In case Q choose a private vertex $x$ on an edge $uv$. If a maximum
matching avoids $uv$, use that matching. When both roots are matched
on distinct edges their joint selection probability is
\[
 \frac{13\cdot12}{17\cdot16}\cdot\frac14=\frac{39}{272}.
\]
If one root is unmatched, it is $(13/32)(13/34)$, which is larger.
Omitting $x$ when both roots are selected therefore gives
\[
 B-39/272=143461501/8500000<17.
\]

If every maximum matching contains $uv$, invoke
Lemma~\ref{paired:essential}. Conditional on focal inactivity, a
third vertex $w$ on another matching edge is selected with probability
$13/(2\cdot16)=13/32$. If $w$ is unmatched, its independently chosen
probability is the same, by design. The special net repair cost is
therefore $(4/17)(19/32)=19/136$. Replacing the generic focal-pair
potential in~\eqref{paired:B35} gives
\[
 B-\lambda_3-\min_a\lambda_a+19/136
   =285690523/17000000<17.
\]
Every branch constructs a 2-dominating set with expected size below
$|L|+17$, contradicting criticality. No order bound on $G$ was used.
\end{proof}

Propositions~\ref{paired:large44}, \ref{paired:h40}, \ref{paired:h39},
\ref{paired:even-sampler}, \ref{paired:h37}, and~\ref{paired:h35},
together with the exclusions of 41--43, now exclude every $h\ge34$.

\subsection{Leaf trimming and matching resilience}
\label{paired:trimming-section}

\begin{lemma}[One-leaf representative]\label{paired:trim}
Every counterexample has a connected counterexample representative
with the same high-degree order and the same set of triangle counts,
and with at most one leaf. If $n_2$ is the number of degree-two
vertices in that representative, then
\begin{equation}\label{paired:rough-order}
 n_2\le\ceil{h/2}+1,
 \qquad n\le h+\ceil{h/2}+2.
\end{equation}
\end{lemma}
\begin{proof}
Deleting a leaf cannot increase $k$: every omission set in the smaller
graph lifts to the larger by leaving the deleted leaf selected.
Therefore a leaf deletion preserving the high set and $T(G)$
preserves counterexample status. The cut bound and critical equality
also force $k$ to remain unchanged.

Every high vertex with triangle count at least two has at least
three nonleaf neighbors. If $h$ is odd, all high vertices have this
property, so their incident leaves may be deleted freely. There is
at most one leaf adjacent to a degree-two vertex, because $G[L]$
has at most one edge. Retain such a leaf if it exists. Otherwise
retain one high-adjacent leaf only when needed to preserve the
triangle count zero. This leaves at most one leaf, preserves every
triangle count, and does not disconnect the remaining graph.

For even $h$, there is at most one exceptional high vertex $v$
without the count-at-least-two guarantee. By even stability every
leaf is high-adjacent. Delete leaves at the other high vertices,
retaining one if needed for the value zero. At $v$, retain only
enough leaves to keep its degree at least three; any retained leaf
there already preserves zero. Since $h\ge11$ by
Proposition~\ref{paired:small}, connectedness gives $v$ at least one
nonleaf neighbor. The only possible need for two retained leaves
is that $v$ has exactly one nonleaf neighbor. Then $v$ lies in no
triangle. Delete $v$ and its two leaves. All remaining $h-1$ high
vertices still have their at least two triangles and remain high.
A maximum cut in the remaining graph, which need not be connected,
gives an omission set of at least $\ceil{(h-1)/2}=h/2$ high vertices.
Lift it and also omit $v$, retaining its two leaves. Existing omitted
vertices retain their supporting neighbors from the smaller graph,
and $v$ has the two selected leaves. This gives $h/2+1>k$ omissions,
a contradiction. Thus at most one leaf is needed in even order too.

At least $n_2-1$ degree-two vertices are independent, by stability
of $G[L]$. An independent set of degree-two vertices is omittable,
so $n_2-1\le k=\ceil{h/2}$. Adding at most one leaf proves
\eqref{paired:rough-order}.
\end{proof}

\begin{lemma}[Matching resilience]\label{paired:resilience}
In a critical counterexample, let $F\subseteq E(H)$ satisfy
$d_G(v)-d_F(v)\ge3$ for every high vertex. Then
$\mu(H-F)=\mu(H)$.
\end{lemma}
\begin{proof}
The high and low sets of $G-F$ are unchanged. Edge deletion cannot
decrease the 2-domination number. Using the equality in
Lemma~\ref{paired:matching-bound},
\[
 |L|+\mu(H)=\gamma_2(G)\le\gamma_2(G-F)
 \le |L|+\mu(H-F)\le |L|+\mu(H).
\]
All inequalities must be equalities. The matching bound applies even
when deletion disconnects the graph.
\end{proof}

\subsection{Simultaneous low omissions and the order-37 representative}
\label{paired:order37-section}

It remains to prove the representative bound. Apply
Lemma~\ref{paired:trim} and work in the resulting connected
counterexample. The completed high-order exclusions give
$11\le h\le33$. If $h\le23$, then~\eqref{paired:rough-order} already
gives $n\le37$. Assume for contradiction that $n\ge38$; since
there is at most one leaf,
\begin{equation}\label{paired:n2-lower}
 n_2\ge37-h.
\end{equation}

\begin{lemma}[Uniform-sample exclusion for $24\le h\le32$]
\label{paired:order37-uniform}
A trimmed counterexample of order at least 38 cannot have
$24\le h\le32$.
\end{lemma}
\begin{proof}
Use the full-degree rank profile~\eqref{paired:profile-definition},
the fixed-size paired alteration~\eqref{paired:uniform-bound} with
$s=13$, and the exact weights in the supplied
\nolinkurl{R21/order37_uniform_input.json}. They are rational weights
for each degree type; an odd-order row has its unmatched correction
as the last entry. For $h=24,25$ all weights are zero. For every
other row they are the indicated exact weights reused from the
earlier order-39 certificate. Verify each inequality
\eqref{paired:uniform-dual} for every type pair and overlap; there
are 1,270 such inequalities over the nine orders.

Put $q_h=156/[h(h-1)]$ and $p_h=13/h$. If $G[L]$ is independent,
omit every degree-two vertex when both its roots are initially
sampled, earning $n_2q_h$ in expectation. If $G[L]$ has one edge,
choose a degree-two endpoint $x$ and keep its low neighbor selected.
Omit $x$ when its one high neighbor is sampled, earning $p_h$.
At most two degree-two vertices are on that low edge, and all
others are ordinary and supply $q_h$ each. Since this case has
odd $h\ge25$, $p_h\ge2q_h$. Therefore the total credit is at least
$n_2q_h$ in both cases. All omissions affect only high roots that
were initially selected, so they can be taken simultaneously after
the pair repairs.

Using~\eqref{paired:n2-lower}, subtract $(37-h)q_h$ from the high
bound. The exact positive margins below $\floor{h/2}$ are
\begin{center}\small
\begin{tabular}{rl@{\qquad}rl}
\toprule
$h$&Margin&$h$&Margin\\\midrule
24&$1968779/1248072$&25&$1274117/1300075$\\
26&$10702053/10400600$&27&$11190077/20058300$\\
28&$235339/367080$&29&$2013666/7540435$\\
30&$2452616/6653325$&31&$2038426/29464725$\\
32&$57723899/347373600$&&\\\bottomrule
\end{tabular}
\end{center}
For example the $h=24$ high bound is $17593393/1248072$, and its
low credit is $169/46$. The checker reconstructs every rank profile,
evaluates all pair inequalities, and compares every displayed margin
as a fraction; it does not require a graph catalogue or an optimizer.
Every row contradicts criticality.
\end{proof}

Only $h=33$ remains. An ordinary degree-two vertex means one with
two high neighbors; these neighbors need not be adjacent. Let $Q$
be the graph of those endpoint pairs which are edges of $H$.
By odd stability, $Q$ is a simple matching with at most one two-edge
fork. If an ordinary vertex's endpoint pair is not in the chosen
matching, the odd sampler of Lemma~\ref{paired:odd-table} omits it
with probability $q^2=9/64$ when both roots are initially selected.
The purpose of resilience is to choose one matching for all these
credits at once.

\begin{lemma}[A fork with two preselected neighbors]
\label{paired:prepaid-fork}
In the odd independent-edge sampler, suppose a full-degree-four high
vertex $u$ is incident to two edges of $Q$. Its profile type is three,
and its two private low neighbors allow the high cost bound $B_h$
to be reduced by at least
\[
 \Delta=321442183/10^9.
\]
At least one of those two private vertices simultaneously provides
an additional omission credit $q^2=9/64$.
\end{lemma}
\begin{proof}
The four neighbors of $u$ are its two private degree-two vertices
and two distinct high roots. Its triangle count is two or three,
so distinctness puts it at rank one or two and gives profile type
three for the orders under consideration. While $u$ is unselected,
both private low neighbors remain selected, and no repair at $u$
is necessary.

If $u$ is matched to a vertex of type $a$, let $Y$ be the partner's
outside artificial support. The focal pair costs at most
$\tfrac14\Pr(Y\le1)+\tfrac38\Pr(Y=0)$. Among its $a-1$ outside
artificial neighbors, let $j$ external matching edges contain two
marked vertices and $m=a-1-2j$ contain a single marked vertex;
a possible marked singleton has the same selection probability
$q=3/8$ as a singly marked edge. With $p=3/4$ put
\[
 P_0=(1-p)^j(1-q)^m,
 \quad P_1=P_0\left(\frac{jp}{1-p}+\frac{mq}{1-q}\right),
\]
\[
 C_a=\max_{0\le j\le\floor{(a-1)/2}}
       \left[\tfrac14(P_0+P_1)+\tfrac38P_0\right].
\]
These are exact support-zero and support-one probabilities. Direct
evaluation for $3\le a\le10$ proves
$\lambda_3+\lambda_a-C_a\ge\Delta$. If $u$ is unmatched, its
whole repair potential $\lambda_3+\zeta$ can be removed, an even
larger saving. At most one of its two incident Q edges belongs
to the matching, so at least one private vertex has roots on distinct
sampling units and earns $q^2$. It is omitted only when $u$ was
initially selected, so this omission is compatible with the use of
both low neighbors to support an unselected $u$.
\end{proof}

\begin{lemma}[The last order case]\label{paired:order37-h33}
A trimmed counterexample of order at least 38 cannot have $h=33$.
\end{lemma}
\begin{proof}
The odd independent-edge certificate gives the exact high bound
$B_{33}=16.430458235$. If the configuration in
Lemma~\ref{paired:prepaid-fork} occurs, its upper bound becomes
\[
 B_{33}-\Delta-q^2=15.968391052<16,
\]
which already contradicts criticality.

Suppose there is no such full-degree-four fork root. A cubic vertex
incident with a Q edge has exactly two triangles, by the same
neighborhood argument as in Lemma~\ref{paired:essential}. There is
at most one such cubic vertex, and it is incident with only one
Q edge $e=uv$. Indeed a second incident private attachment would
force two low vertices with the same endpoint pair, contrary to
stability. If no such cubic vertex exists, there is no exceptional
edge $e$.

Delete every edge of $Q$ except the possible $e$. An endpoint losing
one edge has full degree at least four. A fork center losing two
edges has full degree at least five, because its degree-four case
was already excluded. All high vertices remain high. By
Lemma~\ref{paired:resilience}, a near-perfect matching exists that
avoids all these deleted edges.

If this matching also avoids $e$, every ordinary degree-two vertex
earns $q^2$. If a low--low edge exists, omit one of its degree-two
endpoints when its high neighbor is sampled, earning $q$. At most
two degree-two vertices lie on the edge, and $q\ge2q^2$. Thus the
simultaneous total low credit is at least $n_2q^2$, including the
low-edge case. Using $n_2\ge4$ gives
\[
 B_{33}-4q^2=15.867958235<16.
\]

If the matching contains $e$, use the cubic repair of
Lemma~\ref{paired:essential}. The lemma's local repair only requires
the cubic marked configuration; its essential-edge hypothesis was
used to force that configuration, which is already known here.
Its expected net cost is $5/32$. Replace the pair potential and
do not count its private vertex again. All other ordinary vertices
still earn $q^2$, and the possible low edge is handled as above.
Other low omissions affect only initially selected high roots, so
they are compatible with the special repair. The resulting bound is
\[
 B_{33}-\lambda_3-\min_a\lambda_a+5/32-3q^2
 =15.790655766<16.
\]
All branches contradict the critical 2-domination number.
\end{proof}

\begin{proof}[Completion of Theorem~\ref{paired:terminal}]
The small-core certificate excludes $h\le10$. The independent
alteration, boundary rank arguments, uniform matching repair,
independent-edge sampling, and fixed-active sampling exclude every
$h\ge34$. Thus a counterexample has $11\le h\le33$ and satisfies
the critical equalities. Trim it using Lemma~\ref{paired:trim}. For
$h\le23$ its order is already at most 37. If its order exceeded
37, Lemmas~\ref{paired:order37-uniform} and
\ref{paired:order37-h33} would exclude every remaining possible
$h$. The trimmed graph is therefore the required connected
representative, with unchanged $h$ and triangle-count set.
\end{proof}

\begin{remark}[What has and has not been closed]\label{paired:scope}
The quantified conclusion is that existence of any counterexample
implies existence of a connected counterexample on at most 37 vertices.
The argument does not bound the order of every counterexample. It
does not enumerate all graphs on at most 37 vertices, and it does
not prove~\eqref{paired:target} in the remaining range $11\le h\le33$.
The connectivity conclusion at $h=11$ and the even cubic filters
are additional necessary conditions within that range.
\end{remark}

\subsection{Certificate structure, reproduction, and provenance}
\label{paired:repro}

The companion directory is \nolinkurl{verification/paired_384b/};
its subdirectories \nolinkurl{R08}, \nolinkurl{R09}, and so on retain
the original round numbers solely to identify provenance~\cite{paired:archive}. The proofs
above form one dependency chain; a theorem is not counted again
when a later round strengthens its numerical boundary. The main
sources are the 384b manuscripts of Rounds 8--11 and 14--22, together
with the initial Round-5 reduction and the Round-10 cubic supplement.
The consolidated audit JSON records the exact source paths.

The small-core data are masks and relabeling certificates. The seven-
vertex catalogue is certified by enumerating all permutations and
marking all $2^{21}$ masks. The eight-vertex coverage file assigns
an explicit relabeling to each seven-mask extension, and the nine-
vertex coverage file does the same for each eight-mask extension.
The eight-vertex records use four-byte little-endian words with three
bits per vertex image; the nine-vertex records use eight-byte
little-endian words with four bits per vertex image. A coverage verifier checks that
each record encodes a permutation, reconstructs the relabeled edge
mask, and tests membership in the target list. It also checks record
counts and exhaustion of the file. Coverage therefore does not
depend on trusting a graph-isomorphism or canonical-labeling package.
Duplicates in the extension families are permitted and counted.

The probability certificates have three distinct forms. The rank
certificates evaluate finite binomial sums and all explicitly listed
rank placements. The fixed-size pair certificates check every
type-pair and overlap inequality and then sum the resulting rational
weights. The matching-edge certificates enumerate the six color-pair
counts and, in odd order, all singleton colors; they multiply capped
integer generating functions. The fixed-active certificate also
imposes the actual 16-edge external capacity and conditions on the
proper active-edge count. The extra physical controls verify selected
repair implementations on sampled finite graphs, but are not used
as exhaustive graph coverage. In particular, the 600 controls and
46,800 weighted outcomes in
\nolinkurl{R21/check_joint_low_repairs.py} supplement the universal
compatibility proofs above.

For this consolidation 21 checkers were rerun, all with exit status
zero. They include all finite-kernel checks through $h=10$, all
three coverage checks, both rank exclusions, the uniform matching
certificates, all three matching-edge distributions, the order-37
certificate, and the simultaneous-repair controls. The runner
\nolinkurl{run_verification.py} compiles the C++ verifiers with
\nolinkurl{g++ -O3 -std=c++17}, runs Python standard-library checkers,
and saves each command, exit status, elapsed time, and output log in
\nolinkurl{rerun_results.json} and \nolinkurl{rerun_logs/}.
These runs use copies and do not alter archival source files.
They reproduce exact finite certificates; discovery sweeps and
optimization searches are not verification dependencies.

The original series successively obtained larger representative
bounds, including 217, 67, 64, 43, 42, 40, and 39 before the bound
37. Those numerical consequences are subsumed by
Theorem~\ref{paired:terminal}. Their reusable arguments---leaf
trimming, low omission credit, rank constraints, paired repair, and
matching resilience---are proved here once. The order-37 certificate
also contains its original $h=35$ calculation, with the three
conditional bounds $16.993517969$, $16.9162155$, and $16.812700786$;
these exclude $h=35$ representatives of order at least 38. The
stronger Proposition~\ref{paired:h35} now excludes that high order
without a total-order assumption.

The following full SHA-256 digests identify the data needed for
catalogue coverage and the only external weight input to the final
order reduction. The verifier digests and the manifest binding the
remaining certificate outputs are listed immediately afterward. The complete source manifest
also records byte sizes and archival paths, while the consolidated
audit records which statements are finite certificates and which
are universal proofs.

\begingroup\footnotesize
\setlength{\tabcolsep}{4pt}
\begin{longtable}{p{0.40\textwidth}p{0.55\textwidth}}
\caption{384b certificate input and verifier digests. Paths are relative to
\nolinkurl{verification/paired_384b/}. Each two-line hexadecimal string
is one complete SHA-256 digest.}\label{paired:hash-table}\\
\toprule
File & SHA-256\\\midrule\endfirsthead
\toprule File & SHA-256\\\midrule\endhead
\bottomrule\endfoot
\nolinkurl{R11/core7_masks.txt} & \texttt{de8140ae11b1d575d725febccb6140ae}\newline\texttt{faea77bbce7880d78283bb63b9e5486a}\\[3pt]
\nolinkurl{R11/core8_masks.txt} & \texttt{865fd27c2bb07b21c11f57a5d71d0e20}\newline\texttt{ebcdd6beea937a043751f70cc89f456f}\\[3pt]
\nolinkurl{R11/core9_masks.txt} & \texttt{018a48bfc1548601c0c71c539d4fe25d}\newline\texttt{809e05c6280ee9d3c91c1f50ada24d01}\\[3pt]
\nolinkurl{R11/core8_cover_permutations.bin} & \texttt{ec105831fab7638d77fae72b4d9d2e18}\newline\texttt{1246ba3d73e4f852924e9e6521a1c781}\\[3pt]
\nolinkurl{R11/core9_cover_permutations.bin} & \texttt{2c6f017c01a45645f93b71fa8f13dca8}\newline\texttt{c4eccc5fd22cca76b20e10b3b3cf7549}\\[3pt]
\nolinkurl{R21/order37_uniform_input.json} & \texttt{f97ee255bb05f749a9f2aeef7a83ae6c}\newline\texttt{a99b9114f3e97355c7adc311c387cd81}\\[3pt]
\nolinkurl{R08/check_384b_alteration.py} & \texttt{11d05965d783e89fd2d15fce3488486d}\newline\texttt{64b53db3c001aa38b1627a46ee16dac6}\\[3pt]
\nolinkurl{R08/check_384b_hypergeometric.py} & \texttt{796f885ac58b0e52e6c2db05f13b895f}\newline\texttt{ca5276096029fdfecea0c6c51542438a}\\[3pt]
\nolinkurl{R08/check_384b_kernels.py} & \texttt{92dafc0831fd9f5cb3aae50739593b0d}\newline\texttt{9ea19dd4da6e61a6441b76cbc32f5017}\\[3pt]
\nolinkurl{R08/verify_384b_h7.py} & \texttt{74d894388f8fe04a1478d07a6fc33427}\newline\texttt{51a654e902f432c772c34bd8762eaf51}\\[3pt]
\nolinkurl{R08/verify_384b_h8.cpp} & \texttt{d674b73d21333315aefb6602c5c4f948}\newline\texttt{4d60355dbfb1060f8c8a02651b1a07f6}\\[3pt]
\nolinkurl{R09/verify_384b_h9.cpp} & \texttt{d5bb333c69cd65d778744da1ac0c176d}\newline\texttt{10b9a1714f350da6a8eb2832ea16c3a7}\\[3pt]
\nolinkurl{R10/check_low_omission_credit.py} & \texttt{ebe2632c4606e40093dbe11e2bac3725}\newline\texttt{239cb24ae1eeaa517f4f42bf2ba1d95f}\\[3pt]
\nolinkurl{R10/check_order43.py} & \texttt{578320b635bc69f4d7cd4a0f1791db08}\newline\texttt{fafcf0c4229860fdfcf56e69dd18a33b}\\[3pt]
\nolinkurl{R11/certify_core7_coverage.cpp} & \texttt{6824967f10e2d451189561b79c02ed86}\newline\texttt{0c224ed209c741136cbfe3b296973f43}\\[3pt]
\nolinkurl{R11/verify_384b_h10.cpp} & \texttt{5fc5a41e4ff03ffde85bff00d1cffaa5}\newline\texttt{08064274de7ea95735184ae36f8e68a6}\\[3pt]
\nolinkurl{R11/verify_core8_cover.py} & \texttt{4aae049fd7cc29b83ffef992b2e9bc34}\newline\texttt{571736ef2b18f2633f0254bf023aec36}\\[3pt]
\nolinkurl{R11/verify_core9_cover.cpp} & \texttt{72cc50a60d005a531f82c888be066d98}\newline\texttt{a238bff750f11b93de385492a1693415}\\[3pt]
\nolinkurl{R15/check_rank_h42.py} & \texttt{f441532e3e0cfafc9b8ac7991339dc86}\newline\texttt{7a49e9e6a39e0309eb733645ba44ea5e}\\[3pt]
\nolinkurl{R16/check_rank_h41.py} & \texttt{04fc4d94918fa4065e8fcb476c3c277f}\newline\texttt{9ce5a63cf72c26cdb8ec5337e703b0bf}\\[3pt]
\nolinkurl{R17/check_matching_h40.py} & \texttt{4c53641c6efce77949522a5052eea980}\newline\texttt{883933ce4feff214965a5aba0c80198e}\\[3pt]
\nolinkurl{R18/check_matching_h39.py} & \texttt{b462f1b60db5d65594178c9ee2e33fb9}\newline\texttt{125434c1be96d9b32a84df6d67c0009d}\\[3pt]
\nolinkurl{R19/check_paired_sampling_even.py} & \texttt{69c7bb48ec2a368f0249149211df33ca}\newline\texttt{a7ab8ba1a509d17f1c256ed69934e05e}\\[3pt]
\nolinkurl{R20/check_paired_sampling_h37.py} & \texttt{de00f742fd056f53e3130759243f15c7}\newline\texttt{687885fa0557e8eab20bcb43a09e1d78}\\[3pt]
\nolinkurl{R21/check_joint_low_repairs.py} & \texttt{249a18c5a7f1819bf3f6614922058bf8}\newline\texttt{41ed5e722b22057fd88341c0c4227112}\\[3pt]
\nolinkurl{R21/check_order37.py} & \texttt{65770c371aa08d4182fffe66312e3131}\newline\texttt{6a2c68d5efc036242aa836968cd79e70}\\[3pt]
\nolinkurl{R22/check_fixed_active_h35.py} & \texttt{3501301c351b1e86732ee7f902502c4e}\newline\texttt{b701d65dcd9d049f2300016557d76baa}\\[3pt]
\nolinkurl{artifact_manifest.json} & \texttt{9cf90736538db5a592aaf9ac3ff27612}\newline\texttt{3415bb521ed089ec954b27bf6dc95b0a}\\[3pt]
\nolinkurl{source_manifest.json} & \texttt{f57cca0e02f5956b406969c91418bbb3}\newline\texttt{045a8c300ec35ca483571c92a7ab65a5}\\[3pt]
\nolinkurl{rerun_results.json} & \texttt{928bf4bf7366956b19c2558827229320}\newline\texttt{9aba2bf23d8f5660f6205fb8d80e5c27}\\[3pt]
\nolinkurl{run_verification.py} & \texttt{b04bb7f1e1480cde9411ffb490faf2dd}\newline\texttt{db68e2cd8ae16dca4a25ffb53ca2207f}\\[3pt]
\end{longtable}
\endgroup

The artifact manifest identifies every bundled source, data, and exact
output JSON by full digest. The source manifest separately identifies
the unchanged archival inputs, so a regenerated output can be compared
with its source without conflating the two. Both are bound by their
full digests in Table~\ref{paired:hash-table}. The audit file
\nolinkurl{paired_384b_audit.json} maps each theorem to its mathematical
and finite dependencies, and explicitly distinguishes this rerun from
independent-audit records already present in the supplied archive.


% Integrated source fragment: split_199.tex
\section{Traceability in split graphs: Conjecture 199}
\label{sec:split-small}
This and the following two sections consolidate the split-graph results
of Rounds 17--22. All path conclusions refer to a Hamiltonian path.
The maximum induced-tree order is written $\tau$ here to distinguish
it from the maximum induced-forest order used elsewhere.


\subsection{Definitions and the precise scope}
All graphs are finite and simple. A graph is \emph{traceable} if it has a
Hamiltonian path, that is, a path containing every vertex. Write $\alpha(G)$
for the independence number and $\kappa(G)$ for vertex connectivity.
The order of a graph is its number of vertices. In particular,
$\tau(G)$ is the maximum \emph{vertex order} of a connected acyclic
\emph{induced} subgraph; it is not an induced-forest parameter.

A \emph{split graph} has a partition $V(G)=C\mathbin{\dot\cup}I$ in which
$C$ is a clique and $I$ is independent. WOWII Conjecture 199, dated
January 12, 2006 in the archived list \cite{src:catalogue}, asserts that every
connected graph of order greater than one satisfying
\begin{equation}\label{split17:eq:original}
 \tau(G)-2\le\kappa(G)
\end{equation}
is traceable. We prove the following restricted statement.

\begin{theorem}\label{split17:thm:main}
Let $G$ be a connected split graph of order greater than one satisfying
\eqref{split17:eq:original}. Then $G$ is traceable in either of these cases:
\begin{enumerate}
\item $\kappa(G)\le3$;
\item $G$ is noncomplete, $k=\kappa(G)\ge2$, and some split partition
$V(G)=C\mathbin{\dot\cup}I$ with $I$ maximum satisfies
\[
 k\ge2m-2,\qquad m=\max_{c\in C}|N(c)\cap I|.
\]
\end{enumerate}
Complete graphs are traceable without either additional restriction.
\end{theorem}

The proof of the second part is in Section~\ref{split17:sec:small-m}. The first
part follows from Sections~\ref{split17:sec:k1}--\ref{split17:sec:k3}. Neither part asserts
that an arbitrary capacity-two factor is acyclic or can be patched into
a spanning path.

\subsection{Split partitions and elementary tools}
\begin{lemma}\label{split17:lem:setup}
A split graph has a split partition $C\mathbin{\dot\cup}I$ with $I$ a
maximum independent set. In such a partition every $c\in C$ has a
neighbor in $I$. If the graph is connected and noncomplete, then
\[
 \kappa(G)=\min_{i\in I}d_C(i).
\]
Writing $k=\kappa(G)$ and $A_c=N(c)\cap I$, condition
\eqref{split17:eq:original} implies
\begin{equation}\label{split17:eq:columns}
 d_C(i)\ge k,\qquad 1\le|A_c|\le k+1,\qquad
 |A_c\triangle A_d|\le k\quad(c\ne d).
\end{equation}
\end{lemma}
\begin{proof}
Start with any split partition. If a clique vertex has no neighbor in the
independent side, move it to that side. Every remaining clique vertex is
adjacent to the moved vertex, so no further move is needed. When each
clique vertex has an independent-side neighbor, an independent set using
one clique vertex omits at least one independent-side vertex; it therefore
has size at most $|I|$. An independent set uses at most one clique vertex,
so $I$ is maximum.

For a connected noncomplete graph, $|I|\ge2$. Put
$k_0=\min_{i\in I}d_C(i)$. Removing the neighbors of a vertex attaining this
minimum isolates it from another vertex of $I$, so $\kappa(G)\le k_0$.
Deleting fewer than $k_0$ vertices leaves a nonempty clique, and every
remaining independent vertex retains a neighbor in it. The remaining
graph is connected, proving equality.

The induced star on $\{c\}\cup A_c$ has order $1+|A_c|$. For distinct
$c,d\in C$, the induced graph on
$\{c,d\}\cup(A_c\triangle A_d)$ is a double star of order
$2+|A_c\triangle A_d|$. Common independent neighbors are omitted. Applying
\eqref{split17:eq:original} to these two trees proves \eqref{split17:eq:columns}.
\end{proof}

We use the following precise classical input. Chv\'atal and Erd\H{o}s
\cite[Theorem 2, p.~112]{src:ce} proved that an $s$-connected graph with no
independent set of size $s+2$ has a Hamiltonian path. Thus
\begin{equation}\label{split17:eq:CE}
 \alpha(G)\le\kappa(G)+1\quad\Longrightarrow\quad
 G\text{ is traceable}.
\end{equation}
Their proof adjoins a universal vertex, applies their Hamiltonian-cycle
theorem, and then deletes that vertex. We invoke \eqref{split17:eq:CE} only when
the size of our maximum independent side is explicitly at most $k+1$.

\begin{lemma}\label{split17:lem:insertions}
Let $C\mathbin{\dot\cup}I$ be a split partition with $I$ maximum.
The following path operations are valid.
\begin{enumerate}
\item A path containing every vertex of $I$ can be extended to a spanning
path by inserting the unused clique vertices.
\item If a path has an edge $ac$, with $a\in I$, $c\in C$, and $x\in I$
is outside the path, let $p,q$ be distinct unused neighbors of $x$ in $C$.
If $ap$ is an edge, replace $ac$ by $a,p,x,q,c$.
\item If the path has an edge $cd$ within $C$, and $p,q$ are distinct unused
clique neighbors of a vertex $x\in I$ outside the path, replace $cd$ by
$c,p,x,q,d$. The resulting path still has an edge within $C$.
\end{enumerate}
\end{lemma}
\begin{proof}
For the first operation, an unused $d\in C$ has a neighbor $a\in I$ already
on the path. Replace an incident path edge $ac$, $c\in C$, by $a,d,c$.
The paths to which we apply this operation have an edge incident with
each independent vertex. If the path initially consists of a single
independent vertex, first append one of its clique neighbors.
The other operations follow by checking the displayed incidences and the
clique edges. All inserted vertices are unused, so simplicity is preserved.
The edge $cp$ in the third operation is still a clique edge.
\end{proof}

\begin{lemma}[Classical path closure]\label{split17:lem:closure}
If $u,v$ are nonadjacent vertices of an $n$-vertex graph $F$ and
$d_F(u)+d_F(v)\ge n-1$, then $F$ is traceable if and only if $F+uv$ is
traceable.
\end{lemma}
\begin{proof}
Only the reverse implication needs proof. Add a universal vertex $w$ to
$F$, obtaining $H$ of order $N=n+1$. A spanning path of $F+uv$ gives a
Hamiltonian cycle of $H+uv$, and $d_H(u)+d_H(v)\ge N$. If this cycle does
not use $uv$, delete $w$ and finish. Otherwise deleting $uv$ gives a
spanning path $u=v_1,v_2,\ldots,v_N=v$ in $H$.
Among the $N-1$ indices $j$, the sets specified by $uv_{j+1}\in E(H)$ and
$v_jv\in E(H)$ have total size at least $N$. Some index belongs to both.
Then
\[
 v_1,\ldots,v_j,v_N,v_{N-1},\ldots,v_{j+1},v_1
\]
is a Hamiltonian cycle in $H$. Removing $w$ gives a spanning path in $F$.
This is the standard Bondy--Chv\'atal closure argument \cite{src:bc}.
\end{proof}

\subsection{Small maximum independent neighborhoods}\label{split17:sec:small-m}
\begin{theorem}\label{split17:thm:small-m}
Under the hypotheses and notation of Lemma~\ref{split17:lem:setup}, suppose
$k\ge2$ and $k\ge2m-2$, where $m=\max_{c\in C}|A_c|$. Then $G$ is traceable.
\end{theorem}
\begin{proof}
If $m=1$, the clique-neighbor sets of different independent vertices are
disjoint. Choose two neighbors of each independent vertex, form disjoint
three-vertex paths, and concatenate them and the unused clique vertices
using clique edges. Hence assume $m\ge2$.

Choose $c_0$ with $A=A_{c_0}$ of size $m$, and put $X=I\setminus A$.
Maximality of $m$ and \eqref{split17:eq:columns} give
\[
 2|A_c\setminus A|=|A_c\triangle A|+|A_c|-|A|\le k.
\]
Put $r=\lfloor k/2\rfloor\ge1$. For $J\subseteq X$, count incidences
between $J$ and its clique neighborhood to obtain
\[
 k|J|\le r|N_C(J)|,\qquad |N_C(J)|\ge2|J|.
\]
Hall's theorem \cite{src:hall}, applied to two copies of each vertex of $X$,
therefore supplies two distinct clique neighbors per vertex of $X$, with
all chosen clique vertices distinct. Indeed, a subset of the copied side
with underlying set $J$ has at most $2|J|$ vertices and neighborhood
$N_C(J)$, so it satisfies Hall's condition.

Form these disjoint three-vertex paths and add every unselected clique
vertex as a singleton path. This covers $C\cup X$ by paths whose endpoints
are in $C$. The vertex $c_0$ is a singleton, since it has no neighbors in
$X$. List $A=\{a_1,\ldots,a_m\}$.

Choose a neighbor $p\ne c_0$ of $a_1$, which exists because $d(a_1)\ge k\ge2$.
Join the singleton $c_0$ and the component containing $p$ through $a_1$.
At most one clique vertex becomes internal: $p$ might become internal,
but $c_0$, previously isolated, becomes an endpoint.

After $t$ such insertions joining distinct clique-ended components,
at most $2t-1$ clique vertices are internal. Each subsequent join makes
at most two former endpoints internal. Every other clique vertex is an
endpoint of one component, with a singleton counted as one endpoint vertex.
For $1\le t\le m-2$, the next vertex $a_{t+1}$ has at least
\[
 k-(2t-1)\ge2m-2-2t+1\ge3
\]
clique neighbors that are endpoints. A component has at most two endpoint
vertices, so two such neighbors belong to distinct components. Join those
components through $a_{t+1}$, maintaining the invariant.

After $m-1$ core vertices have been inserted, the final vertex $a_m$ has
at least
\[
 k-(2m-3)\ge1
\]
endpoint neighbor. Append $a_m$ to that component. The resulting path has
one independent endpoint and one clique endpoint. Concatenate every
remaining clique-ended component at its clique endpoint, using clique
edges. This is a spanning path. When $m=2$, the intermediate insertion
loop is empty and the same final step applies.
\end{proof}

\subsection{Connectivity one}\label{split17:sec:k1}
Assume $k=1$. By \eqref{split17:eq:columns}, the nonempty columns have size at most
two and pairwise symmetric difference at most one. If a two-element
column occurs, every two-element column is identical to it and every
singleton belongs to it. If none occurs, all singleton columns are
identical. Every independent vertex has a neighbor, so $|I|\le2=k+1$.
Equation~\eqref{split17:eq:CE} proves traceability.

\subsection{Connectivity two}\label{split17:sec:k2}
Assume $k=2$. The columns satisfy
\begin{equation}\label{split17:eq:k2}
 1\le|A_c|\le3,\qquad |A_c\triangle A_d|\le2,
 \qquad d_C(i)\ge2.
\end{equation}
If $|I|\le3$, apply \eqref{split17:eq:CE}. Henceforth $|I|\ge4$.
The case $m=1$ was handled in Theorem~\ref{split17:thm:small-m}.

\subsubsection{Two-element maximum columns}
Two two-element columns must intersect. A pairwise-intersecting family of
two-element sets either has a common element or is contained in a
three-element ground set: from distinct sets $\{a,b\},\{a,c\}$, a member
not containing $a$ must be $\{b,c\}$; its presence prevents any new ground
element. Every singleton column belongs to every two-element column, by
\eqref{split17:eq:k2}. A family with one distinct two-element member, or a family
confined to three ground elements, would have $|I|\le3$.
Thus all two-element columns are $\{a,x\}$ for a fixed $a$, and their
common intersection is exactly $\{a\}$. Singleton columns can only be $\{a\}$.

Write $I\setminus\{a\}=\{x_1,\ldots,x_t\}$. Each $x_j$ has two distinct
clique neighbors $p_j,q_j$ with column $\{a,x_j\}$, and these groups are
disjoint for different $j$. The path
\[
 a,p_1,x_1,q_1,p_2,x_2,q_2,\ldots,p_t,x_t,q_t
\]
contains all independent vertices and ends in $C$. Append all unused
clique vertices to obtain a spanning path.

\subsubsection{Three-element maximum columns}
Distinct three-element columns intersect in at least two elements.
Such a family either has a common two-element subset or its union lies
in a four-element set. Indeed, take distinct members
$\{a,b,c\},\{a,b,d\}$. A member not containing both $a,b$ must be
$\{a,c,d\}$ or $\{b,c,d\}$. Once such a member occurs, the intersection
condition prevents any fifth element in any member.
Every two-element column is a subset of every three-element column;
every singleton belongs to every three-element column. A single distinct
three-element type would force $|I|\le3$.

Suppose first that all three-element columns have a common pair $\{a,b\}$.
Their intersection is exactly this pair, since there are at least two
distinct types. Every smaller column is a nonempty subset of $\{a,b\}$.
For $I\setminus\{a,b\}=\{x_1,\ldots,x_t\}$, where $t\ge2$, choose distinct
$p_j,q_j$ with column $\{a,b,x_j\}$. The groups are disjoint, and
\[
 a,p_1,x_1,q_1,b,p_2,x_2,q_2,p_3,x_3,q_3,\ldots,p_t,x_t,q_t
\]
is a path containing $I$. Append all remaining clique vertices.

It remains to handle the four-element alternative, where $|I|=4$.
There are at least three clique vertices with three-element columns.
Otherwise the two distinct column types would each have a private
independent vertex occurring in no smaller column and of degree one.
Choose distinct $t_1,t_2,t_3$ with three-element columns not all identical.
Their union is $I$ and their intersection is nonempty.

For any specified $x$ in that intersection there is an alternating
seven-vertex path with endpoint $x$. To verify this, delete $x$ from the
three columns. On the remaining three points, up to relabelling and
reordering, the multiset is either
\[
 \{ab,ab,bc\}\qquad\text{or}\qquad\{ab,bc,ca\}.
\]
Order the columns respectively as $ab,ab,bc$ or $ca,ab,bc$. Then
\begin{equation}\label{split17:eq:seven}
 x,t_1,a,t_2,b,t_3,c
\end{equation}
is the required alternating path.

If an unused clique vertex $d$ has a three-element column, it sees at
least one of the two independent endpoints of \eqref{split17:eq:seven}. Orient
the path to end there, append $d$, and then append all other unused clique
vertices. If all unused clique vertices have smaller columns, choose one
of them, $d$. Its nonempty column is contained in the intersection of all
three-element columns. Choose $x\in A_d$, construct \eqref{split17:eq:seven} with
endpoint $x$, orient it to end at $x$, and append $d$ and the remaining
clique vertices. If there are no unused vertices, \eqref{split17:eq:seven} itself
is spanning. This completes the connectivity-two case.

\subsection{Connectivity three}\label{split17:sec:k3}
Assume $k=3$. We have
\begin{equation}\label{split17:eq:k3}
 1\le|A_c|\le4,\qquad |A_c\triangle A_d|\le3,
 \qquad d_C(i)\ge3.
\end{equation}
If $|I|\le4$, use \eqref{split17:eq:CE}. Assume henceforth $|I|\ge5$.
Choose a maximum column $A=A_{c_0}$ of size $m$, and put $X=I\setminus A$.
The identity
\[
 2|A_c\setminus A|=|A_c\triangle A|+|A_c|-|A|\le3
\]
shows that each clique vertex has at most one neighbor in $X$.
Consequently the groups $C_x=N_C(x)$, $x\in X$, are pairwise disjoint,
and each has at least three vertices. Put
$C_0=C\setminus\bigcup_{x\in X}C_x$; in particular $c_0\in C_0$.
For $c\in C_x$, let $S_c=A_c\cap A$. Comparison with $A$ gives
\begin{equation}\label{split17:eq:core-bound}
 m+1-|S_c|=|A_c\triangle A|\le3,
 \qquad |S_c|\ge m-2.
\end{equation}

If $m\ge3$, it is enough to find a path through all of $A$ using at most
one clique vertex from each $C_x$. At least two unused vertices remain
in every group. Choose distinct unused $p,q\in C_x$; by
\eqref{split17:eq:core-bound}, $p$ has a core neighbor already on the path.
The second operation of Lemma~\ref{split17:lem:insertions} inserts $x$.
Different groups are disjoint, so this works successively for every
$x\in X$. The first operation then inserts all unused clique vertices.

\subsubsection{Maximum column size at most two}
The case $m=1$ was already handled. For $m=2$, write $A=\{a,b\}$ and
choose $r\in N_C(a)\setminus\{c_0\}$. The path $a,r,c_0,b$ has a clique
edge and reserves at most one vertex in each group. Use the third
operation of Lemma~\ref{split17:lem:insertions} with two remaining neighbors
of each $x\in X$. A clique edge remains after each insertion. Finally
insert all unused clique vertices by its first operation.

\subsubsection{Maximum column size three}
Here $|X|\ge2$. If some $d\ne c_0$ has two distinct core neighbors $a,b$,
write $A\setminus\{a,b\}=\{c\}$. The path
\[
 c,c_0,a,d,b
\]
contains $A$ and reserves at most one vertex in a private group, so the
absorption procedure applies.

Otherwise every clique vertex except $c_0$ has at most one core neighbor.
Equation~\eqref{split17:eq:core-bound} says every private-group column is
$\{a,x\}$ for some $a\in A$. Columns $\{a,x\},\{b,y\}$ from different
groups would have symmetric difference four if $a\ne b$. Since there
are at least two groups, comparison with a fixed column in another group
shows that all private-group columns use one common core vertex $a$.
Write $A=\{a,b,c\}$. Both $b,c$ have all their neighbors in $C_0$.
Choose $p\in N_C(b)\setminus\{c_0\}$ and
$q\in N_C(c)\setminus\{c_0\}$. They are distinct, because no column
other than $c_0$ contains two core vertices. The path
\[
 a,c_0,b,p,q,c
\]
contains $A$ and reserves no private-group vertex. Absorb the groups at
$a$ and then insert the remaining clique vertices.

\subsubsection{Four core vertices and at least two outside vertices}
Assume $m=4$ and $|X|\ge2$. Each private column has at least two core
neighbors. Suppose there are $p\in C_x$ and $q\in C_y$, $x\ne y$, such
that $|S_p\cup S_q|\ge3$.
If $S_p,S_q$ are disjoint, they are two-element sets, say
$S_p=\{a,b\}$ and $S_q=\{c,d\}$; use
\[
 a,p,b,c_0,c,q,d.
\]
Otherwise choose distinct $a,b,c$ with $a,b\in S_p$ and $b,c\in S_q$;
such a choice is possible because both sets have size at least two and
their union has size at least three. For the remaining core vertex $d$,
use
\[
 d,c_0,a,p,b,q,c.
\]
Both paths reserve one vertex in each of two different groups. Absorption
therefore applies.

If no such pair exists, all private core neighborhoods equal one
and the same pair $B=\{a,b\}$. Indeed, comparison of columns from different
groups forces equality of their core sets, each of size at least two;
comparison with a fixed column in another group covers every column.
Write $A\setminus B=\{c,d\}$. All neighbors of $c,d$ lie in $C_0$.
Choose $r\in N_C(c)\setminus\{c_0\}$ and
$s\in N_C(d)\setminus\{c_0\}$, and choose any private-group vertex $p$.
If $r=s$, use
\[
 a,p,b,c_0,c,r,d;
\]
if $r\ne s$, use
\[
 a,p,b,c_0,c,r,s,d.
\]
These paths contain $A$ and reserve only the private-group vertex $p$.
Again the absorption procedure completes the proof.

\subsubsection{Four core vertices and one outside vertex}
The final case is $m=4$, $X=\{x\}$, and $|I|=5$.
A singleton column $\{x\}$ would differ from $A$ in five elements,
contrary to \eqref{split17:eq:k3}. Thus all of the at least three clique neighbors
of $x$ have columns of size at least two. Together with $c_0$, this gives
at least four such columns.

Put $q=|C|$ and $n=q+5$. For a missing incidence $ic$ with $|A_c|\ge2$,
\[
 d(i)+d(c)\ge3+(q-1+2)=n-1.
\]
Apply Lemma~\ref{split17:lem:closure} successively to make all these columns
adjacent to all independent vertices. Every independent vertex then has
degree at least four. Each remaining column is a nonempty singleton,
so all of its missing incidences also satisfy
\[
 d(i)+d(c)\ge4+(q-1+1)=n-1.
\]
Fill them using the same closure rule. The final complete split graph has
five independent vertices and at least four clique vertices, and hence
has a spanning path. Reverse the closure steps to obtain one in $G$.
This completes the proof of Theorem~\ref{split17:thm:main}.



\section{Balanced incidence and the clique-size bound}
\label{sec:split-bounds}
Continue with $V(G)=C\mathbin{\dot\cup}I$, $A_c=N(c)\cap I$,
$q=|C|$, $h=|I|$, row degrees $r_i$, and column sizes $s_c=|A_c|$.
The conditions in~\eqref{split17:eq:columns} may also be imposed with
an integer lower bound $k$ that is not necessarily the minimum row
degree. We use that slightly more general incidence formulation below.


\begin{theorem}\label{split19:thm:graph}
Let $G$ be a connected noncomplete split graph satisfying
\eqref{split17:eq:original}. Choose a split partition with $I$ maximum, and put
$k=\kappa(G)$, $q=|C|$, and $h=|I|$. If $q\le2k$, then $G$ is traceable.
Moreover, $h=q$ always implies $q\le2k$, so every equal-size instance
is traceable.
\end{theorem}
Complete graphs are traceable without these restrictions.



For a missing incidence $ic$, the full degree of the
clique vertex is $q-1+s_c$. The path-closure rule of
Lemma~\ref{split17:lem:closure} therefore permits its addition whenever
\begin{equation}\label{split19:eq:threshold}
 r_i+s_c\ge h.
\end{equation}
Only incidences are added, and degrees never decrease.


\subsection{The main incidence theorem}
\begin{theorem}\label{split19:thm:matrix}
Let a split graph have nonempty clique columns and side sizes
$1\le h\le q\le2k$, where $k\ge1$ is an integer. Suppose every independent
vertex has degree at least $k$ and every clique column has size at most
$k+1$. If $q=2k$, suppose additionally that pairwise column symmetric
differences have size at most $k$. Then the graph is traceable.
For $q<2k$ the symmetric-difference hypothesis is unnecessary.
\end{theorem}

\subsubsection{A sufficient descending order}
Sort the original column sizes as
\[
 1\le s_1\le s_2\le\cdots\le s_q,
 \qquad a=q-h\ge0.
\]
If $h\le k+1$, every missing incidence satisfies
$r_i+s_c\ge k+1\ge h$, so closure immediately completes the graph.
Assume henceforth $h\ge k+2$. In particular $q-k\ge a+2$.
A sufficient condition for traceability is
\begin{equation}\label{split19:eq:parking}
 s_j\ge j-a\qquad(1\le j\le q-k).
\end{equation}
For $j\le a+1$ this follows from nonempty columns.

To verify sufficiency, complete the columns in descending index order.
For $j\ge q-k$, their sizes are at least
$s_{q-k}\ge q-k-a=h-k$, so \eqref{split19:eq:threshold} follows from $r_i\ge k$.
For each later index $j<q-k$, the $q-j$ larger-index columns have already
been completed. Every independent vertex has degree at least $q-j$,
and \eqref{split19:eq:parking} gives
\[
 r_i+s_j\ge(q-j)+(j-a)=h.
\]
All closure additions are valid. A complete split graph with $h\le q$
is traceable, and reversing the additions proves the criterion.

\subsubsection{The deficit budget}
Define
\[
 E=\sum_{i\in I}(r_i-k)\ge0,
 \qquad W=\sum_{c\in C}(k+1-s_c)\ge0.
\]
Counting incidences yields
\begin{equation}\label{split19:eq:budget}
 W=q(k+1)-kh-E=q+ka-E\le B:=q+ka.
\end{equation}
If \eqref{split19:eq:parking} fails at $j$, then
$a+2\le j\le q-k$ and $s_j\le j-a-1$. The first $j$ columns each have
deficit at least $k+a+2-j$, so
\begin{equation}\label{split19:eq:quadratic}
 W\ge f(j):=j(k+a+2-j).
\end{equation}
The quadratic $f$ is concave, and its minimum on $[a+2,q-k]$ occurs at
an endpoint. Direct expansion gives
\begin{align}
 f(a+2)-B&=2k-q,\label{split19:eq:endpoint-one}\\
 f(q-k)-B&=(h-k-1)(2k-q).\label{split19:eq:endpoint-two}
\end{align}
If $q<2k$, both differences are positive because $h\ge k+2$.
Equations \eqref{split19:eq:budget} and \eqref{split19:eq:quadratic} then contradict each
other. This proves Theorem~\ref{split19:thm:matrix} in the strict case.

\subsubsection{The boundary and its two extremal patterns}
Let $q=2k$. Both endpoint differences are zero. If $E>0$, then $W<B$
while $f(j)\ge B$, again ruling out any failure of \eqref{split19:eq:parking}.
Suppose that condition fails and $E=0$. Every independent vertex has
original degree exactly $k$, and equality holds throughout
\eqref{split19:eq:quadratic}. Strict concavity forces a failing index to be an
endpoint of $[a+2,k]$, unless the endpoints coincide. Exactly $j$ columns
then have size $j-a-1$, and every other column has size $k+1$; any further
deficit would make $W>B$.

\paragraph{Pattern A.}
For $j=a+2$, there are $a+2$ singleton columns and
$2k-a-2=h-2$ columns of size $k+1$. Since $h\ge k+2$, there are at least
$k$ of the latter. Choose a singleton $\{v\}$. Its symmetric difference
with a size-$(k+1)$ column would have size $k+2$ unless that column
contained $v$. Thus $v$ sees every high column and its own singleton
column, giving original degree at least $k+1$. This contradicts $E=0$.
Pattern A is impossible.

\paragraph{Pattern B.}
For $j=k$, put $\ell=k-a-1=h-k-1\ge1$. There are $k$ low columns of
size $\ell$ and $k$ high columns of size $k+1$.
First complete every high column; its missing incidences satisfy
\[
 r_i+(k+1)\ge2k+1\ge h.
\]
Let $U$ be the union of the original low columns. Every $i\in U$ retains
at least one low-column neighbor, in addition to its $k$ now-completed
high-column neighbors. Therefore every missing incidence from $U$ to a
low column satisfies
\[
 r_i+\ell\ge k+1+\ell=h.
\]
Complete all low columns to all vertices of $U$.

We claim $|U|\ge\ell+1$. Otherwise all low columns equal the same set
$U$ of size $\ell$. Each vertex of $U$ originally had its $k$ neighbors
in those low columns, and had no high-column neighbor because its degree
was exactly $k$. Every high column was therefore disjoint from $U$.
Its symmetric difference with a low column had size
$(k+1)+\ell=h>k$, a contradiction.

Now every low column has size at least $\ell+1$, and every vertex outside
$U$ has degree at least $k$ from the completed high columns. The remaining
missing incidences satisfy
\[
 r_i+s_c\ge k+(\ell+1)=h.
\]
Complete them and reverse closure. If the two endpoint patterns coincide,
Pattern A already rules out that configuration. This proves
Theorem~\ref{split19:thm:matrix}.

\subsection{Consequences for the original split hypothesis}
\begin{lemma}\label{split19:lem:side-size}
Under \eqref{split17:eq:columns}, $h\le q+1$. If $h=q$, then $q\le2k$.
\end{lemma}
\begin{proof}
The incidence bounds give $kh\le q(k+1)$.
Counting unordered column differences by rows gives
\begin{align}
 \sum_{\{c,b\}\subseteq C}|A_c\triangle A_b|
 &=\sum_{i\in I}r_i(q-r_i)\notag\\
 &\ge k\left(qh-\sum_i r_i\right)
 \ge kq(h-k-1).\label{split19:eq:pairs}
\end{align}
The same sum is at most $kq(q-1)/2$.
If $h\ge q+2$, incidence counting forces $q\ge2k$, while
\eqref{split19:eq:pairs} forces $q\le2k-3$, a contradiction. Thus $h\le q+1$.

When $h=q$, equation \eqref{split19:eq:pairs} gives $q\le2k+1$.
Suppose equality holds. Both lower-bound steps in \eqref{split19:eq:pairs} must
then be equalities. Hence every column has size $k+1$, and every row
degree is either $k$ or $q$, the two equality possibilities in
$r_i(q-r_i)\ge k(q-r_i)$. The total row excess is
\[
 E=q(k+1)-kq=q=2k+1.
\]
If $u$ rows have degree $q$, the same excess is
$u(q-k)=u(k+1)$. This would require $(k+1)\mid(2k+1)$, which is impossible
for $k\ge1$. Therefore $q\le2k$.
\end{proof}

The remaining side-size possibility is supplied by the following explicit
dependency. It is the balanced-deficit result; its complete proof is
included in Subsection~\ref{split19:sec:balanced}, so no unstated reduction to that
case is assumed.
\begin{lemma}[Balanced deficit]\label{split19:lem:balanced}
A split incidence graph satisfying \eqref{split17:eq:columns} with $h=q+1$
is traceable.
\end{lemma}

\begin{proof}[Proof of Theorem~\ref{split19:thm:graph}]
Lemma~\ref{split17:lem:setup} supplies \eqref{split17:eq:columns}, and
Lemma~\ref{split19:lem:side-size} gives $h\le q+1$. If $h=q+1$, apply
Lemma~\ref{split19:lem:balanced}. Otherwise $h\le q$, and $q\le2k$ permits
Theorem~\ref{split19:thm:matrix}. If $h=q$ was assumed instead, the second part
of Lemma~\ref{split19:lem:side-size} supplies $q\le2k$ first.
\end{proof}

\subsection{Proof of the balanced deficit dependency}\label{split19:sec:balanced}
Assume \eqref{split17:eq:columns} and $h=q+1$. Counting as in
\eqref{split19:eq:pairs}, but using this value of $h$, gives
\[
 kq(q-k)\le\sum_i r_i(q-r_i)\le\frac{kq(q-1)}2.
\]
Thus $k\le q\le2k-1$. Put
\[
 t=q-k\ge0,\qquad d=2k-q\ge1,\qquad
 E=\sum_i(r_i-k)\ge0.
\]
The nonnegative total column deficit is
\begin{equation}\label{split19:eq:balanced-budget}
 W=\sum_c(k+1-s_c)=q(k+1)-k(q+1)-E=t-E\le t.
\end{equation}

Fix a column $A$ of size $s$ and let $E_A=\sum_{i\in A}(r_i-k)$.
Including its zero self-difference, the sum of its differences with all
columns is
\[
 \sum_b|A\triangle A_b|
 =\sum_{i\in A}(q-r_i)+\sum_{i\notin A}r_i
 =k(q+1)+s(q-2k)+E-2E_A.
\]
It is at most $k(q-1)$, and $0\le E_A\le E$, so
\[
 sd\ge2k+E-2E_A\ge2k-E.
\]
Its own deficit $k+1-s$ is at most $W$, giving
$E\le t-(k+1-s)=s-d-1$. Consequently every column satisfies
\begin{equation}\label{split19:eq:balanced-column}
 (s_c-1)(d+1)\ge2k.
\end{equation}

Call a column high if $s_c\ge t+1=q+1-k$. Its missing incidences satisfy
$r_i+s_c\ge k+t+1=h$, so complete every high column by
Lemma~\ref{split17:lem:closure}. If $\ell$ columns were not high, each contributes
at least $k+1-t=d+1$ to $W$. Hence
\[
 \ell(d+1)\le W\le t.
\]
By \eqref{split19:eq:balanced-column} and $t\le k-1<2k$,
\[
 (s_c-1)(d+1)\ge2k>t\ge\ell(d+1),
\]
so in particular every remaining column has size at least $\ell+1$.
Each independent vertex now sees the $q-\ell$ completed high columns.
Thus every remaining missing incidence satisfies
\[
 r_i+s_c\ge(q-\ell)+(\ell+1)=q+1=h.
\]
Complete all remaining columns. The complete split graph with $q$ clique
and $q+1$ independent vertices has an alternating spanning path.
Reverse closure to prove Lemma~\ref{split19:lem:balanced}. For $k=1$, the bound
$q\le2k-1$ gives $q=1$, so the graph is the three-vertex path.



\subsection{Padding the independent side}
\begin{corollary}\label{split:padding}
Suppose the incidence conditions~\eqref{split17:eq:columns} hold, and let
$m=\max_c|A_c|$. If $0\le b=q+1-h\le k+1-m$, the associated split graph
is traceable.
\end{corollary}
\begin{proof}
Add $b$ independent vertices, each adjacent to every clique vertex.
Every column size becomes at most $m+b\le k+1$, while all pairwise
symmetric differences are unchanged. Old rows retain degree at least
$k$, and new rows have degree $q\ge k$. The new graph has $q+1$ rows,
so Lemma~\ref{split19:lem:balanced} gives a spanning path.
Delete the added independent vertices from that path. Each internal
deletion is spliced through the edge between its two clique neighbors;
an endpoint deletion shortens the path. Added rows are pairwise
nonadjacent, so these operations leave a simple path through every
original vertex exactly once.
\end{proof}


\section{Hall equality, exact factors, and alternating reach}
\label{sec:split-capacity}
We now identify the residual incidence obstruction after all the
traceability cases above. An induced tree of a split graph contains
at most two clique vertices: with one it is a star; with two it may
use only their exclusive independent neighbors, and a maximum such
tree uses all of them.
Consequently
\[
 \tau(G)=\max\bigl(\{1\}\cup\{1+|A_c|:c\in C\}
 \cup\{2+|A_c\triangle A_d|:c\ne d\}\bigr),
\]
omitting an empty collection of pairs. Thus the original premise is
captured exactly by the following incidence bounds when
$k=\kappa(G)=\min_i r_i$; the matrix results permit any lower bound $k$:
\begin{equation}\label{split21:H}
\begin{gathered}
 k\ge1,\quad \varnothing\ne A_c\subseteq I,\quad
 r_i:=|\{c:i\in A_c\}|\ge k,\\
 |A_c|\le k+1,\qquad |A_c\triangle A_d|\le k\quad(c\ne d).
\end{gathered}
\end{equation}
For $S\subseteq I$, write $N(S)=\{c:A_c\cap S\ne\varnothing\}$.


\subsection{Hall inequalities and equality}

The first counting reduction is included for completeness; it is not
presented as a new matching theorem.

\begin{lemma}\label{split21:almost}
Under \eqref{split21:H}, every $S\subseteq I$ satisfies $|N(S)|\geq|S|-1$.
If $N(S)\ne C$, then $|N(S)|\geq|S|$.
Consequently every proper subset of $I$ satisfies ordinary Hall, and a
deficient set can only be $S=I$ with $h=q+1$.
\end{lemma}
\begin{proof}
Fix nonempty $S$ and put $p=|S|$, $t=|N(S)|$. Restrict to rows $S$ and
columns $N(S)$, writing $B_c=A_c\cap S$. Every row retains all its neighbors;
the nonempty restricted columns still have size at most $k+1$ and pairwise
distance at most $k$. If $M=\sum_{i\in S}r_i$, then
\[
 kp\leq M\leq(k+1)t.
\]
Double counting pairs of columns separated by a row gives
\begin{align*}
\sum_{c<d}|B_c\triangle B_d|
 &=\sum_{i\in S}r_i(t-r_i)\\
 &\geq k\sum_{i\in S}(t-r_i)
 =k(tp-M)\geq kt(p-k-1).
\end{align*}
The left side is at most $kt(t-1)/2$. If $p\geq t+2$, the incidence
bound gives $t\geq2k$, whereas the pair count gives $t\leq2k-3$.
This contradiction proves $p\leq t+1$.

If $N(S)\ne C$, choose $d\notin N(S)$. Then $A_d\cap S=\varnothing$,
and for every $c\in N(S)$,
\[
 |A_c\cap S|\leq|A_c\triangle A_d|\leq k.
\]
Hence $k|S|\leq\sum_{c\in N(S)}|A_c\cap S|\leq k|N(S)|$.
If instead $N(S)=C$ and $|S|>|N(S)|$, applying the first conclusion to
both $S$ and $I$ forces $|S|=q+1=h$, so $S=I$.
The empty-set cases are immediate.
\end{proof}

By Hall's marriage theorem \cite{src:hall}, when $h\leq q$ there is a matching
covering all of $I$. When $h=q+1$, adjoining a new clique-side vertex
adjacent to every row and applying Hall gives a matching missing at most
one original independent vertex. Neither conclusion alone is a path theorem.

\begin{theorem}[Tight proper-neighborhood sets]\label{split21:tight}
Suppose \eqref{split21:H} holds and a nonempty $S\subseteq I$ satisfies
\[
 |N(S)|=|S|=p,\qquad Q:=N(S)\ne C.
\]
Then there is a single vertex $x\in I\setminus S$ such that
\begin{equation}\label{split21:structure}
\begin{gathered}
 I=S\mathbin{\dot\cup}\{x\},\\
 A_d=\{x\}\quad(d\notin Q),\qquad
 A_c=\{x\}\cup X_c\quad(c\in Q),\qquad |X_c|=k,\quad X_c\subseteq S,
\end{gathered}
\end{equation}
and every row of $S$ belongs to exactly $k$ of the sets $X_c$.
Moreover $p\leq2k-1$.
\end{theorem}
\begin{proof}
Equality must hold at every step of the final incidence count in
Lemma~\ref{split21:almost}. Therefore $r_i=k$ for $i\in S$ and
$|A_c\cap S|=k$ for $c\in Q$.
For any $d\notin Q$ and $c\in Q$,
\[
 |A_c\triangle A_d|
 = k+|(A_c\setminus S)\triangle A_d|\leq k.
\]
Thus $A_c\setminus S=A_d$ for every such pair. All outside columns equal
one common nonempty set $B$, and every inside column is $B$ together with
$k$ rows of $S$. The rank bound gives $|B|\leq1$, so $B=\{x\}$.
Every row outside $S\cup\{x\}$ would have degree zero, proving
\eqref{split21:structure}. The remaining row and column counts are both $k$.

In this square $p$-by-$p$ block, counting differences gives
\[
 pk(p-k)=\sum_{c<d\in Q}|X_c\triangle X_d|
 \leq kp(p-1)/2.
\]
Since $p\geq k\geq1$, division yields $p\leq2k-1$.
\end{proof}

Conversely, a system of the form \eqref{split21:structure} whose square block is
$k$-regular and has pairwise column distances at most $k$ satisfies
\eqref{split21:H}. A block column and an outside singleton differ in exactly $k$
rows. The theorem therefore describes the equality structure exactly.



\subsection{Eliminating every nonempty Hall-equality set}
\begin{proposition}\label{split21:inputs}
Under~\eqref{split21:H}, the split graph is traceable if $h\le q\le2k$
or $h=q+1$. When $h=q$, the hypotheses imply $q\le2k$.
\end{proposition}
\begin{proof}
These are Theorem~\ref{split19:thm:matrix} and
Lemmas~\ref{split19:lem:balanced} and~\ref{split19:lem:side-size},
whose complete proofs appear above.
\end{proof}


\begin{corollary}\label{split21:tighttrace}
Under \eqref{split21:H}, if some nonempty $S\subseteq I$ satisfies
$|N(S)|=|S|$, then the split graph is traceable.
\end{corollary}
\begin{proof}
First suppose $Q=N(S)\ne C$. Use Theorem~\ref{split21:tight} and retain the $p$
columns of $Q$ together with one outside column $d$. The retained graph has
\[
 |C'|=p+1=|I|\leq2k.
\]
Every row in $S$ retains its $k$ neighbors, and $x$ has $p+1\geq k$
neighbors. The other bounds are inherited, so Proposition~\ref{split21:inputs}
gives a Hamiltonian path in the retained graph.

The column $d$ is adjacent to $x$. Choose an edge incident with $d$ on
that path. Its other endpoint is either $x$ or a clique vertex. Every
omitted column is adjacent to both endpoints, and all omitted columns
are mutually adjacent. Replace that edge by a path through all omitted
columns.

If $N(S)=C$, then $|S|=q$ and Lemma~\ref{split21:almost} gives $h\leq q+1$.
The cases $h=q$ and $h=q+1$ follow from the corresponding explicit inputs
in Proposition~\ref{split21:inputs}.
\end{proof}

\begin{corollary}\label{split21:strict}
Any nontraceable split graph satisfying the original WOWII 199 premise,
with the maximum-independent-set partition fixed above, would satisfy
\[
 |N(S)|\geq|S|+1\qquad(\varnothing\ne S\subseteq I).
\]
\end{corollary}
\begin{proof}
The balanced case $h=q+1$ is traceable by Proposition~\ref{split21:inputs}.
Otherwise Lemma~\ref{split21:almost} gives ordinary Hall on every set, and
Corollary~\ref{split21:tighttrace} excludes equality on any nonempty set.
\end{proof}

\subsection{Strict Hall and an exact degree-two factor}

\begin{theorem}\label{split21:factor}
Assume \eqref{split21:H}, $k\geq3$, and strict Hall expansion on every nonempty
row subset. Then the bipartite incidence graph has a spanning subgraph $F$
such that
\[
 d_F(i)=2\quad(i\in I),\qquad 0\leq d_F(c)\leq2\quad(c\in C).
\]
Every nontrivial component of $F$ is an alternating cycle or a path with
both endpoints in $C$; isolated clique vertices may also occur.
\end{theorem}
\begin{proof}
We first prove the exact capacity inequalities
\begin{equation}\label{split21:capacity}
 \operatorname{cap}_2(S):=\sum_{c\in C}\min(2,|A_c\cap S|)\geq2|S|
 \qquad(S\subseteq I).
\end{equation}
The empty-set case is immediate. For a singleton the left side is its
degree, at least $k\geq3$. Suppose $|S|\geq2$. If no restricted column
has size one, the left side is $2|N(S)|\geq2|S|+2$.

Otherwise choose $d$ with $A_d\cap S=\{v\}$, put $S'=S\setminus\{v\}$,
and write $r'_c=|A_c\cap S'|$. Restriction does not increase symmetric
difference. Comparing a restricted column with $\{v\}$ gives
$0\leq r'_c\leq k$ for every $c$. Hence
\begin{equation}\label{split21:chain}
 \operatorname{cap}_2(S')=\sum_c\min(2,r'_c)
 \geq\frac2k\sum_c r'_c\geq2|S'|=2|S|-2.
\end{equation}
Adding $v$ increases the contribution of column $d$ from zero to one,
so $\operatorname{cap}_2(S)\geq\operatorname{cap}_2(S')+1$. If \eqref{split21:capacity} failed, integrality
would force equality throughout \eqref{split21:chain}. For $k\geq3$, the inequality
$\min(2,r)\geq2r/k$ is strict at every integer $1\leq r\leq k-1$.
Consequently every positive $r'_c$ equals $k$, and every row of $S'$ has
degree exactly $k$. Thus
\[
 k|N(S')|=\sum_c r'_c=k|S'|,
\]
contradicting strict Hall on the nonempty set $S'$.

To construct the factor, use a network with source-to-row arcs of capacity
two, incidence arcs from rows to columns of capacity one, and column-to-sink
arcs of capacity two. Fixing the source-side row set $S$, the minimum cut
capacity after choosing the source-side columns is
\[
 2(|I|-|S|)+\sum_c\min(2,|A_c\cap S|)\geq2|I|.
\]
The integral max-flow min-cut theorem \cite{src:fordfulkerson} gives a flow
of value $2|I|$. Its selected incidence edges have exactly the stated
degrees. Since the selected graph has maximum degree two, its nontrivial
components are paths or cycles. No independent vertex can be a path
endpoint because all its selected degrees equal two.
\end{proof}

The lower bound $k\geq3$ cannot be dropped from this factor statement.
For $k=2$, take rows $I=\{v,a,b\}$ and columns
\[
 \{v\},\quad\{v,a,b\},\quad\{a\},\quad\{b\}.
\]
The row minimum is two, the maximum column size is three, the maximum
pairwise column distance is two, and strict Hall holds. But
$\operatorname{cap}_2(I)=1+2+1+1=5<6$, so an exact factor of the stated kind does not
exist. This split graph is nevertheless traceable: labeling its columns
$c_0,c_1,c_2,c_3$ in the displayed order, one spanning path is
\[
 a,c_2,c_0,v,c_1,b,c_3.
\]
It is not a counterexample to the original conjecture.



\subsection{Hypotheses and explicit inputs}

Let $I$ and $C$ be finite disjoint nonempty sets, with $h=|I|$ and $q=|C|$.
For each $c\in C$, let $\varnothing\ne A_c\subseteq I$. The associated split
graph makes $C$ a clique and $I$ independent, with incidence neighborhoods
$A_c$. Assume an integer $k\geq3$ satisfies
\begin{equation}\label{split22:H}
 r_i:=|\{c:i\in A_c\}|\geq k,\qquad
 |A_c|\leq k+1,\qquad |A_c\triangle A_d|\leq k\quad(c\ne d).
\end{equation}
Write $N(S)=\{c:A_c\cap S\ne\varnothing\}$, and assume strict Hall:
\begin{equation}\label{split22:Hall}
 |N(S)|\geq|S|+1\qquad(\varnothing\ne S\subseteq I).
\end{equation}
In particular $q\geq h+1$. Define
\[
 \operatorname{cap}_2(S)=\sum_{c\in C}\min(2,|A_c\cap S|).
\]

The degree-two factor theorem (Theorem~\ref{split21:factor} and Corollary~\ref{split21:strict}) establishes from
\eqref{split22:H}--\eqref{split22:Hall} that $\operatorname{cap}_2(S)\geq2|S|$ for every $S\subseteq I$,
and that the incidence graph has a spanning subgraph $F$ satisfying
\begin{equation}\label{split22:factor}
 d_F(i)=2\quad(i\in I),\qquad d_F(c)\leq2\quad(c\in C).
\end{equation}
We use this exact factor theorem as an explicit input. A second
input is that \eqref{split22:H} implies traceability when $h\leq q\leq2k$
(Theorem~\ref{split19:thm:matrix}). No factor is silently regarded as a Hamiltonian path.

For connection with the original problem, let $\tau(G)$ be the maximum
\emph{vertex order of a connected induced tree} and $\kappa(G)$ the vertex
connectivity. WOWII 199 asks whether a connected graph of order greater than
one with $\tau(G)\leq\kappa(G)+2$ must be traceable \cite{src:catalogue}. In a
noncomplete split graph, choosing $I$ maximum makes all columns nonempty and
gives $k=\kappa(G)=\min_{i\in I}d_C(i)$. Induced stars and double stars then
give \eqref{split22:H}. Corollary~\ref{split21:strict} shows
that an unresolved nontraceable split instance must satisfy \eqref{split22:Hall}.
The results below address this residual incidence formulation only.

\subsection{Capacity equality}

\begin{lemma}\label{split22:proper}
Under \eqref{split22:H}--\eqref{split22:Hall}, if a nonempty $S\subseteq I$ satisfies
$\operatorname{cap}_2(S)=2|S|$, then $S=I$. Moreover, some whole column of size $k+1$
is contained in $S$.
\end{lemma}
\begin{proof}
There must be a restricted singleton column. Otherwise every nonzero
$|A_c\cap S|$ is at least two and
$\operatorname{cap}_2(S)=2|N(S)|\geq2|S|+2$.
If all restricted columns had size at most $k$, then
$\min(2,r)\geq2r/k$, with strict inequality at $r=1$ since $k\geq3$,
would give
\[
 \operatorname{cap}_2(S)>\frac2k\sum_c|A_c\cap S|\geq2|S|,
\]
a contradiction. Thus some restricted column has size $k+1$. The global
rank bound makes it a whole column $A$ contained in $S$.

Every column $B$ meets $S$, because otherwise $|A\triangle B|>k$.
If $B\cap S=\{v\}$, comparison with $A$ first forces $v\in A$ and then
forces $B=\{v\}$: an extra vertex outside $S$ would increase the difference
beyond $k$.

The exact factor \eqref{split22:factor} uses $2|S|$ incidences at $S$, equal to the
sum of their column capacities. Hence every column meeting $S$ at least
twice uses both selected slots on $S$. A column meeting $S$ once is a global
singleton and has no outside neighbor. All columns meet $S$, so no row
outside $S$ could have selected degree two. Therefore $S=I$.
\end{proof}

\begin{theorem}\label{split22:global}
Under \eqref{split22:H}--\eqref{split22:Hall}, global capacity equality $\operatorname{cap}_2(I)=2h$
implies $q\leq2k$.
\end{theorem}
\begin{proof}
By Lemma~\ref{split22:proper}, there is a column $A$ of size $k+1$.
Let $L$ be the number of singleton columns and $H$ the number of other
columns. Put $d=q-h\geq1$. From $L+2H=2h$ and $L+H=q$ we obtain
\begin{equation}\label{split22:LH}
 L=2d,\qquad H=h-d.
\end{equation}
Every singleton label lies in $A$, by comparison with its $k+1$ vertices.
Let $t=h-k-1$ count the rows outside $A$. For any column $B$,
\[
 2|B\setminus A|=|B\triangle A|+|B|-|A|\leq k.
\]
The outside rows have no singleton neighbors, and one of the $H$
nonsingleton columns is $A$ itself. Thus
\begin{equation}\label{split22:outside}
 kt\leq(H-1)\lfloor k/2\rfloor\leq(H-1)k/2.
\end{equation}
Substituting $H=k+1+t-d$ gives $t\leq k-d$, hence
\begin{equation}\label{split22:bound}
 q=k+1+t+d\leq2k+1.
\end{equation}

Suppose $q=2k+1$. Then $t=k-d$ and $H-1=2t$. If $t=0$, there is only
one nonsingleton column and $L=2k$. The total incidence count is then
$3k+1<k(k+1)$ for $k\geq3$, impossible. Thus $t>0$. Equality throughout
\eqref{split22:outside} forces $k$ even and every nonsingleton column other than
$A$ to have $k/2$ rows outside $A$. For each such column $B$, the distance
condition gives
\[
 |B\cap A|\geq |B\setminus A|+1=k/2+1.
\]
The rank bound forces $|B|=k+1$. Therefore all $H$ nonsingleton columns
have size $k+1$. The total row excess over the degree lower bound is
\begin{equation}\label{split22:excess}
 E=(k+1)H+L-kh=2k+1-kd\geq0.
\end{equation}
Since $k$ is even and at least four, $d\leq2$.

Let $U$ be the set of singleton labels, with $u=|U|$. Comparison with each
singleton shows that every nonsingleton column contains all of $U$.
Also $U\subseteq A$, so $t>0$ implies $I\setminus U\ne\varnothing$.
The latter set has neighbors only among the $H$ nonsingleton columns.
Strict Hall gives
\[
 h-u+1\leq H=h-d,\qquad\text{hence }u\geq d+1.
\]
The rows of $U$ alone contribute at least $u(H-k)+L$ to the total excess.
Here $H-k=k+1-2d>0$. Consequently
\[
 E\geq(d+1)(k+1-2d)+2d.
\]
The right side exceeds \eqref{split22:excess} by $(2d-1)(k-d)>0$, a contradiction.
Equality in \eqref{split22:bound} is impossible, so $q\leq2k$.
\end{proof}

\begin{corollary}\label{split22:strictcap}
Under \eqref{split22:H}--\eqref{split22:Hall}, if $q>2k$, then
\[
 \operatorname{cap}_2(S)\geq2|S|+1\qquad(\varnothing\ne S\subseteq I).
\]
\end{corollary}
\begin{proof}
The earlier exact factor theorem gives the non-strict bound. A proper
nonempty set cannot attain equality by Lemma~\ref{split22:proper}, and the whole
set cannot attain equality by Theorem~\ref{split22:global}. Use integrality.
\end{proof}

\subsection{Directed alternating reach}

\begin{theorem}\label{split22:reach}
Assume \eqref{split22:H}--\eqref{split22:Hall} and $q>2k$, and fix any factor satisfying
\eqref{split22:factor}. Orient nonfactor incidence edges from $I$ to $C$, and
factor edges from $C$ to $I$. From the vertices of any factor cycle, a
directed path reaches a clique vertex of selected degree less than two.
\end{theorem}
\begin{proof}
Call a clique vertex \emph{free} if its selected degree is less than two.
Start reachability from all vertices of a factor cycle. Let $X\subseteq I$
and $Y\subseteq C$ be the reached sets, and suppose no free clique vertex
is reached. Every $c\in Y$ then has selected degree two, with both selected
neighbors in $X$. Every graph edge from $X$ to a clique vertex outside $Y$
must be a factor edge, or that vertex would be reached. Therefore, for
every $c\in C$,
\[
 \min(2,|A_c\cap X|)=d_F(c,X).
\]
For $c\in Y$, both sides are two; outside $Y$, all incidences to $X$ are
selected and there are at most two. Summing yields
$\operatorname{cap}_2(X)=\sum_c d_F(c,X)=2|X|$, contradicting
Corollary~\ref{split22:strictcap}, since $X\ne\varnothing$.
\end{proof}

The argument uses strict capacity, not ordinary strict Hall. For example,
with rows $\{a,b,c,d\}$ and columns
$\{a\},\{b\},\{c,d\},\{a,b,c,d\},\{a,b,c,d\}$, the hypotheses
\eqref{split22:H}--\eqref{split22:Hall} hold with $k=3$ and $q=5$. A factor can consist
of the cycle on rows $c,d$ and the third and fourth columns, and the path
through the first column, $a$, the fifth column, $b$, and the second column.
The directed reach from this cycle stops short of both free path endpoints.
Here $\operatorname{cap}_2(I)=2h$ and $q\leq2k$, as permitted by Theorem~\ref{split22:global}.

\subsection{A shortest switch need not reduce the cycle count}

Theorem~\ref{split22:reach} is not an absorption theorem. The following example
satisfies even the original induced-tree premise, strict capacity, and
$q>2k$, but an arbitrary shortest switch can create a replacement cycle.

Let $I=\{x,y,a,b,c\}$ and let $C_0,\ldots,C_6$ form a clique, with
\begin{align*}
 A_{C_0}=A_{C_1}&=\{x,y,a\}, & A_{C_2}&=\{a,b\},\\
 A_{C_3}&=\{y,a,b,c\}, & A_{C_4}&=\{x,a,b,c\},\\
 A_{C_5}&=\{a,c\}, & A_{C_6}&=\{a\}.
\end{align*}
Here $k=3$, $q=7$, and $h=5$. The maximum column size is four and the
maximum column distance is three. The split formulas give
$\kappa(G)=3$ and $\tau(G)=5$. Directly checking the 31 nonempty row
subsets gives minimum Hall slack two and minimum capacity slack one.

Choose a factor consisting of
\[
 x-C_0-y-C_1-x,
 \qquad C_2-a-C_3-b-C_4-c-C_5,
 \qquad\text{isolated }C_6.
\]
The free clique vertices are $C_2,C_5,C_6$. Neither cycle row sees one,
so the directed route $x\to C_4\to b\to C_2$ of length three is shortest.
Toggle its edges and open the original cycle by removing $xC_0$.
The resulting factor consists of
\[
 C_2-a-C_3-b-C_2,
 \qquad C_0-y-C_1-x-C_4-c-C_5,
 \qquad\text{isolated }C_6.
\]
All independent degrees remain two and all clique degrees remain at most
two, but there is still one cycle.

This graph is traceable; one spanning path is
\[
 C_0-x-C_1-y-C_3-b-C_4-c-C_5-a-C_6-C_2.
\]
Thus it is not a counterexample to WOWII 199, and the displayed initial
factor is not a minimum-cycle factor. The example refutes only automatic
cycle decrease for an arbitrary shortest alternating switch. It does not
refute a carefully chosen minimum-cycle or multi-switch argument.



\subsection{The precise unresolved step}
Combining the preceding results, a nontraceable split graph satisfying
the original premise would have $k\ge4$, $q\ge2k+1$, $h\le q-1$,
strict Hall on every nonempty subset of $I$, and
\[
 \operatorname{cap}_2(S)\ge2|S|+1
 \qquad(\varnothing\ne S\subseteq I).
\]
It has a factor with degree exactly two on $I$ and at most two on $C$,
and each cycle in such a factor has a directed alternating route to
unused clique-side capacity. The displayed twelve-vertex example
shows why an arbitrary shortest switch cannot be assumed to lower the
number of cycles. A proof that suitable switches eliminate all cycles
has not been supplied. Hence these results resolve explicit split
subclasses and strengthen the residual obstruction; they do not resolve
all split graphs or unrestricted Conjecture 199.


% Integrated source fragment: even_distance_189.tex
\section{Even-distance traceability: reductions and obstructions for Conjecture 189}
\label{sec:even-distance}
Let $G$ be connected of order $n>1$, and let $s=d_2$ be its
second-smallest degree, counting multiplicity. For each vertex $v$, put
\[
 \operatorname{Even}(v)=\{x\in V(G):\operatorname{dist}(v,x)
 \text{ is even}\}.
\]
Distance zero is included, so $v\in\operatorname{Even}(v)$. The archived
Conjecture 189 asks whether the hypothesis
\begin{equation}\label{even:hypothesis}
 |\operatorname{Even}(v)|\le s+1\qquad(v\in V(G))
\end{equation}
implies traceability \cite{src:catalogue}. We first exclude every
independent-set neighborhood obstruction to traceability, except for
an explicitly traceable family. We then give exact examples explaining
why several stronger automatic reductions are unavailable.

\subsection{Saturation and deficient independent sets}
\begin{lemma}[Saturation]\label{even:saturation}
Assume~\eqref{even:hypothesis}. If an independent set $A$ of size $s+1$
is contained in $N(c)$ for some vertex $c$, then $A$ dominates $G$.
Moreover, each $a\in A$ is either complete or anticomplete to each
component of $G-A$.
\end{lemma}
\begin{proof}
For any $a\in A$, all of $A$ is at distance zero or two from $a$.
Thus $\operatorname{Even}(a)=A$ by~\eqref{even:hypothesis}.
Suppose $xy$ is an edge of $G-A$, with $ax$ an edge and $ay$ a
nonedge. Then $y$ has distance two from $a$, contrary to
$\operatorname{Even}(a)=A$. Adjacency to $a$ therefore propagates
through every component of $G-A$. Connectedness ensures that each such
component has a neighbor in $A$, so every vertex outside $A$ is
dominated by $A$.
\end{proof}

For an independent set, $N(I)$ denotes the union of its open
neighborhoods; it is disjoint from $I$.
\begin{theorem}[The only neighborhood deficiency]\label{even:deficiency}
Assume~\eqref{even:hypothesis}, and let $I$ be independent with
$|N(I)|<|I|$. Then
\[
 |I|=s+1,\qquad |N(I)|=s,\qquad n=2s+1.
\]
All but at most one vertex of $I$ are complete to $N(I)$, and $G$ has
an explicit Hamiltonian path.
\end{theorem}
\begin{proof}
Connectedness gives $s\ge1$, and at most one vertex has degree below
$s$. Write $t=|I|$, $q=|N(I)|$. The empty set is not deficient, and
$t=1$ gives $q\ge1$. If $2\le t\le s$, some vertex of $I$ has degree
at least $s$, so $q\ge s\ge t$.

Suppose $t=s+1$. At least $s$ vertices of $I$ have degree at least
$s$. Since $q<t$, this forces $q=s$ and makes all those vertices
complete to $N(I)$. Designate the possible exceptional vertex by $x$;
if there is none, choose any $x\in I$. It has a neighbor $c\in N(I)$.
The vertex $c$ is adjacent to all of $I$, so Lemma~\ref{even:saturation}
shows that $I$ dominates $G$. Therefore $V(G)=I\mathbin{\dot\cup}N(I)$
and $n=2s+1$. Enumerate $N(I)=\{c_1,\ldots,c_s\}$ with $xc_1\in E(G)$,
and $I-\{x\}=\{a_1,\ldots,a_s\}$. The sequence
\[
 x,c_1,a_1,c_2,a_2,\ldots,c_s,a_s
\]
is a Hamiltonian path, regardless of the edges within $N(I)$.

Finally suppose $t\ge s+2$. No vertex can have $s+1$ neighbors in
$I$: such a subset $A$ would dominate by Lemma~\ref{even:saturation},
whereas the independent vertex in $I-A$ has no neighbor in $A$.
Every vertex of $N(I)$ therefore has at most $s$ neighbors in $I$.
Counting incidences and using the positive degree of the possible
exceptional vertex yields
\[
 sq\ge\sum_{a\in I}d(a)\ge s(t-1)+1>s(t-1).
\]
Integrality gives $q\ge t$, a contradiction. All cases are covered.
\end{proof}

\begin{corollary}\label{even:factor}
Every graph satisfying~\eqref{even:hypothesis} has
$\alpha(G)\le\lceil n/2\rceil$. If $n\ne2s+1$, or if $G$ is a
nontraceable graph satisfying the hypothesis, then $G$ has a perfect
2-matching: a spanning subgraph with components that are single edges
or odd cycles. In particular it has a perfect fractional matching.
\end{corollary}
\begin{proof}
An independent set larger than $\lceil n/2\rceil$ is deficient, since
its neighborhood is contained in its complement. Theorem~\ref{even:deficiency}
would instead give its size $(n+1)/2$, a contradiction.
Under either extra hypothesis, every independent set $J$ satisfies
$|N(J)|\ge|J|$. For an arbitrary $S\subseteq V(G)$, here let $N(S)$
mean the union of open neighborhoods, possibly intersecting $S$.
Put $J=S\setminus N(S)$. Then $J$ is independent and $N(J)$ is
disjoint from $S$, whence
\[
 |N(S)|\ge |S\cap N(S)|+|N(J)|
 \ge |S\cap N(S)|+|J|=|S|.
\]
The bipartite double cover of $G$ has left and right copies of
$V(G)$, with $u_Lv_R$ an edge exactly when $uv\in E(G)$.
The displayed inequality is Hall's condition, so Hall's theorem
\cite{src:hall} gives a perfect matching. It defines a permutation
$\pi$ with $v\pi(v)\in E(G)$ and no fixed points. A two-cycle gives
a single edge, an odd cycle is retained, and an even cycle of length
at least four is replaced by alternating edges. These components form
the claimed perfect 2-matching. Give weight one to each single edge
and weight one-half to each odd-cycle edge to obtain the fractional
matching. This standard factor implication is included with its proof;
it is not a new matching characterization.
\end{proof}

\subsection{A tight family beyond a connectivity--independence shortcut}
\begin{proposition}\label{even:js}
For every integer $s\ge2$, there is a traceable graph $J_s$ satisfying
\eqref{even:hypothesis} with equality at every vertex and
\[
 n=3s+1,\quad d_2=s,\quad \kappa=s-1,\quad \alpha=s+1.
\]
Thus the premise of Conjecture 189 does not imply the
Chv\'atal--Erd\H{o}s sufficient condition $\alpha\le\kappa+1$.
\end{proposition}
\begin{proof}
Take disjoint cliques $A=\{a_0,\ldots,a_s\}$ and
$B=\{b_0,\ldots,b_s\}$, and an independent set
$Z=\{z_1,\ldots,z_{s-1}\}$. Join every vertex of $Z$ to
$A-\{a_0\}$ and $B-\{b_0\}$, with no edges between $A$ and $B$
and no other added edges.
The two special vertices have degree $s$, each ordinary clique
vertex has degree $2s-1$, and each vertex of $Z$ has degree $2s$.
Thus $d_2=s$.
For $a_0$, the even-distance set is $\{a_0,b_0\}\cup Z$, with
$b_0$ at distance four and $Z$ at distance two. For an ordinary
$A$ vertex, it is that vertex together with the $s$ ordinary
$B$ vertices. For a $Z$ vertex, it is itself, the other $s-2$
vertices of $Z$, and $a_0,b_0$. The symmetric descriptions hold on
$B$. Every count is $s+1$.

Deleting $Z$ disconnects the two cliques. Deleting fewer than $s-1$
vertices leaves a member of $Z$ and an ordinary vertex in each clique;
all surviving vertices remain connected through these. Hence
$\kappa=s-1$. The independent set $Z\cup\{a_0,b_0\}$ has size $s+1$.
Any independent set meeting $Z$ uses at most the two special clique
vertices, while one disjoint from $Z$ uses at most one vertex of each
clique. Thus $\alpha=s+1$.
For a Hamiltonian path, start at $a_0$, append $a_i,z_i$ for
$i=1,\ldots,s-2$, then append
\[
 a_{s-1},a_s,z_{s-1},b_1,b_0,b_2,\ldots,b_s.
\]
The initial loop is empty when $s=2$. All displayed consecutive pairs
are edges and every vertex is used once.
\end{proof}
The family was obtained by generalizing a graph in a previously public
certificate list for $d_2=4$, credited in \cite{even:pr70}. Its members
are examples satisfying Conjecture 189, not counterexamples to it.

\subsection{A Hamiltonian family excluding a linear-order extrapolation}
\begin{proposition}\label{even:gt}
For every integer $t\ge2$, a Hamiltonian graph $G_t$ satisfies
\eqref{even:hypothesis} with $s=3t-1$ and $n=10t-3>3s+1$.
\end{proposition}
\begin{proof}
Replace the vertices of a seven-cycle, in cyclic order, by cliques
of sizes
\[
 (w_0,\ldots,w_6)=(t,t,t,2t-1,2t-1,2t-1,t),
\]
joining bags completely exactly along edges of the seven-cycle.
The order is $10t-3$. The degrees in the successive bags are
\[
 (3t-1,3t-1,4t-2,5t-3,6t-4,5t-3,4t-2).
\]
The first two bags each have at least two vertices, so $d_2=3t-1$.
Distances between different bags equal the quotient-cycle distances;
within a bag, distinct vertices are adjacent. A vertex in bag $i$
therefore has $1+w_{i-2}+w_{i+2}$ even-distance vertices. The seven
counts are
\[
 (3t,3t,3t,3t,2t+1,3t,3t),
\]
all at most $s+1=3t$. Nevertheless
$n-(3s+1)=t-1>0$. Listing each bag consecutively in cyclic bag order
gives a Hamiltonian cycle.
\end{proof}
This disproves an auxiliary extrapolation from small-$s$ data, not
the original conjecture or the quadratic order bound in the credited
prior archive \cite{even:pr70}. No computational solver is needed for
the all-parameter statement above.

\subsection{An explicit obstruction to two automatic sufficient-condition routes}
The following twelve-vertex graph is taken from the same credited
certificate list \cite{even:pr70}; its graph6 string is
\texttt{K?BELqwptrvN}. To specify it without external data, the table
lists for each vertex only its neighbors with larger labels.
\begin{center}
\begin{tabular}{cl@{\qquad}cl}
\toprule
Vertex&Larger neighbors&Vertex&Larger neighbors\\\midrule
0&5,6,7,8,9,10,11&6&9\\
1&5,6,9&7&9,10,11\\
2&7,8,10,11&8&10,11\\
3&7,8,10,11&9&10,11\\
4&7,8,10,11&10&11\\
5&6,9&11&none\\\bottomrule
\end{tabular}
\end{center}
Its second-smallest degree is four. Breadth-first distance counts from
vertices $0,\ldots,11$ give even-distance cardinalities
$(5,5,5,5,5,5,5,5,5,5,4,4)$, so the premise holds.
The set $\{0,1,2,3,4\}$ is independent and $\{0,9\}$ is a vertex
cut. Thus $\alpha\ge5>\kappa+1$, since $\kappa\le2$.
For the independent triple $\{2,3,4\}$, all three neighborhoods equal
$\{7,8,10,11\}$, and hence
\[
 d(2)+d(3)+d(4)-|N(2)\cap N(3)\cap N(4)|=8<11=n-1.
\]
In particular neither the connectivity--independence threshold nor this
modified triple-degree threshold follows automatically from the original
premise. The path
\[
 1,5,6,0,8,2,7,3,10,4,11,9
\]
visits all twelve vertices along edges. The graph is therefore
explicitly traceable. The accompanying standard-library checker
reconstructs this table and verifies the distances, cuts, independent
sets, and spanning path; no completeness claim about a broader census
is needed.

\subsection{Scope and prior work}
Easwar and Johari prove traceability under several additional
connectivity, order, and diameter conditions \cite{even:ej}.
Those attributed results are not added to this paper's resolution
count. The reduction proved here says that any counterexample to the
original conjecture must have a perfect 2-matching and cannot arise
from an independent set with a deficient neighborhood. Joining the
factor components into a spanning path remains unproved. The example
families identify limitations of proposed routes and do not refute
Conjecture 189.


% Integrated source fragment: appendices.tex
\section{Further partial results and a formulation-dependent counterexample}
\label{revision:research}
This section records the further results obtained while revising the
consolidation. The original-item ledger remains at fourteen conclusions.
The new theorems enlarge proved subclasses of 40, 61, 165, and 199,
restrict the residual profile of 384b, and refine a connected-domination
equality case. For 389a a complete counterexample is obtained only after
specifying the lower-rank quartile convention; it is kept separate from
an unqualified resolution of the archived statement. None of the
restricted-class results is asserted to settle its unrestricted entry.

\subsection{A graph of supports with subdivided parallel links}
\label{r12support:section}
The following extension transfers the simultaneous degree-switch argument
of Section~\ref{late:geodesic} to an arbitrary graph of supports. It gives
a further subclass of statement 61 without a triangle-free hypothesis;
it does not settle the unrestricted conjecture.

Let $H$ be any connected nonempty simple graph on $k$ vertices, called
roots. For each edge $uv$ of $H$, choose exactly one replacement:
retain the edge $uv$, or replace it by $c_{uv}\ge1$ internally
vertex-disjoint paths of length two. Internal vertices are distinct
between replacements. Attach at least one private pendant vertex to
every root, and add no other vertices or edges. Write $G$ for the resulting
graph, and let $Q$ be the spanning subgraph of $H$ consisting of the
edges replaced by at least two length-two paths.

\begin{theorem}
\label{r12support:main}
For every graph constructed above,
\[
 \nu(G)=k,\qquad
 f(G)\ge\res(G)+f(H)
       \ge\res(G)+\left\lceil\frac{\diam(G)}2\right\rceil.
\]
Thus the original inequality in WOWII 61 holds on this class, with the
stronger half-diameter term.
\end{theorem}
\begin{proof}
The roots form a vertex cover of size $k$. Matching each root $u$ to
one private pendant neighbor $a_u$ attains this bound, so $\nu(G)=k$.
Choose a maximum induced forest $H[W]$, put $S=V(H)\setminus W$, and
let $t=\nu(Q[W])$. Since $Q[W]$ is a forest, its minimum vertex-cover
size equals $t$.

Delete $S$ and a vertex cover of size $t$ from $Q[W]$. The surviving
root skeleton is a subforest of $H[W]$. Every surviving link with at
least two internal vertices loses one of its roots. Thus it reduces
to a star or isolated vertices. A retained edge or a single subdivided
edge cannot introduce a cycle over a forest skeleton. Pendant vertices
also introduce none. Consequently
\begin{equation}
 f(G)\ge n-|S|-t,\qquad n=|V(G)|.
 \label{r12support:forest}
\end{equation}

Choose a matching of $t$ edges in $Q[W]$. For each chosen $uv$, select
one internal vertex $x_{uv}$ of its link; delete $ua_u$ and $va_v$,
and insert $uv$ and $a_ua_v$. Both inserted edges were absent, because
$uv$ was replaced rather than retained and the partners are private
leaves. The five-vertex sets $\{u,v,x_{uv},a_u,a_v\}$ are pairwise
disjoint. All switches may therefore be performed simultaneously,
yielding a simple graph $G'$ with the same degree at every vertex.

In $G'$ there are $t$ triangles $ux_{uv}v$, $t$ partner edges
$a_ua_v$, and $k-2t$ unchanged root-partner edges. These subgraphs are
mutually vertex-disjoint, and additional edges between them do not
invalidate the ensuing vertex-cover lower bound. Every vertex cover
of $G'$ has size at least
\[
 2t+t+(k-2t)=k+t.
\]
The classical residue inequality $\res(F)\le\alpha(F)$ and dependence
of residue only on the degree sequence give
\begin{equation}
 \res(G)=\res(G')\le\alpha(G')\le n-k-t.
 \label{r12support:residue}
\end{equation}
Combining \eqref{r12support:forest} and \eqref{r12support:residue},
\[
 \res(G)+f(H)=\res(G)+k-|S|
       \le n-|S|-t\le f(G).
\]

Finally every vertex of $G$ is a root or has a root neighbor. Each
edge of $H$ is represented by a path of length at most two, so
$\diam(G)\le2\diam(H)+2$. A diametral geodesic in $H$ is an induced
forest on $\diam(H)+1$ vertices. Therefore
$f(H)\ge\diam(H)+1\ge\lceil\diam(G)/2\rceil$.
\end{proof}

\begin{corollary}[Exact feedback number for a support tree]
\label{r12support:tree}
If $H=T$ is a tree and $t=\nu(Q)$, then
\[
 \operatorname{fvs}(G)=t,\qquad f(G)=n-t,\qquad
 \res(G)\le n-k-t,
 \qquad f(G)\ge\res(G)+\nu(G).
\]
\end{corollary}
\begin{proof}
Take $W=V(T)$ in the preceding proof. This gives a feedback vertex
set of size $t$. Conversely, each edge in a maximum matching of $Q$
supplies a four-cycle by selecting two internal vertices in its link.
These $t$ four-cycles are vertex-disjoint, so every feedback vertex
set has size at least $t$. The residue bound is
\eqref{r12support:residue}, and $f(T)=k=\nu(G)$.
\end{proof}

\begin{example}[Beyond the earlier diameter and matching boundaries]
\label{r12support:beyond}
For $r\ge2$, take the support tree on
$z,u_1,v_1,\ldots,u_r,v_r$ with edges $zu_j,u_jv_j$.
Retain each $zu_j$, replace each $u_jv_j$ by two length-two paths,
and add exactly one private leaf at each root. Then
\[
 n=6r+2,\qquad \operatorname{fvs}(G)=r,\qquad
 \nu(G)=2r+1,\qquad \diam(G)=8.
\]
Leaves at distinct $v_j$ roots attain the diameter. For every $r\ge2$
the feedback-vertex number exceeds one, the diameter exceeds six,
and $\nu(G)>\lceil\diam(G)/2\rceil=4$. Thus the theorem reaches
beyond all three earlier hypotheses. The feedback-vertex number and
the excess matching number are unbounded.
\end{example}

The theorem uses the established degree-equivalent-realization method
and the earlier simultaneous-switch mechanism. Its new content here
is the transfer through a maximum induced forest of the arbitrary
support graph $H$. This is an additional restricted-class theorem,
not another fully resolved catalogue number or a claim of historical
priority. Exact finite checks of the construction and the switch,
forest, matching, and residue certificates accompany this section.


\begin{proposition}[A pendant-neighbor subclass of statement 40]
\label{r12support:forty}
Let $G$ be connected and have a nonempty vertex cover $C$ of size $k$ such that
every vertex of $C$ has at least two pendant neighbors. Then
\[
 p(G)=n(G)-2k,\qquad
 f(G)\ge\left\lceil\frac{p(G)+b(G)+1}{2}\right\rceil.
\]
Here $p$ is the minimum number of vertex-disjoint paths covering all
vertices, and $b$ is the largest order of an induced bipartite subgraph.
In particular the conclusion holds for the preceding support-graph
construction whenever each root receives at least two private leaves.
\end{proposition}
\begin{proof}
A covering by $p$ disjoint paths is a spanning linear forest with
$n-p$ edges. Every edge meets $C$, and each member of $C$ has degree
at most two in a linear forest. Hence $n-p\le2k$. Conversely, use
two pendant neighbors at each member of $C$ to form $k$ disjoint
three-vertex paths, and make all remaining vertices singleton paths.
This gives $2k$ edges and proves $p=n-2k$.

The set $V(G)\setminus C$ is independent. Adding any one vertex of
$C$ induces a star together with isolated vertices, so
$f(G)\ge n-k+1$. Since $b(G)\le n$,
\[
 \left\lceil\frac{p(G)+b(G)+1}{2}\right\rceil
 \le\left\lceil\frac{2n-2k+1}{2}\right\rceil
 =n-k+1\le f(G).
\]
\end{proof}
This elementary sufficient condition is not an unrestricted resolution
of statement 40 and is not presented as a historically new path-cover
principle.


\subsection{A sharp partial result for the maximum-degree boundary in 165}
\label{r12leaf:sec165}
For a connected graph $G$, let $L(G)$ be its maximum spanning-tree
leaf number, let $M=\{v:d_G(v)=\Delta(G)\}$, and put
$t=|N_G(M)\setminus M|$. The archived statement 165 asks whether
\[
 L(G)\geq\left\lceil\sqrt{2t}\right\rceil.
\]
The following result establishes a class not explicitly treated in the
supplied consolidated manuscript. It does not settle the unrestricted
statement. Its forest-of-stars construction extends the full-star
criterion used in the proof of Theorem~\ref{late:leaves:main}.

\begin{lemma}\label{r12leaf:multigraph}
Every finite loopless multigraph $S$ of maximum degree at most three,
with $h$ vertices and $m$ edges, has a feedback vertex set $A$ with
\[
 |A|\leq\min\{\lfloor h/2\rfloor,\lfloor m/2\rfloor\}.
\]
Parallel edges are counted with multiplicity, and two parallel edges
form a cycle.
\end{lemma}
\begin{proof}
Choose a cut maximizing the number of crossing edges. Moving one
vertex cannot improve it, so each vertex has at most one neighbor
of its own color, with multiplicity. Each color therefore induces
a matching together with isolated vertices. Deleting the smaller
color class leaves a forest and costs at most $\lfloor h/2\rfloor$.
For a second construction, repeatedly delete a vertex on a cycle
until the remaining multigraph is a forest. Each deletion removes
at least two edges that have not previously been removed. Thus this
construction costs at most $\lfloor m/2\rfloor$. Take the smaller
of the two resulting deletion sets.
\end{proof}

\begin{theorem}\label{r12leaf:165class}
Let $G$ be connected with $\Delta(G)=3$, and suppose its set $M$ of
degree-three vertices is independent. Then
\[
 L(G)\geq2+\frac{|N_G(M)\setminus M|}{4}.
\]
Consequently, the original inequality in statement 165 holds for $G$.
\end{theorem}
\begin{proof}
Write $h=|M|$. Form a loopless multigraph $S$ on $M$ by introducing
one edge $uv$ for each vertex of $G$ adjacent to both $u,v\in M$.
Such a vertex has degree two, since $M$ is the entire degree-three
set. Distinct such vertices may yield parallel edges. Every vertex
of $S$ has degree at most three. If $m=|E(S)|$, counting incidences
from $M$ to its external neighborhood gives exactly
\[
 t=3h-m:
\]
each external neighbor is counted once, except that the $m$ vertices
represented by edges of $S$ are counted twice.

Take the feedback vertex set $A$ from Lemma~\ref{r12leaf:multigraph}.
The union of all stars centered at $M\setminus A$, including every
edge incident with these centers, is a forest. Indeed, $M$ is
independent and every vertex outside $M$ has degree at most two;
a cycle in this union would therefore project precisely to a cycle
of the multigraph $S[M\setminus A]$. Add uncovered vertices as
isolated vertices and extend this forest to a spanning tree $T$ of
$G$. Every center in $M\setminus A$ retains all three incident
edges and has degree three in $T$.

Every nontrivial subcubic tree has exactly two more leaves than
vertices of degree three. Hence
\[
 L(G)\geq L(T)\geq h-|A|+2
 \geq 2+\max\{h/2,h-m/2\}
 \geq2+\frac{3h-m}{4}=2+\frac t4.
\]
Finally $t\leq4(L(G)-2)$ yields
\[
 L(G)^2-2t\geq L(G)^2-8L(G)+16=(L(G)-4)^2\geq0.
\]
Since $L(G)$ is an integer, this is the required ceiling inequality.
\end{proof}

Connected graphs of maximum degree at most two also satisfy 165:
paths have two spanning-tree leaves and at most two vertices in
$N(M)\setminus M$, while cycles and $K_2$ have $M=V(G)$ and hence
empty external boundary. The singleton case has both sides zero
under the degree-one leaf convention.

\begin{proposition}\label{r12leaf:165higher}
The same conclusion $L(G)\geq2+t/4$, and hence statement 165, holds
when $\Delta(G)\in\{4,5\}$, the maximum-degree set $M$ is independent,
and every vertex outside $M$ has degree at most two.
\end{proposition}
\begin{proof}
Write $D=\Delta(G)$ and construct the auxiliary multigraph $S$ as
above. Then $t=Dh-m$, and $\Delta(S)\leq D\leq5$. Color $S$ with
three colors so as to minimize the number of monochromatic edges,
counted with multiplicity. If a vertex had at least two same-color
neighbors, local optimality would force at least two neighbors of
each other color, contradicting its degree being at most five.
Each color class therefore induces a matching and isolated vertices.
One of them is an induced forest of order at least $h/3$.
The greedy cycle-deletion argument also leaves an induced forest of
order at least $h-m/2$. Choose the larger such vertex set $U$.
The union of its full stars is a forest and extends to a spanning
tree retaining degree $D$ at every vertex of $U$. The tree identity
$L(T)=2+\sum_{v:d_T(v)\geq3}(d_T(v)-2)$ gives
\[
 L(G)\geq2+(D-2)\max\{h/3,h-m/2\}.
\]
For $D=4$, a convex combination with weights $3/4,1/4$ shows
\[
 2\max\{h/3,h-m/2\}\geq h-m/4=(4h-m)/4.
\]
For $D=5$, the weights $5/6,1/6$ give
\[
 3\max\{h/3,h-m/2\}\geq4h/3-m/4
 \geq(5h-m)/4.
\]
Thus $L(G)\geq2+t/4$ in both cases, and the preceding square
identity establishes 165.
\end{proof}

\begin{remark}[Sharpness and scope]
The archived annotation of 165 already credits C. Foster with examples
having $L=k+2$ and $t=4k$. The linear coefficient above is sharp even
within the independent-maximum-degree class. An explicit realization
uses vertices $a_i,b_i$ for $1\leq i\leq k$, two internally disjoint
length-two $a_i$--$b_i$ paths for each $i$, and length-three paths
$b_i$--$a_{i+1}$ with cyclic indices. For $k=1$ this is the theta
graph with path lengths $2,2,3$. All internal path vertices are
distinct. The degree-three set is independent, $h=2k$, and $t=4k$.
Every spanning tree must break each of the $k$ disjoint four-cycles,
so at least one of $a_i,b_i$ loses degree three for each $i$.
The subcubic tree identity gives $L\leq k+2$, while the theorem gives
the reverse inequality. Thus $L=k+2$. The case $k=2$ also gives
equality $L^2=2t$ in 165.
This sharpness pattern is consistent with the previously credited
Foster parameters; it is not counted as a new discovery of that
extremal parameter family. The unrestricted cases of 165, including
adjacent maximum-degree vertices in general, remain unresolved here.
\end{remark}


\subsection{One further independence--connectivity boundary}
\label{r12:199:gap-two}
The closure operation also removes one more residual boundary. This is
an additional restricted case of statement 199, not a resolution of
its unrestricted form.

\begin{lemma}[A near-balanced row-degree criterion]
\label{r12:199:row-degree}
Let $G$ be a split graph with partition $C\mathbin{\dot\cup}I$,
where $h=|I|\ge4$, every vertex of $C$ has a neighbor in $I$, and
\[
 |C|\ge h-1,\qquad d_C(i)\ge h-2\quad(i\in I).
\]
Suppose also that $|N_C(S)|\ge |S|$ for every nonempty proper
subset $S\subsetneq I$. Then $G$ is traceable.
\end{lemma}
\begin{proof}
Repeatedly add a missing incidence edge $ic$ whenever
$d_C(i)+|N(c)\cap I|\ge h$. The total degree sum of its endpoints
is then at least $h+|C|-1=|V(G)|-1$. Thus each addition preserves
traceability in both directions by the path-closure lemma
(Lemma~\ref{split17:lem:closure}). The split partition is unchanged;
all the stated degree and neighborhood inequalities persist.

In the resulting closed graph, any column with at least two
independent neighbors is adjacent to all of $I$: otherwise its
missing incidence satisfies $d_C(i)+|N(c)\cap I|\ge(h-2)+2=h$.
Consequently every column is either universal to $I$ or has exactly
one independent neighbor. Let $U$ be the set of universal columns,
put $u=|U|$, and let $L_i$ be the set of singleton columns adjacent
to $i$. These sets partition $C$, and
$u+|L_i|\ge h-2$ for all $i\in I$.

If $u\le h-4$, every $L_i$ has at least two vertices. Choose distinct
$a_i,b_i\in L_i$ for each $i$, and concatenate the disjoint paths
$a_i,i,b_i$ by clique edges. This gives a path through every vertex
of $I$.

If $u=h-3$, each $L_i$ is nonempty. Label four distinct rows
$i_1,i_2,i_3,i_4$ and choose $c_j\in L_{i_j}$ for $1\le j\le4$.
In an ordering $i_1,\ldots,i_h$ of the rows, realize the first
and third gaps by
\[
 i_1,c_1,c_2,i_2\qquad\hbox{and}\qquad
 i_3,c_3,c_4,i_4.
\]
Realize each of the remaining $h-3$ gaps with one distinct vertex
of $U$. This is again a simple path through all rows.

If $u=h-2$, at least two rows have nonempty singleton sets.
Indeed $|C|\ge h-1$ ensures at least one such row, and if it were
unique, say $x$, then $S=I\setminus\{x\}$ would have
$|N_C(S)|=u=h-2<|S|$, a contradiction. Choose distinct such rows
$x,y$ and columns $c_x\in L_x,c_y\in L_y$. The $h-2$ universal
columns give an alternating path starting at $y$ through all $h-1$
rows of $I\setminus\{x\}$. Prepend $x,c_x,c_y$ to this path.

Finally, if $u\ge h-1$, alternate all $h$ rows with $h-1$ distinct
universal columns. In every case the resulting path contains $I$.
Each unused clique vertex has a neighbor $i$ on this path and can
be inserted into a path edge $ic$ by replacing it with $i,d,c$.
Repeating this operation includes every unused clique vertex.
The closed graph is traceable, so the original graph is traceable.
\end{proof}

\begin{corollary}[The independence gap two case of 199]
\label{r12:199:gap-two-cor}
Let $G$ be a connected split graph of order greater than one,
satisfying the original premise $\tau(G)\le\kappa(G)+2$.
If $\alpha(G)\le\kappa(G)+2$, then $G$ is traceable.
Consequently any remaining nontraceable split instance must have
\[
 \alpha(G)\ge\kappa(G)+3.
\]
\end{corollary}
\begin{proof}
Complete graphs and the cases $\kappa(G)\le3$ are already covered
above. Put $k=\kappa(G)\ge4$, choose $I$ maximum, and write
$h=|I|=\alpha(G)$. The Chv\'atal--Erd\H{o}s path theorem covers
$h\le k+1$, so suppose $h=k+2$. The incidence conditions
\eqref{split21:H} hold, and Lemma~\ref{split21:almost} gives
$|C|\ge h-1$ and ordinary Hall on every proper row subset.
The row-degree lower bound is $k=h-2$. Apply
Lemma~\ref{r12:199:row-degree}.
\end{proof}

This boundary is not subsumed by the previously listed small-connectivity,
small-column, clique-size, balanced, padding, or Hall-equality cases.
For example, for any $k\ge4$, take $h=k+2$ rows and $q=2k+1$
columns, each consisting of all rows except one. Let $k+1$ columns
omit one fixed row $x$, and let the other $k$ columns omit distinct
rows different from $x$. Then the minimum row degree is $k$,
the maximum column size is $k+1$, and pairwise column distances
are at most two. Strict Hall holds, $q>2k$, and the padding allowance
$k+1-\max_c|A_c|$ is zero. The corollary applies to this infinite
family. It supplies an additional sufficient condition, not a claim
that this easy family was previously of unknown traceability.


\subsection{Extending the boundary by the Tan--Hung classification}
\label{r12:199:gap-three}
The following improvement uses a published Hamiltonicity theorem. Its
contribution here is an application to the residual incidence system of
statement 199; we do not claim a new general classification theorem.

\begin{lemma}[A consequence of Tan--Hung]
\label{r12:199:strict-degree}
Let $G$ be a split graph with independent part $I$, where $h=|I|\ge6$.
If every $i\in I$ has degree at least $h-3$ and
\[
 |N_C(S)|\ge |S|+1\qquad(\varnothing\ne S\subseteq I),
\]
then $G$ is traceable.
\end{lemma}
\begin{proof}
Add a vertex $w$ adjacent to every vertex and put $F=G+w$, with split
partition $I\mathbin{\dot\cup}(C\cup\{w\})$. Strict Hall on $I$ gives
$|C|\ge h+1$. Hence every clique vertex of $F$ has degree at least
$h+1$, every independent vertex has degree at least $h-2$, and
\[
 |N_F(S)|\ge |S|+2\qquad(\varnothing\ne S\subseteq I).
\]

We apply Tan and Hung's Theorem~1\cite{TanHung2004}. For an independent
part of size $h\ge6$, their theorem says that a split graph of minimum
degree at least $h-2$ satisfying the stated strict neighborhood
inequalities is Hamiltonian unless it is a spanning subgraph of their
exceptional graph $G^h_t$, where $t$ is the clique size. In that graph,
the $h-3$ rows $u_4,\ldots,u_h$ have all their neighbors in the same
set of $h-2$ clique vertices (Table~1 of that paper). Such a spanning
subgraph has those $h-3$ rows of degree at most $h-2$. Under any
exceptional embedding they must therefore lie in our independent part,
since our clique vertices have degree at least $h+1$. It cannot satisfy
the displayed inequality on
$S=\{u_4,\ldots,u_h\}$: it would require at least $h-1$ neighbors.
Thus $F$ is Hamiltonian. Deleting $w$ from a Hamilton cycle of $F$
gives a Hamiltonian path of $G$.
\end{proof}

\begin{corollary}[The independence gap three case of 199]
\label{r12:199:gap-three-cor}
Let $G$ be a connected split graph of order greater than one satisfying
the original WOWII 199 premise $\tau(G)\le\kappa(G)+2$.
If $\alpha(G)\le\kappa(G)+3$, then $G$ is traceable.
Consequently every remaining nontraceable split instance must satisfy
\[
 \alpha(G)\ge\kappa(G)+4.
\]
\end{corollary}
\begin{proof}
Complete graphs and the cases $\kappa(G)\le3$ are already covered.
Write $k=\kappa(G)\ge4$ and fix a maximum independent part
$I$ of size $h=\alpha(G)$. Corollary~\ref{r12:199:gap-two-cor}
handles $h\le k+2$, so assume $h=k+3\ge7$.
If $G$ were nontraceable, Corollary~\ref{split21:strict} would give
strict Hall on every nonempty row subset. The minimum row degree is
$k=h-3$. Lemma~\ref{r12:199:strict-degree} is then a contradiction.
\end{proof}

This extends a restricted sufficient condition from independence gap two
to gap three. It does not close statement 199 for arbitrary split graphs,
and it does not address the nonsplit residual case.

To see that the new boundary is not contained in the earlier listed
regimes, fix $k\ge4$ and let $I=\{0,1,\ldots,k+2\}$. Take $2k+1$
columns whose omitted row pairs are
\[
 \{0,j\}\ (1\le j\le k+1),\qquad
 \{1,j\}\ (2\le j\le k+1).
\]
Every column has size $k+1$, the minimum row degree is $k$, and
two columns have symmetric difference at most four. Thus the original
premise holds, with $\alpha=k+3$ and $\tau=k+2$. Strict Hall holds:
a singleton has at least $k$ neighbors; a pair misses at most one
column; and a set of at least three rows meets every column.
Here $q=2k+1$, while the padding allowance is zero. In particular this
family lies beyond the gap-two, small-clique, small-column, balanced,
and Hall-equality cases. Its explicit traceability is not itself claimed
to be previously unknown.



\subsection{A further restriction at high-degree order thirty-three}
\label{r12:384h33}
The fixed-active method yields an additional structural restriction without
resolving statement 384b. Degrees in the following proposition are full
degrees in $G$, not degrees in the induced high-vertex graph $H$.

\begin{proposition}\label{r12:384h33-prop}
Every counterexample to 384b with exactly $h=33$ vertices of degree at
least three has exactly two vertices of degree three. Their triangle
counts are two and three. In particular, the latter vertex and its three
neighbors induce a $K_4$. Moreover, no pendant triangle occurs, and the
high-root edge of every ordinary low triangle is avoidable by a
maximum matching of $H$.
\end{proposition}
\begin{proof}
Use the critical profile and matching lemmas from
Section~\ref{paired:section}. The high graph has a near-perfect matching
of sixteen edges, the forbidden high contribution to a 2-dominating
set is sixteen, and all thirty-three high triangle counts are distinct
and at least two. The rank-profile multiplicities for types $3$ through
$10$ are
\[
 (2,3,4,5,6,6,5,2).
\]
Choose exactly twelve matching edges uniformly, select one endpoint
uniformly on each active edge, and independently select the unmatched
vertex with probability $2/5$. Thus a matched vertex has marginal
$3/8$, and the expected initial high selection count is $62/5$.

Use the artificial neighborhoods and pair repair of
Lemma~\ref{paired:artificial}. Conditional on focal inactivity, the
external matching has twelve active edges out of fifteen; conditional
on focal activity, it has eleven active edges out of fifteen. If $X,Y$
are capped outside support counts, the expected focal repair is
\[
 W=\frac14\mathbb E_{12}g(X,Y)
     +\frac38\bigl(\Pr_{11}(X=0)+\Pr_{11}(Y=0)\bigr),
\]
with the same $g$ as in \eqref{paired:gpair} and with the singleton
mixed at probabilities $3/5,2/5$.
The finite-capacity coefficient computation in
\eqref{paired:fixed-poly}--\eqref{paired:fixed-coef}, now with fifteen
external edges, proves $W\le\lambda_a+\lambda_b$ for
\begingroup\setlength{\arraycolsep}{2pt}
\[
\begin{array}{c|rrrrrrrr}
a&3&4&5&6&7&8&9&10\\\hline
10^9\lambda_a&346108345&246668957&171497253&116878435&
77326580&50571772&34135474&24984333
\end{array}
\]
\endgroup
and the unmatched repair is at most $\lambda_a+\zeta$, where
$10^9\zeta=88891656$. The exact verifier checks all $36,985$ feasible
color-pair patterns, all thirty-six unordered type pairs, and all
unmatched double-neighbor placements. Its enumeration uses precisely
the six-count classification in \eqref{paired:pair-patterns}; it fills
the unused positions to a total of fifteen external edges, rather
than adding arbitrarily many zeros. The displayed rounded-up rational
potentials are checked directly, so no numerical optimizer is a proof
dependency. The resulting high-cost bound is
\[
 B=\frac{62}{5}+\sum_a n_a\lambda_a+\zeta
   =\frac{2022441569}{125000000}=16.179532552.
\]

A low triangle must occur, since the critical profile includes the
low count one. If it is a pendant triangle, omit one private vertex
when its root is initially selected. This gives the upper bound
\[
 B-\frac38=15.804532552<16,
\]
so no pendant triangle can occur.

Consider an edge attachment on roots $u,v$. If $uv$ is in every
maximum matching of $H$, use the cubic repair in
Lemma~\ref{paired:essential}. Conditional on the focal pair being
inactive, its third high neighbor has selection probability $2/5$,
whether it is on another matching edge or is unmatched. The special
net repair is therefore $(1/4)(3/5)=3/20$. Replacing the generic
pair potential gives
\[
 B-\lambda_3-\min_a\lambda_a+\frac3{20}
 =15.958439874<16.
\]
Hence every such edge can be avoided by a maximum matching.

Finally suppose there is at most one full-degree-three high vertex.
A cubic vertex has at most three incident triangles. Since all high
counts are distinct and at least two, the only two rank positions
that can be cubic are the first two. At least one of their profile
types can therefore be increased from three to four. This decreases
the expected bound by
$\lambda_3-\lambda_4=0.099439388$.
Choose a maximum matching avoiding the root edge $uv$ of any low
triangle. If both roots are matched, their joint selection probability
is
\[
 \frac{12\cdot11}{16\cdot15}\cdot\frac14=\frac{11}{80}.
\]
If one is unmatched, it is $(2/5)(3/8)=3/20$, which is larger.
Omit the private low vertex whenever both roots are initially selected.
This is compatible with every pair repair and with the profile upgrade,
since only already selected high roots lose a low neighbor. The bound is
\[
 B-\frac{11}{80}-(\lambda_3-\lambda_4)
 =\frac{3985648291}{250000000}=15.942593164<16,
\]
a contradiction. There are consequently exactly two cubic high
vertices, of distinct counts two and three. Three triangles through
a cubic vertex mean its three neighbors are pairwise adjacent, yielding
the stated $K_4$.
\end{proof}

This is a restricted-case exclusion and a necessary profile for any
remaining $h=33$ counterexample. It does not remove $h=33$ entirely,
change the order-$37$ representative bound, or add a resolved original
conjecture. The accompanying standalone
\nolinkurl{verify_h33_restriction.py} checks the rational certificate
and compares small fixed-active cases with direct enumeration.


\subsection{The first-quartile-one case of 389a}
\label{revision:389q1}
Retain the notation and source distinction of Section~\ref{late:389a}.
The following bounds treat the small quartile value analytically.
\begin{proposition}[The first-quartile-one boundary]
Let $G$ be a connected graph of order $n\ge4$, let $\ell$ be its number
of leaves, let $a$ count the vertices of degree at most two, and let
$\mu$ be the largest degree frequency. If $\ell\ge1$ and
$n\le4\ell+1$, then
\[
 \gamma_2(G)\le\mu(a+1)-1.
\]
Consequently, WOWII389a holds when its first-quartile degree is one
and that quartile is taken at rank $\lceil n/4\rceil$. In particular
this proves the quartile-one case for every unambiguous source order
$4\mid n$.
\end{proposition}
\begin{proof}
Use the cut bound of Lemma~\ref{late:389a:cut}
$\gamma_2\le\lfloor(n+a)/2\rfloor$ and the inequalities
$a\ge\ell$ and $\mu\ge\max\{2,\ell\}$.
If $\ell=1$, then
$\lfloor(n+a)/2\rfloor\le\lfloor(5+a)/2\rfloor\le2a+1$
for every $a\ge1$.
If $\ell=2$, then
$\lfloor(n+a)/2\rfloor\le\lfloor(9+a)/2\rfloor\le2a+1$
for every $a\ge2$.
If $\ell\ge3$, then
\[
 \mu(a+1)-1-\frac{4\ell+1+a}{2}
 \ge\left(\ell-\frac12\right)a-\ell-\frac32>0,
\]
since $a\ge\ell\ge3$. Integrality completes all cases.
At quartile rank $\lceil n/4\rceil$, the value one implies
$\ell\ge\lceil n/4\rceil$, and hence $n\le4\ell$.
\end{proof}

\begin{corollary}[Residual orders under the weaker floor-rank condition]
Suppose instead that the quartile-one condition is interpreted only as
$\ell\ge\lfloor n/4\rfloor$. Any failure of 389a in this case must
have one of the profiles
\[
 (n,\ell,a,\mu)=(7,1,1,2),\qquad(10,2,2,2),\qquad(11,2,2,2).
\]
This is a finite parameter reduction, not an asserted resolution of the
remaining three profiles.
\end{corollary}
\begin{proof}
Now $n\le4\ell+3$. If $\ell\ge3$, the same cut bound works, since
\[
 \mu(a+1)-1-\frac{4\ell+3+a}{2}
 \ge\left(\ell-\frac12\right)a-\ell-\frac52>0
\]
for $a\ge\ell\ge3$.
For $\ell=1$, $n\le7$ and the bound
$\lfloor(7+a)/2\rfloor\le2a+1$ holds for $a\ge2$.
If $a=1$ and $\mu\ge3$, the right side is at least five and the
cut bound is at most four. For $a=1$, $\mu=2$, and $n\le6$,
the cut bound is at most three, exactly the required threshold.
For $\ell=2$, $n\le11$ and
$\lfloor(11+a)/2\rfloor\le2a+1$ for $a\ge3$.
If $a=2$ and $\mu\ge3$, the right side is at least eight and the
cut bound is at most six. When $a=2$, $\mu=2$, and $n\le9$,
the cut bound is at most five. The displayed profiles are all that remain.
\end{proof}

\begin{proposition}[Exclusion of the seven-vertex residual profile]
No counterexample occurs in the profile $(n,\ell,a,\mu)=(7,1,1,2)$.
Consequently the floor-rank quartile-one case can fail only in the
profiles $(10,2,2,2)$ and $(11,2,2,2)$.
\end{proposition}
\begin{proof}
Let $p$ be the sole leaf and let $v$ be its neighbor. Every other
vertex has degree at least three, and each degree appears at most twice.
Any degree-six vertex must be $v$, because it is adjacent to $p$.
If there were no degree-six vertex, the other six vertices would have
degrees $3,3,4,4,5,5$; including $p$ would give odd degree sum $25$,
which is impossible. Thus $v$ has degree six. Among the remaining
five vertices there must be a degree-five vertex $w$, since degrees
three and four each appear at most twice. The vertex $w$ is adjacent
to all vertices except $p$. The set $\{p,v,w\}$ is therefore
2-dominating, and its size is the required value $2(1+1)-1=3$.
\end{proof}


\subsection{The lower-rank quartile interpretation fails}
\label{r12:389a:rounding}
The quartile convention is not a harmless presentational choice. The
following example refutes the interpretation
$q=d_{\lfloor n/4\rfloor}$, while not refuting the upper-rank
interpretation $q=d_{\lceil n/4\rceil}$. Since the archived source
specifies rank $n/4$ without resolving nonintegral ranks, this is a
\emph{formulation-qualified result}, not an additional unqualified
resolution of the original numbered entry.

\begin{theorem}[An eleven-vertex lower-rank counterexample]
\label{r12:389a:counter}
There is a connected graph $G$ on eleven vertices with
\[
 \gamma_2(G)=6,\qquad a=2,\qquad \mu=2,\qquad
 d_{\lfloor n/4\rfloor}=1.
\]
Consequently the lower-rank reading of 389a would require $6\le5$.
\end{theorem}
\begin{proof}
Use vertices $0,\ldots,10$. The following table lists each vertex's
neighbors with larger labels and specifies all edges.
\begin{center}
\begin{tabular}{cl@{\qquad}cl}
\toprule
Vertex&Larger neighbors&Vertex&Larger neighbors\\\midrule
0&2,3,5,6,7,8,9,10&6&7,8\\
1&4,6,8&7&8\\
2&5,7&8&none\\
3&4,6,8&9&none\\
4&5,7&10&none\\
5&7,8&&\\\bottomrule
\end{tabular}
\end{center}
The graph is connected and has degree multiset
\[
 (1,1,3,3,4,4,5,5,6,6,8).
\]
Hence $a=2$, $\mu=2$, and the lower quartile is $d_2=1$.
The leaves $9,10$ force both of themselves into every two-dominating
set. The set $\{0,1,2,3,9,10\}$ is two-dominating, so it remains
to rule out a set of size five. It would consist of both leaves
and three vertices of the core $\{0,\ldots,8\}$.

First suppose vertex $0$ is selected and let $T$ be the other two
selected core vertices. The vertices $1,4$ are the only core vertices
not adjacent to $0$; their neighborhoods
$N(1)=\{4,6,8\}$ and $N(4)=\{1,3,5,7\}$ are disjoint.
If neither belongs to $T$, both would require all of $T$ in their
disjoint neighborhoods, an impossibility. If both belong to $T$,
vertex $2$ has only one selected neighbor, namely $0$.
If $1\in T$ but $4\notin T$, the other member of $T$ must be in
$\{3,5,7\}$ to dominate $4$ twice. Choosing $3$ leaves vertex $2$
with only the neighbor $0$; choosing $5$ or $7$ leaves vertex $3$
with only the neighbor $0$. If $4\in T$ but $1\notin T$, the other
member must be $6$ or $8$, and vertex $2$ again has only the neighbor
$0$. All cases fail.

Now suppose $0$ is omitted. The three selected core vertices would
have to two-dominate $K=G[\{1,\ldots,8\}]$, since the two leaves
have no neighbors in $K$. In $K$,
\[
 N_K(2)=\{5,7\},\qquad N_K(1)=N_K(3)=\{4,6,8\}.
\]
If $2$ is omitted, $5$ and $7$ must both be selected, leaving at
most one selected neighbor of whichever of $1,3$ is omitted.
Thus $2$ is selected. Selecting exactly one of $1,3$ leaves the
other without two selected neighbors; selecting both leaves $5$
with only the selected neighbor $2$. If neither is selected, the
remaining two selections must be a pair from $\{4,6,8\}$.
The pairs $\{4,6\}$, $\{4,8\}$, and $\{6,8\}$ fail respectively
at vertices $8$, $6$, and $4$. Thus no three-set two-dominates $K$.
This proves $\gamma_2(G)=6$, and
$\mu(q+a)-1=2(1+2)-1=5$.
\end{proof}

For the same graph, the upper-rank quartile is $d_3=3$, giving the
nonviolated proposed upper bound $2(3+2)-1=9$. Its 2,048 vertex
subsets were independently checked by the standard-library program
\texttt{verify\_389a\_rounding.py}; the construction and case proof
above do not depend on the SAT solver used to discover the graph.
This graph is not claimed to have minimum possible order.


\subsection{Strict neighborhood containment forces positive slack}
\label{post22:strict-containment}
Here $b(G)$ is the maximum order of an induced bipartite subgraph and
$\gamma_c(G)$ is the connected domination number.

\begin{proposition}[A forbidden inclusion at equality]
\label{post22:strict-neighborhood}
Let $G$ be a finite connected graph. If
$N_G[u]\subsetneq N_G[v]$ for two vertices $u,v$, then
\[
 b(G)\ge\gamma_c(G)+2.
\]
Consequently, in an equality graph $b(G)=\gamma_c(G)+1$, any two
closed neighborhoods that are comparable by inclusion are equal.
\end{proposition}
\begin{proof}
The vertices $u,v$ are adjacent. Choose
$x\in N_G[v]\setminus N_G[u]$; then $u,v,x$ induce a path.
Use this path to initiate the rooted bipartite-growth construction
of Subsection~\ref{late:claw:rooted-connected-domination}, and denote
its first maximal connected induced bipartite piece by $B_0$.
Every neighbor of $u$ other than $v$ is also adjacent to $v$.
Such a vertex cannot belong to $B_0$, since it would form a triangle
with $u,v$. Thus $u$ has degree one in $G[B_0]$.

Take a spanning tree of $G[B_0]$. Its vertex $u$ is a leaf, while
$v$ is not a leaf: it joins $u$ to the nonempty remainder of the tree.
Choose another leaf $w$. Deleting $u,w$ leaves a nonempty connected
set that dominates $B_0$. It also dominates every vertex outside
$B_0$ having a neighbor in $B_0$. Indeed, maximality implies that
such a vertex sees both color classes of $B_0$. If all its
$B_0$-neighbors belonged to $\{u,w\}$, it would have to be adjacent
to $u$, and therefore also to the retained vertex $v$, a contradiction.

Continue the cited construction with pairwise anticomplete pieces
$B_1,\ldots,B_k$ and at most $k$ connector vertices. From each later
nonsingleton piece delete one spanning-tree leaf different from its
attachment root; delete a later singleton piece entirely. Keep all
connectors. Each connector retains a neighbor in the first earlier
piece it meets: for $B_0$ this was just proved, and for a later piece
it follows from the two-color property after deletion of only one
vertex. A singleton piece cannot be that first piece. Every retained
piece is therefore connected to the initial retained piece.
The same neighbor property dominates all uncolored vertices, deleted
tree leaves have retained neighbors, and a deleted singleton root is
dominated by its connector. We obtain a connected dominating set of
size at most
\[
 \sum_{j=0}^k |B_j|+k-(2+k)=|B|-2\le b(G)-2,
\]
since the union $B$ of the pairwise anticomplete bipartite pieces is
itself induced bipartite. This proves the inequality.
\end{proof}
This is a refinement of the credited DeLaVi\~na--Waller construction.
It supplies an additional structural restriction at equality; it is
not a complete resolution of any further numbered statement.



\appendix
\section{Archive provenance and coverage}
\label{app:provenance}
The five supplied outer archives were opened recursively through their
twenty-two round archives. Their extracted files were compared with the
original member bytes by SHA-256. All 1,333 compared files matched; all
51 main \LaTeX{} manuscripts were among the matched files. No original
source archive was edited. The main source corpus contains 15,798 lines
of \LaTeX{}, before consolidation. Supporting proof notes and supplemental
manuscripts are included separately in the source map.

\subsection{Outer-archive identity}
The following full SHA-256 values identify the exact supplied bundles.
The final ``(1)'' in each local upload name is only the supplied filename
suffix. Each digest is printed on two lines, with no character omitted.
\begin{longtable}{rp{0.20\linewidth}p{0.61\linewidth}}
\toprule
Rounds&Bytes&SHA-256\\\midrule
\endfirsthead
\toprule Rounds&Bytes&SHA-256\\\midrule
\endhead

1--5&2,196,897&\texttt{9277d9ea942b9dd21d43c947c87de211}\newline\texttt{00363890548a68f3076071611eb57878}\\[5pt]

6--10&3,273,667&\texttt{2779810d2cce8e8e0a5a488d347c825f}\newline\texttt{a2bf81973a10ae30125c4c562e574dff}\\[5pt]

11--15&11,544,658&\texttt{31c39ef6c8ea7b0992ea09e865b26bad}\newline\texttt{1dae5603441f0249ef074d9300bbadc7}\\[5pt]

16--20&10,857,175&\texttt{cbb57fdfbc2852d30428b45cc2611f86}\newline\texttt{b3ead07431f605cf6e771e47f6a444a4}\\[5pt]

21--22&1,367,329&\texttt{8f86fb2135deadecd620fc047ca2813d}\newline\texttt{14963b8cfe795a128ff1a81bf1418a0b}\\[5pt]

\bottomrule
\end{longtable}

\subsection{Round-by-round mathematical coverage}
The table maps all rounds to retained content. Repeated improvements
of the same item are integrated into their strongest terminal statement;
their reusable lemmas and exact certificate dependencies are retained.
\begin{longtable}{rrp{0.70\linewidth}}
\toprule
Round&Main papers&Retained subject matter\\\midrule
\endfirsthead
\toprule Round&Main papers&Retained subject matter\\\midrule
\endhead

1&2&434d proof; 394 counterexamples\\[3pt]

2&1&392e/f counterexamples; auxiliary 133 and 269 reductions\\[3pt]

3&1&395a counterexamples\\[3pt]

4&2&434e proof; 382c counterexamples\\[3pt]

5&1&269 counterexamples; auxiliary 189 and 384b reductions\\[3pt]

6&1&302 counterexamples; auxiliary 389a sufficient conditions\\[3pt]

7&1&298 and the prescribed $4k$-cycle theorem\\[3pt]

8&1&384b critical profiles, kernels, and alteration bounds\\[3pt]

9&2&133 density and equality bounds; 384b refinement; 189/199 notes\\[3pt]

10&2&177/180 claw-free cases; 384b refinement; matching padding\\[3pt]

11&3&133 path restrictions; 171 counterexamples; 384b small kernels\\[3pt]

12&3&133 radius four; 382e counterexamples; forest and residue results\\[3pt]

13&2&133 dominating-five-cycle branch; 154 counterexamples\\[3pt]

14&3&133 nondominating branch; historical 354 refutation; 384b core connectivity\\[3pt]

15&2&133 global $P_7$-free density; 384b rank refinement\\[3pt]

16&4&133 even-cycle classification; 384b refinements; unigraph inheritance\\[3pt]

17&4&155 low-leaf case; 199 split cases; 384b refinement; 61 triangle-free equality\\[3pt]

18&4&133 $P_8$ cores and radius five; 199 balanced case; 384b refinement\\[3pt]

19&5&133 $P_9$ cores; 199 clique bound; 384b sampling/order; 61 apex forests\\[3pt]

20&3&384b sampling/order refinements; 61 matching--geodesic case\\[3pt]

21&2&199 Hall equality and exact factor; 384b order-37 representative\\[3pt]

22&2&199 strict capacity and reach; 384b exclusion of $h=35$\\[3pt]

\bottomrule
\end{longtable}
The early low-average result for 133 is contained in the later
unrounded-average theorem. The Round 11--14 five-, six-, and seven-cycle
arguments are included in the full global $P_7$-free proof. The separate
radius-four attachment proof is replaced by the already proved general
even-cycle theorem at $k=3$, after the complete radius-four reduction;
this substitution introduces no circular dependency. The earlier 384b
representative bounds are superseded by 37, while their structural and
sampling lemmas are proved in the combined chain. The balanced 199 proof
is supplied once and then reused for padding and Hall equality.

\subsection{Main-manuscript index}
Each filename below is relative to the indicated round's original
\texttt{paper/} or \texttt{papers/} directory. A twelve-hex-character
digest prefix is printed for convenience; the full hashes and exact
paths are in \texttt{provenance/source\_manifest.json} in the package.
The original source files are preserved there for comparison.
\begin{longtable}{rp{0.64\linewidth}l}
\toprule
Round&Main manuscript&SHA-256 prefix\\\midrule
\endfirsthead
\toprule Round&Main manuscript&SHA-256 prefix\\\midrule
\endhead

1&\nolinkurl{KEIO_WOWII_R01_394.tex}&\texttt{31685100757c}\\

1&\nolinkurl{KEIO_WOWII_R01_434d.tex}&\texttt{e3efdd2d0a62}\\

2&\nolinkurl{KEIO_WOWII_R02_392ef.tex}&\texttt{48be15dff345}\\

3&\nolinkurl{KEIO_WOWII_R03_395a.tex}&\texttt{2eb5c2ccc475}\\

4&\nolinkurl{KEIO_WOWII_R04_382c.tex}&\texttt{4bd120d314ed}\\

4&\nolinkurl{KEIO_WOWII_R04_434e.tex}&\texttt{b69324b11f4a}\\

5&\nolinkurl{KEIO_WOWII_R05_269.tex}&\texttt{047f4f179ca7}\\

6&\nolinkurl{KEIO_WOWII_R06_302.tex}&\texttt{00dd7a59883d}\\

7&\nolinkurl{KEIO_WOWII_R07_298.tex}&\texttt{a61253f072c7}\\

8&\nolinkurl{KEIO_WOWII_R08_384b.tex}&\texttt{369efb33c34a}\\

9&\nolinkurl{KEIO_WOWII_R09_133.tex}&\texttt{3214b6b7ddd1}\\

9&\nolinkurl{KEIO_WOWII_R09_384b.tex}&\texttt{49bab377de00}\\

10&\nolinkurl{KEIO_WOWII_R10_177_180.tex}&\texttt{9f0b3c608b36}\\

10&\nolinkurl{KEIO_WOWII_R10_384b.tex}&\texttt{fd9b6aa4811a}\\

11&\nolinkurl{KEIO_WOWII_R11_133.tex}&\texttt{bb181f8b0267}\\

11&\nolinkurl{KEIO_WOWII_R11_171.tex}&\texttt{4034f71ff698}\\

11&\nolinkurl{KEIO_WOWII_R11_384b.tex}&\texttt{fb46660351c3}\\

12&\nolinkurl{KEIO_WOWII_R12_133.tex}&\texttt{1e2154125f3d}\\

12&\nolinkurl{KEIO_WOWII_R12_382e.tex}&\texttt{1716684c54d7}\\

12&\nolinkurl{KEIO_WOWII_R12_Forest.tex}&\texttt{7329132992b0}\\

13&\nolinkurl{KEIO_WOWII_R13_133.tex}&\texttt{1a0d09cd05e4}\\

13&\nolinkurl{KEIO_WOWII_R13_154.tex}&\texttt{8e47357e3726}\\

14&\nolinkurl{KEIO_WOWII_R14_133.tex}&\texttt{8f7f56bd93c3}\\

14&\nolinkurl{KEIO_WOWII_R14_354.tex}&\texttt{2943dcaf2beb}\\

14&\nolinkurl{KEIO_WOWII_R14_384b.tex}&\texttt{22355a2d2fee}\\

15&\nolinkurl{KEIO_WOWII_R15_133.tex}&\texttt{7471dcc46945}\\

15&\nolinkurl{KEIO_WOWII_R15_384b.tex}&\texttt{61e6911853ce}\\

16&\nolinkurl{KEIO_WOWII_R16_133_cycles.tex}&\texttt{be2fa44cd46d}\\

16&\nolinkurl{KEIO_WOWII_R16_384b_h41.tex}&\texttt{0ca444303059}\\

16&\nolinkurl{KEIO_WOWII_R16_384b_order42.tex}&\texttt{a78dcbeccd5a}\\

16&\nolinkurl{KEIO_WOWII_R16_unigraph_auxiliary.tex}&\texttt{6677f073c5da}\\

17&\nolinkurl{KEIO_WOWII_R17_155_low_leaves.tex}&\texttt{a62b9f6d70a7}\\

17&\nolinkurl{KEIO_WOWII_R17_199_split.tex}&\texttt{4b759b011451}\\

17&\nolinkurl{KEIO_WOWII_R17_384b_h40.tex}&\texttt{3322baf6c089}\\

17&\nolinkurl{KEIO_WOWII_R17_61_trianglefree.tex}&\texttt{de1deae838ad}\\

18&\nolinkurl{133_p8.tex}&\texttt{7ba353f8879f}\\

18&\nolinkurl{133_radius5.tex}&\texttt{0fc11b81a579}\\

18&\nolinkurl{199_balanced.tex}&\texttt{a015fec418cd}\\

18&\nolinkurl{384b_h39.tex}&\texttt{3c8e13f187ee}\\

19&\nolinkurl{133_p9.tex}&\texttt{bab20ad65bf5}\\

19&\nolinkurl{199_clique_bound.tex}&\texttt{2ac3670001a1}\\

19&\nolinkurl{384b_order40.tex}&\texttt{25b8fe84a922}\\

19&\nolinkurl{384b_paired.tex}&\texttt{22b519f94531}\\

19&\nolinkurl{61_apex_forests.tex}&\texttt{4df8b15b1e8f}\\

20&\nolinkurl{384b_h37.tex}&\texttt{790789191ef2}\\

20&\nolinkurl{384b_order39.tex}&\texttt{3e7b6113c1c0}\\

20&\nolinkurl{61_matching_geodesic.tex}&\texttt{0fa44736f3dd}\\

21&\nolinkurl{199_hall_equality.tex}&\texttt{650065c59273}\\

21&\nolinkurl{384b_order37.tex}&\texttt{fe7dea9c3166}\\

22&\nolinkurl{199_strict_capacity.tex}&\texttt{1054d26069b7}\\

22&\nolinkurl{384b_h35.tex}&\texttt{3ff9cc3e34b0}\\

\bottomrule
\end{longtable}

\section{Verification and reproducibility}
\label{app:verification}
The accompanying package contains the standalone manuscript, the
unmodified source materials retained for the proof record, the exact
certificate data, and rerun outputs from this consolidation. The
machine-readable manifests distinguish original source hashes from
outputs newly generated by a verifier. Exploratory search histories
do not substitute for the proved extension, reduction, or coverage
arguments in the paper.

\subsection{What was rerun}
The verification work is organized into the following groups. A
program count is an execution count within the specified group, not
a count of solved conjectures or independent human reviews.
\begin{center}
\begin{tabular}{p{0.22\linewidth}rp{0.60\linewidth}}
\toprule
Group&Programs&Scope\\\midrule
Early conclusions&10&434d/e checks; exact 394, 392e/f, 395a, 382c,
269, and 302 constructions; 298 coloring and the credited antecedent
chain's finite checks\\[4pt]
Later structural work&6&171, 382e, 154, and 354 witnesses; exact 389a
parameter inequalities; matching-padding and weak-majorization checks\\[4pt]
384b&21&Small kernels and their independent input coverage; exact rank,
sampling, matching-repair, and representative-order certificates\\[4pt]
133, 189, and 199&16&The two closed core certificates; the general
even-cycle path checks; four 189 reduction/example guards; nine 199
incidence, closure, factor, capacity, and warning-example checks\\\bottomrule
\end{tabular}
\end{center}
All listed final runs passed. The original mathematical arguments,
rather than the breadth of these finite tests, justify the analytic
theorems. The finite proof dependencies are identified in the body.

The $P_8$ certificate has 192 graph states, maximum state order 25,
and 93,895 required one-vertex extensions. The $P_9$ certificate has
13,930 graph states, maximum state order 45, and 442,204 contact-tree
nodes. Both checks account for every permitted branch, with zero
unexplained frontier. The maximum state orders are observed properties
of closed certificates, not assumptions on target graph order. Their
standard-library verifiers check paths and permutations directly.

For 384b, the 21 reruns include coverage of the small graph kernels,
not merely arithmetic on a list of unspecified completeness. Exact
integer and rational checks verify the potentials used for the later
matching repairs. The $h=35$ computation uses exactly thirteen active
matching edges, with the correct conditioned probabilities; it does
not replace them by independent-edge probabilities. The final
representative theorem additionally uses the analytic low-vertex
trimming and simultaneous-omission lemmas.

\subsection{Restored dependencies and explicitly limited reruns}
Two packaging issues required particular care. The Round 17
$\kappa=3$ checker expected a companion family JSON file that was
absent from its main verification directory but present in that
round's independent-audit snapshot. That original file was restored
to the rerun directory, with its hash recorded; the checker was
unchanged and then passed.

The Round 7 independent checker completed its mathematical tests but
then attempted to hash two prior manuscript snapshots absent from
the supplied archives. In a separately recorded copy, only those
final metadata hash operations were omitted. The mathematical tests
were unchanged. The original and adapted checker hashes and this
limited adaptation are recorded in the audit. That limited adaptation belongs to the original consolidation record.
For the revised edition, the primary manuscript archive was independently
retrieved and the written proof checked. The archive digest is
\begin{center}\small\texttt{d067e88bdbf41f4fd4405b61e7cc71acfc9e4a6a6e64d7ca8d55f35b289d2afc}\end{center}
The source URL, retrieval record, and dependency audit are retained in
\nolinkurl{provenance/revision/298_primary_audit.md}. The classical
transversal bound and all author attributions remain unchanged.

No external Lean certificate was rebuilt. No claim is made that every
historical exploratory enumeration or discovery program was rerun.
In particular, a historical output's older residual range or
``pending'' label is not used as the terminal theorem's status when
the subsequent audited proof gives a stronger conclusion.

\subsection{Checks added for the revised edition}
The original 53-program verification record remains a historical record
of the supplied consolidation. In this edition the additional 394 checker
was rerun on the triangle-free 15- and 19-vertex witnesses, enumerating
all $2^{15}$ and $2^{19}$ vertex subsets respectively and independently
checking maximum matching, local independence, and odd-cycle-transversal
numbers. The new results have direct written proofs for arbitrary family
size. Additional independent checks and research-specific certificates are
identified under the \texttt{extensions} verification group. Original
proof-certificate data and source texts are preserved.

\subsection{Using the package}
The root-level \texttt{README.txt} identifies the paper, source archive,
verification groups, and exact commands. The companion
\texttt{verify.py} lists or runs the recorded checks and compiles the
small C++ verifiers when needed. Python code uses either the standard
library or the explicitly listed graph-library dependency; C++ checks
use the documented compiler options. The paper itself is a single
\LaTeX{} source file with its bibliography inline and no external
figure or input-file dependency. It can be rebuilt by running
\texttt{pdflatex} until the cross-references stabilize.

\section{Attribution and consolidation corrections}
\label{app:attribution}
The following distinctions are maintained throughout the paper.
The total-domination--annihilation antecedents are credited to the
research record attributed to Gebendorfer; the classical transversal
bound is credited to Chv\'atal and McDiarmid. The half-coefficient
354 counterexample modifies a construction credited to Feit. The
precursor for the Maxine mechanism is credited to Cao and Eppstein.
The common-neighbor observation used in the induced-forest discussion
is credited to Gunderson and collaborators, as corrected in the
later source audit. The triangle-free $P_8$ classification is credited
to Chudnovsky and Stacho, and the $Y$--$\Delta$--Heawood graph is a
known graph. Hall's theorem, Hamiltonian closure, elimination-sequence
majorization, and the published matching bounds retain their original
attributions.

The consolidated 133 statements explicitly restore the connectedness
assumption of the original deficit and equality theorems. The phrase
``distance-two dominating'' is used for distance at most two from a
cycle, to distinguish it from two-neighbor domination. The local triangle
count is denoted consistently by $t(v)$, while the total local
independence is $\Lambda$. Graph6 data strings were compared with the
source strings and preserved. The 199 incidence formulation states
that an induced tree with two clique vertices may use only exclusive
neighbors, and a maximum such tree uses all of them. These are
clarifications of scope and notation during integration, not new
conjecture conclusions.

The revised manuscript preserves every complete original-item conclusion
and the proved partial programs through Round 22. The subsequent structural
extensions are identified separately. No unproved cycle-elimination claim,
finite failure to find a counterexample, or exploratory search is promoted
to a full resolution. The 392e/f family statement now explicitly restricts
``fails'' to positive-deficit parameters: $r\ge3$ and $r\ge4$ respectively
for $J_r$; the original principal witnesses are unchanged.

\clearpage
\addcontentsline{toc}{section}{References}
\begin{thebibliography}{99}
\bibitem{TanHung2004}
N. D. Tan and L. X. Hung,
Hamilton cycles in split graphs with large minimum degree,
\emph{Discussiones Mathematicae Graph Theory} 24 (2004), 23--40.
\url{https://doi.org/10.7151/dmgt.1210}.

\bibitem{src:catalogue}
E. DeLaVi\~na,
\emph{Written on the Wall II: Conjectures of Graffiti.pc}, original catalogue.
\url{http://cms.uhd.edu/faculty/delavinae/research/wowII/all.html}.
The source-qualified numbered formulations used in this paper are preserved
in the supplied round archives and in the pinned HTML-row transcription,
revision \texttt{3be380b4222f974bedbefb228c8e75fbd9a7b3ae}:
\url{https://gitlab.com/ai-village-agents/village/graffiti-verification/-/blob/3be380b4222f974bedbefb228c8e75fbd9a7b3ae/data/wow2_rows_html.json}.
These references identify the formulas, not a current curator decision.

\bibitem{src:definitions}
E. DeLaVi\~na,
\emph{Written on the Wall II: Original invariant definitions},
\texttt{wowIIdefs.js}, Internet Archive snapshot of 12 May 2024,
13:36:36 UTC.
\url{https://web.archive.org/web/20240512133636id_/http://cms.uhd.edu/faculty/delavinae/research/wowII/wowIIdefs.js}.

\bibitem{early:dpw-independent}
E. DeLaVi\~na, R. Pepper, and B. Waller,
\emph{Graffiti.pc on the Independent-Domination Number of a Graph},
author-hosted manuscript, p.~8, Conjecture~3.2 and the preceding
degeneracy and set-eccentricity definitions.
\url{https://www.uhd.edu/documents/academics/sciences/dpw-cn2011.pdf}.

\bibitem{early:baseline}
KEIO, \emph{Induced Trees, Domination, and Clique Structure in Graffiti.pc:
A consolidated treatment of forty-four WOWII statements}, seventh
expanded edition, supplied unpublished manuscript, 26 October 2026,
Proposition~38.1, Corollary~38.2, Lemma~40.6, and Theorem~40.7.
The equality-quotient and zero-slack results used here are credited
antecedents in this research record.

\bibitem{early:glossary}
\emph{WoWII invariant glossary}, archived transcription,
\texttt{INVARIANT-GLOSSARY.md}, revision
\texttt{2455820ef9b925c0746cfb72561bca61105b1aa2}.
\url{https://github.com/Kuberwastaken/c5-k4/blob/2455820ef9b925c0746cfb72561bca61105b1aa2/data/INVARIANT-GLOSSARY.md}.
Used together with the original definitions in the source audits.

\bibitem{early:dpw-two}
E. DeLaVi\~na, C. E. Larson, R. Pepper, and B. Waller,
Graffiti.pc on the 2-domination number of a graph,
\emph{Congressus Numerantium} \textbf{203} (2010), 15--32.
\url{https://math1um.github.io/Research/r07.pdf}.

\bibitem{early:dlpww-total}
E. DeLaVi\~na, Q. Liu, R. Pepper, B. Waller, and D. B. West,
Some conjectures of Graffiti.pc on total domination,
\emph{Congressus Numerantium} \textbf{185} (2007), 81--95.
\url{https://www.uhd.edu/documents/academics/sciences/total-dom-2007.pdf}.

\bibitem{early:gebendorfer}
J. J. Gebendorfer,
\emph{A mixed transversal bound and the total-domination--annihilation
conjecture}, manuscript dated 11 October 2026, as attributed and
mathematically audited in the supplied Round~7 package, especially
Lemmas~2, 3, 7, 8, Proposition~4, Corollary~5, and Theorem~9.
The archived attribution identifies the attachment to
\url{https://github.com/google-deepmind/formal-conjectures/issues/7242}
and the archive SHA256
\path{d067e88bdbf41f4fd4405b61e7cc71acfc9e4a6a6e64d7ca8d55f35b289d2afc}.
The prior manuscript snapshots are not included in the supplied ZIPs;
the present text reconstructs the required arguments from
\cite{early:298-audit}.  The linked issue and associated Lean record
were not independently fetched or rebuilt for this consolidation.

\bibitem{early:298-audit}
\emph{Independent mathematical audit: WOWII 298 and the prescribed
$4k$-cycle theorem}, supplied research record,
\texttt{KEIO\_WOWII\_Round\_07/independent\_audit/AUDIT.md},
26 October 2026.  Includes the detailed proof audit of the credited
antecedents and of the prescribed-cycle deduction, together with
an independent finite checker and its archived output.

\bibitem{early:chvatal-mcdiarmid}
V. Chv\'atal and C. McDiarmid,
Small transversals in hypergraphs,
\emph{Combinatorica} \textbf{12} (1992), 19--26.
\url{https://doi.org/10.1007/BF01191201}.

\bibitem{early:henning-lowenstein}
M. A. Henning and C. L\"owenstein,
\emph{A characterization of hypergraphs that achieve equality in the
Chv\'atal--McDiarmid Theorem}, arXiv:1401.4851v1, 20 January 2014,
Section~1.
\url{https://arxiv.org/html/1401.4851v1}.

\bibitem{late:382e:ce}
Y. Cao and D. Eppstein, comments of 6 January 2018 on
\emph{Will core decomposition get a maximal clique?}, Theoretical Computer
Science Stack Exchange.
\href{https://cstheory.stackexchange.com/questions/39903/will-core-decomposition-get-a-maximal-clique}{Question 39903, comments and discussion}.

\bibitem{late:totaldom}
E. DeLaVi\~na, C. E. Larson, R. Pepper and B. Waller,
\emph{Graffiti.pc on the total domination number of a tree},
Congressus Numerantium \textbf{195} (2009), 5--18.
Feit's construction is credited on p.~7 of the
\href{https://www.uhd.edu/documents/academics/sciences/boca2009-dlpw.pdf}{author manuscript}.

\bibitem{late:DW}
E. DeLaVi\~na and B. Waller,
\emph{Spanning Trees with Many Leaves and Average Distance},
Electronic Journal of Combinatorics \textbf{15} (2008), R33.
\url{https://doi.org/10.37236/757}.

\bibitem{late:Minty}
G. J. Minty,
\emph{On maximal independent sets of vertices in claw-free graphs},
Journal of Combinatorial Theory, Series B \textbf{28} (1980), 284--304.
\href{https://doi.org/10.1016/0095-8956(80)90074-X}{doi:10.1016/0095-8956(80)90074-X}.
Cited for the classical alternating-path method, not as an attribution
of the exact intermediate lemma proved here.

\bibitem{late:forest:Hakimi}
S.~L. Hakimi,
\emph{On realizability of a set of integers as degrees of the vertices of a
linear graph. I},
Journal of the Society for Industrial and Applied Mathematics
\textbf{10} (1962), 496--506.
\href{https://doi.org/10.1137/0110037}{doi:10.1137/0110037}.

\bibitem{late:forest:WOW61}
E. DeLaVi\~na,
\emph{Written on the Wall II}, Conjecture 61.
Recorded mathematical formulation in the Formal Conjectures project,
\texttt{GraphConjecture61}, consulted 26 October 2026:
\href{https://google-deepmind.github.io/formal-conjectures/src/FormalConjectures/WrittenOnTheWallII/GraphConjecture61/}{Formal Conjectures: GraphConjecture61}.

\bibitem{late:forest:Triesch}
E. Triesch,
\emph{Degree sequences of graphs and dominance order},
Journal of Graph Theory \textbf{22} (1996), 89--93.
The precise elimination-sequence theorem used here is restated with attribution
in \cite[Theorem 1.6]{late:forest:Hiller}.

\bibitem{late:forest:Hiller}
M. Hiller,
\emph{On the domination order among elimination sequences},
arXiv:2206.05820 (2022).
\url{https://arxiv.org/abs/2206.05820}.

\bibitem{late:forest:GK}
J.~R. Griggs and D.~J. Kleitman,
\emph{Independence and the Havel--Hakimi residue},
Discrete Mathematics \textbf{127} (1994), 209--212.
\url{https://www.sciencedirect.com/science/article/pii/0012365X9200479B}.

\bibitem{late:forest:ADP}
D. Amos, R. Davila, and R. Pepper,
\emph{On the $k$-residue of disjoint unions of graphs with applications to
$k$-independence},
Discrete Mathematics \textbf{321} (2014), 24--34.
\href{https://doi.org/10.1016/j.disc.2013.12.013}{doi:10.1016/j.disc.2013.12.013}.

\bibitem{late:FMS}
O. Favaron, M. Mah\'eo and J.-F. Sacl\'e,
\emph{On the residue of a graph},
Journal of Graph Theory \textbf{15} (1991), 39--64.
\href{https://doi.org/10.1002/jgt.3190150107}{doi:10.1002/jgt.3190150107}.

\bibitem{late:unigraph:Barrus}
M.~D. Barrus,
\emph{Havel--Hakimi residues of unigraphs},
Information Processing Letters \textbf{112} (2012), 44--48.
\href{https://doi.org/10.1016/j.ipl.2011.10.011}{doi:10.1016/j.ipl.2011.10.011}.
Author manuscript:
\url{https://math.uri.edu/~barrus/papers/UniRes_v1.pdf}.

\bibitem{late:unigraph:Tyshkevich}
R. Tyshkevich,
\emph{Decomposition of graphical sequences and unigraphs},
Discrete Mathematics \textbf{220} (2000), 201--238.
\href{https://doi.org/10.1016/S0012-365X(99)00381-7}{doi:10.1016/S0012-365X(99)00381-7}.

\bibitem{late:Gunderson}
K. Gunderson, K. Meagher, J. Morris and V. R. T. Pantangi,
\emph{Induced forests in some distance-regular graphs},
arXiv:2301.05207; author manuscript of 12 September 2023, Section 5.
\url{https://www.cs.uleth.ca/~morris/Research/CRG-1.pdf}.
\url{https://doi.org/10.1016/j.dam.2023.12.003}.

\bibitem{late:equality:LM}
V.~E. Levit and E. Mandrescu,
\emph{Matrices and $\alpha$-stable bipartite graphs},
arXiv:math/9912120v1 (1999), Theorem~2.5.
\url{https://arxiv.org/abs/math/9912120}.

\bibitem{late:apex:GHHL}
W. Goddard, S.~M. Hedetniemi, S.~T. Hedetniemi and R. Laskar,
\emph{Generalized subgraph-restricted matchings in graphs},
Discrete Mathematics \textbf{293} (2005), 129--138.
\href{https://doi.org/10.1016/j.disc.2004.08.027}{doi:10.1016/j.disc.2004.08.027}.

\bibitem{late:apex:BR}
J. Baste and D. Rautenbach,
\emph{Degenerate Matchings and Edge Colorings},
arXiv:1702.02358 (2017).
\url{https://arxiv.org/abs/1702.02358}.

\bibitem{late:apex:Fajtlowicz}
S. Fajtlowicz, \emph{Written on the Wall}, July 2004 version, p.~37.
\url{https://independencenumber.wordpress.com/wp-content/uploads/2012/08/wow-july2004.pdf}.

\bibitem{late:apex:FR}
M. F\"urst and D. Rautenbach,
\emph{On some hard and some tractable cases of the maximum acyclic
matching problem}, Annals of Operations Research \textbf{279} (2019),
291--300.
\href{https://doi.org/10.1007/s10479-019-03311-1}{doi:10.1007/s10479-019-03311-1}.
Preprint arXiv:1710.08236 (2017), Conclusion:
\url{https://arxiv.org/abs/1710.08236}.

\bibitem{src:bh}
A. Brandst\"adt and C. T. Ho\`ang,
On clique separators, nearly chordal graphs, and the Maximum Weight Stable
Set Problem, \emph{Theoretical Computer Science} \textbf{389} (2007),
295--306. \url{https://doi.org/10.1016/j.tcs.2007.09.031}.

\bibitem{src:ghp}
S. Gaspers, S. Huang, and D. Paulusma,
Colouring Square-Free Graphs without Long Induced Paths,
\emph{STACS 2018}, article 35, 35:1--35:15.
\url{https://doi.org/10.4230/LIPIcs.STACS.2018.35}.

\bibitem{src:cwxpre}
R. Chen, D. Wu, and B. Xu,
\emph{Structure and coloring of some $(P_7,C_4)$-free graphs},
arXiv:2302.11115v1, 22 February 2023, Theorem 1.3.
\url{https://arxiv.org/abs/2302.11115v1}.

\bibitem{src:cwx}
R. Chen, D. Wu, and B. Xu,
Structure of some $(P_7,C_4)$-free graphs with application to colorings,
\emph{Discrete Applied Mathematics} \textbf{357} (2024), 14--23.
\url{https://doi.org/10.1016/j.dam.2024.05.038}.

\bibitem{src:ess}
P. Erd\H{o}s, M. Saks, and V. T. S\'os,
Maximum induced trees in graphs,
\emph{Journal of Combinatorial Theory, Series B} \textbf{41} (1986),
61--79. The radius argument is credited there to Fan Chung.
\url{https://users.renyi.hu/~sos/1986_Maximum_Induced_Trees_in_Graphs.pdf}.

\bibitem{src:hof}
P. van 't Hof and D. Paulusma,
A new characterization of $P_6$-free graphs,
\emph{Discrete Applied Mathematics} \textbf{158} (2010), 731--740,
Theorem 6. Primary author manuscript:
\url{https://pimvanthof.github.io/p6free.pdf}.

\bibitem{src:camby}
E. Camby and O. Schaudt,
A new characterization of $P_k$-free graphs,
\emph{Algorithmica} \textbf{75} (2016), 205--217.
\url{https://doi.org/10.1007/s00453-015-9989-6}.
Exact version used: \url{https://arxiv.org/abs/1402.7213v1}, Theorem~5.

\bibitem{src:fajtlowicz}
S. Fajtlowicz,
A characterization of radius-critical graphs,
\emph{J. Graph Theory} \textbf{12}(4) (1988), 529--532.
\url{https://doi.org/10.1002/jgt.3190120409}.

\bibitem{src:dlpw}
E. DeLaVi\~na, C. E. Larson, R. Pepper, and B. Waller,
A characterization of graphs where the independence number equals the
radius, \emph{Graphs Combin.} \textbf{28} (2012), 315--332.
\url{https://doi.org/10.1007/s00373-011-1054-x}.

\bibitem{src:heawoodfamily}
R. Hanaki, R. Nikkuni, K. Taniyama, and A. Yamazaki,
On intrinsically knotted or completely 3-linked graphs,
\emph{Pacific J. Math.} \textbf{252} (2011), 407--425.
\url{https://doi.org/10.2140/pjm.2011.252.407}.
Author version: \url{https://arxiv.org/abs/1006.0698}.

\bibitem{src:cs}
M. Chudnovsky and J. Stacho,
3-colorable subclasses of $P_8$-free graphs,
\emph{SIAM J. Discrete Math.} \textbf{32}(2) (2018), 1111--1138.
\url{https://doi.org/10.1137/16M1104858}.
Author-hosted primary text:
\url{https://web.math.princeton.edu/~mchudnov/p8free1.pdf}.

\bibitem{paired:DGH}
E.~DeLaVi\~na, W.~Goddard, M.~A.~Henning, R.~Pepper, and E.~R.~Vaughan,
\emph{Bounds on the $k$-domination number of a graph},
Applied Mathematics Letters \textbf{24} (2011), 996--998.
\href{https://doi.org/10.1016/j.aml.2011.01.013}{doi:10.1016/j.aml.2011.01.013}.
Theorems 3 and 5; the specialization of Theorem 3 used here is proved
in Lemma~\ref{paired:matching-bound}.

\bibitem{paired:archive}
KEIO, \emph{WoWII research rounds 1--22: supplied manuscripts,
verification programs, certificates, and independent-audit records},
unpublished research archive, supplied for this consolidation (2026).
For 384b the proof sources are Round 5, partial-results reduction;
Rounds 8--11 and 14--22, 384b manuscripts; and Round 10, cubic-filter
supplement. Exact file paths, SHA-256 digests, dependencies, and rerun
records are provided in \nolinkurl{paired_384b_audit.json} and the
companion verification directory.

\bibitem{src:ce}
V.~Chv\'atal and P.~Erd\H{o}s,
A note on Hamiltonian circuits,
\emph{Discrete Mathematics} \textbf{2} (1972), 111--113.
\url{https://doi.org/10.1016/0012-365X(72)90079-9}.
Author archive scan:
\url{https://combinatorica.hu/~p_erdos/1972-02.pdf}.

\bibitem{src:bc}
J.~A.~Bondy and V.~Chv\'atal,
A method in graph theory,
\emph{Discrete Mathematics} \textbf{15} (1976), 111--135.
\url{https://doi.org/10.1016/0012-365X(76)90078-9}.

\bibitem{src:hall}
P.~Hall, On representatives of subsets,
\emph{Journal of the London Mathematical Society}
\textbf{10} (1935), 26--30.
\url{https://doi.org/10.1112/jlms/s1-10.37.26}.

\bibitem{src:fordfulkerson}
L. R. Ford, Jr. and D. R. Fulkerson, Maximal flow through a network,
\emph{Canadian Journal of Mathematics} \textbf{8} (1956), 399--404.
\url{https://doi.org/10.4153/CJM-1956-045-5}.

\bibitem{even:pr70}
Public verification materials for WOWII Conjecture 189, PR~70,
\emph{agentic-conjectures}, pinned commit
\texttt{8c4288df412f8aecadfb6998d0182abb42c21aea}.
The archived $d_2=4$ certificate list is credited as the source of the
twelve-vertex example and the discovery input for $J_s$.
\url{https://github.com/umaia1234/agentic-conjectures/pull/70}.
\url{https://raw.githubusercontent.com/umaia1234/agentic-conjectures/8c4288df412f8aecadfb6998d0182abb42c21aea/problems/wowii-graph-conjecture-189/sigma4_certificates.txt}.

\bibitem{even:ej}
R. Easwar and S. Johari,
\emph{Conditions for traceability under a bound on the size of
even-distance sets}, arXiv:2609.11719v1, 10 September 2026.
\url{https://arxiv.org/abs/2609.11719v1}.
\end{thebibliography}
\end{document}
